% 本文件由 LaTeX 排版 Agent 根据有效 translation.json 自动生成。
% author=Augustine of Hippo; work_id=1602; title=福音的和谐

\chapter{福音的和谐}

% structure: p0001 source=h1 target=omitted reason=page-title-already-rendered-by-parent

% structure: p0002 source=h2 target=section reason=relative-heading-rank; base=h2

\section{卷一}

\Emph{权威、数目、次序和福音的计划。对攻击福音的回应。}\par

\Emph{第一章 论福音的权威。}\newline{}\Emph{第二章 论圣史家们的次序以及他们写作的原则。}\newline{}\Emph{第三章 论玛窦与马尔谷特别着眼基督君王性而路加处理司祭性这一事实。}\newline{}\Emph{第四章 论若望着手阐述基督的天主性这一事实。}\newline{}\Emph{第五章 论两种美德：若望默观性，其他圣史家行动性。}\newline{}\Emph{第六章 论\WorkTitle{默示录}中的四种活物被一些人用作、另一些人用作不同的，作为四位圣史家的合适形象。}\newline{}\Emph{第七章 陈述圣奥思定着手这部\WorkTitle{福音的和谐}工作的理由，以及他如何应对那些声称基督本人未写任何东西、而祂的门徒在宣认祂为天主时作了无根据肯定的人的例子。}\newline{}\Emph{第八章 论为何若基督因普通叙事报告的见证而被相信为最智慧的人，祂就不应当因更优越的宣讲报告的见证而被相信为天主这一问题。}\newline{}\Emph{第九章 论某些假装基督写过关于巫术之书的人。}\newline{}\Emph{第十章 论一些疯狂到假设这些书卷以伯多禄和保禄之名题署的人。}\newline{}\Emph{第十一章 反对那些愚昧地幻想基督用巫术使人皈依自己的。}\newline{}\Emph{第十二章 论犹太人的天主，在该民族被征服后，仍未被罗马人接纳，因为祂的诫命是唯独祂当受崇拜、偶像当被摧毁这一事实。}\newline{}\Emph{第十三章 论天主为何容许犹太人被降服的问题。}\newline{}\Emph{第十四章 论希伯来人的天主，尽管该民族被征服，却藉着推翻偶像并引导所有外邦人归向祂自己的服务而证明自己是不被征服的这一事实。}\newline{}\Emph{第十五章 论异教徒被迫赞美基督时，却向祂的门徒发泄侮辱这一事实。}\newline{}\Emph{第十六章 论宗徒们在摧毁偶像的主题上所教导的与基督或先知所教导的并无不同这一事实。}\newline{}\Emph{第十七章 反对那些唯独拒绝以色列天主的罗马人。}\newline{}\Emph{第十八章 论希伯来人的天主不被罗马人接纳，因为祂的旨意是唯独祂应受崇拜这一事实。}\newline{}\Emph{第十九章 证明此天主是真天主。}\newline{}\Emph{第二十章 论未发现异教徒的先知预言任何反对希伯来天主的事这一事实。}\newline{}\Emph{第二十一章 一个论证尽力唯独崇拜此天主——祂在禁止崇拜其他神祇的同时，自己并不被其他神祇禁止被崇拜。}\newline{}\Emph{第二十二章 论外邦人对我们天主的看法。}\newline{}\Emph{第二十三章 论异教徒关于朱庇特和萨图恩所沉溺的愚行。}\newline{}\Emph{第二十四章 论那些拒绝以色列天主的人，结果未能崇拜所有神祇；另一方面，那些崇拜其他神祇的人也未能崇拜祂这一事实。}\newline{}\Emph{第二十五章 论假神并不禁止其他神与他们一起被崇拜这一事实。以色列的天主是真天主，由祂在预言和应验中的工作所证明。}\newline{}\Emph{第二十六章 论偶像崇拜被基督的名字和基督徒的信德按照预言所颠覆这一事实。}\newline{}\Emph{第二十七章 一个敦促残余偶像崇拜者最终成为此真天主仆人的论证，祂正在各处推翻偶像。}\newline{}\Emph{第二十八章 论预言中的偶像被拒弃。}\newline{}\Emph{第二十九章 论为何异教徒应当拒绝崇拜以色列的天主，即使他们认为祂只是掌管元素的神？}\newline{}\Emph{第三十章 论随着预言已应验，以色列的天主现已处处被人认识这一事实。}\newline{}\Emph{第三十一章 论关于基督的预言的应验。}\newline{}\Emph{第三十二章 用预言的话语为宗徒的教义辩护反对偶像崇拜的陈述。}\newline{}\Emph{第三十三章 反对那些抱怨人类生命的幸福因基督徒时代的来临而受损之人的陈述。}\newline{}\Emph{第三十四章 以上各章跋语。}\newline{}\Emph{第三十五章 论中保的奥秘藉着预言向古代人启示，正如现今在福音中向我们宣告这一事实。}\par

% structure: p0005 source=h2 target=section reason=relative-heading-rank; base=h2

\section{卷二}

\Emph{\BibleBook{玛窦}{福音}与其他福音的和谐，直至最后晚餐。}\par

\Emph{序言} 序言。\newline{}\Emph{第一章} 说明为何基督的祖先名录一直列到若瑟，而基督并非由那人的种子所生，而是由童贞玛利亚所生。\newline{}\Emph{第二章} 解释基督是达味之子的意义，尽管他并非由达味之子若瑟按普通生育方式所生。\newline{}\Emph{第三章} 说明为何玛窦为基督列举了一系祖先，而路加列举了另一系。\newline{}\Emph{第四章} 说明为何在玛窦福音中（不计基督本人）有四十代，尽管他将这些世代分为三组，每组十四代。\newline{}\Emph{第五章} 说明路加是如何在有关基督的受孕、婴孩期或童年的那些被一位记述而另一位省略的事上，证明与玛窦和谐一致。\newline{}\Emph{第六章} 论所有四位福音书作者对洗者若翰宣讲的记法位置。\newline{}\Emph{第七章} 论两位黑落德。\newline{}\Emph{第八章} 解释玛窦的陈述，即若瑟因害怕阿赫劳斯而不敢带婴孩基督去耶路撒冷，却不怕去加里肋亚，那里是王子之弟黑落德做分封侯的地方。\newline{}\Emph{第九章} 解释玛窦说若瑟带婴孩基督去加里肋亚的原因是他害怕阿赫劳斯（那时在耶路撒冷继位），而路加告诉我们去加里肋亚的原因是他们的城纳匝肋在那里。\newline{}\Emph{第十章} 说明为什么路加告诉我们，\Quote{他的父母每年逾越节都上耶路撒冷去}，与孩子一起，而玛窦暗示他们因害怕阿赫劳斯而从埃及回来后不敢去那里。\newline{}\Emph{第十一章} 探讨这样一个问题：如果玛窦正确地记载了黑落德已从贤士那里得知基督诞生，并且在他搜寻时杀害了许多孩子，那么按照路加的说法，他们怎么可能在基督母亲的取洁日届满时，带着他上耶路撒冷圣殿去履行常规礼仪呢？\newline{}\Emph{第十二章} 关于四位福音书作者各自归给若翰的话。\newline{}\Emph{第十三章} 论耶稣的圣洗。\newline{}\Emph{第十四章} 论耶稣受洗后来自天上的声音或话语。\newline{}\Emph{第十五章} 解释为何按若望福音，洗者若翰说，\Quote{我不认识他；}，而按其他福音书作者，他确实已经认识他。\newline{}\Emph{第十六章} 论耶稣的受试探。\newline{}\Emph{第十七章} 论宗徒们在捕鱼时蒙召。\newline{}\Emph{第十八章} 论他前往加里肋亚的日期。\newline{}\Emph{第十九章} 论玛窦所记载的他在山上所作的长篇讲道。\newline{}\Emph{第二十章} 解释玛窦告诉我们百夫长亲自来求耶稣医治他的仆人，而路加则说百夫长打发朋友去求耶稣。\newline{}\Emph{第二十一章} 论有关伯多禄岳母的叙述是如何引出的。\newline{}\Emph{第二十二章} 论本章之后所记载的事件顺序，以及玛窦、马尔谷和路加在这些事上是否一致。\newline{}\Emph{第二十三章} 论那人对主说，\Quote{不论你往哪里去，我要跟随你；}的事，以及与之相关的其他事情，和玛窦与路加记载这些事的顺序。\newline{}\Emph{第二十四章} 论主在那次在船上睡觉时渡过湖去，并驱逐那些被他准许进入猪群的魔鬼；以及玛窦、马尔谷和路加对这些事所行所说的记载的一致性。\newline{}\Emph{第二十五章} 论那个瘫痪的人，主对他说，\Quote{你的罪赦了，}和\Quote{拿起你的床来；}以及特别是关于玛窦和马尔谷在记载这事发生的地点是否一致的问题，因为玛窦说发生在\Quote{他自己的城}，而马尔谷说在葛法翁。\newline{}\Emph{第二十六章} 论玛窦的蒙召，以及玛窦自己的叙述是否与马尔谷和路加提到肋未（阿尔斐的儿子）的记载一致。\newline{}\Emph{第二十七章} 论那个宴席上同时有人指责基督与罪人一同吃饭，以及他的门徒不禁食；论到这些指责是由谁提出的，福音书作者似乎给出不同的叙述；以及玛窦、马尔谷和路加在这些人所说的话和主的答复的报告上是否也彼此和谐一致。\newline{}\Emph{第二十八章} 论会堂主管的女儿复活和摸他衣边的妇人；以及这些事被记载的顺序是否在任何记录这些事的作者中显出矛盾；特别是会堂主管向主提出请求时所用的言辞。\newline{}\Emph{第二十九章} 论那两个盲人和哑巴附魔者，这些故事仅有玛窦记载。\newline{}\Emph{第三十章} 论那段记载，他怜悯群众，派遣门徒，赐给他们治病的能力，并给他们许多教训，指导他们如何生活；并论到玛窦在此与马尔谷和路加和谐一致的证据，特别是关于手杖的事，玛窦说主告诉他们不要带，而马尔谷说那是唯一要带的东西；以及关于穿鞋和穿外衣的事。\newline{}\Emph{第三十一章} 论玛窦和路加记载的洗者若翰在狱中打发门徒去见主的事。\newline{}\Emph{第三十二章} 论他责备那些不悔改的城邑的事，路加和玛窦都记载了；以及玛窦在叙述顺序上与路加和谐一致的问题。\newline{}\Emph{第三十三章} 论他呼召他们负他的轭和担子，以及玛窦和路加在叙述顺序上没有不一致的问题。\newline{}\Emph{第三十四章} 论门徒掐麦穗吃的那段经文；以及玛窦、马尔谷和路加在叙述顺序上如何彼此和谐一致。\newline{}\Emph{第三十五章} 论那个枯干了一只手的人在安息日被治好；以及玛窦的记述如何能与马尔谷和路加的相和谐，无论是在事件顺序上，还是在主和犹太人所说的话的报告中。\newline{}\Emph{第三十六章} 论另一个需要思考的问题，即这三位福音书作者在从枯干手被治好的人的故事转向下一主题时，是否在叙述顺序上产生了矛盾。\newline{}\Emph{第三十七章} 论玛窦和路加关于那个又哑又瞎的附魔者的记载的一致性。\newline{}\Emph{第三十八章} 论有人对他说他凭贝耳则步驱魔，以及由此引出的关于亵渎圣神和两棵树的言论；以及在这些段落中玛窦与其他两位福音书作者之间，特别是玛窦与路加之间，是否存在不一致。\newline{}\Emph{第三十九章} 论玛窦与路加在主回答那些求征兆的人的话（说到约纳先知、尼尼微人、南方女王，以及那污灵出去后又返回的比喻）中如何一致。\newline{}\Emph{第四十章} 论当有人向他报告他的母亲和兄弟时，玛窦与马尔谷和路加在记载此事的顺序上是否存在不一致。\newline{}\Emph{第四十一章} 论他在船上所说的关于撒种的比喻，种子撒下后落在路旁等；以及关于那人在麦子中撒下莠子；关于芥菜种和面酵；还有他在家里所说的关于藏在地里的宝贝、珍珠和撒在海里的网，以及那从库里拿出新旧东西的人；并证明玛窦与马尔谷和路加在共同记载的事上以及在叙述顺序上如何和谐一致。\newline{}\Emph{第四十二章} 论他来到自己的家乡，人们因他的教训而惊奇，藐视他的出身；论玛窦在此与马尔谷和路加和谐一致；特别是第一位福音书作者所呈现的叙述顺序是否与其他两位不完全一致。\newline{}\Emph{第四十三章} 论玛窦、马尔谷和路加所记载的黑落德听到主的奇妙作为后所说的话彼此一致，以及他们在叙述顺序上的和谐。\newline{}\Emph{第四十四章} 论若翰被监禁和死亡的故事在这三位福音书作者中的记载顺序。\newline{}\Emph{第四十五章} 论四位福音书作者在记述五饼奇蹟时的顺序和方法。\newline{}\Emph{第四十六章} 论四位福音书作者在同一个五饼奇蹟的主题上如何彼此和谐。\newline{}\Emph{第四十七章} 论他在水面上行走，以及关于记载这一幕的福音书作者之间的和谐问题，以及他们如何从用五个饼喂饱群众的段落过渡。\newline{}\Emph{第四十八章} 论玛窦和马尔谷一方与若望一方，在到达湖的另一边后所发生的事的记载中没有任何矛盾。\newline{}\Emph{第四十九章} 论那个客纳罕妇人说，\Quote{可是小狗也吃主人桌子上掉下来的碎屑}，以及玛窦和路加的记载之间的和谐。\newline{}\Emph{第五十章} 论他用七个饼喂饱群众的事，以及玛窦和马尔谷在记载这个奇蹟上如何和谐。\newline{}\Emph{第五十一章} 论玛窦说离开那些地方后他来到玛加丹境内；以及他在这点与马尔谷的一致，并关于约拿的言论，那再次作为对求征兆之人的回答。\newline{}\Emph{第五十二章} 论玛窦与马尔谷在关于法利塞人的酵母的陈述上的一致，无论在主题本身还是在叙述顺序上。\newline{}\Emph{第五十三章} 论他问门徒人们说他是谁的事；并且论到无论是主题还是顺序，玛窦、马尔谷和路加之间是否有任何不一致。\newline{}\Emph{第五十四章} 论他向门徒预言即将来临的受难，以及玛窦、马尔谷和路加在记载此事上的一致程度。\newline{}\Emph{第五十五章} 论三位福音书作者在后来记载主命令那愿意跟随他的人来跟从他的方式上的一致。\newline{}\Emph{第五十六章} 论主与梅瑟和厄里亚一起在山上的显现；以及前三位福音书作者在这件事的顺序和情况上的一致性问题；特别是天数，玛窦和马尔谷说六天后，路加说八天后。\newline{}\Emph{第五十七章} 论玛窦和马尔谷在记载主对门徒谈论厄里亚的来临时的一致。\newline{}\Emph{第五十八章} 论那人带他的儿子来见耶稣，门徒们不能治好他；以及这三位福音书作者在此处的叙述顺序上的一致。\newline{}\Emph{第五十九章} 论当主向他们谈论他的受难时，门徒们极其忧愁的事，三位福音书作者以同样的顺序记载。\newline{}\Emph{第六十章} 论他从鱼口中拿出钱来交税，这事只有玛窦记载。\newline{}\Emph{第六十一章} 论他叫一个小孩来站在他们中间作为效法的榜样，以及世界的绊脚石；论那引人犯罪的肢体；论小孩子的天使常见天父的面；论那一百只羊中的一只；论私下劝戒弟兄；论释放和束缚罪；论两人的同心合意和三人的聚集；论饶恕七十个七次；论那仆人自己的大债被豁免却不肯豁免同伴的小债；并论玛窦与其他福音书作者在这些主题上的和谐。\newline{}\Emph{第六十二章} 论玛窦和马尔谷在记载有人问他是否可以休妻时的和谐，特别是关于主和犹太人之间具体的问答，而福音书作者似乎有些微分歧。\newline{}\Emph{第六十三章} 论他按手在小孩子身上；论那富有人，他对他说，\Quote{变卖你所有的一切；}论在不同时辰雇工人的葡萄园；并论玛窦与其他两位福音书作者在这些主题上没有不一致。\newline{}\Emph{第六十四章} 论他私下向门徒预言他的受难；以及载伯德儿子的母亲带着儿子来请求让一个坐在他右边，一个坐在他左边；以及玛窦与其他两位福音书作者在这些主题上没有不一致。\newline{}\Emph{第六十五章} 论玛窦与马尔谷之间，或玛窦与路加之间，在记载耶里哥的瞎子得看见的事上没有对立。\newline{}\Emph{第六十六章} 论玛窦提到的小驴驹，以及他的记载与其他只提到驴的福音书作者的一致。\newline{}\Emph{第六十七章} 论驱逐圣殿中买卖之人的事，以及前三位福音书作者与若望之间的和谐，若望在一个非常不同的关联中记载了同一件事。\newline{}\Emph{第六十八章} 论无花果树枯萎，以及玛窦与其他福音书作者在这事的记载上没有任何矛盾，以及与它相关的事；特别是玛窦与马尔谷在叙述顺序上的一致。\newline{}\Emph{第六十九章} 论前三位福音书作者在犹太人问主凭什么权柄做这些事时的一致。\newline{}\Emph{第七十章} 论父亲命令两个儿子去葡萄园做工，以及葡萄园租给其他园户；论玛窦对这些段落的版本与其他两位保持同样顺序的福音书作者的一致性；特别论到那三个福音书作者都记载的关于租出葡萄园的比喻的和谐，以及特别论到听比喻的人的回答，玛窦似乎与其他有些不同。\newline{}\Emph{第七十一章} 论王的婚宴，许多人被邀请；以及玛窦引入这段的次序与路加的比较，路加在另一个关联中给了我们一个类似的叙述。\newline{}\Emph{第七十二章} 论这三位福音书作者关于将带凯撒像的钱币归还凯撒，以及那嫁过七个兄弟的妇人的记载中的和谐特征。\newline{}\Emph{第七十三章} 论那被嘱咐爱天主和爱邻舍两条诫命的人；以及玛窦和马尔谷遵守的叙述顺序，以及他们与路加没有不一致的问题。\newline{}\Emph{第七十四章} 论犹太人被问他们认为基督是谁的儿子的段落；以及玛窦与其他两位福音书作者之间是否有不一致，他提出问题为\Quote{你们对基督的看法如何？他是谁的儿子？}并告诉我们他们回答\Quote{达味的儿子}，而其他人则问\Quote{经师们怎么说基督是达味的儿子？}\newline{}\Emph{第七十五章} 论那些坐在梅瑟位子上的法利塞人，他们吩咐人却不行，以及主对这些法利塞人的其他斥责；论玛窦的叙述是否在此与其他两位福音书作者一致，特别是与路加，他引入了一段类似的经文，尽管不是按这个顺序，而是在另一个关联中。\newline{}\Emph{第七十六章} 论玛窦与其他两位福音书作者在预言圣殿毁灭的事的叙述顺序上的一致。\newline{}\Emph{第七十七章} 论三位福音书作者在记载主在橄榄山上门徒问何时发生终局时的讲论中的和谐。\newline{}\Emph{第七十八章} 论玛窦和马尔谷一方与若望一方之间是否有矛盾，因为前者说逾越节是两天后，然后说他是在伯达尼，而后者给出了伯达尼发生的平行叙述，但提到那是在逾越节前六天。\newline{}\Emph{第七十九章} 论玛窦、马尔谷和若望在记载伯达尼的晚餐上的一致，那妇人在晚餐中把珍贵的香膏倒在主身上，以及这些记载如何与路加在另一个时期记载的类似事件调和。\newline{}\Emph{第八十章} 论玛窦、马尔谷和路加所记载的，他打发门徒去预备吃逾越节的地方时的一致特征。\par

% structure: p0008 source=h2 target=section reason=relative-heading-rank; base=h2

\section{第三卷}

\Emph{玛窦福音与其他福音的协调，从最后晚餐开始}\par

\Emph{序言} 序言。\newline{}\Emph{第一章} 论四位福音作者在记述主的晚餐及指出出卖祂的人上显示一致的方法。\newline{}\Emph{第二章} 论在提及预言伯多禄否认的记载中证明他们没有任何分歧。\newline{}\Emph{第三章} 论如何可以证明他们在记述主所说的话上（直至祂离开他们晚餐的房屋时）之间不存在分歧。\newline{}\Emph{第四章} 论晚餐后离开房屋来到的一块地或园子里所发生的事；以及论在若望对此事沉默的情况下，如何可以证明其他三位福音作者——即玛窦、马尔谷和路加——之间真正的和谐。\newline{}\Emph{第五章} 论四位福音作者就祂被捕时所行所言的一切所作的记述；以及论证明这些不同叙述并无真正分歧。\newline{}\Emph{第六章} 论这些福音作者记述主被带到大司祭家中所发生的事的和谐，以及论祂在夜间被带到那屋里后所发生的事，特别是伯多禄否认的事件。\newline{}\Emph{第七章} 论福音作者们在清晨、把耶稣解送比拉多之前所发生的事的不同记述中的完全和谐；以及论关于为主所定的价格而被引用的那段经文的问题，玛窦将其归于耶肋米亚，但那位先知的著作中却找不到这样的段落。\newline{}\Emph{第八章} 论福音作者们记述比拉多面前所发生的事中没有分歧。\newline{}\Emph{第九章} 论祂在比拉多的卫队手中所受的嘲弄，以及论报告那场景的三位福音作者——即玛窦、马尔谷和若望——之间的和谐。\newline{}\Emph{第十章} 论我们如何能够调和玛窦、马尔谷和路加所作的声明（即另一个人被强逼来背耶稣的十字架）与若望的声明（他说耶稣自己背着十字架）。\newline{}\Emph{第十一章} 论玛窦的版本与马尔谷的版本在记述给祂喝的苦酒方面的一致性，这记述出现在祂被钉十字架的叙述之前。\newline{}\Emph{第十二章} 论所有四位福音作者在分祂衣服的事上保存的一致。\newline{}\Emph{第十三章} 论主受难的时辰，以及论关于马尔谷和若望在\Quote{第三}时辰和\Quote{第六}时辰这一点上没有分歧的问题。\newline{}\Emph{第十四章} 论所有福音作者在与祂同钉的两个强盗的事上保存的和谐。\newline{}\Emph{第十五章} 论玛窦、马尔谷和路加在侮辱主的人的事上记述的一致性。\newline{}\Emph{第十六章} 论归给强盗的讥讽，以及论关于玛窦和马尔谷一方与路加之间没有分歧的问题，因为路加说其中一个强盗讥笑祂，另一个相信了祂。\newline{}\Emph{第十七章} 论四位福音作者在醋的记述上的和谐。\newline{}\Emph{第十八章} 论主临终时连续的话语；以及论玛窦和马尔谷是否与路加在这些话语的报告上和谐，以及这三位福音作者是否与若望和谐的问题。\newline{}\Emph{第十九章} 论圣所帐幔的撕裂，以及论玛窦和马尔谷是否真正与路加就那事件发生的顺序和谐的问题。\newline{}\Emph{第二十章} 论玛窦、马尔谷和路加关于百夫长和与他在一起的人感到惊奇的各种记述的一致性问题。\newline{}\Emph{第二十一章} 论站在那里的一些妇女，以及论玛窦、马尔谷和路加（他们说她们远远站着）是否与若望（他说其中一位站在十字架旁）矛盾的问题。\newline{}\Emph{第二十二章} 论福音作者们在关于向比拉多求主身体的若瑟的叙述上是否完全一致，以及若望的版本是否包含互相矛盾的陈述的问题。\newline{}\Emph{第二十三章} 论前三位福音作者是否与若望在祂安葬的记述上完全和谐的问题。\newline{}\Emph{第二十四章} 论四位福音作者在关于主复活前后所发生的事件的叙述中没有分歧。\newline{}\Emph{第二十五章} 论基督后来向门徒们的显现，以及论当四位不同福音作者的记载以及宗徒保禄和宗徒大事录所呈现的记载放在一起比较时，是否能够建立完全和谐的问题。\par

% structure: p0011 source=h2 target=section reason=relative-heading-rank; base=h2

\section{第四卷}

\Emph{马尔谷、路加或若望独有的段落}\par

\Emph{序言} 序言。\newline{}\Emph{第一章} 论有关证明马尔谷福音与他所叙述的其他内容和谐一致的问题（忽略他与玛窦共有的部分），从其开始直到记载说\Quote{他们去了\Place{葛法翁}，安息日他立刻教导他们}的地方；路加也报告了这一事件。\newline{}\Emph{第二章} 论从那人身上赶出折磨他的邪魔，以及关于\Person{马尔谷}的版本是否与\Person{路加}的版本完全一致的问题，\Person{路加}在报告这一事件上与他一致。\newline{}\Emph{第三章} 论\Person{马尔谷}对多次提到\Person{伯多禄}之名的记载是否与\Person{若望}告诉我们的、宗徒得到那个名字的具体时间不一致的问题。\newline{}\Emph{第四章} 论这些话：\Quote{他越吩咐他们不要告诉任何人，他们越发传扬出去；}以及关于这句话是否与他的预知（福音中提醒我们注意的）不一致的问题。\newline{}\Emph{第五章} 论\Person{若望}关于那个赶出魔鬼但并非门徒圈子中的人所作的陈述；以及主的回答：\Quote{不要禁止他们，因为不反对你们的，就是帮助你们的；}以及关于这个回答是否与另一句话，他说：\Quote{不与我一起的，就是反对我。}相矛盾的问题。\newline{}\Emph{第六章} 论\Person{马尔谷}记录了比\Person{路加}更多的、主就那以基督之名赶魔鬼（尽管他不跟随门徒）之人所说的话；以及关于如何证明这些附加的话与基督禁止阻止那奉祂名行奇迹之人的用意真正相关的问题。\newline{}\Emph{第七章} 论从这一点到主的晚餐（四位福音书作者的叙述的共同讨论从这一事件开始）之间，马尔谷福音没有提出需要特别考察的问题。\newline{}\Emph{第八章} 论路加福音，尤其是其开头与\WorkTitle{宗徒大事录}开头之间的和谐。\newline{}\Emph{第九章} 论如何证明路加所记载的鱼获事件不与\Person{若望}在主复活后记载的明显相似的事件视为同一；以及从这一点到主的晚餐（从这一事件到结束，所有福音书作者的记载已被联合考察）之间，路加福音和马尔谷福音一样，没有出现需要特别考虑的困难。\newline{}\Emph{第十章} 论福音书作者\Person{若望}，以及他与另外三位之间的区别。\par

\SourceRecord{Translated by S.D.F. Salmond.；Nicene and Post-Nicene Fathers, First Series；Vol. 6.；Edited by Philip Schaff.；Buffalo, NY: Christian Literature Publishing Co.,；1888.；<http://www.newadvent.org/fathers/1602.htm>.}

% structure: p0001 source=h1 target=section reason=page-title

\section{\WorkTitle{福音的和谐，卷一}}

\EditorialNote{本书以一段简短的陈述开始，讨论福音作者的权威、他们的数目、顺序以及他们叙述的不同计划。\Person{奥古斯丁}随后准备讨论关于福音和谐的问题，他在本书中与那些提出以下困难的人进行辩论：\Person{基督}没有留下自己的著作，或者有人错误地声称某些关于魔法艺术的书是由祂写的。他也回应了那些反对福音教导的人的异议，他们断言\Person{基督}的门徒立即将他们的师傅抬高到超出祂实际所是，当他们肯定祂是天主时，并教导他们未曾从祂领受的教训，当他们禁止敬拜众神时。针对这些对手，他通过引用先知的话，并证明以色列的天主应是唯一受敬拜的对象，来为宗徒的教导辩护。这位天主，虽然从前\Place{罗马}人拒绝接纳祂为唯一神祇，且正因祂禁止他们与祂一起敬拜别神，但如今最终使\Place{罗马}帝国服从于祂的名，并在万民中通过福音的宣讲粉碎了他们的偶像，正如祂通过祂的先知所应许的那样。}\par

% structure: p0003 source=h2 target=subsection reason=relative-heading-rank; base=h2

\subsection{第一章 论福音的权威}

1. 在圣经所包含的全部神圣记载中，福音理应居于首位。因为法律和先知从前预言必将成就的事，在福音中得以实现和应验。这福音的首批宣讲者是宗徒，他们亲眼看见了我们的主、救主耶稣基督，当时祂尚在肉身中。他们不仅记着从祂口中听到的话语和亲眼所见的事迹，而且当宣传福音的职责加于他们时，他们还精心让人类了解那些在祂与门徒建立门徒关系之前发生的神圣且值得纪念的事件，这些事件涉及祂的诞生、幼年或少年时期。关于这些事，他们能够进行精确查问，并从祂本人或祂的父母或其他当事人那里，基于最可靠的提示和最可信的见证获得信息。他们中的某些人——即\Person{玛窦}和\Person{若望}——在他们各自的书中，将关于祂的一切必要记录留给世界。\par

2. 为防止有人以为，在理解和宣讲福音方面，若那些当主在肉身中显现并与门徒同在时实际跟随过祂的人，与那些通过前者以可信方式得知事实的人所宣告的，有何区别，天主眷顾藉圣神的作为安排使那些只是首批宗徒追随者的人也被授予权柄，不仅宣讲福音，而且也将其书面记载下来。我指的是\Person{马尔谷}和\Person{路加}。然而，所有其他那些曾尝试或敢于写出主或宗徒事迹书面记录的人，却未能在他们自己的时代使人相信他们是那种能使教会信任并接纳其作品为圣经正典权威的人。这不仅因为他们无权在叙述中要求人们相信，也因为他们以欺骗性方式在其作品中引入了某些被公教会与宗徒信德准则及纯正教义所谴责的内容。\par

\SourceRecord{Translated by S.D.F. Salmond.；Nicene and Post-Nicene Fathers, First Series；Vol. 6.；Edited by Philip Schaff.；Buffalo, NY: Christian Literature Publishing Co.,；1888.；<http://www.newadvent.org/fathers/1602101.htm>.}

% structure: p0001 source=h1 target=section reason=page-title

\section{\WorkTitle{福音的和谐 卷一}}

% structure: p0002 source=h2 target=subsection reason=relative-heading-rank; base=h2

\subsection{第二章  论福音作者的次序及他们写作的原则。}

3.  现在，那四位福音作者的名字已传遍全世界，他们的人数被定为四——或许仅仅是因为这个世界有四部分，而他们以人数如同一种神秘的标记，表明基督教会正在扩展并遍及整个世界——据信是按以下次序写作：首先是\Person{玛窦}，其次是\Person{马尔谷}，第三是\Person{路加}，最后是\Person{若望}。因此，[似乎]他们之间有一个确定的次序，关于他们个人所知和宣讲（福音）之事，但关于撰写书面叙述的任务，则有不同的次序。确实，就他们获取自己的知识和受命宣讲而言，那些在主带着肉身在世时实际追随祂、亲耳听祂说话、亲眼见祂行动的人，无疑是首先的；他们从主口中接受了使命，并被派遣去宣讲福音。但就撰写那应当被接受为神圣权威所定的福音记录而言，在主逾越节前所选的人中，只有两位获得了一席之地——即第一位和最末位。按次序，第一位是\Person{玛窦}，最末位是\Person{若望}。因此，其余两位，虽不属于上述之列，但已成为在他人中说话之基督的门徒，他们如同要被拥抱的儿子，被这两位所支持，并被安置在两者之间。\par

4.  在这四位中，确实只有\Person{玛窦}被认为用希伯来文写作；其余的人用希腊文。尽管他们每人似乎都保持了自己特定的叙述次序，但这当然不应理解为每位作者选择在不知其前任所做之事的情况下写作，或遗漏了那些别人显然已记录的事件，仿佛这些事没有信息。事实是，正如他们各人领受了灵感的恩赐，他们避免在其各自的劳作中添加任何多余的合作作品。因为\Person{玛窦}被认为着手按照王室谱系来建构主降生成人的记录，并讲述祂大部分言行，这些言行涉及人类今世的生活。\Person{马尔谷}紧随其后，他的著作看起来像是\Person{玛窦}的随从和缩影。因为在他的叙述中，他没有与\Person{若望}一同记载任何其他作者没有记载的事；他单独记载的很少；他与\Person{路加}一起而不同于其余作者记载的更少；但与\Person{玛窦}一致时，他有大量的段落。他叙述时，许多地方几乎用了与\Person{玛窦}相同的字词和数目，这些一致之处或仅与那位福音作者，或与他与其他作者一起。另一方面，\Person{路加}似乎更多专注于主的司铎的谱系和身份。因为他虽然按自己的方式将谱系追溯到\Person{达味}，但他遵循的不是王室谱系，而是非王者的谱系。他也将谱系带到\Person{达味}之子\Person{纳堂}那里，此人同样不是国王。然而\Person{玛窦}并非如此。因为他在谱系中一路沿着\Person{撒罗满}王，严格而规律地追溯其他君王的继承；在列举这些时，他也保持了那神秘的数目，我们以后会谈到。\par

\SourceRecord{Translated by S.D.F. Salmond.；Nicene and Post-Nicene Fathers, First Series；Vol. 6.；Edited by Philip Schaff.；Buffalo, NY: Christian Literature Publishing Co.,；1888.；<http://www.newadvent.org/fathers/1602102.htm>.}

% structure: p0001 source=h1 target=section reason=page-title

\section{福音的和谐，卷一}

% structure: p0002 source=h2 target=subsection reason=relative-heading-rank; base=h2

\subsection{第三章 论玛窦与马尔谷特别注目基督的君王特质，而路加则论及司铎特质。}

5. 因为主耶稣基督，是唯一真正的君王和唯一真正的司铎，前者为治理我们，后者为替我们赎罪，向我们显示了祂自己的形象如何同时承载了这两部分，而这在先祖中只是分别被推荐[注意]的。这一点若（例如）我们查看钉在祂十字架上的题字时就会显明——\BibleQuote{犹太人的君王：}与此相连，并出于一种隐秘的本能，比拉多回答说：\BibleQuote{我所写的，我已经写了。}因为在圣经中早有预言：\BibleQuote{不要毁坏标题的文字。}同样，就司铎的部分而言，如果我们留意祂所教导我们关于奉献和领受的道理，这一点也是明显的。因为祂就是这样预先向我们发出关于祂自己的预言：\BibleQuote{你照默基瑟德的品位，永远做司铎。}在圣经的许多其他见证中，基督都同时显现为君王和司铎。因此，甚至达味自己——祂被宣告为达味之子，这比被称为亚巴郎之子更频繁，且玛窦和路加都一致如此认为：前者视他为从撒罗满而下可以追溯祂世系的人，后者则视他为从纳堂而上可以追溯祂族谱的人——当他吃陈设饼时，也代表了司铎的部分。因为除了司铎，任何人吃那饼都是不合法的。此外，只有路加提到天使如何发现玛利亚，以及她与司祭匝加利亚的妻子依撒伯尔有亲属关系。同一福音作者记述了这位匝加利亚的妻子是亚郎的女儿之一，也就是说她属于司铎支派。\par

6. 因此，玛窦着眼于君王特质，路加着眼于司铎特质，他们同时都突出地展示了基督的人性：因为基督是按照祂的人性才成为君王和司铎的。天主也将祂祖先达味的宝座赐给祂，使祂的王国没有终结。这样做是为了使天主与人之间有一位中保，就是降生成人的基督耶稣，为我们转求。路加则没有人与他相连，作为他的概括者，就像马尔谷依附于玛窦那样。这或许并非没有某种庄严的意义。因为君王有权拥有随从的顺服跟随；因此，这位致力于记述基督君王特质的福音作者，就有一个人作为他的同伴，以某种方式跟随他的脚步。但由于司铎必须独自进入至圣所，按照这一原则，路加——他的目标着眼于基督的司铎职务——就没有任何人作为同工跟随他，以某种方式作为他叙述的简缩者。\par

\SourceRecord{Translated by S.D.F. Salmond.；Nicene and Post-Nicene Fathers, First Series；Vol. 6.；Edited by Philip Schaff.；Buffalo, NY: Christian Literature Publishing Co.,；1888.；<http://www.newadvent.org/fathers/1602103.htm>.}

% structure: p0001 source=h1 target=section reason=page-title

\section{《福音的和谐》，卷一}

% structure: p0002 source=h2 target=subsection reason=relative-heading-rank; base=h2

\subsection{第四章 论若望致力于阐释基督的天主性}

7. 然而，这三位福音作者所关注的大部分内容，是基督藉着人的肉身、按时间方式所行的事。但若望则着眼于主真正的那天主性，即祂与圣父同等的那一位，并竭尽全力在他的福音中，以他认为足以满足人类需要和认知的方式，阐述这神圣的本性。因此他被提升到更高的境界，将其他三位远远抛在身后；以至于，你在他们身上看到的是某种与尘世之人基督交往的人，而在若望身上，你看到的是已经穿透环绕大地的云雾，达到清澈天空的人，他能以最清晰、最稳定的心灵之眼，仰望那在起初就与天主同在、万有藉祂而造的圣言。在那里，他也认出了那成为血肉、以便居住在我们中间的那一位；[我们所说的那位圣言]，祂取了肉身，并非变成了肉身。因为如果这肉身的取用不是以同时保存不变的天主性的方式完成的，就绝不可能说出这样的话：\BibleQuote{我与父原是一体。}因为父与肉身当然不是一体。同一位若望也记下了主论及自己的见证：\BibleQuote{看见了我，就是看见了父；}以及\BibleQuote{我在父内，父在我内；} \BibleQuote{使他们合而为一，如同我们合而为一；} \BibleQuote{凡父所做的，子也照样做。}还有其他类似的陈述，旨在向那些拥有正确理解的人表明基督的天主性，即祂与父平等，我们几乎可以说，这些陈述得以进入福音叙事，全靠若望一人。他像是一个比别人更丰富、更亲切地汲取了祂天主性奥秘的人，仿佛是从他惯常在主坐席时靠着的那个主怀中汲取的。\par

\SourceRecord{Translated by S.D.F. Salmond.；Nicene and Post-Nicene Fathers, First Series；Vol. 6.；Edited by Philip Schaff.；Buffalo, NY: Christian Literature Publishing Co.,；1888.；<http://www.newadvent.org/fathers/1602104.htm>.}

% structure: p0001 source=h1 target=section reason=page-title

\section{\WorkTitle{福音的和谐，第一部}}

% structure: p0002 source=h2 target=subsection reason=relative-heading-rank; base=h2

\subsection{第五章　论两种德性：若望从事于默观性，其他福音作者从事于行动性}

8. 再者，有两个不同的德性（或才能）已呈现于人的心灵。其中，一是行动性的，一是默观性的：一个乃是藉以行走道路者，另一个则是藉以到达目标者；一个是人劳苦以求净化心灵而能看见天主者，另一个则是人解脱羁绊而看见天主者。如此，这两个德性中，前者专注于有关正确操练现世生活的规诫，而后者则处理那永生生命的教义。同样地，一个运作，一个安息；因为前者在涤除罪中找到它的领域，后者则在已被涤除之光中活动。并且再次，在此有死生命中，一个从事于良好行为的工程；而另一个则更凭信德而存在，且仅在极少数人身上可见，如同在镜中模糊地观看，且仅部分地得见不变真理的某种景象。现在，这两个德性被理解为象征性地呈现在雅各伯的两位妻子身上。关于这些，我已经在我力所能及的范围内，并尽可能充分地针对我的任务，在我所写的反对摩尼教徒福斯图斯的著作中讨论过了。因为肋阿，按解释，意思是\Quote{劳苦}，而辣黑耳则意味\Quote{所看见的第一原则}。藉此，我们得以明白——只要人细心留意此事——那三位以卓越的丰富性处理了主的现世作为以及祂那些主要旨在塑造现世生活习俗的言论的福音作者，他们专注于那行动性的德性；而另一方面，若望，他叙述的主的事迹远少得多，却以更大的仔细和更丰富的细节记录了祂所说的话，尤其是那些旨在引领我们认识天主圣三的统一与永生福乐的讲论，他制定了他的计划并构思了他的陈述，以推荐默观性的德性让我们重视。\par

\SourceRecord{Translated by S.D.F. Salmond.；Nicene and Post-Nicene Fathers, First Series；Vol. 6.；Edited by Philip Schaff.；Buffalo, NY: Christian Literature Publishing Co.,；1888.；<http://www.newadvent.org/fathers/1602105.htm>.}

% structure: p0001 source=h1 target=section reason=page-title

\section{\WorkTitle{福音的和谐，第一卷}}

% structure: p0002 source=h2 target=subsection reason=relative-heading-rank; base=h2

\subsection{第六章 论\BibleBook{}{默示录}中的四种活物，有人以此喻指四福音作者，或以他法喻指，皆相宜。}

9. 基于这些理由，我也认为，在那些将\BibleBook{}{默示录}中的活物解释为象征四福音作者的不同群体中，凡以狮子指\Person{玛窦}、以人指\Person{马尔谷}、以牛犊指\Person{路加}、以鹰指\Person{若望}者，对图形的运用比那些将人配给\Person{玛窦}、将鹰配给\Person{马尔谷}、将狮子配给\Person{若望}者更为合理。因为，后者在形成他们对这事的特定见解时，选择只着眼于各福音书的开端，而未顾及几位福音作者各自设计的全貌，而这正是最应当详加考察之事。因为，那位最充分地使我们注意到基督王权特性的作者，自然更应该由狮子来代表。因此，我们也在\BibleBook{}{默示录}那段提及王室支派本身时读到狮子，那里说：\BibleQuote{犹大支派的狮子已得胜了。}因为在玛窦的记述中，贤士们从东方来寻找那位君王，并朝拜祂——祂的诞生由星象向他们显示。同样，\Person{黑落德}自己也是一位君王，[据说]对这位王室婴儿感到恐惧，并下令杀害如此多的小孩子，以确保那一个孩子可能被杀。再者，路加以牛犊为象征，指向司祭所奉献的卓越祭献，这一点两派解经家均无异议。因为在这部福音中，叙述者的记述从\Person{匝加利亚}司祭开始。其中也提到\Person{玛利亚}与\Person{依撒伯尔}的亲戚关系。同样，它也记载了为婴孩基督履行了最初司祭职务的适当礼仪；仔细审视便可在本福音中发现许多其他事项，由此明显可见，\Person{路加}的目的是处理司祭的角色。由此可进一步推论，\Person{马尔谷}既不自命记述王室谱系，也不阐述司祭职的任何独特之处，无论关于亲属关系或祝圣方面，同时他呈现在眼前，处理基督作为人所行的事，因此他似乎简单地由四种活物中的人像所指代。再者，那三种活物——狮子、人或牛犊——均在地上行动；同样，这三位福音作者主要专注于基督在肉身中所行的事，以及祂向人们传授的诫命，这些人也背负肉身的重担，为在此有死生命中正当实践而受训诲。而\Person{若望}则如鹰般翱翔于人性软弱的云层之上，用心中最锐利最坚定的双眼凝视不变真理之光。\par

\SourceRecord{Translated by S.D.F. Salmond.；Nicene and Post-Nicene Fathers, First Series；Vol. 6.；Edited by Philip Schaff.；Buffalo, NY: Christian Literature Publishing Co.,；1888.；<http://www.newadvent.org/fathers/1602106.htm>.}

% structure: p0001 source=h1 target=section reason=page-title

\section{\WorkTitle{福音的和谐，第一卷}}

% structure: p0002 source=h2 target=subsection reason=relative-heading-rank; base=h2

\subsection{第七章 圣奥古斯丁从事这部福音和谐著作之理由的陈述，以及他以何种方式回应那些妄称基督本人未著一字、而祂的门徒宣告祂为天主实属不当之人的实例。}

10. 然而，主那些神圣的战车——祂乘着它们巡行大地，使万民负上祂轻松的轭与轻省的重担——却被某些人恶言攻击。这些人以不敬的虚妄或愚昧的鲁莽，妄图剥夺那些导师作为信实史家的信誉；而正是藉着这些导师的传递，基督信仰遍布世界，并结出如此丰硕的果实，以致不信者现今几乎不敢私下相互诽谤，因外邦人的信仰和万民的虔敬而受抑制。尽管如此，他们仍以其恶意的争辩试图阻止一些人认识信仰、从而成为信徒，又竭力搅扰那些已获信仰的人，使其困恼；此外，也有一些弟兄，在不损害自身信德的前提下，渴望知道如何回答此类问题，或为增长学识，或为驳斥敌人的虚妄言论。因此，在主我们天主的启迪与助佑下——但愿这一切能有益于这些人的得救——我们在此著作中，旨在显明那些自以为能对四位圣史所写的四部福音书提出微妙指控之人的错误或鲁莽。为成功完成此计划，我们必须证明这些作者彼此间并无矛盾。因为那些仇敌惯于以此作为他们一切虚妄异议中最主要的理由，即福音作者彼此不和谐。\par

11. 但我们必须首先讨论一个常使某些人心生困惑的问题：为何主自己未曾写作任何东西，而让我们必须接受他人为其历史所作的见证？这尤其是外教人责难之处，当他们缺乏勇气直接指控或亵渎主耶稣基督自己，并承认祂——尽管只是作为一个人——拥有最卓越的智慧时，他们便如此说。在作出这一让步的同时，他们断言门徒为他们的师傅所宣称的远超过祂的实际；甚至称祂为天主子、天主的圣言，万物藉祂而造，并宣称祂与天主原是一体。同样，他们也以此打发宗徒书信中其他类似段落；按照这些书信的教导，我们被训示当与圣父同等地朝拜祂为唯一的天主。因为他们认为，祂固然应被尊为最智慧的人，但否认祂应被当作天主来朝拜。\par

12. 因此，当他们质问他为何没有亲自写作时，似乎暗示：若祂写下任何关于自己的事，他们便会相信；至于他人凭自己的判断所传于世的关于祂生平的事，他们则不必接受。但我倒要反问他们：为什么对于他们自己某些最杰出的哲学家，他们却接受其门徒在著作中留下的陈述，尽管那些哲人自己并未留下任何关于自己生平的文字？例如，\Person{毕达哥拉斯}——当时希腊在默观德行方面无人比他更显赫——据信绝对没有写过任何东西，无论是关于他个人史还是其他主题。至于\Person{苏格拉底}——另一方面，他们判定他在陶铸道德生活的实践德行上超越所有其他人，甚至毫不迟疑地声称，连他们的神祇\Person{阿波罗}也作证称他为最智慧的人——他确实用一些简短的韵文演绎了\Person{伊索}的寓言，并用自己的言辞和韵律来转述他人的主题。但也仅此而已。他本人绝无意愿写任何东西，以至于宣称就连那一点也是被迫于他的精灵的帝王意志，正如他所有门徒中最杰出的\Person{柏拉图}告诉我们的。那部作品也并非旨在优美地表达他自己的思想，而更多地是表达他人的观点。因此，他们有何合理理由相信那些哲人的门徒所记录的一切关于其生平的事，却拒绝相信在基督身上，祂的门徒所写的关于祂生平的事？尤其当我们看到他们承认基督在智慧上超越所有其他人，尽管否认祂是天主时，就更应如此主张。难道那些他们毫不犹豫地承认远比基督低劣的人，竟能使他们门徒在叙述其生平的一切上值得信赖，而基督却没有这种能力？但如果这种说法荒谬至极，那么对于这位他们赋予智慧荣誉者，他们理应相信的不仅是符合自己想法的内容，还应当相信那些从这位智者本人学习到各种事实、并将其关于祂生平的叙述留传后世之人的记载。\par

\SourceRecord{Translated by S.D.F. Salmond.；Nicene and Post-Nicene Fathers, First Series；Vol. 6.；Edited by Philip Schaff.；Buffalo, NY: Christian Literature Publishing Co.,；1888.；<http://www.newadvent.org/fathers/1602107.htm>.}

% structure: p0001 source=h1 target=section reason=page-title

\section{\WorkTitle{福音的和谐}，第一卷}

% structure: p0002 source=h2 target=subsection reason=relative-heading-rank; base=h2

\subsection{第八章 论为何若因普遍的传闻而相信基督是最有智慧的人，就不应因宣讲的卓越报道而相信祂是天主。}

13. 此外，他们应当告诉我们，他们是通过什么途径认识到祂是最有智慧的人这一事实，或者这消息是如何传到他们耳中的。如果仅仅是通过当前的传闻得知，那么，普通的传闻比起祂的门徒们（他们四处宣讲祂，将同一传闻如同芬芳之气散播至全世界）所传的，是否更为可信？总之，他们应当比较这两种传闻，并相信其中更为优越的那一种关于祂生平的记述。事实上，从那个天主教会（它的扩展遍及全球，令此类人惊讶不已）以奇妙清晰度传播开来的这则传闻，以无与伦比的方式胜过了像他们这样的人所谈论的虚浮谣言。此外，这则传闻具有如此分量和流通性，以至于他们在恐惧中只能在内心喃喃自语，焦虑而微弱地提出些微的反对意见，仿佛他们现在更害怕被听到，而不是希望获得信任；这传闻宣告基督是天主的独生子，祂自己就是天主，万物由祂而造成。因此，如果他们选择传闻作为见证，为何不特别选择这一则极为清晰、卓越非凡的传闻呢？如果他们渴望文字的证据，为何不采纳那些在权威上远超其他一切的福音书呢？我们这一方确实接受了那些关于他们神祇的陈述，这些陈述同时出现在他们最古老的著作和最普遍的传闻中。但如果这些神祇应被视为敬拜的对象，为何在剧场中他们又成为笑柄？反之，若他们本该引人发笑，那么当他们在剧场中被崇拜时，这笑料就更大了。剩下的就是我们应将那些人视为有意为基督作见证者，他们因谈论自己不知道的事，而失去了知道自己所谈之事的功绩。另外，如果他们声称拥有任何据称由基督撰写的书籍，他们应当将它们拿出供我们查阅。因为这些书（若真有的话）必定是最有益、最健全的，因为它们出自一位他们认为是最有智慧的人之手。然而，如果他们不敢拿出来，那必是因为它们内容邪恶；但若内容邪恶，那么最有智慧的人就不可能写过它们。然而，他们承认基督是最有智慧的人，因此基督不可能写过任何这类东西。\par

\SourceRecord{Translated by S.D.F. Salmond.；Nicene and Post-Nicene Fathers, First Series；Vol. 6.；Edited by Philip Schaff.；Buffalo, NY: Christian Literature Publishing Co.,；1888.；<http://www.newadvent.org/fathers/1602108.htm>.}

% structure: p0001 source=h1 target=section reason=page-title

\section{福音的和谐，卷一}

% structure: p0002 source=h2 target=subsection reason=relative-heading-rank; base=h2

\subsection{第九章 论某些人妄称基督撰写关于魔术之书}

14. 然而，这些人的愚昧竟达到如此地步，竟敢声称他们认为由祂所写的书卷含有那些技艺，他们认为祂正是藉着这些技艺施行了那些奇迹，而这些奇迹的名声已传遍各地。而他们这种幻想暴露了他们真正喜爱的是什么，以及他们的真正目的是什么。因为确实如此，他们向我们表明，他们持有这种意见——基督是最智慧的人——仅仅是因为祂拥有某种我不知道的非法技艺的知识，而这些技艺不仅受到基督宗教纪律的公正谴责，甚至连地上的政府行政也予以谴责。而且，实话实说，如果有人断言他们读过这类由基督所著的书卷，那么他们为何不利用从这些书中获得的知识，亲手施行一些类似的工作，就像那些由祂所行、令他们如此惊叹的奇迹一样呢？\par

\SourceRecord{Translated by S.D.F. Salmond.；Nicene and Post-Nicene Fathers, First Series；Vol. 6.；Edited by Philip Schaff.；Buffalo, NY: Christian Literature Publishing Co.,；1888.；<http://www.newadvent.org/fathers/1602109.htm>.}

% structure: p0001 source=h1 target=section reason=page-title

\section{\WorkTitle{福音的和谐，卷一}}

% structure: p0002 source=h2 target=subsection reason=relative-heading-rank; base=h2

\subsection{第十章　　论某些人竟狂妄到以为这些书卷上题有\Person{伯多禄}与\Person{保禄}之名}

15. 不仅如此，由于天主的审判，有些人或相信或希望别人相信\Person{基督}曾写过此类文字，他们甚至谬误到宣称这些书卷的封面上，以书信题署的形式，写着致\Person{伯多禄}与\Person{保禄}的称呼。完全可能的是，\Person{基督}之名的敌人，或是某些以为借此可赋予这等可憎技艺以出自如此荣耀之名的权威的人，曾以\Person{基督}和宗徒的名义写过这类东西。但在这种极其欺诈的狂妄中，他们如此盲目，以至于使自己成为即便那些对基督教文献仅具孩童般认识、仅居于诵读等级的青年人也觉得可笑的对象。\par

16. 因为当他们决定想象\Person{基督}曾以这样的笔调写信给祂的门徒时，他们想到了那些最容易被当作\Person{基督}很可能写信之人的追随者，即那些与祂保持最亲密友谊关系的人。于是\Person{伯多禄}和\Person{保禄}出现在他们脑中，我相信，仅仅是因为他们在许多地方偶然看到这两位宗徒在画像中与祂同在。因为\Place{罗马}以一种特别光荣而隆重的方式推崇\Person{伯多禄}与\Person{保禄}的功绩，原因之一便是他们于同一日殉道。因此，那些不在圣经中而在彩绘墙上寻求\Person{基督}及其宗徒的人，完全陷入错误乃是应得的报应。这些虚构故事的作者被画家误导，也不足为奇。因为在\Person{基督}与祂的门徒在可朽肉身中共同生活的整个时期，\Person{保禄}从未成为祂的门徒。只有在祂受难、复活、升天、圣神自天降临、许多犹太人悔改并表现出奇妙信德、执事兼殉道者\Person{斯德望}被石头砸死之后，且当\Person{保禄}仍名为\Person{扫禄}、正残酷迫害那些信从\Person{基督}的人时，\Person{基督}才从天上呼召那人，使他成为祂的门徒和宗徒。那么，\Person{基督}怎么可能在祂死前就写了那些他们希望别人相信是祂所写、且是写给那些在门徒中与祂最亲近的\Person{伯多禄}和\Person{保禄}的书卷呢？因为直到那时，\Person{保禄}根本尚未成为祂的门徒。\par

\SourceRecord{Translated by S.D.F. Salmond.；Nicene and Post-Nicene Fathers, First Series；Vol. 6.；Edited by Philip Schaff.；Buffalo, NY: Christian Literature Publishing Co.,；1888.；<http://www.newadvent.org/fathers/1602110.htm>.}

% structure: p0001 source=h1 target=section reason=page-title

\section{《\WorkTitle{福音的和谐}》卷一}

% structure: p0002 source=h2 target=subsection reason=relative-heading-rank; base=h2

\subsection{第十一章　驳斥那些愚昧地认为基督藉邪术使人民归向自己的人}

17. 此外，让那些疯狂幻想祂是藉邪术才能成就大事，并且认为祂是藉施行这类仪式，才使祂的名在那些将要归向祂的人民中成为神圣的，注意这个问题：即，祂是否可以在尚未降生于世之前，便藉邪术使那些大先知们充满圣神，使他们预先把那些关于祂的、注定要发生的事宣告出来？现在我们可以在福音中读到这些事的应验，并在世界中看到它们的实现。诚然，即使祂是藉邪术为自己赢得敬拜，而且还是死后的事，但祂在出生之前却不可能是一个术士。相反，有一个民族被特别委派来担任预言祂来临的职分，这个国家的整个管理制度都被设定为对这将要来临的君王的一个预言，祂要建立一个由万民中召选出来的天上国度。\par

\SourceRecord{Translated by S.D.F. Salmond.；Nicene and Post-Nicene Fathers, First Series；Vol. 6.；Edited by Philip Schaff.；Buffalo, NY: Christian Literature Publishing Co.,；1888.；<http://www.newadvent.org/fathers/1602111.htm>.}

% structure: p0001 source=h1 target=section reason=page-title

\section{\WorkTitle{福音的和谐，卷一}}

% structure: p0002 source=h2 target=subsection reason=relative-heading-rank; base=h2

\subsection{第十二章　论犹太人的天主在那一民族被征服后仍未被罗马人接纳，因为祂的诫命是唯独祂应受敬拜，且要消灭偶像。}

18. 此外，那希伯来民族，正如我所说，受命预言基督，除了一位天主，即创造天地及其中万物的真天主外，别无其他神祇。他们常在祂的震怒下被交在敌人手中。如今，确实因他们将基督置于死地这极其严重的罪，他们已从他们王国的首都耶路撒冷本身被彻底根除，并臣服于罗马帝国。罗马人习惯通过亲自敬拜他们所征服民族的神祇来安抚它们，并惯于承担其神圣礼仪的职责。但他们在对希伯来民族的天主一事上拒绝依此原则行事，无论是进攻之时，还是征服该民族之后。我相信他们意识到，若他们接纳这位天主的敬拜——祂的诫命是唯独祂应受敬拜，且要消灭偶像——他们就必须舍弃一切从前他们曾从事宗教服务之物，并相信其帝国因敬拜这些偶像而壮大。但在这一点上，他们魔鬼的虚伪大大欺骗了他们。因为他们本该认识到这一事实：惟有凭真天主那隐藏的旨意——万有之最高权柄都掌握在祂手中——帝国才被赐予他们并得以壮大，而他们的地位并非由于那些神祇的眷顾；那些神祇若在此事上能施加任何影响，宁可保护自己的子民不受罗马人压制，或者使罗马人完全臣服于他们。\par

19. 当然，他们绝不可能断言，他们自身所表现的那种虔敬与行为方式，竟成了他们所征服民族的神祇所喜爱和选择的对象。他们只要回忆起自己早期的起源——那为被遗弃罪犯设立的庇护所，以及\Person{罗慕路斯}的杀弟行为——便绝不会如此断言。因为当\Person{雷慕斯}和\Person{罗慕路斯}设立他们的庇护所时，意图是凡逃到那里的人，无论所犯何罪，其行为均可免于惩罚；他们并未颁布任何补赎的规诫，以引导这些可怜之人的心灵回归正途。难道不是借着这免罚的贿赂，他们武装了聚集而来的恐惧逃犯，对抗他们原本所属的城邦，以及他们所畏惧的法律吗？或者，当\Person{罗慕路斯}杀死并未对他作恶的弟弟时，难道他的心思是专注在伸张正义，而非攫取绝对权力吗？神祇们果真喜欢这样的行为，仿佛它们本身是各自城邦的敌人，竟偏爱那些与这些社群为敌的人吗？不，相反，它们既未因离弃一方而加害于他们，也未因投靠另一方而给予任何帮助。因为他们无力给予或剥夺王权。那些事是由唯一真天主按其隐秘的旨意成就的。祂并不意欲使那些祂赐予世上王权的人必定蒙福，或使那些祂剥夺了那地位的人必定不幸。祂却因其他理由、藉其他方式使人蒙福或不幸，并或出于许可，或作为实际恩赐，按祂所喜悦的任何人和祂所选择的任何时期，依照预定时代的秩序，分赐短暂尘世的王国。\par

\SourceRecord{Translated by S.D.F. Salmond.；Nicene and Post-Nicene Fathers, First Series；Vol. 6.；Edited by Philip Schaff.；Buffalo, NY: Christian Literature Publishing Co.,；1888.；<http://www.newadvent.org/fathers/1602112.htm>.}

% structure: p0001 source=h1 target=section reason=page-title

\section{\WorkTitle{福音合参} 第一卷}

% structure: p0002 source=h2 target=subsection reason=relative-heading-rank; base=h2

\subsection{第十三章 论天主为何容许犹太人被征服}

20. 因此，他们也不能这样公平地追问我们：为什么你们宣称是至高无上真天主的希伯来人天主，不仅没有使罗马人屈服于他们权下，甚至未能保护这些希伯来人自己免遭罗马人的征服？因为他们公开的罪行先行，先知们早已预言这事必将临到他们；而最重要的原因在于，他们以不敬的狂怒处死了基督，在犯此罪时，因其他隐秘过犯的应得惩罚，他们被蒙蔽（看不清自己罪行的罪责）。祂的受难也将有益于外邦人，这同样由先知预言所预告。从另一角度看，以下事实同样清楚：那民族的王国、圣殿、司祭职、祭献制度，以及那在希腊文中称为\OriginalTerm{傅油}{\GreekText{χρῖσμα}}的神秘傅油——\Emph{基督}之名明显由此而来，那民族也惯于称他们的君王为受傅者——这一切被设立，其明确目的是预表基督；同样清楚的是，在被害的基督复活后，开始向信从的外邦人宣讲时，所有这些事物都到达了终点，无论是通过胜利的罗马人，还是通过被征服的犹太人，它们都以这种方式悄然结束，无人认清这一情况。\par

\SourceRecord{Translated by S.D.F. Salmond.；Nicene and Post-Nicene Fathers, First Series；Vol. 6.；Edited by Philip Schaff.；Buffalo, NY: Christian Literature Publishing Co.,；1888.；<http://www.newadvent.org/fathers/1602113.htm>.}

% structure: p0001 source=h1 target=section reason=page-title

\section{\WorkTitle{福音的和谐}，卷一}

% structure: p0002 source=h2 target=subsection reason=relative-heading-rank; base=h2

\subsection{第十四章　论希伯来人的天主虽其人民被征服，却借推翻偶像并使众外邦归服于祂，显明祂是不可征服的}

21. 这确实是一个奇妙的实事，却被那些至今仍为外教人的少数人所忽略——即这位曾被征服者所冒犯，又被征服者拒绝接纳的希伯来人的天主，如今却在全国万民中被宣讲和敬拜。这就是那位先知早在很久以前所预言的天主，当他这样对天主子民说：\BibleQuote{那领你们出来的、以色列的天主，将被称为全地的（天主）。} 这预言是通过基督的名而实现的，基督以那希伯来种族始祖亚伯拉罕之孙——以色列的后裔形式来到人间。因为对这同一以色列也曾说过：\BibleQuote{地上万族都要因你的后裔蒙福。} 由此表明，以色列的天主，那位创造天地、公正而仁慈地管理人间事务的真天主，既不因祂的仁慈而排除公义，也不因公义而妨碍仁慈，当祂允许那民族的王国和司祭职被罗马人夺取和推翻时，祂自身并未因祂的希伯来人民遭受覆灭而被击败。因为如今，藉着这福音——其来临曾由那王国和司祭职所预表——的真正君王和司祭基督的大能，以色列的天主亲自在各地摧毁列邦的偶像。事实上，正是为了防止这种摧毁，罗马人才拒绝像接纳他们所征服的其他民族的神明那样，接纳这位天主的敬礼。因此，祂从那预言的民族中挪走了王国和司祭职，因为藉着那民族向人类所应许的那一位已经来到。而藉着君王基督，祂使那曾击败该民族的罗马帝国臣服于祂自己的名下；并藉着基督徒信德的力量和虔诚，将其转变，以至于推翻了那些偶像，那些偶像所受到的尊崇曾阻碍对祂的敬拜得以进入。\par

22. 我认为，基督在诞生于世之前，并非藉魔术的技艺，使得那些注定要在祂历史中发生的事，由如此多的先知预先宣告，并藉着在一个特定民族中设立的王国和司祭职而预先描绘。因为那与现已废止的王国相关联的民众，在天主奇妙的眷顾下分散在各地，始终没有领受真正君王和司祭的傅油——在这傅油中，基督之名的含义得以显明。但尽管如此，他们仍然保留了他们某些礼仪的残余；另一方面，即使在被推翻和征服的状态下，他们也没有接受那些与偶像崇拜相关的罗马仪式。因此，他们仍保留先知书作为基督的见证；这样，在祂敌人的文献中，我们找到了证明这预言之基督真理的证据。那么，这些不幸的人通过他们不配称的方式赞美基督的名，究竟显明自己是什么？若有什么与魔术实践相关的东西是以祂的名写下的，而基督的教义却如此剧烈地反对这些技艺，那么这些人倒应藉此事实，领会这名的伟大——连那些违反祂诫命生活的人，也试图藉着加上这名来尊崇他们的邪恶行径。因为正如在人类各种谬误的进程中，许多人以祂的名作掩护，建立了反对真理的各式异端一样，就连基督的敌人也认为，为了使他们所提出反对基督教义的意见得到接纳，除非有基督之名，否则他们便无任何权威可供利用。\par

\SourceRecord{Translated by S.D.F. Salmond.；Nicene and Post-Nicene Fathers, First Series；Vol. 6.；Edited by Philip Schaff.；Buffalo, NY: Christian Literature Publishing Co.,；1888.；<http://www.newadvent.org/fathers/1602114.htm>.}

% structure: p0001 source=h1 target=section reason=page-title

\section{\WorkTitle{福音的和谐 卷一}}

% structure: p0002 source=h2 target=subsection reason=relative-heading-rank; base=h2

\subsection{第十五章 论外教人被迫赞美基督时，却将侮辱倾泻于祂的门徒这一事实。}

23. 但对此该说什么呢？如果那些徒然赞美基督者以及那些曲解基督教信仰的诽谤者，竟不敢亵渎基督，其原因正是他们中有些哲学家——如\Place{西西里}的\Person{波斐利}在其著作中告知我们的——曾就关于基督的【宣称】求问他们的神祇，并被自己的神谕所迫而赞美基督呢？这也不足为奇。因为我们在\WorkTitle{福音}中也读到，魔鬼曾承认祂；在我们的\WorkTitle{先知书}中也这样记载：\BibleQuote{外邦的神祇都是魔鬼。}于是，为了避免与自身神祇的回答相抵触，他们便避开对基督的亵渎，转而将其倾泻于祂的门徒身上。然而在我看来，这些外邦人的神祇——外教哲学家或许也曾求问过他们——如果被问及要对基督的门徒以及基督本人作出评判时，也必会同样被迫赞美他们。\par

\SourceRecord{Translated by S.D.F. Salmond.；Nicene and Post-Nicene Fathers, First Series；Vol. 6.；Edited by Philip Schaff.；Buffalo, NY: Christian Literature Publishing Co.,；1888.；<http://www.newadvent.org/fathers/1602115.htm>.}

% structure: p0001 source=h1 target=section reason=page-title

\section{福音的和谐，卷一}

% structure: p0002 source=h2 target=subsection reason=relative-heading-rank; base=h2

\subsection{第十六章：论及在摧毁偶像一事上，宗徒们所教导的与基督或先知所教导的并无不同。}

24. 然而，这些人仍然争辩说，这种拆毁庙宇、谴责祭献、粉碎一切偶像的行为，并非藉基督本人的教导而实现，而只是藉祂的宗徒们的手，他们声称宗徒们教导了与基督不同的东西。他们以为藉此手段，在尊敬和赞美基督的同时，就能撕裂基督信仰。因为至少有一点是真的：正是基督的门徒们使基督的言行得以传扬，基督教信仰正是建立在这些言行之上，而目前只有极少数这样性格的人仍与之对立，他们现在虽不公开攻击，但直至今日仍在暗中抱怨。但如果他们拒绝相信基督是按所说的那样教导的，就让他们阅读先知们，他们不仅命令彻底摧毁偶像迷信，而且还预言这种颠覆将在基督教时代发生。如果这些先知说了谎话，为什么他们的话以如此强大的明证应验了？但如果他们说了真话，为什么还要抗拒如此神圣的能力？\par

\SourceRecord{Translated by S.D.F. Salmond.；Nicene and Post-Nicene Fathers, First Series；Vol. 6.；Edited by Philip Schaff.；Buffalo, NY: Christian Literature Publishing Co.,；1888.；<http://www.newadvent.org/fathers/1602116.htm>.}

% structure: p0001 source=h1 target=section reason=page-title

\section{\WorkTitle{福音的和谐，第一卷}}

% structure: p0002 source=h2 target=subsection reason=relative-heading-rank; base=h2

\subsection{第十七章 驳斥那些唯独拒绝以色列的天主的罗马人。}

25. 然而，有一件事应当受到他们更加审慎的考虑——即，他们认为以色列的天主是什么，以及为什么他们没有像对待其他屈服于罗马帝国权力的民族的众神那样，在崇拜的尊荣中承认祂？当我们看到他们持守这样的见解：智者应当崇拜所有的神祇时，这个问题就更加需要回答。那么，为什么祂被排除在这众神之外？如果祂非常强大，为什么唯有祂不被他们崇拜？如果祂力量很小或没有力量，为什么其他神的偶像被万邦打碎，而祂如今几乎是唯一在那些民族中被崇拜的天主？那些既崇拜所认为的神祇中的大尊与小尊，同时又拒绝崇拜这位比他们所有侍奉的神祇更强有力的天主的人，绝无法从这问题的困境中解脱出来。如果祂是[一位]大能的天主，为什么祂只被视为该当被拒绝？如果祂是[一位]力量很小或没有力量的天主，为什么祂虽被拒绝，却能成就如此之多？如果祂是善的，为什么祂是唯一与其它善神分离的？如果祂是恶的，为什么这位孤立无援的天主不被那么多善神征服？如果祂是诚实的，为什么祂的诫命被藐视？如果祂是说谎的，为什么祂的预言得以应验？\par

\SourceRecord{Translated by S.D.F. Salmond.；Nicene and Post-Nicene Fathers, First Series；Vol. 6.；Edited by Philip Schaff.；Buffalo, NY: Christian Literature Publishing Co.,；1888.；<http://www.newadvent.org/fathers/1602117.htm>.}

% structure: p0001 source=h1 target=section reason=page-title

\section{\WorkTitle{福音的和谐，卷一}}

% structure: p0002 source=h2 target=subsection reason=relative-heading-rank; base=h2

\subsection{第十八章　论希伯来人的天主不为罗马人所接纳，乃因祂願意唯独祂被敬拜}

26. 总之，他们可以随意思想祂。然而，我们仍可问：是否那些为苍白之神和热病之神建造了庙宇、并命令既应当祈求善灵又应当平息恶灵的罗马人，也拒绝把邪神作为崇拜的对象？那么，无论他们对祂抱有何种见解，问题仍然在于：为什么只有祂是他们认为既不值得呼求帮助、也不值得平息的神祇？这位神是谁？是与如此众多的神祇中唯一尚未被发现的那样一位未知者，还是如今被如此多人唯独敬拜的那样一位已知者？所以，他们除了说祂的旨意是唯独祂应被敬拜、祂的命令是当时他们所敬拜的外邦人的那些神祇应停止被敬拜之外，再也无法提出任何理由来解释他们拒绝接纳对这位天主的敬拜。但是，更应当向他们要求回答的却是另一个问题，即：他们认为这位禁止其他神祇（他们曾为之建造庙宇和偶像）的敬拜、并且拥有如此巨大的权能以致其意志在摧毁他们的偶像方面比他们的意志在阻止祂的敬拜方面更占上风的神祇，是何等或何种神祇？事实上，他们那位哲学家——甚至根据他们自己的神谕，他们视其为所有人中最有智慧者——的意见是明明白白的：\Person{苏格拉底}的意见是，每一位神祇都应当以他可能规定的那样被敬拜。因此，对他们而言，拒绝敬拜希伯来人的天主便成了最为必然的事。因为如果他们有意以一种不同于祂所宣称应当被敬拜的方式来敬拜祂，那么他们必定不是在敬拜这位天主本身，而是在敬拜他们自己的虚构。另一方面，如果他们愿意按照祂所指示的方式来敬拜祂，那么他们不能不明了，他们不得敬拜那些祂所禁止敬拜的其他神祇。因此，他们拒绝了唯一真天主的服务，因为他们害怕得罪那许多假神。他们认为那些神祇的愤怒对他们的伤害，会大于这位天主的善意对他们的益处。\par

\SourceRecord{Translated by S.D.F. Salmond.；Nicene and Post-Nicene Fathers, First Series；Vol. 6.；Edited by Philip Schaff.；Buffalo, NY: Christian Literature Publishing Co.,；1888.；<http://www.newadvent.org/fathers/1602118.htm>.}

% structure: p0001 source=h1 target=section reason=page-title

\section{\WorkTitle{福音的和谐 第一卷}}

% structure: p0002 source=h2 target=subsection reason=relative-heading-rank; base=h2

\subsection{第十九章 证明此天主是真天主}

27. 但那必是一种徒劳的必要性和可笑的胆怯。我们现在问：那些以所有神祇都当受敬拜为乐的人，对这位天主持何看法？因为如果祂不当受敬拜，当祂不受敬拜时，如何能说所有神祇都受敬拜呢？而如果祂当受敬拜，那么就不可能所有其他神祇与祂一同受敬拜。因为除非唯有祂受敬拜，否则祂实在完全不受敬拜。或者也许是这样的情况：他们断言祂根本不是天主，而称那些神祇为神——按照我们的信仰，这些神祇除非凭借祂的审判所赐予的许可，否则毫无能力行任何事——不仅不能对任何人行善，甚至也不能对任何人加害，除非对那些被祂（拥有全部权能的那位）判定为理应受害的人。但是，正如他们自己被迫承认的，那些神祇所显示的能力比祂所行的要小。因为如果那些神祇被认为是神，他们的先知被人咨询时给出回应——我不说它们是虚假的，但至少最有利于他们私人利益——那么，祂为何不被视为天主？祂的先知不仅在被咨询的特殊时间对相关主题给出了适当的回答，而且对于那些他们未被咨询、却关乎全人类和万民的主题，在事件发生前很久就预言了那些我们现在读到、并确实亲眼目睹之事。如果他们向那位受其感召而歌咏罗马命运的\Person{西比尔}赋予神的名号，那么祂为何不当被称为天主？祂依照从前的预告，向我们显明了罗马人和万民如何通过基督的福音而相信祂是唯一的天主，并拆毁他们祖先的一切偶像。最后，如果他们称那些从未敢通过他们的先知对这神说过任何反对之话的神祇为神，那么祂为何不当被称为天主？祂不仅通过祂先知的口命令毁灭他们的偶像，还预言在外邦人中，它们将被那些受命放弃偶像、唯独敬拜祂的人所毁灭，而这些人接受了这些命令，就成为祂的仆人。\par

\SourceRecord{Translated by S.D.F. Salmond.；Nicene and Post-Nicene Fathers, First Series；Vol. 6.；Edited by Philip Schaff.；Buffalo, NY: Christian Literature Publishing Co.,；1888.；<http://www.newadvent.org/fathers/1602119.htm>.}

% structure: p0001 source=h1 target=section reason=page-title

\section{《福音的和谐》，卷一}

% structure: p0002 source=h2 target=subsection reason=relative-heading-rank; base=h2

\subsection{第二十章 论未发现异教先知曾预言任何反对希伯来天主之事}

28. 或让他们说明——若他们能够——他们的某位\OriginalTerm{西卜拉}{Sibyl}，或他们其他先知中的任何一位，是否早已预言过，希伯来人的天主、以色列的天主将受万民敬拜，同时宣称此前崇拜他神者曾正当地拒绝了祂；并且，祂的先知们的著作将享有如此崇高的权威，以致罗马政府本身将因服从它们而命令摧毁偶像，而所说的预言者们同时警告不可遵行这些法令——我指的是，若他们能够，让他们从他们先知的任何书中读出类似的话语。因为我并不停下来陈述，我们在他们书中读到的东西，如何为我们的宗教——即基督宗教——重复了证言，这些证言他们本可听自圣天使和我们自己的先知；正如魔鬼本身，在基督居于肉身时，也被迫承认基督。但我略过这些事；当我们提出它们时，他们的争辩是，这些事是由我们的人捏造的。然而，最肯定的是，他们自己可能被催逼提出任何由他们自己神明的预言者们所预言的、反对希伯来天主的事；而就我们这边，我们可以指出在我们先知书中记载的、反对他们神明的宣言，其数目之众多、分量之重大，实为卓著；在其中我们既能留意到命令，又能背诵预言，并证明其应验。面对这些事实的实现，那数量较小的、仍为异教徒的人，更倾向于悲伤，却不愿承认那位有能力预言这些事终将实现的天主；而他们在对待自己虚假的神明——实为真正的魔鬼——时，再没有比通过其神谕得知与自己相关的未来之事更珍视的了。\par

\SourceRecord{Translated by S.D.F. Salmond.；Nicene and Post-Nicene Fathers, First Series；Vol. 6.；Edited by Philip Schaff.；Buffalo, NY: Christian Literature Publishing Co.,；1888.；<http://www.newadvent.org/fathers/1602120.htm>.}

% structure: p0001 source=h1 target=section reason=page-title

\section{\WorkTitle{福音的和谐}，第一卷}

% structure: p0002 source=h2 target=subsection reason=relative-heading-rank; base=h2

\subsection{第二十一章  论证唯独应敬拜这位天主，祂禁止敬拜其他神祇，而其他神祇并不禁止人敬拜祂。}

29. 既然如此，这些不幸的人为何不领悟：这位天主是真天主，他们看出祂被置于与他们自己的神祇完全分开的地位，以至于虽然他们不得不承认祂是天主，但那些宣称当敬拜一切神祇的人，本身却不得与其余神祇一起敬拜祂？如今，既然这些神祇和这位天主不能一同被敬拜，为何不选择那位禁止敬拜其他神祇的祂？为何不放弃那些不禁止敬拜祂的神祇呢？或者，如果他们确实禁止敬拜祂，就请读出禁令。有什么比这更应当在他们的神庙中向民众宣读的呢？然而在那里从未听到过这样的声音。事实上，如此众多的（神祇）反对一位（神祇）的禁令，应当比一位（神祇）反对如此众多（神祇）的禁令更显著、更有力。因为，如果敬拜这位天主是不敬，那么那些不阻止人做这种不敬之事的神祇就是无益的；但若敬拜这位天主是虔诚，那么，既然在那种敬拜中命令不得敬拜这些神祇，敬拜他们就是不敬。再者，如果这些神祇禁止敬拜祂，却如此畏怯，以致他们宁可担心被人听见，而不敢禁止，那么谁还会如此不明智，不从这一事实得出结论？谁看不出应当选择这位天主，祂如此公开地禁止敬拜他们，命令摧毁他们的偶像，预言了他们被摧毁，并且亲自实现了这摧毁；而宁可选择那些神祇？我们不知道他们曾命令禁绝敬拜祂，没有读到他们预言过这样的事，也看不出他们有足够的能力使其实现。我问这个问题，让他们回答：这位天主是谁，祂如此扰乱所有外邦人的神祇，如此揭露他们的一切神圣仪式，如此使他们归于灭绝？\par

\SourceRecord{Translated by S.D.F. Salmond.；Nicene and Post-Nicene Fathers, First Series；Vol. 6.；Edited by Philip Schaff.；Buffalo, NY: Christian Literature Publishing Co.,；1888.；<http://www.newadvent.org/fathers/1602121.htm>.}

% structure: p0001 source=h1 target=section reason=page-title

\section{\WorkTitle{福音的和谐，卷一}}

% structure: p0002 source=h2 target=subsection reason=relative-heading-rank; base=h2

\subsection{第22章 论外邦人对我们天主的看法}

30. 然而，我为何要质问那些天生才智已丧失的人，他们无法回答这位天主是谁的问题？有人说祂是\OriginalTerm{萨图尔努斯}{Saturn}。我想，原因在于安息日的圣化，因为那些人将那一天归给\OriginalTerm{萨图尔努斯}{Saturn}。但他们自己的\Person{瓦罗}——在他们中间，他们不能指出有谁比他更有学问——认为犹太人的天主就是\OriginalTerm{朱庇特}{Jupiter}，并且他判断，无论用什么名字，只要所指的对象相同，就无关紧要；我认为，从这一点我们可以看出，他是如何被天主的至高无上所征服的。因为，罗马人习惯崇拜的对象没有比\OriginalTerm{朱庇特}{Jupiter}更崇高的，他们的卡比托利欧山便是明显而充分的证据，并视他为众神之王；当他观察到犹太人崇拜至高天主时，他无法想到任何其他对象配得上那个称号，除了\OriginalTerm{朱庇特}{Jupiter}本身。但是，无论人们将希伯来人的天主称为\OriginalTerm{萨图尔努斯}{Saturn}，还是宣称祂是\OriginalTerm{朱庇特}{Jupiter}，让他们告诉我们，\OriginalTerm{萨图尔努斯}{Saturn}何时敢于禁止崇拜另一位神祇。他甚至不敢禁止崇拜这位据说是将他赶出王国的\OriginalTerm{朱庇特}{Jupiter}——儿子驱逐父亲。如果\OriginalTerm{朱庇特}{Jupiter}作为更强大的神祇和征服者，已被他的崇拜者接受，那么他们就不应该崇拜被征服和被驱逐的\OriginalTerm{萨图尔努斯}{Saturn}。但另一方面，\OriginalTerm{朱庇特}{Jove}也没有禁止对他的崇拜。不，那位他有能力战胜的神祇，他仍然容许其继续为神。\par

\SourceRecord{Translated by S.D.F. Salmond.；Nicene and Post-Nicene Fathers, First Series；Vol. 6.；Edited by Philip Schaff.；Buffalo, NY: Christian Literature Publishing Co.,；1888.；<http://www.newadvent.org/fathers/1602122.htm>.}

% structure: p0001 source=h1 target=section reason=page-title

\section{\WorkTitle{福音的和谐 第一卷}}

% structure: p0002 source=h2 target=subsection reason=relative-heading-rank; base=h2

\subsection{第二十三章 论外教人在\Person{朱庇特}和\Person{萨图尔努斯}方面所沉溺的愚行}

31. 他们说：你们这些叙述不过是寓言，需要智者来解释，否则只配被嘲笑；但我们尊敬\Person{玛罗}所说的那位\Person{朱庇特}，他说：\par

\Quote{万物充满了\Person{朱庇特}，}\newline{}——\Person{维吉尔}的\WorkTitle{《牧歌》}，iii. v. 60\par

这就是说，是赋予万物生机的生命之灵。因此，\Person{瓦罗}认为\Person{朱庇特}受到犹太人的崇拜，并非毫无道理；因为犹太人的天主通过祂的先知说：\BibleQuote{我充满了天地。}但同一诗人所命名的以太是什么意思？他们如何理解这个术语？因为他这样说道：\par

然后全能的父亲以太，携带着丰沛的雨水，\newline{}降入他那多产的配偶的怀中。\newline{}——\Person{维吉尔}的\WorkTitle{《农事诗》}，ii. 325\par

他们确实说，这个以太不是精神，而是一个高耸的物体，天空在其上延伸于空气之上。诗人是否被允许一会儿用\Person{柏拉图}追随者的语言说话，仿佛天主不是身体，而是精神，一会儿又用斯多葛学派的语言说话，仿佛天主是一个身体？那么，他们在\Place{卡皮托尔}崇拜的是什么呢？如果它是一个精神，或者，简言之，如果它就是有形的天空本身，那么他们称为埃癸斯的\Person{朱庇特}盾牌又是怎么回事？这个名字的起源，确实是由一头母山羊在\Person{朱庇特}被他母亲藏匿时哺育了他这一情况来解释的。难道这是诗人的虚构吗？那么，罗马的\Place{卡皮托尔}宫也是诗人的纯粹创造吗？你们根据哲学家的观念在书中寻求你们的神，又按照诗人的观念在你们的寺庙中崇敬他们，这种即使不是诗的，也是显然滑稽的多变性，又是什么意思呢？\par

32. 但是，那位\Person{欧赫美鲁斯}不也是诗人吗？他极其直白地宣称\Person{朱庇特}本人、他的父亲\Person{萨图尔努斯}、以及他的兄弟\Person{普路托}和\Person{涅普顿}都是人，以至于他们的崇拜者更应当感谢诗人，因为诗人的虚构与其说是意在诋毁他们，不如说是意在粉饰他们？尽管\Person{西塞罗}提到，这位\Person{欧赫美鲁斯}被诗人\Person{恩尼乌斯}译成了拉丁文。或者，\Person{西塞罗}本人是诗人吗？他在\WorkTitle{《图斯库卢姆论辩》}中与论辩对手商议时，称他为知晓秘密事情的人，说了以下的话：\Quote{诚然，我若试图探究古事，从那里拿出希腊作家呈给世人的题材，就会发现，即使是那些被认为是高级神祇的神，也已从我们这里升到了天上。问问在\Place{希腊}有哪些坟墓被指出来；回想一下，既然你已入教，那些奥秘中传授的事物；那么，无疑你就会明白这种信念多么广泛。}这位作者确实充分承认他们的那些神原本是人的教义。他确实善意地推测他们升到了天上。但他毫不犹豫地公开说，就是众人普遍给予他们的那种荣誉，也是由人授予的，当他这样谈到\Person{罗慕路斯}时：\Quote{凭善意和声誉，我们把那位创建此城的\Person{罗慕路斯}提升到了不朽的神祇之列。}因此，认为更古老的人对\Person{朱庇特}、\Person{萨图尔努斯}和其他神所做的，正如罗马人对\Person{罗慕路斯}所做的，并且事实上他们在更近的时代也考虑过对\Person{凯撒}所做的那样，有什么可惊奇的呢？对于这些，\Person{维吉尔}又加上了谄媚的诗句，说：\par

\Quote{看，\Person{凯撒}的星，\Person{狄俄涅}的后裔，升起了。}\newline{}——\Person{维吉尔}的\WorkTitle{《牧歌》}，ix. 47\par

那么，他们当心，免得历史的真相反而向我们展示在世上为他们伪神竖立的坟墓！并当心，免得诗的虚浮，不是为他们在天上固定星辰，而只是假托星辰！因为，实际上，那颗星不是\Person{朱庇特}的星，这颗星也不是\Person{萨图尔努斯}的星；而简单的事实是，在自创世以来就已安置的这些星辰上，那些人死后被那些想在他们离世后尊他们为神的人加上了他们的名字。关于这些，我们确实可以问：为何贞洁如此不配，而淫逸如此配得，以致\Person{维纳斯}有一颗星，而\Person{弥涅耳瓦}却在那些与日月同行的星辰中得不到一颗？\par

33. 但也许可以说，那位敢于提及他们神的坟墓并加以记载的学院派智者\Person{西塞罗}，比起诗人来更是可疑的权威；虽然他不曾贸然将此论断作为他本人的个人意见，而是把它作为包含在他们自己神圣仪式传统中的说法记录下来。那么，是否也可以说，\Person{瓦罗}要么仅仅表达了像诗人那样的个人创作，要么像学院派那样只是含糊其辞，因为他宣称，在所有这类神祇的例子中，他们崇拜的事宜都源于他们作为人的生平或死亡？或者，那位\Place{埃及}祭司\Person{莱昂}，他给\Place{马其顿}的\Person{亚历山大}解释他们这些神的起源，解释方式虽与希腊人的意见略有不同，但在指出他们的神原本是人这一点上却如此一致，难道他不是一个诗人或学院派吗？\par

34. 但这与我们何干？让他们声称他们敬拜的是朱庇特，而非一个死人；让他们坚持说他们将卡皮托利山奉献给的不是一个死人，而是那赋予万物生命、充满世界的神灵。至于那面用母山羊皮制成的、为纪念他的乳母的盾牌，就让他们随意解释吧。然而关于萨图尔努斯，他们又怎么说？他们以萨图尔努斯的名义敬拜的，难道不是这位神灵？他不是第一个从奥林匹斯山降临到我们中间的神祇吗？诗人唱道：\par

\Quote{然后从奥林匹斯山巅下来\newline{}善良的萨图尔努斯，被他的更强大的继承人\newline{}朱庇特夺走了王冠，流放他乡：\newline{}他首先使原本分散在山顶的\newline{}种族联合起来，\newline{}并给他们律法遵守，\newline{}并意愿他躺卧的土地\newline{}应被称为拉丁姆。}\newline{}——维吉尔的\WorkTitle{埃涅阿斯纪}卷八第320-324行，科宁顿译本\par

他的雕像那被遮盖的头部，不就是把他呈现为一个隐藏者吗？难道不是他向意大利人民传授了农耕技术，这由镰刀所表示吗？不，他们说；因为你可以看到，被如此记载的那位存在是一个人，而且确实是一位特定的国王；然而，我们将萨图尔努斯解释为\Quote{普遍的时间}，这也在他的希腊名称中表明：因为他被称为克洛诺斯，这个词加上送气音后也是表示时间的词语；因此，在拉丁语中他被称为萨图尔努斯，仿佛意味着他\Quote{年岁饱足}。但现在，我不知道该如何看待这些人，他们竭力给他们的神祇的名称和形象赋予更有利的意义，却承认他们的主神和众神之父就是时间。这除了暴露他们所有的神祇都是暂时的，因为他们的根本之父被证明是时间，还能说明什么呢？\par

35. 因此，他们中较晚近的柏拉图学派哲学家，在基督徒时代兴盛起来，却耻于此类幻想，并试图以另一种方式解释萨图尔努斯，声称他得名\OriginalTerm{克洛诺斯}{\GreekText{Χρόνος}}，仿佛是为了表示理智的饱满；他们的解释是：希腊语中\Quote{饱满}由\OriginalTerm{饱满}{\GreekText{χόρος}}表示，\Quote{理智}或\Quote{心灵}由\OriginalTerm{理智}{\GreekText{νοῦς}}表示；这一词源似乎也受到拉丁语名称的支持，假设这个单词的前半部分（\OriginalTerm{萨图尔努斯}{Saturnus}）来自拉丁语，后半部分来自希腊语：因此他被称为萨图尔努斯，相当于\Quote{饱满的理智}。因为他们看到，将朱庇特视为时间之子是多么荒谬——朱庇特要么被他们视为永恒的神祇，要么他们希望他被视为永恒的神祇。此外，根据这种新奇解释——如果他们的古代权威真的持有这种看法，西塞罗和瓦罗竟未注意到它，这令人惊奇——他们称朱庇特为萨图尔努斯之子，从而可能表示他是从那个最高理智发出的精神——这个精神他们选择视为世界的灵魂，可以说，充满了所有天上和地上的身体。由此也产生了马罗的那句话，我稍早引用过，即\Quote{万物充满朱庇特}。那么，如果他们有能力，他们难道不应该改变人们所沉迷的迷信，就像他们改变自己的解释一样吗？或者根本不树立雕像，或者至少为萨图尔努斯而不是朱庇特建造卡皮托利山？因为他们也坚持认为，除非分有他那至高不变的上智，否则不可能产生出具有智慧的理性灵魂；这一论断不仅针对人的灵魂，甚至也针对他们称之为朱庇特的同一个世界灵魂。现在，我们不仅承认，而且特别宣讲，确实有某种天主的上智，任何真正有智慧的灵魂通过分有它而获得智慧。但那个被称为世界的普遍物质体是否具有一种灵魂，或所谓它自己的灵魂，即一种理性生命，通过它能够支配自己的运动，就像每种动物那样，这是一个既宏大又模糊的问题。这种意见不应被肯定，除非其真实性被明确确证；也不应被否定，除非其虚假性被明确确证。即使这个问题永远无法解决，对人又有什么要紧呢？因为无论如何，没有一个灵魂是凭借从任何其他灵魂，而是只能从那唯一至高不变的天主上智汲取而成为智慧或蒙福的。\par

36. 然而，那些在卡皮托利山上敬拜朱庇特而非萨图尔努斯的罗马人，以及那些认为朱庇特应当高于其他神祇受到敬拜的其他民族，他们当然不同意前面提到的那些人的观点；那些按照他们疯狂看法，若在这些事上有能力，宁愿将他们的最高堡壘奉献给萨图尔努斯，并且最特别地想要消灭那些占星术士和占星家——这些人为萨图尔努斯（对手们称其为智者的创造者）在其他星辰中赋予了凶神星相。然而，这种看法在人类心中如此强烈地反对他们，以至于人们甚至不愿提这个神的名字，而称他为\Quote{古神}而非萨图尔努斯；而且出于一种如此可怕的迷信精神，迦太基人现在几乎更改他们城镇的名称，更多时候称其为\Quote{古神之城}，而非\Quote{萨图尔努斯之城}。\par

\SourceRecord{Translated by S.D.F. Salmond.；Nicene and Post-Nicene Fathers, First Series；Vol. 6.；Edited by Philip Schaff.；Buffalo, NY: Christian Literature Publishing Co.,；1888.；<http://www.newadvent.org/fathers/1602123.htm>.}

% structure: p0001 source=h1 target=section reason=page-title

\section{\WorkTitle{福音合参，卷一}}

% structure: p0002 source=h2 target=subsection reason=relative-heading-rank; base=h2

\subsection{第二十四章 论拒绝以色列的天主之人事实上无法敬拜所有神；反之，敬拜其他神之人也无法敬拜祂}

因此，这些崇拜偶像之人实际上所尊崇的是什么，以及他们试图掩饰的是什么，是显而易见的。但是，就连这些萨图尔努斯的新解释者也必须告诉我们他们对希伯来人的天主有什么看法。因为在他们看来，敬拜众神也是合理的，正如外邦人所行的那样，因为他们的骄傲使他们耻于在基督面前谦卑自己以获得赦罪。那么，他们对以色列的天主持什么看法呢？因为如果他们不敬拜祂，那么他们就没有敬拜众神；而如果他们敬拜祂，他们也不是以祂所规定的方式敬拜祂，因为他们还敬拜其他祂所禁止敬拜的神。针对这些做法，祂通过同样的那些先知颁布了禁令，也正是通过这些先知，祂预先宣告了那些注定要发生的事，即他们的偶像如今正从基督徒手中遭受的事。无论解释如何，要么是天使被派往那些先知，以预象的方式，并通过可见对象的恰当形式，向他们显示那唯一真天主、万物的创造者，宇宙万物都服从于祂，并指示祂所命令的敬拜应如何进行；要么是他们中有些人的心智被圣神如此有力地提升，以致他们能够在天使自身目睹对象的那类神视中看见那些事物：无论如何，这是一个无可辩驳的事实，即他们确实事奉了那位禁止敬拜其他神的天主；而且，同样确定的是，他们以虔诚的忠信，在君王和司铎的职务上，同时为国家的福祉和那些神圣典章的利益而服务，这些典章预示着基督作为真君王和真司铎的来临。\par

\SourceRecord{Translated by S.D.F. Salmond.；Nicene and Post-Nicene Fathers, First Series；Vol. 6.；Edited by Philip Schaff.；Buffalo, NY: Christian Literature Publishing Co.,；1888.；<http://www.newadvent.org/fathers/1602124.htm>.}

% structure: p0001 source=h1 target=section reason=page-title

\section{\WorkTitle{福音的和谐}，卷一}

% structure: p0002 source=h2 target=subsection reason=relative-heading-rank; base=h2

\subsection{第二十五章  论假神不禁人同时敬拜别的神；以色列的天主是真天主，由其预言与应验之事证明}

38. 然而，在外邦人的诸神中（他们愿意敬拜这些神，表明他们不愿敬拜那位不能与他们一起被敬拜的天主），请他们告诉我们，为何在众多神明中，没有一位曾想到要禁止敬拜另一位？他们却赋予诸神不同的职位与职能，并认为他们各自掌管专属领域的事物？因为如果 \Person{朱庇特} 不禁止敬拜 \Person{萨图尔努斯}，那是因为他不仅被视为一个人，即驱逐其父王的人，而应被视为苍天之体，或充满天地之神明，因此他不能阻止那从他而出的至圣心灵被敬拜；同样，\Person{萨图尔努斯} 也不能禁止敬拜 \Person{朱庇特}，因为他不仅被视为被反叛者征服的人——如同一位同名人被某个所谓的 \Person{朱庇特} 击败而逃往意大利——而且因为原初心灵偏爱出自它的心灵。但至少 \Person{伏尔甘} 应该禁止敬拜 \Person{玛尔斯}，其妻的情夫；\Person{赫拉克勒斯} 也应禁止敬拜迫害他的 \Person{朱诺}。他们之间竟有如此污秽的默契，以至于连贞洁的 \Person{狄安娜} 也不禁止敬拜——我不说 \Person{维纳斯}，甚至 \Person{普里阿普斯}？因为同一个人若想同时做猎人与农夫，就必须服侍这两位神明；然而他甚至连为他们并排建造庙宇也会感到羞愧。但他们可能辩解，说 \Person{狄安娜} 按寓意代表某种美德，随便他们怎么说；他们或许告诉我们，\Person{普里阿普斯} 实际上象征丰饶之神——至少使 \Person{朱诺} 因有这样的帮助者来使妇女多产而感到羞愧。他们可以随意说，可以凭其智慧作任何解释：然而尽管如此，\Place{以色列} 的天主必推翻他们所有的论证。因为祂禁止敬拜所有这些神明，而祂的敬拜却不受任何神明阻碍，并且祂立刻命令、预告并实现毁灭他们的偶像与圣礼，这已足够清楚地表明他们是虚假说谎的神明，而祂自己是唯一真实的天主。\par

39. 此外，谁不感到奇怪：那些如今人数已少的众多假神的敬拜者，竟拒绝向祂致敬——当被问及这是哪位神明时，无论他们如何回答，至少不敢否认祂是神？因为若他们否认祂的神性，很容易被祂在预言与应验中的作为所驳斥。我不提那些他们自以为可以不信的作为，如起初祂创造天地及其中万物。我也不列举那些追溯至远古的事件，如 \Person{哈诺客} 被接升天、洪水毁灭不敬虔者、义人 \Person{诺厄} 及其全家藉木舟得救。我以 \Person{亚巴郎} 开始陈述祂在人间的事迹。确实，这位男子曾藉天使的神谕得到一个明确的应许，我们现在看到它实现了。因为对他说：\Quote{天下万民都要因你的后裔蒙受祝福。}从他的后裔中，产生了 \Place{以色列} 民族，从中出现了 \Person{童贞玛利亚}，她是基督的母亲；而万民都在祂内蒙受祝福，如果他们能，现在就大胆地否认吧。这同一应许也给了 \Person{亚巴郎} 的儿子 \Person{依撒格}。又给了 \Person{亚巴郎} 的孙子 \Person{雅各伯}。\Person{雅各伯} 又称为 \Place{以色列}，整个民族都从他得名与血统，因此这民族的天主被称为 \Place{以色列} 的天主：并非祂不是外邦人的天主，无论他们不认识祂或现在认识祂；而是祂愿意在这民族中更明显地彰显祂应许的力量。因为这民族起初在 \Place{埃及} 繁衍，后来藉 \Person{梅瑟} 之手，以许多征兆与奇事从那里脱离奴役，看见大多数外邦民族被制服于其下，并获得了应许之地，在那地由出生于 \Person{犹大} 支派的本族君王统治。这位 \Person{犹大} 也是 \Place{以色列} 十二子之一，\Person{亚巴郎} 的孙子。从他产生了称为犹太人的民族，他们藉天主自己帮助做出伟大事业，也当祂惩罚他们时，因他们的罪遭受许多苦难，直到那应许的后裔来临——万民要因祂蒙受祝福，并自愿打碎他们祖先的偶像。\par

\SourceRecord{Translated by S.D.F. Salmond.；Nicene and Post-Nicene Fathers, First Series；Vol. 6.；Edited by Philip Schaff.；Buffalo, NY: Christian Literature Publishing Co.,；1888.；<http://www.newadvent.org/fathers/1602125.htm>.}

% structure: p0001 source=h1 target=section reason=page-title

\section{《福音的和谐，第一卷》}

% structure: p0002 source=h2 target=subsection reason=relative-heading-rank; base=h2

\subsection{第二十六章　论基督之名与基督徒的信德如何按预言推翻了偶像崇拜}

40. 因为基督徒所成就的，并非仅属于基督教时代的事，而是早已预言的。那些至今仍与基督之名作对的犹太人——关于他们注定的背信，这些预言书并未缄默——他们自己拥有并诵读那位先知的话：\BibleQuote{我的主天主，我患难之日的避难所，外邦人要从地极来到你面前，说：\InnerQuote{我们的祖先诚然拜了虚妄的偶像，它们毫无益处。}}看，这事如今正在成就；看，外邦人正从地极来到基督面前，说着这样的话，打碎他们的偶像！天主为其普世教会所行的这事意义重大：那应得惩罚的犹太民族被分散到各地，到处携带我们预言的记录，使人不能以为这些预言是我们编造的；这样，我们信德的仇敌便成了真理的见证。那么，基督的门徒怎么可能教导他们未从基督学到的东西，如同那些愚昧人凭其荒唐的幻想所指责的那样，意图推翻对异教神明和偶像的迷信崇拜？难道那些至今仍在基督仇敌的书中被人诵读的预言，也是基督门徒的杜撰吗？\par

41. 那么，除以色列的天主外，谁废除了这些制度？因为正是向这百姓颁布了那些对梅瑟所说的神圣诫命：\BibleQuote{以色列，请听：主你的天主是唯一的天主。}\BibleQuote{你不可为你制造任何雕像，或任何天上、地上之物的肖像。}并且，为使这百姓在所到之处有能力除灭这些事物，又向他们颁下诫命：\BibleQuote{你不可向他们的神祇跪拜，也不可事奉他们；不可效法他们的作为，而要彻底推翻它们，打碎他们的雕像。}但谁能说基督和基督徒与以色列无关呢？以色列原是亚巴郎的孙子，而亚巴郎、他的儿子依撒格、及他的孙子以色列自己，先後领受了那项我已提到的应许：\BibleQuote{地上万民都要因你的后裔蒙受祝福。}我们如今看见这预言在基督身上应验。因为童贞女出自这一支系，以色列子民和以色列天主的先知曾这样歌颂她：\BibleQuote{看，童贞女将怀孕生子，人要称他的名为厄玛奴耳。}因为厄玛奴耳解释出来就是\BibleQuote{天主与我们同在。}因此，这位禁止崇拜他神、禁止制造偶像、命令销毁偶像、并藉先知预言外邦人从地极要说\BibleQuote{我们的祖先诚然拜了虚妄的偶像，毫无益处}的以色列天主，正是那藉基督之名和基督徒的信德，计划、应许并实现了推翻这一切迷信的天主。因此，这些不幸的人——他们明知连自己的神祇（即畏惧基督之名的魔鬼）也禁止他们亵渎基督之名——却徒劳地试图证明这种教义与基督无关，而基督徒正是凭这教义的力量反驳偶像，并在有机会时根除所有虚假宗教。\par

\SourceRecord{Translated by S.D.F. Salmond.；Nicene and Post-Nicene Fathers, First Series；Vol. 6.；Edited by Philip Schaff.；Buffalo, NY: Christian Literature Publishing Co.,；1888.；<http://www.newadvent.org/fathers/1602126.htm>.}

% structure: p0001 source=h1 target=section reason=page-title

\section{\WorkTitle{福音的和谐}，第一卷}

% structure: p0002 source=h2 target=subsection reason=relative-heading-rank; base=h2

\subsection{第二十七章 一项说服残余的偶像崇拜者最终成为这位处处摧毁偶像的真天主的仆人的论证。}

42. 现在让他们就以色列的天主做出回答吧。这位天主教导并命令这样的事，不仅有基督徒的书籍为他作证，犹太人的书籍也为他作证。关于他，让他们从自己的神祇那里寻求意见，这些神祇阻止了对基督的亵渎。关于以色列的天主，如果他们胆敢，就让他们给出一个无礼的回答吧。但是，现在他们要向谁请教呢？或者他们要在哪里寻求意见呢？让他们细读自己权威著作的书籍吧。如果他们认为以色列的天主是朱庇特，正如\Person{瓦罗}所写的那样（为了暂时附和他们的思维方式，我姑且这样说），那么他们为什么不信由朱庇特来摧毁偶像呢？如果他们认为他是撒杜恩，那么他们为什么不敬拜他呢？或者，他们为什么不以那位天主通过先知的声音所规定的方式敬拜他呢？这些先知所言之事，天主已使其应验。他们为什么不信由他来摧毁偶像，并禁止敬拜其他神祇呢？如果他不是朱庇特，也不是撒杜恩（而且，如果他真是其中之一，他肯定不会如此强烈地抨击针对他们的朱庇特和撒杜恩的宗教仪式），那么这位天主是谁呢？在所有其他神祇都受到敬拜的情况下，唯有他未受敬拜。然而，他如此明显地促成自己成为唯一敬拜的对象，推翻所有其他神祇，并使一切自以为高傲、自高自大、与基督为敌（因偶像而迫害和杀害基督徒）的事物蒙羞。但是，事实上，人们现在在问，当这些崇拜者想要献祭时，他们退避到什么样的秘密角落？或者，他们把自己的这些神祇又退回至何等的黑暗区域，以防止它们被基督徒发现并打碎？这种处理方式从何而来，如果不是出于对那些法律和统治者的恐惧，天主藉着这些法律和统治者展示他的能力，而他们如今已服从了基督的名？天主早在古时就已应许这事，当他藉着先知说：\BibleQuote{普世的君王都要崇拜他，万民都要事奉他。}\par

\SourceRecord{Translated by S.D.F. Salmond.；Nicene and Post-Nicene Fathers, First Series；Vol. 6.；Edited by Philip Schaff.；Buffalo, NY: Christian Literature Publishing Co.,；1888.；<http://www.newadvent.org/fathers/1602127.htm>.}

% structure: p0001 source=h1 target=section reason=page-title

\section{\WorkTitle{福音的和谐，第一卷}}

% structure: p0002 source=h2 target=subsection reason=relative-heading-rank; base=h2

\subsection{\Emph{第二十八章 论对偶像摈弃的预言}}

43. 祂的先知们在不同时期所预言的，如今正在应验，这是无可置疑的——即祂将抛弃祂那不信神的子民（其实并不是整个民族，因为连以色列人中也有许多人相信了基督；因为祂的宗徒们本身就出自该民族），并要贬抑一切傲慢和有害的人，使祂自己独自被高举，也就是说，唯独祂自己向人显明为崇高而大能的；直到偶像被信徒抛弃，被不信者隐藏；那时大地因对主的畏惧而崩裂，即世上的人因畏惧而屈服，也就是因畏惧祂的法律，或畏惧那些既信祂的名又作万邦统治者的人的法律，这些人将禁止此类亵渎的作为。\par

44. 因为以上我所简要陈述的这些事情，为了便于理解，先知是这样表述的：\BibleQuote{雅各伯家啊，来罢，让我们行在主的光明中。因为祂已抛弃了祂的子民以色列家，因为那地从头就充满了他们和异邦人的占卜，又为他们生了许多异族的孩子。因为那地充满了金银，他们的财宝不可胜数；那地也充满了马匹，他们的战车不可胜数；那地也充满了他们亲手所造的可憎之物，他们崇拜自己手指所做的。卑贱人屈膝，尊贵人降卑；我必不宽恕他们。现在你们要进入磐石，藏在尘土中，躲避主的威吓和祂威严的荣光，当祂起来震动大地时：因为主的眼目崇高，而人卑微；人的骄傲必被抑低，惟独主在那日被高举。因为万军之主的日子要临到一切有害和骄傲的人，临到一切自高自大而被贬抑的人，他们必被降卑；又临到黎巴嫩的一切高大香柏树和一切橡树，临到巴珊的一切树木，临到一切山岳，一切高冈，临到海中的一切船只，以及一切华丽的船帆。人的侮慢必被抑低而倾倒，惟独主在那日被高举；凡人手所造的一切，他们必藏在洞穴、磐石缝和地窟中，躲避主的威吓和祂威严的荣光，当祂起来粉碎大地时：因为在那日，人必抛弃金银所做的可憎之物，就是他们为敬拜而制造的虚妄邪恶之物，以便进入坚石之罅隙和岩石之孔洞，躲避主的威吓和祂威严的荣光，当祂起来将大地打得粉碎。}\par

\SourceRecord{Translated by S.D.F. Salmond.；Nicene and Post-Nicene Fathers, First Series；Vol. 6.；Edited by Philip Schaff.；Buffalo, NY: Christian Literature Publishing Co.,；1888.；<http://www.newadvent.org/fathers/1602128.htm>.}

% structure: p0001 source=h1 target=section reason=page-title

\section{\WorkTitle{福音的和谐，卷一}}

% structure: p0002 source=h2 target=subsection reason=relative-heading-rank; base=h2

\subsection{\OriginalTerm{第二十九章}{Chapter 29} 论外邦人为何拒绝崇拜以色列的天主，即使他们视祂仅为元素的主神？}

45. 关于这位\OriginalTerm{万军}{Sabaoth}的天主，他们的看法如何？这一词按解释意指众能力或众军的天主，因为天使的众能力和众军都事奉祂。关于这位以色列的天主，他们看法如何？因为祂是那民族的天主，从该民族而来那将要祝福万民的苗裔。为什么那些主张所有神祇都应当受到崇拜的人，偏偏将这位神祇排除在崇拜之外？为什么他们拒绝相信那既证明其他神祇为假神又推翻它们的天主？我曾听见他们中一位说，他在某位哲学家那里读到，从犹太人在其神圣仪式中所做的事，他认出了他们所崇拜的是哪一位天主。他说：\Quote{祂是那位神明，掌管构成这个可见物质宇宙的元素。}然而在祂先知的圣经中清楚显示，以色列民族被命令崇拜那创造天地、一切真智慧之源的天主。但对此问题何需进一步争辩，因为对我当前目的而言，指出他们是如何对那位他们仍无法否认是神的天主怀有各种专断见解，已足够。如果祂真是掌管构成这个世界之元素的神明，为什么祂不被优先崇拜，胜过只掌管海洋的\OriginalTerm{尼普顿}{Neptune}？为什么又不被优先崇拜，胜过只掌管田野和森林的\OriginalTerm{西尔瓦努斯}{Silvanus}？为什么不被优先崇拜，胜过只掌管白昼的太阳，或祂也统治整个天体的热力？为什么不被优先崇拜，胜过只掌管黑夜的月亮，或祂也因掌管湿气而显耀？为什么不被优先崇拜，胜过被认为只掌管空气的\OriginalTerm{尤诺}{Juno}？因为那些掌管部分的神明，无论他们是谁，必然低于那位掌管所有元素和整个宇宙的神明。但这位神明禁止崇拜所有那些神明。那么为何这些人，违背那位大于他们的神明的命令，不仅选择崇拜他们，而且还为了这些神明而拒绝崇拜祂？他们尚未发现任何关于这位以色列天主的恒定而明智的判断；他们也将永远找不到，直到他们发现唯独祂是真天主，万有由祂而造。\par

\SourceRecord{Translated by S.D.F. Salmond.；Nicene and Post-Nicene Fathers, First Series；Vol. 6.；Edited by Philip Schaff.；Buffalo, NY: Christian Literature Publishing Co.,；1888.；<http://www.newadvent.org/fathers/1602129.htm>.}

% structure: p0001 source=h1 target=section reason=page-title

\section{\WorkTitle{福音的和谐，第一卷}}

% structure: p0002 source=h2 target=subsection reason=relative-heading-rank; base=h2

\subsection{第三十章  论事实证明，随着预言的应验，以色列的天主如今已到处被认识。}

46. 有一个人名叫\Person{卢坎}，是他们中一位伟大的韵文演说家。我深信，他长久以来，或凭自己的思考，或藉阅读本国同胞的著作，试图查明犹太人的天主是谁；但因未能以虔诚的方式进行探讨，他没有成功。然而，他宁愿称祂为祂不曾找着的\OriginalTerm{不确定的天主}{uncertain God}，也不绝对否认那神的天主称号，因为他察觉到其存在有如此巨大的证据。因为他说道：\par

\Quote{犹大，信奉崇拜\newline{}一位不确定的天主。\newline{}——\Person{卢坎}，卷二，结尾处。}\par

然而，这位天主，以色列圣而真实的天主，在\Person{卢坎}的时代之后直到我们今日，尚未藉基督的名字在万民中成就如此伟大的工程，正如后来所成就的。但如今，谁还如此顽固而不被感动？谁还如此迟钝而不被点燃？因为圣经的话已经应验：\BibleQuote{因为没有什么能躲避它的热气。}；而且同一篇圣咏（我从其中引用了一小节）中如此久远以前预言的其他事情，也已在应验中最清楚地展现出来。因为在这\Quote{诸天}的名称下，指的是耶稣基督的宗徒们，因为天主要在他们当中临在，以宣扬福音。所以，如今诸天已彰显了天主的荣耀，穹苍已宣告了他手中的化工。日复一日传递言论，夜复一夜显示知识。如今没有言语，没有语言，听不到他们的声音。他们的声音传遍普世，他们的话语直到地极。如今祂在太阳中支搭了祂的帐幕，即在显明之中；这帐幕就是祂的教会。因为为这样做（正如经文接下来所说），祂像新郎一样从自己的洞房出来；这就是说，\OriginalTerm{圣言}{Word}与人的血肉结合，从童贞女的子宫中出来。如今祂像勇士一样欢跃，奔跑了祂的赛程。如今祂的出发从高天之上，祂的回归直到高天之上。因此，最恰当地，接下来就是我已提到的那节经文：\BibleQuote{没有什么能躲避它的热气（或：祂的热气）。}。然而，这些人选择他们那些渺小、软弱、饶舌的反对意见，它们在那火中像糠秕一样化为灰烬，而不像金子那样被炼净；同时，他们假神的虚假纪念碑已经化为乌有，而那\OriginalTerm{不确定的天主}{uncertain God}的真实许诺已被证明是确实的。\par

\SourceRecord{Translated by S.D.F. Salmond.；Nicene and Post-Nicene Fathers, First Series；Vol. 6.；Edited by Philip Schaff.；Buffalo, NY: Christian Literature Publishing Co.,；1888.；<http://www.newadvent.org/fathers/1602130.htm>.}

% structure: p0001 source=h1 target=section reason=page-title

\section{\WorkTitle{福音的和谐，第一卷}}

% structure: p0002 source=h2 target=subsection reason=relative-heading-rank; base=h2

\subsection{第三十一章　论及基督的预言的应验}

47. 因此，让那些赞美基督却拒绝成为基督徒的恶人，不要再指控基督没有教导要放弃他们的神祇、打碎他们的像。因为关于以色列的天主，早已预言祂要被称为全地的天主，如今确实被称为全地的天主。祂藉先知的口预言此事必将成就，并藉基督在适当的时候最终成就了。无疑，如果以色列的天主如今被称为全地的天主，那么祂所命令的必须实现；因为颁布诫命的那位如今已为人所知。此外，祂藉基督并在基督内被显明，为使祂的教会扩展至普世，并藉此使以色列的天主被称为全地的天主，凡愿意的人可在同一位先知稍前的地方读到。那段话我也可以引用。它并不长，无需我们跳过。这里有许多关于基督的临在、谦卑和受难，以及关于祂为元首的身体，即祂的教会的话，其中教会被称为不生育的，如同不能怀孕的。因为多年以来，那注定要存在于万民之中、与它的子女（即圣徒）同在的教会，并未显现，因为基督尚未被传福音者向那些先知未曾宣告过祂的人宣告。此外，经上说：被遗弃的那位的儿女，比有丈夫的那位的儿女更多；这里的\Quote{丈夫}指法律，或以色列子民最初接受的君王。因为先知发言时，外邦人尚未接受法律；基督的君王也尚未显现给万民，尽管如今从这些外邦人中已涌现出更丰硕、更众多的圣徒。因此，依撒意亚这样说，从基督的谦卑开始，然后转向对教会说话，直到我们已引用的那句：\BibleQuote{领你们出来的那位，就是以色列的天主，要被称为全地的天主。看，他说：我的仆人将行事明智，必受尊崇，被高举，且受极崇敬。许多人将因你而惊愕，以致你的容貌被毁，然而众人必看见你的尊荣，世人必看见你的荣耀。因为如此，许多国家将因他而惊愕，君王将闭口无言。因为他们看见了未曾有人告诉过他们的事，听到了未曾听闻过的事。上主，谁相信了我们的报道？上主的手臂向谁显露了？我们在祂面前如同嫩芽，像根出于干地；祂没有俊美，也没有华美。我们看见祂，祂没有美貌，也无威仪；祂的面容被轻视，祂的处境被众人厌弃：一个受苦的人，熟悉病痛；因此祂的面容被转开，受到伤害，不被重视。祂承担了我们的病患，为我们而忧伤。我们以为祂是受责罚，被击打，被折磨。但祂是为我们的过犯被刺透，为我们的罪恶被压伤；祂受的惩罚使我们得到平安，祂受的鞭伤使我们痊愈。我们都像羊一样迷了路，上主却把祂献为赎罪祭。祂受虐待，却仍不开口；祂像被牵去宰杀的羔羊，又像在剪毛者面前默不作声的羊，祂不开口。祂受审判时被剥夺了公义。谁能述说祂的世代？因为祂的生命从地上被剪除，因我百姓的罪过祂被引至死亡。因此，我要使恶人作祂的坟墓，使富人作祂的葬地，因为祂没有行过强暴，口中也没有诡诈。上主喜悦以祂的击打来洁净祂。若你们为你们的过犯献上灵魂，你们必看见长寿的后裔。上主喜悦从忧伤中取回祂的灵魂，使祂看见光明，并将祂显明，使那服事众人的义人成义；祂要担当他们的罪过。因此，祂必获得许多人为产业，且分得强者的战利品；因为祂的灵魂被交于死亡，祂被列于罪犯之中，祂担当了许多人的罪，并为他们的过犯而被交出。不生育的，不怀孕的，你应当欢呼！未曾经过产痛的，你应当高声歌唱！因为被遗弃的那位的儿女，比有丈夫的那位的儿女更多。因为上主说：拓展你帐幕的空间，固定你的庭院；不必吝惜，拉长你的绳索，坚固你的橛子。是的，你要向左向右扩张。因为你的后裔将继承外邦人，你将居住荒废的城市。你不要害怕。因为你必得胜，不至于羞愧，也不至于蒙羞。因为你必忘记你的羞耻，不再记念你寡居的羞辱，因为造你的是你的丈夫——万军的上主是祂的名；领你们出来的那位，就是以色列的天主，要被称为全地的天主。}\par

对于这些证据，这些既是预言又已应验之事的表现，还有什么可说的呢？如果他们以为祂的门徒对基督的神性作了虚假的见证，他们也会怀疑基督的受难吗？不：他们不习惯于相信祂从死者中复活了；但同时，他们却十分愿意相信祂承受了人通常所承受的一切，因为他们愿意祂只被当作一个普通人，而非其他。照此，祂被牵到宰杀之地，如同羔羊；祂被列于罪犯之中；祂为我们的罪受创；因祂的伤痕我们得了痊愈；祂的面容毁损，被轻视，被人掌掴，被唾沫玷污；祂在十字架上的位置被扭曲；祂因以色列人民的罪孽被引至死亡；祂是那位在被人拳打、被荆棘冠冕、被挂在木架上的时候受讥笑、毫无俊美与威仪的人；祂是那位如同羔羊在剪毛者前无声、不开口的人，当那些嘲笑祂的人对祂说：\Quote{基督，给我们说预言吧。}如今，祂确已被高举，如今祂备受尊荣；实在，许多国家都因祂而惊愕。如今君王们闭口不言，他们曾惯于颁布最残酷的反基督徒法律。实在，那些未曾被告知祂的人现在看见了，那些未曾听见的人现在明白了。因为那些先知未曾向他们宣告的外邦民族，如今反倒亲眼看见先知们古时所传的这些事是何等真实；而那些未曾亲耳听见\Person{依撒意亚}本人说话的人，如今从他的著作中明白了他关于祂所说的话。因为在所说的犹太民族中，谁信了先知们的报告呢？主的膀臂，即他们所宣布的这位基督，又向谁启示了呢？可见他们亲手对基督犯下了那些罪行，而这些罪行的施行早已被他们所拥有的先知们预言了。但现在，祂确实借着继承拥有了许多人；祂分得了强者的战利品，因为魔鬼和邪魔如今已被驱逐和放弃，他们一度占有的产业已被祂分给祂的教会的建筑物和其他必要的服务。\par

\SourceRecord{Translated by S.D.F. Salmond.；Nicene and Post-Nicene Fathers, First Series；Vol. 6.；Edited by Philip Schaff.；Buffalo, NY: Christian Literature Publishing Co.,；1888.；<http://www.newadvent.org/fathers/1602131.htm>.}

% structure: p0001 source=h1 target=section reason=page-title

\section{\WorkTitle{福音的和谐，第一卷}}

% structure: p0002 source=h2 target=subsection reason=relative-heading-rank; base=h2

\subsection{第32章 借用预言的话为宗徒的道理辩护，反对偶像崇拜}

49. 那么，这些既悖逆地赞美基督又诽谤基督徒的人，对这些事实作何解释？难道基督用法术使先知们如此久远之前预言了这一切？或是祂的门徒捏造了它们？难道教会虽一度荒芜，如今在外邦民族中扩展，因拥有比那在律法或君王中好像有了丈夫的会堂更多的儿女而欢喜？难道这教会因此被引导扩大她帐幕之地，占据万民万舌，以致如今她把绳索延伸到罗马帝国权力范围之外，甚至远达波斯人、印度人及其他蛮族的领土？或是藉着右手边的真基督徒和左手边的假基督徒，祂的名在如此众多的民族中被传扬？或是祂的后裔要承受外邦人，以至于如今居住在因缺乏对天主的真崇拜和真宗教而荒废的城邑？或是祂的教会，即使在殉道者鲜血覆盖如披紫红袍的时候，也如此不因人的恐吓和暴怒而畏惧，竟能胜过如此众多、暴烈而强大的迫害者？或是她并未像蒙羞之人那样困惑——当成为或成为基督徒是极大的罪行时？或是她得以永远忘记她的羞惭，因为罪在那里增多，恩宠却更加洋溢？或是她被教导不再记念她寡居的耻辱，因为她只是暂时被遗弃、遭受辱骂，如今却重现如此显赫的荣光？最后，难道这只是基督门徒编造的虚构故事——那位造她、将她从魔鬼和邪魔的权势中领出来的主，以色列的天主，如今被称为全地之神？然而所有这些事，先知们——现今在基督仇敌手中的书籍——在基督成为人子之前如此长久地预言了。\par

50. 因此，让他们明白，这件事即便对最迟钝、最愚笨的心思也并非模糊或可疑；我再说，让这些悖逆地赞美基督、却咒骂基督宗教的人明白，基督的门徒所学所教的反对他们神明的，正是基督教义所包含的。因为以色列的天主在先知书中命令：那些这些人执意崇拜的对象应当被憎恶并毁灭；而祂自己如今借着基督和基督的教会，被称作全地之神，正如祂许久以前所应许的。因为，如果他们真以其惊人之愚妄幻想基督曾崇拜他们的神明，且惟有藉着他们祂才能成就如此大事，我们大可以问：以色列的天主——祂如今已借着基督应验了祂关于自己崇拜扩展至万民、以及憎恶并推翻那些异神的应许——是否也曾崇拜他们的神明？他们的神明在哪里？他们的狂热者之预言和先知之占卜在哪里？鸟占、征兆、占卜或邪魔的神谕在哪里？为什么从那些构成此类宗教记录的古书中，拿不出任何形式的告诫或预言来反对基督徒的信仰，或驳斥我们那些如今在万民中已如此为人理解的先知的真理？他们回答说：\Quote{我们得罪了我们的神明，}他们回答说，\Quote{他们因此抛弃了我们；这也解释了为什么基督徒胜过我们，为什么人类生活的福祉耗尽受损，在我们中走向毁灭。}然而，我们挑战他们：拿起他们自己先见者的书，向我们读出任何陈述，指出降临在他们身上的这种结局是由基督徒带来的；不，我们挑战他们背诵任何段落，其中若不是基督（因为他们想把他描绘成自己神明的崇拜者），至少这位以色列的天主——他被承认是其他神明的推翻者——被当作一位要被弃绝、应受憎恶的神明。但他们永远拿不出这样的段落，除非可能是他们自己编造的。如果他们真的引用了任何这样的声明，那不过是他们自己的虚构，这将会在如此严肃之事被引述时的不引人注目中暴露无遗；而在真理中，在预言应验之前，这应当在万国神明的庙宇中被宣告，以便适时预备和警告所有现在愿意成为基督徒的人。\par

\SourceRecord{Translated by S.D.F. Salmond.；Nicene and Post-Nicene Fathers, First Series；Vol. 6.；Edited by Philip Schaff.；Buffalo, NY: Christian Literature Publishing Co.,；1888.；<http://www.newadvent.org/fathers/1602132.htm>.}

% structure: p0001 source=h1 target=section reason=page-title

\section{福音的和谐，卷一}

% structure: p0002 source=h2 target=subsection reason=relative-heading-rank; base=h2

\subsection{第三十三章 驳斥那些抱怨基督徒时代到来损害人生幸福之人的论述}

51. 最后，关于他们抱怨基督徒时代的进入损害了人生幸福，如果他们只是翻阅自己那些哲学家的著作——这些哲学家谴责的正是那些尽管他们不情愿并抱怨却正在被移除的事物——他们就会确实发现，基督的时代是极值得赞美的。因为他们的幸福除了在最卑鄙、最奢侈地滥用（这严重损害了他们的造物主）方面有所减少外，还能在其他方面减少呢？或者，也许恶劣的时代源于这样的情况：在几乎所有的邦国中，剧院正在衰落，随之还有罪恶的巢穴和公开的邪恶行径；是的，就连过去用来崇拜魔鬼的广场和城市也正在倾颓。那么，它们为何倾颓呢？难道不是因为正是那些建造它们时所贪欲地、亵渎地使用的东西的衰败吗？他们自己的\Person{西塞罗}在称赞一位名叫\Person{罗修斯}的演员时，不是称他是一位\Quote{如此聪明的人，以至于他是唯一配得上理应登台的人；而同时，他又是一位如此善良的人，以至于他是唯一配得上理应不接近舞台的人}吗？他这句名言还如此清晰地揭示了什么呢？无非是，舞台在那里是那样卑贱，一个人越是有德行，就越有责任不与它发生关联。然而，他们的神明却喜爱这些他认为只应远离好人的可耻之事。但我们也有同一\Person{西塞罗}的公开承认，他说他不得不通过频繁的庆典来安抚\Person{福罗拉}，这位体育之母；在这些庆典中，通常展示出如此过度的邪恶，以至于相比之下，其他活动倒显得可敬了，然而善良之人仍被禁止参与那些活动。这位母亲\Person{福罗拉}是谁？她是怎样的女神，竟要藉着放纵邪恶的行为（且比平时更频繁、更放荡）来讨好和平息她的怒气？如今，\Person{罗修斯}登上舞台比\Person{西塞罗}崇拜这样的女神要光荣得多！如果外邦人的神明因为为此类庆典而设立的供奉减少而生气，那么，那些喜爱这类事物者的品性也就显而易见了。但另一方面，如果神明自己在愤怒中减少这些供奉，那么他们的愤怒比他们的息怒为我们提供更好的服务。因此，这些人要么驳斥他们自己的哲学家——这些哲学家已经谴责了放荡者方面的相同做法；要么就粉碎他们那些对崇拜者提出如此要求的神明——如果他们还能找到任何这样的神明去粉碎或隐藏的话。但让他们停止他们亵渎的习惯，即指责基督徒时代是导致他们真正繁荣丧失的原因——实际上，他们如此滥用这种繁荣，以至于陷入了所有卑劣和有害的事物中——以免他们反而更加强烈地提醒我们要更充分地赞美基督的大能。\par

\SourceRecord{Translated by S.D.F. Salmond.；Nicene and Post-Nicene Fathers, First Series；Vol. 6.；Edited by Philip Schaff.；Buffalo, NY: Christian Literature Publishing Co.,；1888.；<http://www.newadvent.org/fathers/1602133.htm>.}

% structure: p0001 source=h1 target=section reason=page-title

\section{\WorkTitle{福音的和谐，卷一}}

% structure: p0002 source=h2 target=subsection reason=relative-heading-rank; base=h2

\subsection{第三十四章 前章跋语}

52. 关于这个问题，我本可说得更多，若非我所承担的任务之要求迫使我结束本书，并回归最初设定的目标。事实上，当我着手解决福音书中那些在四福音作者之间似乎彼此不协调的问题（正如某些批评者所见），并尽力阐述他们各自特有的意图时，我首先遇到了一个讨论的必要，即有些人惯常向我们提出的问题：为何我们不能拿出任何由基督亲自撰写的著作？因为他们的目的是要使人相信基督写过某种其他作品——我不知道是何种作品——以迎合他们的倾向，并且相信基督对他们的神明毫无敌意，反而以一种巫术式的崇拜尊敬他们；他们还希望使人相信，基督的门徒不仅虚假地报道了祂，宣称祂是创造万物的天主，而实际上祂不过是一个人——尽管确实是一位具有最高智慧的人——而且就他们这些神明而言，门徒们所教导的与从基督那里学来的不同。正因如此，我们优先以关于以色列天主的论证来驳斥他们；这位天主如今通过基督徒的教会受万民敬拜，祂还在全世界推翻他们亵渎的虚妄，正如祂许久前藉先知的口所宣告的，并且如今藉基督的名实现了这些预言，祂曾应许万民要因基督的名蒙福。从这一切，他们应当明白：基督对于他们的神明，不可能知道或教导任何与以色列天主藉其先知所命令和预言的不同之事；以色列天主藉这些先知应许并最终派遣了这位基督，因着给祖先的应许，万民被宣称为有福，而这位以色列天主也成了全地的天主。他们也应当由此看出，基督的门徒并未偏离他们师傅的教训，因为他们禁止敬拜外邦人的神明，以免我们向无感觉的偶像祈祷，或与魔鬼交往，或以宗教崇拜的敬意事奉受造物而非造物主。\par

\SourceRecord{Translated by S.D.F. Salmond.；Nicene and Post-Nicene Fathers, First Series；Vol. 6.；Edited by Philip Schaff.；Buffalo, NY: Christian Literature Publishing Co.,；1888.；<http://www.newadvent.org/fathers/1602134.htm>.}

% structure: p0001 source=h1 target=section reason=page-title

\section{福音的和谐 第一卷}

% structure: p0002 source=h2 target=subsection reason=relative-heading-rank; base=h2

\subsection{第35章 论中保的奥迹如何藉预言为古代的人所知晓，如今又藉福音向我们宣告}

53. 因此，既然基督本身就是那创造万物的天主的智慧，并且考虑到无论是天使还是人的理性受造物，除非藉参与这一智慧，并藉着那将爱德倾注在我们心中的圣神与这智慧结合（而这天主圣三同时构成唯一的天主），否则不能获得智慧；天主眷顾考虑到必死之人的利益，其时间性的生命被束缚于产生和消亡的事物中，便预定这同一的智慧，在其位格的统一中取了人性，在此人性中他可按照时间条件出生、生活、死亡并复活，说出、实行、承受并维持与我们救恩相宜的事；如此，以典范的方式，同时向地上的人显示回归天堂之路，并向我们之上的天使显示保留其天堂地位之路。因为除非在理性灵魂的本性中，并在时间存在的条件下，有某种新事物产生——也就是说，某种以前不存在的东西开始存在——否则从完全败坏和愚昧的生命过渡到智慧和真正善好的生命绝不可能。所以，正如在默观生活中，真理享有永恒之物，而在信德中，信赖的是被造之物，人通过那处理时间之事的信德被净化，以便能够接受永恒之事的真理。因为他们最卓越的智者之一的哲学家\Person{柏拉图}，在题为\WorkTitle{蒂迈欧篇}的论文中也这样说：\Quote{永恒之于被造之物，如同真理之于信德。}前者属于至高之物——即永恒和真理；后者属于低处之物——即被造之物和信德。因此，为了使我们从最低之物中被召唤出来，并被再次引领到至高之物，并且为了使被造之物能达至永恒，我们必须经由信德到达真理。又因为所有对立面都通过某种中间因素被还原为统一，也因为时间的不义使我们与永恒的正义分离，就需要某种时间性的中保正义；这中保因素在低处之物方面是时间性的，但在高处之物方面是正义的，如此，通过适应前者而不切断与后者的联系，将低处之物带回至高之物。因此，基督被称为天主与人之间的中保，祂站在不朽的天主与有死的人之间，因为祂既是天主又是人，祂使人同天主和好；祂继续是祂先前所是的，但祂也成为了祂先前所不是的。这同一位就是我们所说的在被造之物中的信德和在永恒事物中的真理（的中心）。\par

54. 这伟大而不可言喻的奥迹，这王国与司祭职，藉预言启示给古代的人，如今又藉福音宣讲给他们的后代。因为必需要有一个时期，那曾藉一个民族长期应许的，要在万民中实现。因此，那在祂自己降生之前派遣先知的那一位，也在祂升天之后派遣了宗徒。而且，因着祂所取的人性，祂对所有门徒就如同头对身体肢体一样。因此，当这些门徒写下祂向他们宣告和讲述的事时，绝不能说祂自己什么也没有写；因为事实是，祂的肢体只完成了他们通过头的反复陈述而熟悉的事。因为凡祂愿意让我们关于祂自己的言行而阅读的一切，祂都命令这些门徒去写，祂使用他们如同使用自己的手一样。凡理解这合一的一致性和肢体在这头之下不同职责的和谐服务的人，将以同样的精神接受他从福音中经由门徒所编的叙述而得到的记述，就如同他亲眼看见主本身的手——祂在成为祂自己身体的肢体中所拥有的手——在写字一样。因此，现在让我们转而考察那些批评家认为福音作者给出了矛盾说法的段落（只有那些未能正确理解的人才这样认为）的真正性质；以便当这些问题得到解决时，也能表明那身体中的肢体在身体本身的合一中保持了适当的和谐，不仅在观点上同一，而且在构建与同一性一致的记录上也如此。\par

\SourceRecord{Translated by S.D.F. Salmond.；Nicene and Post-Nicene Fathers, First Series；Vol. 6.；Edited by Philip Schaff.；Buffalo, NY: Christian Literature Publishing Co.,；1888.；<http://www.newadvent.org/fathers/1602135.htm>.}

% structure: p0001 source=h1 target=section reason=page-title

\section{福音的和谐 第二卷}

\EditorialNote{在此书中，奥古斯丁对\BibleBook{玛窦}{福音}进行了系统性的考察，直至最后晚餐的叙述，并将其与\BibleBook{马尔谷}{福音}、\BibleBook{路加}{福音}和\BibleBook{若望}{福音}其他三福音进行比较，旨在证明四位圣史在整个这些段落中完全和谐一致。}\par

% structure: p0003 source=h2 target=subsection reason=relative-heading-rank; base=h2

\subsection{序言。}

1. 我们用了较长篇幅且极为重要的论述，已在一卷书内完成，批驳了那些认为撰写这些福音史的门徒只应受轻蔑对待之人的愚妄，其理由是：我们未能提供任何声称直接出自基督之手的著作——他们虽拒绝敬拜基督为天主，却不否认他远超一切智者的智慧，并认为他配得尊敬；此外，我们也驳斥了那些试图将基督的言论歪曲成符合他们悖逆喜好的样式，而非通过阅读和领受这些言论使人归正或弃邪的人。现在，让我们审视四位圣史关于基督的记载，观察它们如何自洽，并与彼此真正和谐。我们这样做，是希望那些好奇心强于理解力的人，在基督信仰中不会因此类事情感到丝毫冒犯，即便他们认为对福音书的研究——不是肤浅的浏览，而是超出常规的仔细考察——已向他们揭示了某些不恰当且矛盾之处，并且他们认为这些应作为争辩的论据而非审慎思考的内容。\par

\SourceRecord{Translated by S.D.F. Salmond.；Nicene and Post-Nicene Fathers, First Series；Vol. 6.；Edited by Philip Schaff.；Buffalo, NY: Christian Literature Publishing Co.,；1888.；<http://www.newadvent.org/fathers/1602200.htm>.}

% structure: p0001 source=h1 target=section reason=page-title

\section{\WorkTitle{福音的和谐 第二卷}}

% structure: p0002 source=h2 target=subsection reason=relative-heading-rank; base=h2

\subsection{第一章 说明为何基督祖先的谱系追溯至若瑟，而基督并非由若瑟所生，而是由童贞玛利亚所生。}

2. 福音作者玛窦以这样的话开始他的叙述：\BibleQuote{耶稣基督，达味之子，亚巴郎之子的族谱}。他以此开头充分表明，他的任务是按肉身记述基督的谱系。因为按肉身而言，基督是人子——他也常以此自称，从而提醒我们注意，他出于怜悯为我们降卑成了什么。因为那属天且永恒的诞生，藉此他是天主的独生子，在一切受造物之前，因为万物由他而造，是如此不可言说，以致先知的话必须被理解为指着它说的：\BibleQuote{谁能述说他的世代？} 因此，玛窦追溯基督按人身的谱系，从亚巴郎往下提及他的祖先，一直延伸到玛利亚的丈夫若瑟，耶稣由她所生。因为基于玛利亚不是作为若瑟的妻子，而是作为童贞女给基督带来生命这一事实，若瑟被视为与玛利亚结成的婚姻关系毫无关联，这是不被允许的。因为通过这个例子，向已婚的信友们显明了一条杰出的原则：即使他们同意保持节制，这种关系仍可存留，且仍可被称为婚姻，因为虽然身体上没有性结合，但心中有情感的持守；特别是因为在他们身上，我们看到儿子诞生可以不需要那种仅为生育目的而行的肉体交合。此外，仅仅因为他没有通过自己的生育行为生下基督，这并不构成若瑟不能被称为基督之父的充分理由；因为他的确可以适当地成为他未由自己妻子所生、而是从别处收养之人的父亲。\par

3. 基督也确实以另一种方式被当作若瑟的儿子，仿佛他仅是出于那个男人的种子而生。但这种猜想起源于那些忽略玛利亚童贞的人。因为路加说：\BibleQuote{耶稣自己开始年纪大约三十岁，} \Quote{依人看来} \BibleQuote{是若瑟的儿子。} 然而，这位路加不仅没有犹豫只称玛利亚为他唯一的双亲，反而也毫不迟疑地把双方都称为他的双亲，他说：\BibleQuote{孩子渐渐长大，强壮起来，充满智慧，天主的恩宠在他身上；他的父母每年逾越节上耶路撒冷去。} 但恐怕有人会以为这里所说的\Quote{父母}与其说是母亲本人连同玛利亚的血亲，那么对路加同一作品前面的话又该怎么说呢，即：\BibleQuote{他的父亲和母亲因论他所说的话而惊讶。} 既然他也陈述基督并非因若瑟与母亲结合而生，而是仅由童贞玛利亚所生，他又怎能称他为他的父亲呢？除非我们理解他确实是玛利亚的丈夫，虽没有肉体的结合，却因婚姻的真实结合；因此，他在更紧密的关系上也是基督的父亲（只要基督是从他的妻子所生），这胜于若仅是从别人那里收养的情况。这清楚表明，\Quote{依人看来} 这句话是插入的，针对那些以为若瑟像其他人一样生下了基督的人。\par

\SourceRecord{Translated by S.D.F. Salmond.；Nicene and Post-Nicene Fathers, First Series；Vol. 6.；Edited by Philip Schaff.；Buffalo, NY: Christian Literature Publishing Co.,；1888.；<http://www.newadvent.org/fathers/1602201.htm>.}

% structure: p0001 source=h1 target=section reason=page-title

\section{\WorkTitle{福音的和谐} 第二卷}

% structure: p0002 source=h2 target=subsection reason=relative-heading-rank; base=h2

\subsection{第二章 解释基督为何是达味之子，尽管他并非由达味之子若瑟按寻常生育方式而生。}

4. 因此，即使有人能够证明，按血统法则，\Person{玛利亚}不能声称出自\Person{达味}，我们仍有足够理由认定\Person{基督}是\Person{达味}之子，基于那同样的计算方式，即\Person{若瑟}也被称为祂的父亲。但既然宗徒\Person{保禄}明确告诉我们：\BibleQuote{基督按肉体说是达味的后裔}，我们更应毫不犹豫地接受\Person{玛利亚}本人也按血统法则，在某些方面出自\Person{达味}世系的观点。再者，既然这位妇人与司祭家族的关系也并非全无记载，因为\Person{路加}记载了\Person{依撒伯尔}——他记录她出自\Person{亚郎}的女儿们——是她的表亲，我们最应坚定地持守这一事实：基督的血肉出自两条世系，即君王世系和司祭世系；在这两者身上，也曾设立了一种奥妙的傅油，在希伯来人中具有象征意义。换句话说，有一种\OriginalTerm{傅油}{chrism}；这个术语揭示了基督之名的含义，并将其呈现为很久以前就以如此明晰的暗示所指示的事物。\par

\SourceRecord{Translated by S.D.F. Salmond.；Nicene and Post-Nicene Fathers, First Series；Vol. 6.；Edited by Philip Schaff.；Buffalo, NY: Christian Literature Publishing Co.,；1888.；<http://www.newadvent.org/fathers/1602202.htm>.}

% structure: p0001 source=h1 target=section reason=page-title

\section{福音合参，卷二}

% structure: p0002 source=h2 target=subsection reason=relative-heading-rank; base=h2

\subsection{第三章 说明为何玛窦为基督罗列一组祖先而路加列举另一组的原因}

5. 此外，对于那些因玛窦列出一组祖先（从达味往下直到若瑟）而路加列出另一组（从若瑟往上直到达味）而感到困惑的批评者，他们不难看出若瑟可能有两位父亲——一位是生他的父亲，另一位是收养他的父亲。因为在那民族中自古也有收养孩子的习俗，以便将那些并非自己所生的人立为儿子。且不提法郎的女儿收养梅瑟（她是个外邦人），雅各伯本人就以这些明确的话收养了自己的孙子、若瑟的儿子们：\BibleQuote{现在，你那两个在你来埃及以前所生的儿子，归我所有；厄弗辣因和默纳协应归我，如勒乌本和西默盎一样；你以后所生的儿子，就归你自己。}由此也产生了以色列的十二支派，尽管肋未支派没有被计算在内（他们负责在圣殿服务），因为加上那一个支派总共是十三个，而雅各伯的儿子们本是十二个。同样，我们也能理解路加在他福音的族谱中为若瑟指出一位父亲，不是生他的父亲，而是收养他的父亲，并将祖先的名单往上追溯到达味。因为既然两位福音作者（玛窦和路加）都真实叙述，因此必然一位遵循生若瑟的父亲那条线，另一位遵循收养他的父亲那条线。那么，我们难道不应该认为，那位拒绝说若瑟是被那个人所生（却仍称他为若瑟的父亲）的福音作者，更可能保留了收养之父的家谱吗？因为一个人被称为收养他之人的儿子，比被称为不是从他血肉所生之人的儿子更为合适。因此，当玛窦使用了\BibleQuote{亚巴郎生依撒格}、\BibleQuote{依撒格生雅各伯}等语句，始终使用\Quote{生}这个词，直到最后说\BibleQuote{雅各伯生若瑟}，他清楚地让我们知道，他追溯了祖先的次序直到那位不是收养若瑟、而是生养若瑟的父亲。\par

6. 但即使路加说若瑟是赫黎所生，这种说法也不应使我们困扰到以为，两位福音作者中一位提到生父，另一位提到养父，而不是别的。因为说一个人收养了一个儿子，并在爱中生了他（即使不是在肉身上生他），这并非荒谬。那些天主赐予权能成为祂儿子的人，祂并非如对待祂独生子那样，由自己的本性和本体所生，而确实是在祂的爱中收养了我们。宗徒反复使用这一说法，正是为了将我们与那在一切受造物之前的独生子区别开来，万物是藉着祂造的，祂独自从圣父的本体而生；祂按照神性的平等，完全与圣父相同，祂被宣称前来是为了将属于我们本性的种族所固有的血肉赋予自己，以便通过祂参与我们的有死性，藉着祂对我们的爱，使我们借着义子名分分享祂自己的神性。因为宗徒这样说：\BibleQuote{但时期一满，天主就派遣了自己的儿子，生于女人，生于法律之下，为把在法律之下的人赎出来，使我们获得义子的名分。}然而，我们也被称为由天主而生——也就是说，我们这些已经身为人的，获得了权能成为天主的子女，这更是靠着恩宠，而非本性。因为如果我们本性地为儿子，我们就不可能成为别的。但当若望说：\BibleQuote{祂赐给他们权能，成为天主的子女，即那些信祂名的人，}他立刻接着说了这些话：\BibleQuote{他们不是由血，也不是由肉欲，也不是由男欲，而是由天主生的。}因此，对于同一群人，他先说他们接受权能而成为天主的子女，这就是保禄所说的义子名分；其次，他们是由天主生的。并且为了更清楚地显示这是藉着何种恩宠成就的，他继续说：\BibleQuote{圣言成了血肉，居住在我们中间，}——仿佛他要说：那些成为天主子女的人虽然是血肉，有什么可奇怪的呢？因为为了他们的缘故，独生子虽然原是圣言，却成了血肉。然而这两种情况之间有一个巨大的区别：当我们成为天主的子女时，我们变得更好了；但当圣子成为人之子时，祂并非变得较差，而是确实承担了比祂低下的东西。雅各伯也这样说道：\BibleQuote{祂按自己的旨意，用真理之言生了我们，叫我们在祂所造的万物中，成为初熟的果子。}为了不让我们因为使用了\Quote{生}这个词就以为我们成为了祂自己所是的，他在这里非常清楚地指出，因着这种义子名分赐予我们的，是在受造物中的一种首生者地位。\par

7. 因此，即使路加说若瑟是由真正收养他的那个人所生，这也不会偏离真理。实际上，即使他以那种方式生了他，不是生为一个男人，而是生为一个儿子；正如天主生了我们为祂的儿子，而我们先前曾被祂造成人。但祂只生了一位，祂不只是圣子——圣父不是圣子——而且也是天主——圣父同样也是天主。同时，显然如果路加采用了那种措辞，那么对于两位作者中哪一位提到收养的父亲，哪一位提到亲生父亲，就完全是一个疑问了；同样，另一方面，即使两人都没有使用\\Quote{生了}这个词，并且前一位福音作者称他为某人的儿子，后一位称他为另一人的儿子，但仍然会不确定他们中谁提到了他的生父，谁提到了他的养父。然而，实际情况是——一位福音作者说\\Quote{\\Person{雅各伯}生了\\Person{若瑟}}，另一位说\\Quote{\\Person{若瑟}是\\Person{赫里}的儿子}——他们通过各自所使用的表述之间的区别，优雅地指出了他们所处理的不同对象。但肯定的是，正如我所说，一个足够坚定以致认为有必要寻求某种更合理解释、而不只是相信福音作者说了假话的虔诚之人，很容易想到这一点；我重复，这样的人很容易想到去考察一个人可能（被认为）有两个父亲的理由是什么。确实，连那些诽谤者也本可想到这一点，可惜他们偏爱争论而不愿思考。\par

\SourceRecord{Translated by S.D.F. Salmond.；Nicene and Post-Nicene Fathers, First Series；Vol. 6.；Edited by Philip Schaff.；Buffalo, NY: Christian Literature Publishing Co.,；1888.；<http://www.newadvent.org/fathers/1602203.htm>.}

% structure: p0001 source=h1 target=section reason=page-title

\section{\WorkTitle{福音的和谐，第二卷}}

% structure: p0002 source=h2 target=subsection reason=relative-heading-rank; base=h2

\subsection{第四章 论玛窦为何列出四十代（不包括基督本人），尽管他将其分为三组，每组十四代。}

8. 接下来要引入的问题，需要读者以最大的注意和谨慎来正确理解和思考。因为有人敏锐地观察到，玛窦，他立意推崇基督的君王特性，在世代谱系中列出了四十个人，不包括基督本人。这个数字标志着，在此世和此地上，我们应当受基督统治的时期，这统治符合那种痛苦的纪律，正如经上所载：\BibleQuote{天主鞭挞祂所收纳的每一个儿子}；又有一宗徒说：\BibleQuote{我们必须历经许多困难，才能进入天主的国。}这种纪律也由那铁杖所象征，我们在圣咏中读到这样的话：\BibleQuote{祢要用铁杖统治他们}；这话出现于\BibleQuote{但我是由祂在祂的圣山上被立为君王！}之后。因为善人也受铁杖的统治，正如论到他们所说：\BibleQuote{时候到了，审判要从天主的家开始；若是先从我们开始，那些不信从天主福音的人将要有什么结局呢？如果义人仅仅得救，那不虔诚和犯罪的人将出现在哪里呢？}同一批人也适用后面的话：\BibleQuote{祢要摔碎他们，如同陶匠的器皿。}因为善人确实受此纪律的统治，而恶人被它击碎。这里提到这两类不同的人如同相同，因为应用于恶人与善人的符号相同。\par

9. 这个数字作为那劳苦时期的标记，在此期间我们在基督君王纪律下与魔鬼争战，也由律法和先知在梅瑟和厄里亚身上所举行的四十天斋戒——即灵魂的谦卑——所指明，他们各自斋戒了四十天。福音叙述关于主本人四十天的斋戒，在此期间祂也被魔鬼试探，还意味着什么？岂非那属于我们整个现世期间的试探状态，祂降生成人，从我们的有死性中取了肉身来承担？复活后，祂也愿意在地球上与门徒们一同停留不超过四十天，在这段时间里继续以人类交往的方式与他们共同生活，并和他们一起分享凡人所需的食物，尽管祂自己不再死亡；这一切都是为了通过这四十天向他们表明，虽然从他们眼中看祂是隐藏的，祂仍要实现祂所应许的：\BibleQuote{看，我与你们天天在一起，直到世界的终结。}至于这个具体数字为何表示这暂时的尘世生命，目前最直接想到的解释（尽管可能有其他更微妙的方法来解释）是：一年四季也按四个交替的序列旋转，世界本身也被四个方向界定，圣经有时用风的名字称呼它们——东、西、\OriginalTerm{北}{Aquilo}、\OriginalTerm{南}{Meridian}。但四十等于四乘以十。再者，十本身由从一到四的连续数字相加而成。\par

10. 这样，玛窦肩负着呈现基督的纪录的责任，祂是来到这世界的君王，来到这尘世及人类有死的生命中，为要统治我们这些与诱惑斗争的人。于是他始于亚巴郎，数算了四十个人。因为基督是带着血肉从那希伯来民族而来的，为使那民族与其他民族分别，天主曾将亚巴郎从他的本乡和亲族中分别出来。而且，那应许暗示了祂将出于哪一民族，这一点特别使关于祂的预言和宣告更加清晰。因此，福音作者确实在三个部分中各标出十四代，声称从亚巴郎到达味共十四代，从达味到流徙巴比伦又有十四代，再从那个时期到基督诞生又有十四代。但他并没有逐一数算并加总说总共有四十二代。因为在那些先祖中，有一位被数算了两次，即耶苛尼雅，在流徙巴比伦时，似乎偏向外邦民族而发生了转折。此外，当数算从直线的顺序弯曲，形成一个可说是转角，以便采取不同方向时，我们在这个转角遇到的就是被提到两次的那个名字——即在前一系列末尾，以及所指明的转折起点。这也是基督的预象，祂在某种意义上要从割礼转到未割礼，或可说从耶路撒冷到巴比伦，并且成为所有信祂的人（无论是这边还是那边的）的角石。这样，天主当时以预象的方式为后来将要成真的事作了预备。因为耶苛尼雅本人，他的名字预表了我所指的那个转角，按解释是\Quote{天主的准备}。因此，这里实际并没有四十二个不同的先祖，那会是三乘十四的适当总数；但由于其中一个名字被数了两次，这里总共有四十个先祖，同时考虑到基督自己也被数算在其中，祂如同统治这[有意义的]数字四十的君王，掌管我们这暂时的尘世生命。\par

11. 既然玛窦的意图是展示基督下降以与我们分享这有死的状态，他就在福音的开端，以递减的顺序列举了从亚巴郎到若瑟、直到基督自己诞生的那些世代。另一方面，路加并非在福音的开端，而是在基督受洗时详列那些世代，并且不是以递减的顺序，而是以递增的顺序，将祂优先视为在赦罪中身为司铎的角色，正如那从天上来的声音宣告祂，以及若翰本人如此作证说：\BibleQuote{请看，天主的羔羊，除免世罪者！}此外，在追溯家谱的过程中，他越过了亚巴郎，一直上溯到天主，我们经过净化与赎罪后，与祂和好。祂也在自己身上承载了我们被收为义子的起源；因为我们藉着信从天主子，通过义子名分成为天主的子女。而且，为了我们，天主子乐意藉着肉身的诞生成为人子。福音作者清楚表明，他称若瑟为赫里的儿子，并非由于若瑟是赫里所生，而仅是由于他被赫里收养。因为他也称亚当为天主子，严格说来，亚当是天主所造，但也可以说，他因恩宠在乐园中被立为儿子，后来因过犯失去了这恩宠。\par

12. 这样，在玛窦的家谱中，象征的是主基督将我们的罪担负在自己身上；而在路加的家谱中，表达的则是主基督消除我们的罪。根据这些观念，一个列举名字是递减顺序，另一个是递增顺序。因为当宗徒说：\BibleQuote{天主派遣了自己的儿子，带着罪恶肉身的形状}，他指的是基督将我们的罪担负在自己身上。但他接着加上：\BibleQuote{为赎罪祭，在肉身上判定罪过}，则表达了罪的赎免。因此，玛窦通过撒罗满从达味向下追溯，撒罗满的母亲使他犯了罪；而路加通过纳堂向上追溯到达味，天主藉着纳堂先知除去了他的罪。路加采用的数字也确实最清楚地表明了罪的除去。因为在基督身上，祂自己本无罪，确实没有与祂在肉身中所背负的人之不义相连的不义，因此玛窦所用的数字，在基督除外时是四十。相反，藉着祂清除我们一切罪并净化我们，祂使我们与祂自己及圣父的正义建立正确的关系（这样宗徒的话便应验了：\BibleQuote{但那与主结合的，便是与主成为一神}），在路加的数目中，我们包括了基督自己（数算从祂开始）和天主（数算以祂结束），总数成为七十七，这表示所有罪的彻底宽免与消除。关于这数字的奥秘，主自己也曾清楚表明这种完全的罪之消除，当祂说：应该原谅犯罪的人，\BibleQuote{不是七次，而是七十个七次}。\par

仔细查考就会清楚，后一数字被用来指称所有罪的涤除，不是没有理由的。因为，可以说，从法律的十诫来看，十是\OriginalTerm{正义}{justice [righteousness]}之数。此外，罪是违反法律。而十的违反适当地以十一来表示；因此我们也发现指示说，应在会幕中制作十一幅山羊毛幔子；因为谁能怀疑山羊毛织物与罪的表达有关呢？同样，既然所有时间在其流转中都是以七天为周期运行的，我们发现，当十一乘以七时，我们便十分恰当地得到七十七这一数字，作为全部罪的标记。因此，在此列举中，我们遇到了完全赦罪的象征，因为赎罪是由我们司铎的肉身为我们完成的；这数字的计算在这里以祂的名字开始；而修和也是借着圣神为我们与天主达成的；这数字的推算在这里以天主的名字结束；圣神在圣洗时以鸽子的形像显现，所提到的数字就是与那次圣洗相关的。\par

\SourceRecord{Translated by S.D.F. Salmond.；Nicene and Post-Nicene Fathers, First Series；Vol. 6.；Edited by Philip Schaff.；Buffalo, NY: Christian Literature Publishing Co.,；1888.；<http://www.newadvent.org/fathers/1602204.htm>.}

% structure: p0001 source=h1 target=section reason=page-title

\section{福音的和谐，卷二}

% structure: p0002 source=h2 target=subsection reason=relative-heading-rank; base=h2

\subsection{第五章 陈述路加如何被证明与玛窦在有关基督受孕、婴孩或童年之事上和谐一致，这些事被一位省略，而由另一位记载。}

14. 在列举世系之后，玛窦继续写道：\Quote{基督的诞生是这样的：他的母亲玛利亚许配给若瑟，在他们同居以前，她因圣神有孕的事已显示出来。}关于此事如何发生，玛窦在此省略未述，而路加在记述了若翰的受孕之后，加以陈述。他的叙述如下：\BibleQuote{到了第六个月，天使加俾额尔奉天主差遣，往加里肋亚一座名叫纳匝肋的城去，到一位童贞女那里，她已与达味家族中一个名叫若瑟的男子订了婚，童贞女的名字叫玛利亚。天使进去对她说：\InnerQuote{万福！充满恩宠者，主与你同在！你在女人中应受赞颂。}她因这话惊惶不安，便思虑这样的请安有什么意思。天使对她说：\InnerQuote{玛利亚，不要害怕，因为你在天主前获得了宠幸。看，你将怀孕生子，并要给他起名叫耶稣。他将是伟大的，并被称为至高者的儿子；上主天主要把他祖先达味的御座赐给他；他要为王统治雅各伯家，直到永远；他的王权没有终结。}玛利亚便对天使说：\InnerQuote{这事怎能成就？因为我不认识男人。}天使回答她说：\InnerQuote{圣神要临于你，至高者的能力要庇荫你，因此，那要诞生的圣者，将称为天主的儿子。}}然后接着是与此问题无关的事。如今玛窦［简略地］记录了这一切，当他告诉我们玛利亚\BibleQuote{她因圣神有孕的事已显示出来}时。两位福音作者之间并无矛盾，因为路加详述了玛窦省略未提的事；两人都证明玛利亚是因圣神怀孕的。同样，他们之间亦无缺失和谐之处，因为玛窦又接续叙述了一些路加省略的事。玛窦继续给我们如下陈述：\BibleQuote{她的丈夫若瑟，因是义人，不愿公开羞辱她，有意暗中休退她。当他在思虑这事时，看，上主的天使在梦中显现给他说：\InnerQuote{达味之子若瑟，不要怕娶你的妻子玛利亚，因为那在她内受生的，是出于圣神。她要生一个儿子，你要给他起名叫耶稣，因为他要把自己的百姓从他们的罪恶中拯救出来。}这一切事的发生，是为应验主藉先知所说的话：\InnerQuote{看，一位童贞女将怀孕生子，人将称他的名字为厄玛奴耳，意思是：天主与我们同在。}若瑟从睡梦中醒来，就照上主的天使所吩咐的，娶了他的妻子；但没有与她同房，直到她生了头胎儿子，给他起名叫耶稣。}当耶稣在黑落德王时代生于犹太的伯利恒时，等等。\par

15. 关于白冷城，玛窦和路加是一致的。但路加说明了若瑟和玛利亚如何及为何来到那里；而玛窦没有这种说明。另一方面，路加对贤士从东方来的行程保持沉默，而玛窦却提供了一段记述。这段记述紧接他先前所述，构造如下：\BibleQuote{看，有贤士从东方来到耶路撒冷，说：\InnerQuote{那生下来作犹太人之王的在哪里？我们在东方见到了他的星，特来朝拜他。}黑落德王听了这些话，就惊慌起来。}这样，记述一直继续到关于这些贤士所记载的一段话：\BibleQuote{他们在梦中得到指示，不要回到黑落德那里，就由另一条路返回了家乡。}整段内容被路加省略了，正如玛窦也省略了一些路加提及的其他情况：例如，主被放在马槽里；有天使向牧羊人宣布他的诞生；有大队天军同天使一起赞美天主；牧羊人前来看见天使向他们宣布的事果然应验；在他受割损的日子，他接受了他的名字；以及同一路加所报道的在玛利亚取洁的日子满了以后发生的事件——即他们带他到耶路撒冷，西默盎或亚纳在圣殿里论及他的话，当时他们充满圣神，认出了他。关于这一切，玛窦只字未提。\par

16. 因而，有值得探讨的问题涉及这些事件发生的确切时间：哪些事件被\Person{玛窦}省略而由\Person{路加}记载，反之，哪些被路加省略而由玛窦记载。因为玛窦在记述了从东方来的贤士返回本国之后，接着告诉我们天使如何警告\Person{若瑟}带着婴孩逃往\Place{埃及}，以免被\Person{黑落德}杀害；然后黑落德如何找不到他，却杀尽了两岁及以下的婴孩；此后，黑落德死后，若瑟如何从埃及回来，听说\Person{阿赫劳斯}继其父黑落德在\Place{犹太}作王，便带着孩子往\Place{加里肋亚}境内，住在\Place{纳匝肋}城。所有这些事实，路加又都略过不提。然而，仅凭后者记载了前者省略的事，或前者提及了后者忽略的事，并不能证明两位作者之间有任何不和谐。真正的问题是：玛窦连在其叙述中的这些事究竟可能发生在什么时期？即：全家逃往埃及，黑落德死后从埃及返回，以及那时他们住在纳匝肋镇——正是路加告诉我们，他们在\Place{圣殿}中按主的法律为婴孩办完了一切事后，返回的那个地方。因此，在这里我们需要注意一个事实，这一事实也适用于其他类似情况，并将使我们的心灵在后续事例中免于类似的激动或不安。我指的是这一点：每位福音作者都按照某种计划构建他自己的独特叙述，使其看起来像是事件的完整有序记录。因为，对于他不打算记述的那些事件，他只是保持沉默，然后将想要叙述的事与他刚刚讲述的事连接起来，使叙述显得连贯。同时，当其中一位提到另一位没有记载的事实时，若仔细考虑叙述顺序，就会发现那位省略的作者在其记述中跳过的那个点，这样，他将有意引入的事实以这样一种方式附着在前文上，使得看起来像是连续系列，其中一件事紧接着另一件事，中间没有插入任何其他东西。因此，按照这一原则，我们理解到：当玛窦告诉我们贤士如何在梦中受警告不要回黑落德那里，并从另一条路返回本国时，他只是省略了路加所记载的一切关于主在圣殿中所发生的事，以及\Person{西默盎}和\Person{亚纳}所说的话；而另一方面，路加在同一处省略了玛窦所记的逃往埃及的全部经过，并将返回纳匝肋城描述为似乎是紧接着的事。\par

17. 如果谁愿意把这两位福音作者关于基督的诞生、婴孩期或童年所说或未说的所有内容整理成一个完整的叙事，可以按下列顺序安排不同的陈述：— 基督的诞生是这样的。\newline{}\BibleQuote{在犹太王黑落德的时候，阿彼雅班中有一位司铎，名叫匝加利亚，他的妻子是亚郎的后代，名叫依撒伯尔。二人在天主面前是义人，遵行主的一切诫命和典章，无可指摘。他们没有孩子，因为依撒伯尔不能生育，两人又都年纪老迈。及至他在自己班次供职，在天主面前尽司铎的职分时，按照司铎职分的惯例，他抽中了签，进入主的圣殿焚香。众百姓在外面祈祷，到了焚香的时候。有一位上主的天使显现给他，站在香坛右边。匝加利亚看见，就惊慌害怕。天使对他说：\InnerQuote{匝加利亚，不要害怕！因为你的祈祷已蒙应允，你的妻子依撒伯尔要给你生一个儿子，你要给他起名叫若翰。你必要喜乐欢跃，许多人也要因他的诞生而喜乐。因为他在主前将是伟大的，淡酒浓酒都不喝，而且他还在母胎中就要充满圣神。他要使许多以色列子民转向上主，他们的天主。他将以厄里亚的精神和能力，在祂前面先行，使父亲的心转向儿子，使悖逆者转向义人的明智，为上主准备一个完善的百姓。}匝加利亚对天使说：\InnerQuote{我凭什么能知道这事呢？因为我已经老了，我的妻子也年纪老迈了。}天使回答说：\InnerQuote{我是加俾额尔，站在天主面前的，奉命来向你说话，将这些喜讯报告给你。看，你必成为哑巴，不能说话，直到这些事成就的那一天，因为你没有相信我的话；但我的话到了时候必要应验。}百姓等候匝加利亚，都惊讶他在圣殿里滞留。他出来以后，不能和他们说话，他们就知道他在圣殿里见了异象；他向他们打手势，仍成了哑巴。他供职的日子一满，就回了家。这些日子以后，他的妻子依撒伯尔怀了孕，隐藏了五个月，说：\InnerQuote{主在眷顾我的日子，这样待了我，除去了我在人间的羞辱。}到了第六个月，天使加俾额尔奉天主差遣，往加里肋亚一座名叫纳匝肋的城去，到一位童贞女那里，她已与达味家族中一个名叫若瑟的男子订了婚，童贞女的名字叫玛利亚。天使进去向她說：\InnerQuote{万福！充满恩宠者，主与你同在！你在女人中受赞颂。}她因这话惊慌不安，思忖这样的请安有什么意思。天使对她说：\InnerQuote{玛利亚，不要害怕！因为你在天主前获得了宠幸。看，你将怀孕生子，并要给祂起名叫耶稣。祂将是伟大的，并被称为至高者的儿子，上主天主要把祂祖先达味的御座赐给祂。祂要为王统治雅各伯家，直到永远；祂的王权没有终结。}玛利亚便对天使说：\InnerQuote{这事怎能成就？因为我不认识男人。}天使回答她说：\InnerQuote{圣神要临于你，至高者的能力要庇荫你，因此，那要诞生的圣者，将称为天主的儿子。且看，你的亲戚依撒伯尔，她虽在老年，却怀了男胎，本月是第六个月，她原是素称不生育的。因为在天主没有不可能的事。}玛利亚说：\InnerQuote{看，上主的婢女，愿照你的话成就于我吧！}天使便离开她去了。玛利亚就在那几天起身，急速往山区去，来到犹大的一座城。进了匝加利亚的家，问候依撒伯尔。依撒伯尔一听到玛利亚问安，胎儿就在她腹中欢跃。依撒伯尔遂充满了圣神，大声呼喊说：\InnerQuote{你在女人中受赞颂，你的胎儿也受赞颂！吾主的母亲驾临我这里，这是我哪里得来的呢？看，你请安的声音一入我耳，胎儿就在我腹中欢喜踊跃。那信了由上主传于她的话必要完成的，是有福的。}玛利亚说：\InnerQuote{我的灵魂颂扬上主，我的心神欢跃于天主，我的救主。因为祂垂顾了祂婢女的卑微，今后万世万代都要称我有福。因为全能者在我身上行了大事，祂的名字是圣的。祂的仁慈世世代代临于敬畏祂的人。祂施展了手臂，驱散那些心高气傲的人。祂从高座上推下权势者，却举扬了卑微贫困的人。祂使饥饿者饱享美物，却使富有者空手而去。祂扶助了祂的仆人以色列，纪念祂的仁慈，正如祂向我们的祖先所说过的，施恩于亚巴郎和他的子孙，直到永远。}玛利亚同依撒伯尔住了大约三个月，就回家去了。}\newline{}接着叙述如下：—\newline{}\BibleQuote{她因圣神怀了孕。她的丈夫若瑟，因是义人，不愿公开羞辱她，有意暗暗地休退她。当他在思虑这事时，看，上主的天使在梦中显现给他说：\InnerQuote{达味之子若瑟，不要怕娶你的妻子玛利亚，因为那在她内受生的，是出于圣神。她要生一个儿子，你要给祂起名叫耶稣，因为祂要把自己的民族，从他们的罪恶中拯救出来。}这一切事的发生，是为应验上主藉先知所说的话：\InnerQuote{看，一位贞女将怀孕生子，人将称祂的名字为厄玛奴耳，意思是：天主与我们同在。}若瑟从睡梦中醒来，就照上主的天使所吩咐的，娶了他的妻子，但没有与她同房。}\par

\Person{依撒伯尔}产期到了，她生了一个儿子。她的邻居和亲戚听说\Person{主}对她大施仁慈，就与她一同欢乐。到了第八天，他们来给孩子行割损礼，并照他父亲的名字叫他\Person{匝加利亚}。他母亲回答说：\Quote{不，要叫他\Person{若翰}。}他们对她说：\Quote{你亲族中没有叫这个名字的。}他们便向他父亲打手势，看他愿意叫他什么名字。他要了一块写字板，写上：\Quote{他的名字是\Person{若翰}。}众人都惊讶。匝加利亚的口立时开了，舌头也解了，就开口赞美\Person{天主}。所有住在他们四周的人都害怕；这一切话就传遍了\Place{犹太}全山地。凡听见的人都把这些事存在心里，说：\Quote{这孩子将来会怎样呢？因为主的手与他同在。}他父亲\Person{匝加利亚}充满圣神，预言说：\Quote{上主以色列的天主应受赞美，因为祂眷顾救赎了祂的子民，并在祂的仆人\Person{达味}家中，为我们兴起了救恩的角；正如祂借历代先知之口所说的，从创世以来，（赐予）拯救我们脱离仇敌和一切恨我们的人之手；祂向我们的祖先施行仁慈，并记起祂的圣约，就是祂对我们的祖宗\Person{亚巴郎}所起的誓：赐给我们，使我们从敌人手中被救出后，无恐无惧地在祂面前，以圣善和正义事奉祂，终身不渝。你，孩子，要称为至高者的先知，因为你要走在主前面，为祂预备道路，将救恩的知识赐给祂的子民，使他们罪过得赦，这是出于我们天主的仁慈；由此，旭日从高天眷顾了我们，照亮那些坐在黑暗中、死影下的人，引导我们的脚步走向和平之路。}这孩子渐渐长大，心灵强健，住在旷野中，直到他在以色列人前公开出现的日子。那时，\Person{凯撒奥古斯都}出了一道旨意，叫天下人登记户口。这是\Person{季黎诺}作\Place{叙利亚}总督时，初次登记户口。众人各归本城去登记。\Person{若瑟}也从\Place{加里肋亚}的\Place{纳匝肋}城上\Place{犹太}去，到了\Place{达味城}，名叫\Place{白冷}，因为他是\Person{达味}家族的人，与他聘妻\Person{玛利亚}一同登记，那时\Person{玛利亚}已有身孕。他们在那里的时候，\Person{玛利亚}产期到了，生了她的首生子，用襁褓裹起来，放在马槽里，因为客店里没有地方。在那地方有牧人露宿，夜间看守羊群。忽然，\Person{主}的天使站在他们身边，\Person{主}的荣光四面照射他们，他们非常害怕。天使对他们说：\Quote{不要害怕！看，我给你们报告一个全民族的大喜讯：今天在\Place{达味城}中，为你们诞生了一位\OriginalTerm{救主}{Saviour}，就是\OriginalTerm{主基督}{Christ the Lord}。这是给你们的记号：你们要看见一个婴儿，裹着襁褓，躺在马槽里。}忽然有一大队天军，同那天使一起赞美\Person{天主}说：\Quote{\OriginalTerm{天主在天受光荣}{Glory to God in the highest}，\OriginalTerm{主爱的人在世享平安}{on earth peace to men of goodwill}。}众天使离开他们升天去了，牧人们彼此说：\Quote{我们往\Place{白冷}去，看看主所告诉我们的这件事。}他们急忙去了，找到了\Person{玛利亚}和\Person{若瑟}，并那躺在马槽中的婴儿。他们看见后，就明白了天使对他们所说有关这孩子的消息。凡听见的人都惊讶牧人告诉他们的事。\Person{玛利亚}却把这一切事默存在心中，反复思考。牧人回去了，为所听见和看见的一切事荣耀赞美\Person{天主}，正如天使向他们所说的。满了八天，给孩子行割损礼时，祂的名字叫\Person{耶稣}，这是在祂\OriginalTerm{成孕}{conceived}前天使所起的。接着叙述如下：看，有\OriginalTerm{贤士}{wise men}从\Place{东方}来到\Place{耶路撒冷}，说：\Quote{那生来作犹太人的王在哪里？我们在\Place{东方}见到了祂的星，特来朝拜祂。}\Person{黑落德王}听见了，就惊慌不安，全\Place{耶路撒冷}也同他一起惊慌。他召集了众司祭长和民间的经师，问他们基督应当生在哪里。他们回答：\Quote{在\Place{犹太}的\Place{白冷}，因为先知这样记载：\InnerQuote{你，\Place{犹太地}的\Place{白冷}，在\Place{犹太}的郡邑中，决不是最小的，因为将由你出来一位领袖，祂将牧养我的百姓以色列。}}\Person{黑落德}暗中召见了\OriginalTerm{贤士}{wise men}，仔细询问那星出现的时间。然后打发他们往\Place{白冷}去，说：\Quote{你们去仔细寻访那婴儿，找到了，就报告我，好让我也去朝拜祂。}他们听了国王的话，就走了。看，他们在\Place{东方}所见的那星，走在他们前面，直到婴儿所在的地方，就停在上面。他们一见到那星，极其高兴欢喜。进了屋子，看见婴儿和祂的母亲\Person{玛利亚}，就俯伏朝拜了祂；打开宝盒，献上礼物：黄金、\OriginalTerm{乳香}{frankincense}和\OriginalTerm{没药}{myrrh}。他们在梦中得到指示，不要回去见\Person{黑落德}，就从另一条路返回了本国。然后，记述了\OriginalTerm{贤士}{wise men}回去之后，叙述继续如下：按\WorkTitle{梅瑟法律}，她（祂的母亲）\OriginalTerm{取洁}{purification}的日子满了，他们就带祂上\Place{耶路撒冷}，将祂献给\Person{主}（正如\Person{主}的法律上所记载的：\Quote{凡开胎首生的男性，应奉献给\Person{主}}），并依照\Person{主}的法律所规定的，献上祭物：一对斑鸠或两只雏鸽。看，在\Place{耶路撒冷}有一个人，名叫\Person{西默盎}；这人正义虔诚，期待着\OriginalTerm{以色列的安慰}{consolation of Israel}；圣神也在他身上。\par

\Person{西默盎}由圣神得了启示，知道自己未见到主基督以前，不会见死亡。他因圣神的感动进了\Place{圣殿}。当他的父母抱着婴孩\Person{耶稣}进来，要按律法的规矩为他办理时，\Person{西默盎}就把他接在怀里，说：\BibleQuote{主啊！现在可以照你的话，让你的仆人平安去了；因为我亲眼看见了你的救恩，就是你为万民所预备的：为启示\Place{外邦人}的光明，你子民\Place{以色列}的荣耀。}\BibleRef{路 2:29-32} 他的父母就惊奇他所说关于他的这些话。\Person{西默盎}祝福了他们，对孩子的母亲\Person{玛利亚}说：\BibleQuote{看，这孩子被立定，使\Place{以色列}中许多人跌倒和兴起，并成为一个遭反对的记号；而一把剑也要刺透你自己的灵魂，为叫许多人心中的思念显露出来。}\BibleRef{路 2:34-35} 又有一位女先知\Person{亚纳}，是\Person{发努耳}的女儿，属\Place{阿协尔}支派。她年纪老迈，从作处女出嫁后，和丈夫同住了七年，以后便寡居，现年约八十四岁，并不离开\Place{圣殿}，以禁食祈祷昼夜事奉天主。正在那时，她也上前来称谢上主，并向一切盼望\Place{耶路撒冷}得救赎的人讲论这孩子。他们按着上主的法律办完了一切事，看，上主的天使在梦中显现给\Person{若瑟}说：\BibleQuote{起来，带着孩子和他的母亲逃往\Place{埃及}，住在那里，直到我通知你；因为\Person{黑落德}将要寻找这孩子，将他杀害。}\BibleRef{玛 2:13}\Person{若瑟}便起来，连夜带着孩子和他的母亲往\Place{埃及}去，住在那里直到\Person{黑落德}死去；这就应验了上主藉先知所说的话：\BibleQuote{我从\Place{埃及}召回了我的儿子。}\BibleRef{玛 2:15} 那时，\Person{黑落德}见自己受了贤士们的愚弄，就大为震怒，派遣人将\Place{白冷}及其周围境内所有两岁以内的婴孩都杀尽了，依照他由贤士们仔细查问的时间。这就应验了\Person{耶肋米亚}先知所说的话：\BibleQuote{在\Place{辣玛}听到了声音，哀哭和极大的号咷，\Person{辣黑耳}为她的儿女哭泣，不肯受安慰，因为他们不在了。}\BibleRef{玛 2:18} 但\Person{黑落德}死后，看，上主的天使在\Place{埃及}在梦中显现给\Person{若瑟}说：\BibleQuote{起来，带着孩子和他的母亲，往\Place{以色列}地去，因为那些想要孩子性命的已经死了。}\BibleRef{玛 2:20}\Person{若瑟}便起来，带着孩子和他的母亲，进了\Place{以色列}地。只因他一听说\Person{阿尔赫劳}继承他父亲\Person{黑落德}作了\Place{犹太}王，就怕往那里去；并在梦中得到天主的指示，便往\Place{加里肋亚}地区去了，来住在一座名叫\Place{纳匝肋}的城，这就应验了先知们所说的：\BibleQuote{他将被称为纳匝肋人。}\BibleRef{玛 2:23} 孩子渐渐长大，强健起来，充满智慧，天主的恩宠在他身上。他的父母每年在逾越节上\Place{耶路撒冷}去。当他十二岁时，他们照节期的规矩上\Place{耶路撒冷}去。日子满了，他们回去时，孩童\Person{耶稣}仍留在\Place{耶路撒冷}，他的父母并不知道。他们以为他在同行的人中，走了一天的路程，就在亲戚和熟人中寻找他。既找不着，就回\Place{耶路撒冷}去找他。过了三天，他们在\Place{圣殿}里找到了他，坐在教师中间，一面听，一面问。凡听见他的，都希奇他的聪明和他的应对。他们看见他，就非常惊奇。他的母亲对他说：\BibleQuote{孩子，为什么这样对待我们？看，你的父亲和我多么伤心地找你。}\BibleRef{路 2:48}\Person{耶稣}对他们说：\BibleQuote{你们为什么找我？难道你们不知道我必须在我父那里吗？}\BibleRef{路 2:49} 他们不明白他对他们所说的话。他就同他们下去，来到\Place{纳匝肋}，并服从他们；他的母亲把这一切话都存在心里。\Person{耶稣}在智慧和身量上，并在天主和人前的恩爱上，渐渐地增长。\par

\SourceRecord{Translated by S.D.F. Salmond.；Nicene and Post-Nicene Fathers, First Series；Vol. 6.；Edited by Philip Schaff.；Buffalo, NY: Christian Literature Publishing Co.,；1888.；<http://www.newadvent.org/fathers/1602205.htm>.}

% structure: p0001 source=h1 target=section reason=page-title

\section{福音的和谐，第二卷}

% structure: p0002 source=h2 target=subsection reason=relative-heading-rank; base=h2

\subsection{第六章 论四位福音书作者对洗者若翰宣讲位置的安排}

18. 如今在此处，关于若翰宣讲的记述开始了，四位福音书作者都呈现了它。因为在我按叙述顺序最后放置的话——即他引入先知证言的那句话，即\BibleQuote{祂将被称为纳匝肋人}，——之后，\Person{玛窦}立刻给我们作了如下叙述：\BibleQuote{在那些日子，洗者若翰来了，在犹太旷野宣讲，等等。}而\Person{马尔谷}，他对主的诞生、婴孩或少年时期都未作记述，却使他的福音以同一事件开始——也就是说，以若翰的宣讲开始。因为他正是这样开始的：\BibleQuote{耶稣基督、天主子，福音的开始，正如先知依撒意亚所记载的：看，我派遣我的使者在你面前，他要在你前面预备你的道路。有声音在旷野呼喊：预备主的道路，修直他的途径。若翰在旷野施洗，宣讲悔改的洗礼，以得罪之赦，等等。} \Person{路加}则在他所说的\BibleQuote{耶稣在智慧和身量上，并在天主和人前的恩爱上，日益增长}那段话之后，紧接着以这样的措辞谈论若翰的宣讲：\BibleQuote{在提庇留凯撒执政第十五年，般雀比拉多做犹太总督，黑落德做加里肋亚分封侯，他的兄弟斐理伯做依突勒雅和特拉曷尼地区的分封侯，吕撒尼雅做阿彼肋纳的分封侯，亚纳斯和盖法做大司祭时，天主的话在旷野临到匝加利亚的儿子若翰，等等。} 宗徒\Person{若望}，四位福音书作者中最杰出的，在论述了天主的圣言，即圣子，它先于一切受造存在的世代，因为万物都是由他造成的之后，紧接着就引入了他关于若翰宣讲和证言的记述，并这样继续：\BibleQuote{有一个人，从天主那里差来，名叫若翰。} 这就足以立即阐明，四位福音书作者关于\Person{洗者若翰}的叙述彼此并不矛盾。而且，再也没有必要要求或坚持在此事上完全复制我们已在由\Person{玛利亚}所生的\Person{基督}的族谱一事上所做过的事，即证明\Person{玛窦}和\Person{路加}如何彼此和谐，展示我们如何能从两者中建构一个一致的叙述，并且为那些领悟力较弱的人说明，即使其中一位福音书作者提到另一个省略的，或省略另一个提到的，也绝不会因此使得接受另一位的叙述的真实性有任何困难。因为当我们面前有了一个调整一致的方法的例子，无论是按我所呈现的方式，还是按某种可能更令人满意的展示方式，每个人都能理解，在所有其他类似段落中，他在这里所见到的做法都可以再次应用。\par

19. 因此，正如我所说的，现在让我们研究四位福音书作者关于\Person{洗者若翰}的叙述中的和谐一致。\Person{玛窦}这样继续：\BibleQuote{在那些日子，洗者若翰来了，在犹太旷野宣讲}。\Person{马尔谷}没有使用\Quote{在那些日子}这个短语，因为他所给的福音开头，就在这个故事之前，没有叙述任何一连串的事件，所以人们可能理解为，他是用\Quote{在那些日子}这个措辞来指涉这些事件的时间。另一方面，\Person{路加}更加精确地通过政权的纪年来界定若翰宣讲或施洗的那些时间。因为他这样说：\BibleQuote{在提庇留凯撒执政第十五年，般雀比拉多做犹太总督，黑落德做加里肋亚分封侯，他的兄弟斐理伯做依突勒雅和特拉曷尼地区的分封侯，吕撒尼雅做阿彼肋纳的分封侯，亚纳斯和盖法做大司祭时，天主的话在旷野临到匝加利亚的儿子若翰。} 然而，我们不应理解为，当\Person{玛窦}说\Quote{在那些日子}时，他的实际意思仅仅是指这些政权所指的特定时期的若干日子。相反，显然他意图让\Quote{在那些日子}这个措辞所传达的时间记号被理解为指更长的时期。因为他首先向我们叙述了\Person{基督}在黑落德死后从\Place{埃及}返回——这确实发生在祂婴孩或少年时期，并且因此，\Person{路加}关于祂十二岁时在圣殿所遇之事的陈述与之相当一致。然后，紧接着这个婴孩或孩童从埃及被召回的故事之后，\Person{玛窦}按顺序这样继续：\BibleQuote{如今，在那些日子，洗者若翰来了。} 如此，在这个短语下，他当然不仅涵盖了祂的童年时期，而且涵盖了从祂诞生到若翰开始宣讲和施洗的这一段所有时期。而且，在这个时期，\Person{基督}已经达到了成年；因为若翰和祂年龄相同；并且记载说，当祂被若翰施洗时，祂大约三十岁。\par

\SourceRecord{Translated by S.D.F. Salmond.；Nicene and Post-Nicene Fathers, First Series；Vol. 6.；Edited by Philip Schaff.；Buffalo, NY: Christian Literature Publishing Co.,；1888.；<http://www.newadvent.org/fathers/1602206.htm>.}

% structure: p0001 source=h1 target=section reason=page-title

\section{\WorkTitle{福音的和谐，第二卷}}

% structure: p0002 source=h2 target=subsection reason=relative-heading-rank; base=h2

\subsection{第七章　论两位黑落德}

20. 但关于\Person{黑落德}的提及，人们很容易理解，有些人会受到以下情况的影响：\Person{路加}告诉我们，在\Person{若翰}施洗的那些日子，当那时已是成人的主也受洗时，\Person{黑落德}是\Place{加里肋亚}的\OriginalTerm{分封侯}{tetrarch}；而\Person{玛窦}告诉我们，\Person{耶稣}孩童时期在\Person{黑落德}死后从\Place{埃及}回来。现在，这两个记载不能同时为真，除非我们还可以设想有两个不同的\Person{黑落德}。但既然没有人会不知道这是一个完全可能的情况，那么那些急于对福音真理提出虚假指控的人，在他们疯狂的愚行中所表现出的盲目性该有多大；而他们又是多么可怜地轻率，不去思考两个人可能被称作同一个名字？然而，这是一个随处可见的例子。因为后一位\Person{黑落德}被认为是前一位\Person{黑落德}的儿子；就像\Person{阿赫劳斯}一样，\Person{玛窦}说他在父亲死后继承了\Place{犹大}的王位；也像\Person{斐理伯}一样，\Person{路加}将他介绍为\OriginalTerm{分封侯}{tetrarch}\Person{黑落德}的兄弟，他自己是\Place{依突勒雅}的\OriginalTerm{分封侯}{tetrarch}。因为那寻求孩童\Person{基督}性命的\Person{黑落德}是王；而他的儿子，另一位\Person{黑落德}，却不被称为王，而是\OriginalTerm{分封侯}{tetrarch}，这是一个希腊词，词源上意为被派治理王国四分之一的人。\par

\SourceRecord{Translated by S.D.F. Salmond.；Nicene and Post-Nicene Fathers, First Series；Vol. 6.；Edited by Philip Schaff.；Buffalo, NY: Christian Literature Publishing Co.,；1888.；<http://www.newadvent.org/fathers/1602207.htm>.}

% structure: p0001 source=h1 target=section reason=page-title

\section{\WorkTitle{福音合参} 第二卷}

% structure: p0002 source=h2 target=subsection reason=relative-heading-rank; base=h2

\subsection{第八章 关于\Person{玛窦}所述\Person{若瑟}因\Person{阿赫劳}而不敢携带婴孩基督进入\Place{耶路撒冷}，却又毫不害怕进入\Place{加里肋亚}，因为在那里\Person{黑落德}（即那位王子的兄弟）是分封侯的解释}

21. 然而在此又可能发现一个难题：有些人看到\Person{玛窦}告诉我们\Person{若瑟}在回程时不敢带婴孩进入\Place{犹太}，明确理由乃是儿子\Person{阿赫劳}代替其父\Person{黑落德}在那里为王，便可能被引导而问：他如何能够进入\Place{加里肋亚}呢？因为据\Person{路加}证实，那里有那位\Person{黑落德}的另一个儿子，即分封侯\Person{黑落德}。但这种困难只能基于一种臆想，认为当时因婴孩的缘故而有所恐惧的时期，与\Person{路加}现在所处理的时期是同一时期；然而，时代已经改变，这是显而易见的，因为我们在\Place{犹太}再也找不到\Person{阿赫劳}被描写成王，而是代之以\Person{般雀·比辣多}，他也不是犹太人的王，只是他们的总督。在他治下，大\Person{黑落德}的儿子们在\Person{提庇留·凯撒}的统辖下，所拥有的不是王国，而是分封统辖。所有这一切在\Person{若瑟}因害怕当时统治\Place{犹太}的\Person{阿赫劳}而带着婴孩逃往\Place{加里肋亚}（那里也有他的城\Place{纳匝肋}）时，肯定还没有发生。\par

\SourceRecord{Translated by S.D.F. Salmond.；Nicene and Post-Nicene Fathers, First Series；Vol. 6.；Edited by Philip Schaff.；Buffalo, NY: Christian Literature Publishing Co.,；1888.；<http://www.newadvent.org/fathers/1602208.htm>.}

% structure: p0001 source=h1 target=section reason=page-title

\section{\WorkTitle{福音的和谐 第二卷}}

% structure: p0002 source=h2 target=subsection reason=relative-heading-rank; base=h2

\subsection{\Emph{第九章 论及\Person{玛窦}所载，\Person{若瑟}带婴孩\Person{基督}前往\Place{加里肋亚}的原因，是他害怕当时在\Place{耶路撒冷}接替他父亲为王的\Person{阿赫劳斯}；而\Person{路加}告诉我们，前往\Place{加里肋亚}的原因是因为他们的本城\Place{纳匝肋}在那里。}}

22. 或许有人会提出疑问：\Person{玛窦}告诉我们，祂的父母带着男孩\Person{耶稣}去了\Place{加里肋亚}，因为他们由于害怕\Person{阿赫劳斯}而不愿进入\Place{犹太}；然而，\Person{路加}却指出，他们前往\Place{加里肋亚}的原因是因为他们的城市是\Place{加里肋亚}的\Place{纳匝肋}。但我们必须注意到，当天使在\Place{埃及}向\Person{若瑟}梦中说话时，说：\Quote{起来，带着婴孩和祂的母亲，往以色列地去。}若瑟最初理解这句话，使他觉得自己被命令前往\Place{犹太}。因为这是对\Quote{以色列地}这个短语可能有的第一种解释。但是，在查明\Person{黑落德}的儿子\Person{阿赫劳斯}在那里为王后，他拒绝冒这样的危险，因为\Quote{以色列地}这个短语也可以理解为包括\Place{加里肋亚}，因为以色列人也占据着那片领土。同时，这个问题也可以用另一种方式解决。因为\Person{基督}的父母可能认为，他们被召唤与这男孩一同居住，而关于这男孩的信息是通过天使的回应传达给他们的，就在\Place{耶路撒冷}本身，那里有主的圣殿。因此，当他们从\Place{埃及}回来时，他们本会直接去那里并住在那里，若不是因为\Person{阿赫劳斯}的存在而恐惧。当然，他们没有从天上得到任何指示，要他们必须不顾对\Person{阿赫劳斯}的恐惧而住在那里。\par

\SourceRecord{Translated by S.D.F. Salmond.；Nicene and Post-Nicene Fathers, First Series；Vol. 6.；Edited by Philip Schaff.；Buffalo, NY: Christian Literature Publishing Co.,；1888.；<http://www.newadvent.org/fathers/1602209.htm>.}

% structure: p0001 source=h1 target=section reason=page-title

\section{\WorkTitle{福音合参 第二卷}}

% structure: p0002 source=h2 target=subsection reason=relative-heading-rank; base=h2

\subsection{第十章 说明为何\Person{路加}告诉我们\Quote{他的父母每年逾越节同孩子上耶路撒冷}；而\Person{玛窦}暗示他们从埃及回来时因害怕\Person{阿基劳斯}而不敢去那里。}

23. 或者有人向我们提出这个问题：那末，既然\Person{路加}的叙述说基督童年时他的父母每年都上\Place{耶路撒冷}，假如他们因害怕\Person{阿基劳斯}而不敢去那里，这怎么可能呢？嗯，我认为解决这个问题并不十分困难，即使没有一位圣史告诉我们\Person{阿基劳斯}在那里统治了多久。因为可能的情况是，只是为了那一天，并且打算立刻返回，他们在节日那天上去，在聚集的广大群众中不引人注意地前往那座城市，而他们平时却不敢在那里居住。这样，他们既避免了显得如此不虔诚以致忽略了遵守节日的义务，又通过短暂的停留避免了引起注意。而且，尽管所有圣史都忽略了告诉我们\Person{阿基劳斯}统治的年限，我们仍然有一个明显的解释方法，即理解\Person{路加}所指的习惯，他说他们每年去\Place{耶路撒冷}，是在\Person{阿基劳斯}不再是一个可怕的对象的时候进行的。但如果根据任何看起来可信的福音外历史，\Person{阿基劳斯}的统治持续了相当长的时间，那么我在上面指出的考虑仍然足够——即假定孩子的父母虽然害怕在\Place{耶路撒冷}居住，但这种恐惧并不足以使他们忽略遵守他们因敬畏天主而有义务遵守的神圣节日，而且他们很容易以一种不引起公众注意的方式去履行。因为，凭借有利的时机，无论是白天还是小时，人们可以安全地接近那些他们仍然害怕被发现逗留的地方，这并没有什么不可思议的。\par

\SourceRecord{Translated by S.D.F. Salmond.；Nicene and Post-Nicene Fathers, First Series；Vol. 6.；Edited by Philip Schaff.；Buffalo, NY: Christian Literature Publishing Co.,；1888.；<http://www.newadvent.org/fathers/1602210.htm>.}

% structure: p0001 source=h1 target=section reason=page-title

\section{福音的和谐，卷二}

% structure: p0002 source=h2 target=subsection reason=relative-heading-rank; base=h2

\subsection{第十一章　查考一个问题：按照\Person{路加}的记载，当基督之母取洁的日子届满时，他们怎能与他一同上\Place{耶路撒冷}到圣殿去履行惯常的礼仪？如果\Person{玛窦}的记载正确，即\Person{希律}已从贤士那里得知那孩子出生，而他追寻那孩子时，杀死了那么多孩子。}

24. 由此我们也看到另一个问题如何被解决，如果有人确实对此感到困难的话。我指的是这个问题：假定老\Person{希律}已经从贤士那里得到关于犹太人之君王的出生消息而早已焦虑不安，并为此激动，那么当他的母亲取洁的日子届满时，他们怎能安全地与他一同上圣殿去，按照主的法律履行那些事，就是\Person{路加}所具体记载的那些事？因为谁能看不出，这一单独的日子很容易被一个国王忽略，因为他的注意力被众多事务占据？或者，如果看起来不太可能，\Person{希律}正在极度焦虑地等待贤士带回有关那孩子的报告，却长时间没有发现他如何被愚弄，以至于只有在母亲取洁之后，并且在圣殿中为那孩子举行了首生子的礼仪之后，甚至在他们逃往\Place{埃及}之后，他才想到要寻索那孩子的性命并杀害那么多小孩——如果，我说，有人对此感到困难，我就不再列举国王的注意力可能被众多重要事务所占据，并在许多天内要么完全转移了这些念头，要么无法追踪它们。因为无法列举所有可能使这件事完全可能的情况。然而，没有人对人事如此无知，以至于否认或质疑很轻易就有许多这样重要的事情（占据国王的注意力）。因为谁不会想到，许多其他更可怕的消息，无论真假，可能已经呈报给国王，以致那个曾担心某个王族孩童（多年后可能成为他自己或他儿子的对手）的人，可能被某些更迫近的危险的恐惧所激动，以至于他的注意力被强行从早先的焦虑中转移，而去谋划防范其他更紧迫威胁的措施？因此，不具体说明所有这些考虑，我仅斗胆断言：当贤士未能给他带回任何报告时，\Person{希律}可能相信他们被一个欺骗性的星象所误导，并且在未能发现他们以为已经出生的那一位之后，他们羞于回到他那里；就这样，国王的恐惧得以平息，放弃了寻找和迫害那孩子的念头。因此，当他们在他的母亲取洁之后与他一同上\Place{耶路撒冷}，并在圣殿中履行了\Person{路加}所叙述的那些事情时，因为\Person{西默盎}和\Person{亚纳}关于他所说的预言的话，当那些听见的人开始传扬时，可能会把国王的心思再次唤回到他最初的计划，\Person{若瑟}便听从在梦中得到的警告，带着孩子和他的母亲逃往\Place{埃及}。后来，当在圣殿中所行所说的事情完全公开后，\Person{希律}意识到自己被愚弄了；于是，由于渴望置基督于死地，他杀死了大批孩童，正如\Person{玛窦}所记载的。\par

\SourceRecord{Translated by S.D.F. Salmond.；Nicene and Post-Nicene Fathers, First Series；Vol. 6.；Edited by Philip Schaff.；Buffalo, NY: Christian Literature Publishing Co.,；1888.；<http://www.newadvent.org/fathers/1602211.htm>.}

% structure: p0001 source=h1 target=section reason=page-title

\section{福音合参，第二卷}

% structure: p0002 source=h2 target=subsection reason=relative-heading-rank; base=h2

\subsection{第十二章 论四福音作者各自归于若翰的话}

25. 再者，玛窦以下列方式记述若翰的事：—\Quote{那时，有洗者若翰出现在犹太旷野宣讲说：\InnerQuote{你们悔改吧！因为天国临近了。}这就是那藉先知依撒意亚所说的人：\InnerQuote{旷野中有呼号者的声音：你们该预备主的道路，修直祂的途径。}}马尔谷和路加也同意将这依撒意亚的证词视为指向若翰。的确，路加在叙述洗者若翰时，还记下了同一先知的其他话，紧接在已引用的经文之后。福音作者若望也提到洗者若翰亲自将这依撒意亚的证词用在自己身上。与此类似，玛窦在此给我们提供了一些若翰的话，是其他福音作者没有记载的。因为他提到他\Quote{在犹太旷野宣讲说：\InnerQuote{你们悔改吧！因为天国临近了。}}这些话被其他人省略了。但在接下来的内容中，紧接玛窦福音的那一段——即\Quote{旷野中有呼号者的声音：你们该预备主的道路，修直祂的途径}——这位置是含糊的；不清楚这是玛窦以自己口吻述说的，还是继续若翰本人的话，以致让我们理解整段都是再现若翰自己的发言，即：\Quote{你们悔改吧！因为天国临近了；这就是那藉先知依撒意亚所说的人}等等。因为若他说\Quote{我就是那藉先知依撒意亚所说的人}，而是用\Quote{这就是那被说的人}这种说法，这并不构成反对后一种观点的困难。因为这种说话方式是玛窦和若望福音作者习惯于用在自己身上的。例如，玛窦用了\Quote{他看见一个人坐在税关上}而不是\Quote{他看见我}。若望也说：\Quote{这是为这些事作证并写下这些事的门徒，我们知道他的见证是真实的}，而不是\Quote{我是}等等，或\Quote{我的见证是真实的}。是的，我们的主自己也经常使用\Quote{人子}或\Quote{天主子}这些词，而不说\Quote{我}。同样，祂告诉我们\Quote{基督必须受苦，第三天从死者中复活}，而不说\Quote{我必须受苦}。因此，紧接着\Quote{你们悔改吧！因为天国临近了}之后的\Quote{这就是那藉先知依撒意亚所说的人}这句话，完全可能是洗者若翰继续谈论自己；因此，在引用了说话者本人的这些话之后，玛窦自己的叙述才继续，即这样说：\Quote{这若翰身穿骆驼毛的衣服}等等。但若情况如此，那么当他被问及关于自己有什么话说时，他根据福音作者若望的叙述回答\Quote{我是旷野中呼喊者的声音}，也就不足为奇了，因为他早已在吩咐他们悔改时说过同样的话。因此，玛窦接着告诉我们关于他的衣着和生活方式的细节，并这样记述：\Quote{这若翰身穿骆驼毛的衣服，腰束皮带，吃的是蝗虫和野蜜。}马尔谷也几乎以同样的说法给出了同样的陈述。但其他两位福音作者省略了它。\par

26. \Person{玛窦}继续叙述说：那时，\Place{耶路撒冷}、全\Place{犹太}和\Place{约但河}一带的人，都出来到他那里，在\Place{约但河}里受他的洗，承认自己的罪。但他见到许多\Person{法利塞人}和\Person{撒杜塞人}也来受洗，就对他们说：毒蛇的种类！谁指示你们逃避那将要来的忿怒？你们要结出与悔改相称的果实；不要心里说：我们有\Person{亚巴郎}为父。我告诉你们，\Person{天主}能从这些石头中给\Person{亚巴郎}兴起子孙来。现在斧子已经放在树根上，凡不结好果子的树，就砍下来丢在火里。我是用水给你们施洗，叫你们悔改；但那在我以后来的，能力比我更大，我就是给他提鞋也不配，他要用\OriginalTerm{圣神}{Holy Spirit}与火给你们施洗。他手里拿着簸箕，要扬净他的场，把麦子收在仓里，把糠用不灭的火烧尽了。\Person{路加}也记载了这段经文，几乎把同样的话归于\Person{若翰}。而且即使有些措辞上的变化，但实际意义并无偏离。例如，\Person{玛窦}告诉我们\Person{若翰}说：\BibleQuote{你们不要心里说：\InnerQuote{我们有\Person{亚巴郎}为父。}}而\Person{路加}则这样记载：\BibleQuote{你们不要开始说：\InnerQuote{我们有\Person{亚巴郎}为父。}}同样，在前者中有这句话：\BibleQuote{我是用水给你们施洗，叫你们悔改；}而后者则插入了群众问他们该做什么的问题，并记载\Person{若翰}以善行作为悔改的果实来回答他们——所有这些\Person{玛窦}都省略了。所以，当\Person{路加}告诉我们\Person{洗者若翰}回答那些心中默想关于祂、并思索祂是否就是\OriginalTerm{基督}{Christ}的众人时，他只简单地说：\BibleQuote{我是用水给你们施洗，}并没有加上\BibleQuote{叫你们悔改}这句话。进一步，在\Person{玛窦}中，\Person{洗者若翰}说：\BibleQuote{但那在我以后来的，能力比我更大；}而在\Person{路加}中，他被描述为说：\BibleQuote{但有一位能力比我更大的要来。}同样，根据\Person{玛窦}，他说：\BibleQuote{我就是给他提鞋也不配；}而根据另一位，他的话是：\BibleQuote{我就是给他解鞋带也不配。}后面这些说法\Person{马尔谷}也记载了，虽然他未提及那些其他事。因为，在描述了他的衣着和生活方式后，他这样继续：\BibleQuote{他宣讲说：有一位能力比我更大的在我以后来，我就是弯腰给他解鞋带也不配。我用水给你们施洗，但他要用\OriginalTerm{圣神}{Holy Spirit}给你们施洗。}因此，在鞋带的记载上，他与\Person{路加}不同，因为他加上了\BibleQuote{弯腰}一词；在洗礼的记载上，他与前两位都不同，因为他没有说\BibleQuote{与火}，而只说\BibleQuote{用圣神}。因为正如在\Person{玛窦}中，同样在\Person{路加}中，字句相同，且次序相同：\BibleQuote{他要用圣神与火给你们施洗}——唯一例外是\Person{路加}没有加上\BibleQuote{圣}这个形容词，而\Person{玛窦}这样写了：\BibleQuote{用圣神与火。}这三位的叙述得到福音作者\Person{若望}的证实，他说：\BibleQuote{\Person{若翰}为祂作证，呼喊说：\InnerQuote{这就是我曾说过的那一位：在我以后来的，成了在我以前的，因他原先在我以前。}}因为这样他便表明，那些话是\Person{若翰}在其他福音作者记载他说那些话的时候所说的。同样，他也让我们明白，\Person{若翰}是在重复并再次提请人们注意他已经说过的话，当他这样说：\BibleQuote{这就是我曾说过的那一位：在我以后来的。}\par

27. 如果有人问，我们该认为哪一些话最可能是\Person{洗者若翰}实际所说的确切言词，是\Person{玛窦}福音中所记载的，还是\Person{路加}福音中的，或是\Person{马尔谷}在提及他所说的寥寥数语中所引入的——他提及\Person{若翰}说了这几句，却省略了其余一切——那么，任何一个明智地懂得要获得真理的知识真正需要的就是确实把握所指的事物，不论它们碰巧用什么确切言词来表达，都不会认为在这样的事上给自己制造困难是值得的。因为即使一位作者可能保持某种措辞次序，而另一位呈现不同的次序，这当然并不是真正的矛盾。再者，尽管一位说了另一位省略的事，两者之间也不需要有任何对立。因为显然，福音作者们正是按照各自对这些事的记忆来陈述它们，并按照他们各自的偏好促使他们在某些点上采用更简短或更详细的叙述，然而对所叙述的事物本身却给出了相同的描述。\par

28. 如此，更切合当前问题的是，足够明显的是：既然福音真理——藉着那在一切受造之上永恒不变的天主圣言传达，却又藉着暂时性象征和人的口舌传播于普世——已掌握了最崇高的权威，我们就不应认为任何一位作者给出了不可靠的记述，如果当多人回忆某件他们听过或见过的事时，他们没有遵循完全相同的计划，或使用完全相同的词语，却仍描述了完全相同的事实。我们也不应产生这种想法，即使词语顺序有所变化，或某些词被替换为其他意思相同的词，或某些事略而不提（或因记录者未想起，或因从其他陈述中容易理解），或在其他事务中（可能不直接关联其当前目的，但为了叙述并保持时间顺序而决定提及），其中一人引入某事，他认为不必全面详细阐述，而只部分触及；或为了说明含义并使其彻底清晰，被授权撰写叙述者添加了自己的补充——确实不在题材本身，而在表达它的词语上；或虽然对事实本身保留完全可靠的理解，但他可能未能完全成功（尽管以此为目）准确照字面回忆并复述当时听到的话。此外，若有人断言，福音作者本当由圣神的能力获得那种能力，以确保他们彼此之间在词语种类、顺序或数量上毫无变化，此人便未能认识到：恰好是福音作者（在其现有条件下）的权威越卓越，所有其他提供真实事件陈述的人的可信度就越应通过他们得以建立得更坚固；这样，当多人碰巧叙述同一件事时，若一人与他人仅在能以福音作者自身先例为理由辩护的方式上有所不同，便绝不能正当地被指控为不真实。因为我们既不能认为也不能说任何一位福音作者说了虚假的话，同样明显的是，任何其他作者，若在回忆某事以叙述时，其遭遇仅与福音作者被证明所遭遇的方式相似，也同等不承担不真实的指控。正如防范一切虚假属于最高道德，我们更应受如此崇高的权威所支配，以至于当我们发现任何作者的叙述彼此不同，其方式与福音作者记录被证明包含变化的方式一致时，我们不应认为我们遇到了必然虚假的东西。同时，在最严重关乎教义教导的信实性方面，我们也应理解，所要寻求并接受的与其说是单纯的词语，不如说是事实本身的真理；因为对于那些不采用完全相同的陈述方式，但在事实和思想本身不呈现差异的作者，我们视他们在真实性上持有相同地位。\par

29. 那么，关于我在圣史们的几个叙述之间所作的那些比较，这些叙述中有什么必须被视为矛盾的呢？难道我们应当这样看待这一点吗：其中一位圣史给了我们这句话：\Quote{我连提他的鞋也不配}，而其他圣史说到\Quote{解鞋带}？因为在这里，差异似乎既不在单纯的措辞中，也不在措辞的顺序中，也不在任何简单的用语中，而是在实际的事实本身中，当一边提到\Quote{提鞋}，而另一边提到\Quote{解鞋带}时。因此，可以非常合理地问：若翰宣布自己不配做什么——是提鞋，还是解鞋带？因为如果他只说这两句话中的一句，那么那个能够记录所说之话的圣史似乎给出了正确的叙述；而以另一种形式给出这话的圣史，虽然可能确实没有给出[故意地]错误的陈述，但无论如何可能被认为是记忆失误，因而被认为是以一件事代替了另一件事。然而，合理的做法是，不对圣史们提出完全失实的指控，这不仅针对那种通过明确说谎而来的失实，也针对因遗忘而来的失实。因此，如果从\Quote{提鞋}和\Quote{解鞋带}这两句话中引申出不同含义是相关的，那么人们应该怎样正确解释这些事实呢？除非若翰确实说出了这两句话，要么在两个不同场合，要么在同一上下文中。因为他很可能这样表达：\Quote{我连解他的鞋带也不配，我连提他的鞋也不配}；然后一位圣史可能再现了这话的一部分，而其余的圣史再现了另一部分；尽管如此，他们所有人都确实给出了真实的叙述。此外，如果若翰在谈到主的鞋时，仅仅意在传达祂的至高和他自己的卑微，那么，无论他实际上说了哪句话——无论是关于解鞋带还是关于提鞋——任何作者只要用自己的话提到鞋，同时表达了同样的卑微之意，仍然正确地保存了同样的含义，因而没有偏离他所写之人的真实心意。因此，当我们谈论圣史的和谐时，有一个有用的原则，尤其值得牢记：即使他们引入了一些与叙述之人实际所说不同的话，只要他们仍然表达了与那位按字面再现话语的圣史所传达的完全相同的心意，那么就不应认为[有]背离真理的情况。因为这样我们学到有益的教训：我们的目标应该是除了确定说话者的心意和意图之外，别无他物。\par

\SourceRecord{Translated by S.D.F. Salmond.；Nicene and Post-Nicene Fathers, First Series；Vol. 6.；Edited by Philip Schaff.；Buffalo, NY: Christian Literature Publishing Co.,；1888.；<http://www.newadvent.org/fathers/1602212.htm>.}

% structure: p0001 source=h1 target=section reason=page-title

\section{\WorkTitle{福音的和谐，第二卷}}

% structure: p0002 source=h2 target=subsection reason=relative-heading-rank; base=h2

\subsection{第十三章 论耶稣的圣洗}

30. 接着玛窦继续叙述如下：\BibleQuote{那时耶稣由加里肋亚来到约但河若翰那里，为受他的洗；但若翰阻止祂说：\InnerQuote{我需要受祢的洗，而祢却来就我吗？}耶稣回答说：\InnerQuote{你暂且容许吧！因为我们应当这样，以成全全义。}于是若翰就容许了祂。}其他两位也证实了耶稣来到若翰那里。三位也提及祂受了洗。但其他两位完全省略了玛窦所记载的一个情节，即若翰曾对主说话，或主曾回答若翰。\par

\SourceRecord{Translated by S.D.F. Salmond.；Nicene and Post-Nicene Fathers, First Series；Vol. 6.；Edited by Philip Schaff.；Buffalo, NY: Christian Literature Publishing Co.,；1888.；<http://www.newadvent.org/fathers/1602213.htm>.}

% structure: p0001 source=h1 target=section reason=page-title

\section{\WorkTitle{福音合参卷二}}

% structure: p0002 source=h2 target=subsection reason=relative-heading-rank; base=h2

\subsection{第14章 论他受洗后从天上来的话语或声音}

31. 此后，玛窦这样记载：\BibleQuote{耶稣受洗后，立时从水里上来；忽然天为他开了，他看见天主圣神有如鸽子降下，来到他上面；又有声音由天上说：\InnerQuote{这是我的爱子，我所喜悦的。}} 其他两位，即马尔谷和路加，也以类似方式记载了这件事。但同时，他们在忠实原意的情况下，以不同的措辞重现了来自天上的声音。因为虽然玛窦说那话是\BibleQuote{这是我的爱子}，而其他两位则写成\BibleQuote{祢是我的爱子}，但这些不同的表达方式只是传达了同一意义，根据前面讨论过的原则。因为天上的声音只发出了其中一句话；但通过采用\BibleQuote{这是我的爱子}这样的措辞，福音作者意在表明，这话是特别向在场的听众（而不是对耶稣）启示他是天主子的事实。出于这个观点，他选择了将\BibleQuote{祢是我的爱子}这一句转述为\BibleQuote{这是我的爱子}，仿佛是直接对人民说的。因为这话并非要向基督启示他已经知道的事实；而是要让在场的人民听见，为他们的缘故，那声音本身才发出。此外，关于第一位将这话表述为\BibleQuote{我所喜悦的}，第二位表述为\BibleQuote{在祢身上我所喜悦的}，第三位表述为\BibleQuote{在祢身上我得了喜悦}——如果你问这些不同形式中哪一种代表了声音实际说出的话，你可以选择你认为合适的一种，只要明白那些没有重复同样措辞的作者仍然传达了同样的意义。这些措辞的变化也有其用处：它们使我们能比只有一种形式时更充分地理解这句话，并且确保它不被解释为与实际情况不符。因为对于\BibleQuote{我所喜悦的}这句话，如果有人以为它意味着天主在圣子内自我悦纳，那么另一句\BibleQuote{在祢身上我所喜悦的}会教导他谨慎。另一方面，如果单看后一句，有人以为意思是圣父在圣子身上获得了人的喜爱，那么第三句\BibleQuote{在祢身上我得了喜悦}会让他有所领悟。由此充分表明，无论哪一位福音作者为我们保留了天上声音的逐字原话，其他人都只是以更熟悉的措辞来传达同一意义；这样，他们所有人给出的这句话可以被理解为：在祢身上我定了我的喜悦；也就是说，藉着祢来成就我的喜悦。此外，关于路加福音某些抄本中所包含的另一段记载，即天上的声音所说的话是圣咏中写的\BibleQuote{祢是我的儿子，我今日生了祢}；虽然据说在更古老的希腊抄本中找不到，但如果它可以得到任何可信抄本的证实，那么结论岂不是我们假定有两个声音从天上发出，以某种顺序说出？\par

\SourceRecord{Translated by S.D.F. Salmond.；Nicene and Post-Nicene Fathers, First Series；Vol. 6.；Edited by Philip Schaff.；Buffalo, NY: Christian Literature Publishing Co.,；1888.；<http://www.newadvent.org/fathers/1602214.htm>.}

% structure: p0001 source=h1 target=section reason=page-title

\section{福音的和谐 卷二}

% structure: p0002 source=h2 target=subsection reason=relative-heading-rank; base=h2

\subsection{第15章 解释为何根据福音作者若望，若翰洗者说：\Quote{我不认识祂；}而根据其他福音，我们发现他已经认识祂。}

32. 再者，在若望福音中关于鸽子的记载并未提及事件发生的时间，而是记录了若翰洗者报告他所见到的言论。在此处，问题在于如何理解经上所说的：\Quote{我不认识祂；但派遣我来以水施洗的那一位对我说：你看见圣神降下并停留在谁身上，谁就是那以圣神施洗的。}因为如果若翰只是在看见鸽子降在祂身上时才认识祂，那么就要追问：当祂前来受洗时，他怎能对祂说：\Quote{我本应受祢的洗。}因为洗者在鸽子降下之前便对祂说了这话。由此显见，虽然在此以前他（在某种意义下）已经认识祂，——因为当玛利亚拜访依撒伯尔时，他甚至在她腹中跳跃——，但直到此时仍有一件事他不知道，且藉着鸽子的降下而知道，即：祂以属于祂自己的某种神圣权能以圣神施洗；因此，凡从天主领受这洗礼的人，即使他为一些人施洗，也不能说他分施的是自己的东西，或者说圣神是由他赐予的。\par

\SourceRecord{Translated by S.D.F. Salmond.；Nicene and Post-Nicene Fathers, First Series；Vol. 6.；Edited by Philip Schaff.；Buffalo, NY: Christian Literature Publishing Co.,；1888.；<http://www.newadvent.org/fathers/1602215.htm>.}

% structure: p0001 source=h1 target=section reason=page-title

\section{福音的和谐，第二卷}

% structure: p0002 source=h2 target=subsection reason=relative-heading-rank; base=h2

\subsection{第十六章。论耶稣的诱惑。}

玛窦继续叙述如下：\Quote{当时耶稣被圣神领往旷野，为受魔鬼的试探。他禁食了四十昼夜，后来就饿了。试探者前来对他说：\InnerQuote{你若是天主子，就命令这些石头变成饼吧！}耶稣回答说：\InnerQuote{经上记着：\BibleQuote{人生活不只靠饼，也靠天主口中所发的一切言语。}}于是记述继续，直到我们读到这句话：\InnerQuote{于是魔鬼离开了他；看，有天使前来伺候他。}} 这段叙述路加也以类似的方式给出，只是顺序不同。这就使得我们不确定后两次试探中哪一次先发生：究竟是先向祂展示世上的万国，然后祂被带到圣殿的顶端；还是后者先发生，然后前者随之而来。然而，这并不重要，只要清楚所有的事件确实发生了。而且，由于路加用不同的词句陈述了相同的事件和观念，无需注意事实：真理并未因此受损。此外，玛尔谷确实证实了他在旷野中四十昼夜受魔鬼试探的事实，但没有说明对他说的或他回答的话。同时，他没有忽略路加省略的情节，即天使伺候他。然而，若望却略过了整个段落。\par

\SourceRecord{Translated by S.D.F. Salmond.；Nicene and Post-Nicene Fathers, First Series；Vol. 6.；Edited by Philip Schaff.；Buffalo, NY: Christian Literature Publishing Co.,；1888.；<http://www.newadvent.org/fathers/1602216.htm>.}

% structure: p0001 source=h1 target=section reason=page-title

\section{\WorkTitle{福音合参}，第二卷}

% structure: p0002 source=h2 target=subsection reason=relative-heading-rank; base=h2

\subsection{第十七章　论宗徒捕鱼时的蒙召}

34. \Person{玛窦}的叙述如下：\BibleQuote{耶稣听见若翰被下狱，就退往加里肋亚去了。} \Person{马尔谷}也记载同样的事，\Person{路加}亦然，只是路加在本段内没有提到若翰被下狱。圣史\Person{若望}又说，在耶稣往加里肋亚以前，有一天伯多禄和安德肋同祂在一起，那时前者得了伯多禄这个名字，而在此之前他被称为西满。同样，\Person{若望}告诉我们，次日耶稣决意往加里肋亚去时，遇见了斐理伯，并对他说，要他跟随自己。如此，圣史便叙述了关于纳塔乃耳的事。他又告诉我们，第三日在加里肋亚，耶稣在加纳行了变水为酒的奇迹。这些事迹其他圣史都略而不述，他们直接继续叙述耶稣回到加里肋亚的事。因此我们要明白，这里相隔了几天，其间发生了若望在这里插入的门徒们的事迹。这与\Person{玛窦}所说主对伯多禄说的话并不矛盾：\BibleQuote{你是伯多禄，在这磐石上我要建立我的教会。}但我们不应认为那是他第一次得到这个名字；而应认为这件事发生在若望所提到的那一次，即当祂对他说：\BibleQuote{你要称为刻法，意为磐石。}因此主后来可以这样称呼祂，当祂说：\BibleQuote{你是伯多禄。}因为祂当时不是说：\BibleQuote{你要称为伯多禄，}而是说：\BibleQuote{你是伯多禄；}因为先前已经这样对他说过：\BibleQuote{你要称为。}\par

35. 此后，\Person{玛窦}继续叙述说：\BibleQuote{祂离开纳匝肋城，来住在海滨的葛法翁，即在则步隆和纳斐塔里的境内；}等等，直到我们在山上讲道的结尾为止。在这段叙述中，\Person{马尔谷}与他一致，证实了门徒伯多禄和安德肋的蒙召，以及稍后雅各伯和若望的蒙召。但是，虽然玛窦紧接着在他治愈群众、大量人群跟随祂之后，介绍了祂在山上所讲的漫长讲道，马尔谷却在此时插入了其他事项，涉及祂在会堂里教导，以及人们对祂教义的惊奇。然后，他也陈述了玛窦所陈述的，虽然是在那漫长讲道之后才说的，即：\BibleQuote{祂教导他们像有权威的人，不像经师。}他同样给我们讲述了驱逐不洁之神的事，以及随后伯多禄岳母的故事。在这些事上，\Person{路加}也与他一致。但玛窦在这里没有提到邪魔的事。然而，伯多禄岳母的故事他并没有省略，只是在后面才提到。\par

36. 此外，在我们目前所考虑的这段中，同一\Person{玛窦}在叙述了他召唤那些捕鱼的门徒并命令他们跟随祂之后，接着叙述耶稣走遍加里肋亚，在会堂里教导，宣讲福音，并治好各种疾病；当群众聚集到祂周围时，祂上了山，并讲了那段漫长的讲道［已提及］。这样，圣史使我们理解到，那些由\Person{马尔谷}在召叫同一批门徒之后所记录的事迹，发生在祂走遍加里肋亚并在他们的会堂里教导的时期。我们也可以假设，发生在伯多禄岳母身上的事是在这个时候；他后来提到了他这里略过的事，尽管他在那后来之处确实没有直接复述所有在前面省略的事。\par

37. \Person{若望}所记载的门徒蒙召事迹，其内容是这样的：最初，并不在\Place{加里肋亚}，而是在\Place{约但河}附近，\Person{安德肋}同另一位名字未透露的门徒首先成为了主的跟随者；其次，\Person{伯多禄}从他那里得到了这个名字；第三，\Person{斐理伯}蒙召跟随他。而其他三位福音作者，彼此协调一致，尤其是\Person{玛窦}和\Person{马尔谷}在此处高度一致，告诉我们这些人在捕鱼时蒙召。诚然，\Person{路加}并未提\Person{安德肋}的名字。然而，我们可以从\Person{玛窦}和\Person{马尔谷}的叙述中推断他也在同一条船上，他们简要叙述了事情发生的经过。但\Person{路加}却给我们提供了更详尽、更清晰的说明，并且还记载了在那里捕鱼时发生的一个奇迹，以及在此之前主坐在船上向群众讲话的事实。在这方面似乎也有一个不一致之处：\Person{路加}记载了主所说的这句话：\BibleQuote{从今以后，你要得人了。}仿佛这话是主单对\Person{伯多禄}说的，而其他作者则显示这话是对兄弟二人说的。但这很可能的情况是，这些话最初是对\Person{伯多禄}本人说的，当时他因捕获了大量鱼而惊讶不已，这便是\Person{路加}所记的事件；稍后，这些话又对两人一起说了，这【第二次说话】便是另外两位福音作者所提及的。因此，我们提到关于\Person{若望}叙述的这段情况值得仔细考虑；因为它确实可能提出一个相当重要的矛盾。因为情况若是这样：在\Place{约但河}附近，\Person{耶稣}进入\Place{加里肋亚}之前，有两个人听了\Person{洗者若翰}的见证就跟随了\Person{耶稣}；这两个门徒中有一个是\Person{安德肋}，他立即去把自己兄弟\Person{西满}带到\Person{耶稣}那里；当时那兄弟得到了\Person{伯多禄}这个名字，以后他就要被这样称呼——那么，其他福音作者怎能说\Person{耶稣}在\Place{加里肋亚}看见他们捕鱼，并在那里召叫他们做他的门徒呢？这些不同的记载如何才能协调一致呢？除非我们理解为，那些人在\Place{约但河}附近的那次相遇中，并没有对\Person{耶稣}形成一种足以使他们永远跟随他的认识，而只是知道了他是谁，并且在最初对他的惊叹之后，又回到了他们以前的工作。\par

38. 因为【值得注意的是】同样在\Place{加里肋亚的加纳}，在他变水为酒之后，这位\Person{若望}告诉我们他的门徒如何相信了他。那奇迹的叙述如下：\BibleQuote{第三天，在加里肋亚加纳有婚宴，耶稣的母亲也在那里；耶稣和他的门徒也被请去赴宴。}现在，当然，如果正是在这个场合他们相信了他，如福音作者稍后告诉我们的，那么他们在被邀请赴婚宴时还不是他的门徒。但这是一种与我们说\Quote{宗徒\Person{保禄}生于\Place{基里基雅的塔尔索}}时所用的相同表达方式；因为那时他当然还不是宗徒。同样，这里告诉我们\Person{基督}的门徒被邀请赴婚宴，我们要理解的是，他们并非已经是门徒，而只是将要成为他的门徒。因为，当这叙述被整理并记录下来时，他们事实上已经是\Person{基督}的门徒；这就是为什么福音作者，作为过去历史事件的记录者，这样称呼他们。\par

39. 更进一步，关于\Person{若望}的陈述：\BibleQuote{此后，耶稣和他母亲、弟兄及门徒下到\Place{葛法翁}；他们在那里住了不多几天。}不确定到这时这些人是否已经依附于他，特别是\Person{伯多禄}和\Person{安德肋}，以及载伯德的儿子们。因为\Person{玛窦}首先告诉我们，他来到\Place{葛法翁}并住在那里，然后他从船上召叫他们，当时他们正在捕鱼。另一方面，\Person{若望}说他的门徒和他一起来到\Place{葛法翁}。也许\Person{玛窦}在这里只是回顾了一些他未按顺序记载的事情。因为他并没有说：\BibleQuote{此后，耶稣在\Place{加里肋亚海}边行走时，看见两个兄弟，}而是没有任何严格时间顺序的标记，简单地说：\BibleQuote{耶稣在\Place{加里肋亚海}边行走时，看见两个兄弟，}等等：因此，很可能他是在较晚的时候记载了并非实际发生在那个较晚时间的事情，而只是他之前省略未提的事情；这样一来，就可以理解为这些人已经和他一起来到了\Place{葛法翁}，\Person{若望}记载他确实去了那里，他和他母亲及门徒：或者我们是否应该认为这是一群不同的门徒，因为他可能已经有了一个跟随者\Person{斐理伯}，他正是以这种特别的方式召叫了他，对他说：\BibleQuote{跟随我}？因为全部十二位宗徒是如何被召的，从福音作者的叙述中并不清楚。实际上，不仅各种召叫的顺序没有记载，而且甚至并非所有宗徒的召叫事实都被提及，唯一具体说明的召叫是\Person{斐理伯}、\Person{伯多禄}和\Person{安德肋}、载伯德的儿子们，以及税吏\Person{玛窦}，他也被叫作\Person{肋未}。然而，单独从主那里得到一个名字的第一个人是\Person{伯多禄}。因为他没有个别给载伯德的儿子们起名字，而是称他们两个一起为\Quote{雷霆之子}。\par

40. 此外，我们当然应该注意到，\WorkTitle{福音书}和\WorkTitle{圣经}并非将祂\Quote{门徒}这一称呼仅限于那十二人，而是将所有信祂并在祂教导下为天国受培养的人都称为门徒。祂从他们全体中拣选了十二人，祂也称他们为宗徒，正如\Person{路加}所提及的。因为稍后他说：祂同他们下来，站在平地上，祂的门徒和大批群众聚集。如果他仅指十二个人，他肯定不会说\Quote{聚集} [或\Quote{人群}]的门徒。在\WorkTitle{圣经}的其他段落也明显可见，所有在祂指导下学习永生之道的人都被称为祂的门徒。\par

41. 但可能会有人问，祂如何从船上成对地召唤渔夫，即首先召唤\Person{伯多禄}和\Person{安德肋}，然后稍微往前走又召唤另外两人，即\Person{载伯德}的儿子们，根据\Person{玛窦}和\Person{马尔谷}的记载；而\Person{路加}对此事的说法是，他们两只船都装满了大量鱼。他还进一步陈述，\Person{伯多禄}的伙伴，即\Person{载伯德}的儿子\Person{雅各伯}和\Person{若望}，被叫来帮助这两人，因为他们无法拉出拥挤的网，所有在场的人都对捕到的大量鱼感到惊奇；当耶稣对\Person{伯多禄}说：\BibleQuote{不要怕，从此以后你要得人，}虽然这话只对\Person{伯多禄}说，但他们把船靠岸后，都跟随了祂。嗯，我们应当这样理解：\Person{路加}在这里引入的是最先发生的事，这些人此时并未被主召唤，只是向\Person{伯多禄}一人预言他将成为捕人的渔夫。而且，这句话的意思并非他们从此以后不再捕鱼。因为我们读到，甚至在主复活后，他们仍然捕鱼。因此，这些话只是意味从此以后他要得人，并无从此以后不再捕鱼的意思。这样，我们完全可以假定他们按习惯回到捕鱼上；因此，\Person{玛窦}和\Person{马尔谷}所记载的那些事件可能发生在之后。我指的是当祂成对地召唤门徒、并亲自命令他们跟随祂时，首先对\Person{伯多禄}和\Person{安德肋}说话，然后对另两人，即\Person{载伯德}的两个儿子。因为在那次，他们并非在把船拖上岸后、打算返回时跟随祂，而是立即跟随了祂，如同跟随一位召唤并命令他们跟随祂的人。\par

\SourceRecord{Translated by S.D.F. Salmond.；Nicene and Post-Nicene Fathers, First Series；Vol. 6.；Edited by Philip Schaff.；Buffalo, NY: Christian Literature Publishing Co.,；1888.；<http://www.newadvent.org/fathers/1602217.htm>.}

% structure: p0001 source=h1 target=section reason=page-title

\section{\WorkTitle{福音的和谐，第二卷}}

% structure: p0002 source=h2 target=subsection reason=relative-heading-rank; base=h2

\subsection{第十八章　论他前往\Place{加里肋亚}的日期。}

42. 此外，我们必须考虑这个问题：圣史\Person{若望}在提及洗者\Person{若翰}被下狱之前，告诉我们耶稣前往\Place{加里肋亚}。因为，在叙述他如何在\Place{加里肋亚}的\Place{加纳}变水为酒，以及他如何与他的母亲和门徒下到\Place{葛法翁}，并在那里住了不多几天之后，他告诉我们他随后为逾越节上耶路撒冷去；此后他带着门徒来到\Place{犹太}地，在那里同他们停留并施洗；然后在此处接着圣史说：\BibleQuote{若翰也在靠近\Place{撒冷}的\Place{艾农}施洗，因为那里水多；众人前来受洗，因为若翰尚未被下狱。}\BibleRef{若 3:23-24}另一方面，\Person{玛窦}说：\BibleQuote{耶稣听见若翰被下狱，就退往\Place{加里肋亚}去了。}\BibleRef{玛 4:12}同样，\Person{马尔谷}的话是：\BibleQuote{若翰被下狱以后，耶稣来到\Place{加里肋亚}。}\BibleRef{谷 1:14}\Person{路加}对若翰被下狱一事只字未提；但尽管如此，在他记述了基督受洗和受试探之后，他也发表了与其他两位相同的声明，即耶稣前往\Place{加里肋亚}。因为他在这里如此连接他的叙述各部分：\BibleQuote{当一切诱惑结束时，魔鬼离开了他，暂时；耶稣带着圣神的能力回到\Place{加里肋亚}，他的名声传遍了周围各地。}\BibleRef{路 4:13-14}然而，从这一切我们可以看出，并非这三位圣史发表了与圣史\Person{若望}相悖的声明，而只是他们没有记录主在受洗后首次前往\Place{加里肋亚}；在那次，他也曾在\Place{加纳}变水为酒。因为那时\Person{若翰}尚未被下狱。并且我们也应理解，这三位圣史在这些叙述的上下文中插入了另一次他前往\Place{加里肋亚}的记载，这次发生在\Person{若翰}被监禁之后。关于这次返回\Place{加里肋亚}，圣史\Person{若望}本人提供了以下信息：\BibleQuote{因此，耶稣知道法利塞人听见他收门徒并施洗比\Person{若翰}更多（虽然耶稣自己不施洗，而是他的门徒施洗），他就离开\Place{犹太}，又往\Place{加里肋亚}去了。}\BibleRef{若 4:1-3}因此，我们看到那时\Person{若翰}早已被下狱；而且，犹太人听闻他收门徒并施洗比\Person{若翰}更多。\par

\SourceRecord{Translated by S.D.F. Salmond.；Nicene and Post-Nicene Fathers, First Series；Vol. 6.；Edited by Philip Schaff.；Buffalo, NY: Christian Literature Publishing Co.,；1888.；<http://www.newadvent.org/fathers/1602218.htm>.}

% structure: p0001 source=h1 target=section reason=page-title

\section{\WorkTitle{福音合参 第二卷}}

% structure: p0002 source=h2 target=subsection reason=relative-heading-rank; base=h2

\subsection{第十九章 论\Person{玛窦}所载主在山上的长篇讲道}

43. 现在，关于主按照\Person{玛窦}的记载在山上的那篇长篇讲道，让我们先看看其他福音作者是否似乎与它毫无矛盾。\newline{}\Person{马尔谷}确实完全没有记录它，也没有保留任何与之相似的基督的言论，除了某些句子没有连贯地给出，而是分散出现在各处，是主在其他场合重复的。不过，他在叙述的文本中留下了一个空隙，指示我们可以理解为这篇讲道是在那一点讲的，尽管它没有被叙述出来。那就是他说：\Quote{祂在他们的会堂中，并走遍全\Place{加里肋亚}宣讲，并且驱逐魔鬼。}的地方。\newline{}在这段宣讲（他说耶稣在\Place{加里肋亚}各处所做的）的标题下，我们也可以理解那篇由\Person{玛窦}详述的山上讲道也包括在内。因为同一\Person{马尔谷}接着这样记述：\Quote{有一个癞病人来到祂跟前，恳求祂，并跪下向祂说：\InnerQuote{祢若愿意，就能洁净我。}}\newline{}他接着讲述这个癞病人被洁净的其余故事，其方式使我们明白，此人正是\Person{玛窦}所提到的那人，即主在讲道后下山时所治愈的。因为\Person{玛窦}这样记载：\Quote{耶稣从山上下来，有许多群众跟随祂；看，有一个癞病人前来叩拜祂说：\InnerQuote{主！祢若愿意，就能洁净我。}}等等。\par

44. 这个癞病人也被\Person{路加}提及，但未按此次序；而是依照福音作者们惯常的做法，有时在稍后记录先前遗漏的事，有时提前记下后来发生的事，这取决于他们受到上天的感动而想起那些事件——他们之前已经知道这些事，但后来被提示凭记忆将其写成文字。\newline{}然而，同一\Person{路加}也为我们留下了他自己版本的那篇丰富的讲道，从一段与\Person{玛窦}篇起始相同的段落开始。\newline{}因为在后者中，话是这样说的：\Quote{神贫的人是有福的，因为天国是他们的。}而在前者中，则是这样说的：\Quote{你们贫穷的是有福的，因为天主的国是你们的。}\newline{}此外，\Person{路加}叙述中的许多内容也与另一福音的记载相似。最后，两福音中这篇讲道的结尾完全相同——即聪明人把房子盖在磐石上，愚昧人把房子盖在沙土上的比喻；唯一不同的是，\Person{路加}只提到河水冲击房屋，而没有像\Person{玛窦}那样也提到雨和风。因此，很容易相信他在那里引用了与\Person{玛窦}相同的讲道，但同时省略了\Person{玛窦}插入的一些句子；他也加入了\Person{玛窦}未提及的其他言论；并且，他以类似的方式，用稍有不同的措辞表达了某些言论，但并未损害真理的完整性。\par

45. 我们很可以这样设想，正如我所说的，要不是感到难以处理的一个情况，即\Person{玛窦}告诉我们，这篇教诲是主坐在山上讲的；而\Person{路加}说，是主站在平地上讲的。这个差异因此使人觉得前者是指一篇教诲，后者是指另一篇。确实，有什么能阻止我们假定\Person{基督}在别处重复了祂已经说过的一些话，或者再做一次祂在先前某个场合已经做过的事呢？然而，这两篇教诲——一篇由\Person{玛窦}插入，另一篇由\Person{路加}插入——相隔时间不长，这一点很有可能从以下事实推断出来：在这两篇教诲的前后，两位福音书作者都叙述了某些相似或完全一致的事件，因此我们有理由感到，引入这些事件的作者所叙述的是相同的地点和日子。因为\Person{玛窦}的叙述继续如下：\BibleQuote{有许多群众从\Place{加里肋亚}、\Place{十城区}、\Place{耶路撒冷}、\Place{犹太}和\Place{约但河}对岸来跟随祂。祂看见群众，就上了山；祂坐下后，祂的门徒来到祂跟前：祂就开口教训他们，说：\InnerQuote{神贫的人是有福的，因为天国是他们的。}等等}。这里似乎表明，祂想摆脱大群的人，因此祂上了山，仿佛要离开群众，寻找单独与门徒讲话的机会。这一点似乎也被\Person{路加}证实了，他的叙述如下：\BibleQuote{在那些日子，祂出去上山祈祷，整夜向天主祈祷。到了天亮，祂叫祂的门徒来，从他们中拣选了十二人，祂称之为宗徒：\Person{西满}，祂也称之为\Person{伯多禄}，和他的兄弟\Person{安德肋}、\Person{雅各伯}与\Person{若望}、\Person{斐理伯}与\Person{巴尔多禄茂}、\Person{玛窦}与\Person{多默}、\Person{阿耳斐的儿子雅各伯}和号称热诚者的\Person{西满}、\Person{雅各伯的兄弟犹达斯}和\Person{依斯加略人犹达斯}，他就是那负卖者。祂同他们下来，站在平地上，又有祂的门徒和一大群从全\Place{犹太}和\Place{耶路撒冷}，并从\Place{提洛}和\Place{漆冬}的海边来的人，他们来听祂，并求医治他们的疾病；那些被污灵缠磨的人也得了医治。全体群众都寻求触摸祂，因为有能力从祂身上发出，治好了他们所有人。祂举目看着祂的门徒说：\InnerQuote{你们贫穷的是有福的，因为天国是你们的。}等等}。这里的关系让我们明白，在山上从较大的团体中选了十二个门徒（祂称之为宗徒，这件事\Person{玛窦}省略了）之后，祂就讲了\Person{玛窦}所引述的那篇教诲，而\Person{路加}没有记载它——即山上的那篇；然后，当祂下山以后，祂在平地上讲了第二篇与第一篇相似的教诲，对此\Person{玛窦}没有提及，但\Person{路加}详细叙述了；此外，这两篇讲道都以同样的方式结束。\par

46. 但是，关于\Person{玛窦}在那篇教诲结束后接着所陈述的，即：\BibleQuote{\Person{耶稣}讲完了这些话，众人都惊奇祂的教训，}——这可能表明，当时说话的正是那些群众中的门徒，祂就是从他们中拣选了十二人。此外，当福音书作者紧接着用这些话继续时：\BibleQuote{\Person{耶稣}下了山，有许多群众跟随祂；看，有一个癞病人前来朝拜祂，}我们可以假定，那件事发生在两篇教诲之后——不仅是在\Person{玛窦}所记载的那篇之后，也是在\Person{路加}所插入的那篇之后。因为下山后经过了多少时间并不清楚。但\Person{玛窦}的意图仅仅是表明事实本身：下山后，在祂洁净癞病人时，有大量群众与主在一起，而不是指明中间经过了多长时间。这一假定更容易被接受，因为\Person{路加}告诉我们，同样的癞病人在主当时在某城里时被洁净——这是\Person{玛窦}没有费心提及的情况。\par

47. 然而，毕竟也可以提出这样的解释：即起初，主和祂的门徒，没有其他人，是在山上某个更高的地方；在祂停留期间，祂从跟随者中选出了那十二人；然后祂和他们一同下来，并非从山本身下来，而是从山上所说的那个高处下到平地上——就是说，下到山坡上某个能容纳众多人群的平坦地方；此后，祂坐下，祂的门徒在祂旁边站好位置，在这种情况下，祂向他们和在场的其他群众讲了一篇道理，\Person{玛窦}和\Person{路加}都记录了这篇道理，他们叙述的方式确实不同，但在涉及两人都叙述的事实和言论时，两位作者同样忠实地传达了真理。因为我们在探讨之前已经预设了一个立场，这个立场本来应该不言自明，无需任何人事先提醒，即作者之间不一定有矛盾，尽管某人可能省略了另一人提及的内容；同样，尽管一人以某种方式陈述事实，另一人用不同方法，只要关于对象和言论本身的同一真理得以呈现，就没有矛盾。因此，\Person{玛窦}的话：\Quote{祂下了山以后}也可以同时理解为指平地，这平地很可能就在山坡上。之后，\Person{玛窦}讲述洁净癞病人的故事，\Person{马尔谷}和\Person{路加}也以类似方式记述。\par

\SourceRecord{Translated by S.D.F. Salmond.；Nicene and Post-Nicene Fathers, First Series；Vol. 6.；Edited by Philip Schaff.；Buffalo, NY: Christian Literature Publishing Co.,；1888.；<http://www.newadvent.org/fathers/1602219.htm>.}

% structure: p0001 source=h1 target=section reason=page-title

\section{福音的和谐，第二卷}

% structure: p0002 source=h2 target=subsection reason=relative-heading-rank; base=h2

\subsection{第二十章  解释为何玛窦记述百夫长为自己仆人前来见耶稣，而路加却说百夫长差遣朋友到祂那里}

48. 此后，玛窦继续叙述如下：\\BibleQuote{耶稣进了葛法翁，有一个百夫长来到祂跟前，恳求祂说：主，我的仆人瘫痪在家，痛苦难当；}以及以下经文，直到所说\\BibleQuote{他的仆人就在那一刻痊愈了。}这位百夫长仆人的事例，路加也记述了；只是路加不像玛窦那样，将之放在癞病人洁净之后（那个故事是他在后来回忆时记录的），而是在那篇已经讨论过的长篇讲道结束之后才引入。因为他这样连接两段：\\BibleQuote{耶稣在民众面前说完了这一切话，就进了葛法翁；有一个百夫长的仆人，是他所疼爱的，患病快要死了；}等等，直到记载他痊愈的那节经文。因此，我们注意到，基督是在民众面前说完了所有话之后才进入葛法翁；这仅仅意味着祂在结束这些言论之前没有进城；我们不应理解为这是指这些讲论与进入葛法翁之间所隔的时间长短。在这段时间内，那个癞病人被洁净了，玛窦将其记录在适当的位置，但路加只在较后之处提到。\par

49. 因此，让我们来考察玛窦和路加在这仆人的记述上是否一致。玛窦的话是这样的：\\BibleQuote{有一个百夫长来到祂跟前，恳求祂说：我的仆人瘫痪在家。}这似乎与路加提供的版本不一致，路加这样说：\\BibleQuote{他听见了耶稣，就打发犹太人的几个长老到祂那里，恳求祂来医治他的仆人。他们到了耶稣那里，便切切地求祂说：祢为他做这事，是配得的；因为他爱护我们的民族，并且给我们建造了会堂。耶稣就同他们去了。当祂离那家不远时，百夫长打发朋友到祂那里，对祂说：主，不要劳驾；因为祢进我屋顶下，我当不起；为此我也自以为不配去见祢；只要祢说一句话，我的仆人就会痊愈。}如果事情是这样发生的，那么玛窦说\\BibleQuote{有一个百夫长来到祂那里}，如何正确呢？因为那人不是亲自来，而是打发朋友来的。然而，如果我们仔细考察，这个表面的不一致就会消失，并注意到玛窦只是采用了一种非常常见的表达方式。因为我们习惯上说一个人\\Quote{来了}，即使他实际上还没有到达他所要接近的地方；同样，我们也说某人\\Quote{接近}或\\Quote{很接近}他所想要达到的目标；而且，我们也常常把为了达到目的而进行的\\Quote{接近}称为\\Quote{达到}，即使那个被称为\\Quote{达到}另一个人的人可能并没有亲眼见到他所接近的人，因为他可能是通过朋友来接近那个对他有帮助的人。事实上，这种习惯已经如此根深蒂固，以至于在今天的日常用语中，那些在拉选票时通过适当人物的中介接触到那些难以接近的、有影响力的人，被称为\\OriginalTerm{钻营者}{Perventores}。因此，如果\\Quote{接近}本身都可以通过他人来达成，那么通过他人而做出的\\Quote{靠近}（总是比实际\\Quote{接近}要少一些）为什么不能被称为\\Quote{来}呢？因为确实，一个人可以在\\Quote{靠近}方面做很多，但可能没有成功地实际达到他所寻求的目标。因此，玛窦在处理百夫长通过他人向主所作的\\Quote{靠近}时，选择了这种简洁的方式来表达，说\\BibleQuote{有一个百夫长来到祂那里}，这并非不合情理——实际上，任何有理智的人都能理解。\par

50. 然而，同时我们必须足够谨慎，以辨明圣史所采用的措辞中某种奥秘的深度，这措辞与\WorkTitle{圣咏}的这些话语相符：\BibleQuote{你们来就祂，就必得光明。}因为主亲自称赞了百夫长的信德——实际上，百夫长正是藉此信德真实地接近了耶稣——主如此说：\BibleQuote{在以色列我从未见过这么大的信德。}圣史明智地选择述说那人亲自来到耶稣面前，而不是提及那些替他传话的人。而且，\Person{路加}将整个事件如其所发生的那样向我们阐述，其叙述方式迫使我们理解到，另一位同样不可能作任何虚假陈述的作者，是如何可以说那人亲自前来。同样地，那位患血漏的妇人，虽然只是摸了祂衣服的繸头，却比那些拥挤祂的群众更有效地触摸了主。因为正如她因更加恳切地相信而更有效地触摸了主，同样，百夫长因更加彻底地相信而更加真实地来到了主面前。至于这段其余的部分，详细讨论一位作者叙述而另一位省略的各点，将是多余的工作。因为按照一开始所提及的原则，在这些差异中，找不到两位作者之间有任何实际矛盾。\par

\SourceRecord{Translated by S.D.F. Salmond.；Nicene and Post-Nicene Fathers, First Series；Vol. 6.；Edited by Philip Schaff.；Buffalo, NY: Christian Literature Publishing Co.,；1888.；<http://www.newadvent.org/fathers/1602220.htm>.}

% structure: p0001 source=h1 target=section reason=page-title

\section{福音的和谐，第二卷}

% structure: p0002 source=h2 target=subsection reason=relative-heading-rank; base=h2

\subsection{第二十一章 论伯多禄岳母事迹插入的顺序。}

51. 玛窦继续写道：\Quote{耶稣来到伯多禄家里，看见他的岳母躺着发烧，就摸了她的手，热症就退了；她便起来伺候他们。}玛窦没有指明此事的时间；也就是说，他没有说明它发生在何事之前或之后。因为我们当然不必认为，因为记载在某事件之后，就必然在事实上也发生在该事件之后。毫无疑问，此处我们应理解为，他在这里记下了先前省略的内容。因为马尔谷将这段叙述放在他记述癞病人洁净之前，而后者似乎被他置于宣讲山中圣训之后；然而，那篇讲道他却未提及。同样，路加插入伯多禄岳母的故事，也跟在马尔谷版本中同样的事件之后，但却在那篇较长的讲道之前；那篇讲道被他重述，并可能等同于玛窦所说在山中宣讲的讲道。因为他们在何处叙述，有何重要？是按正确顺序插入，或是在某处提前插入先前省略之事，或是提前提及后来真正发生之事，只要在叙述相同或不同事实时既不与自己矛盾，也不与其他作者矛盾，又有何区别？因为一个人虽然对事实有卓越可靠的知识，但决定他回忆的顺序并不在自己的能力之内（因为一件事进入脑海的先后并不是随我们意愿，而只是按所赐予的），所以完全可以假定，每位福音作者都认为自己的责任是按照天主乐意启示给他回忆的顺序，来叙述他所记录的事。至少对于那些顺序问题并不减损福音权威和真理的事迹，这一点可以成立。\par

52. 至于圣神为何随意分施给各人，并无疑是为了使他们的书建立在如此卓越的权威之上，也治理和支配圣人们的心思，在启示他们记录的事项方面，却允许一位史家以这种方式、另一位以那种方式建构叙述，这是一个可以虔诚探究的问题，并可能借天主助佑找到答案。但这并非我们目前着手的工作的目标。我们提议的任务只是要证明，无论每位福音作者在记述基督言行时采用何种具体顺序（是出于能力还是偏好），他既不与自己矛盾，也不与其他史家矛盾；这既适用于与他人相同的事，也适用于不同的事。因此，当时间顺序不明显时，我们不应以为任何一位所采用的顺序有何重要。但每当顺序明显时，如果福音作者呈现了似乎与自己或他人不一致的内容，我们当然应当考虑这段经文，并尽力澄清困难。\par

\SourceRecord{Translated by S.D.F. Salmond.；Nicene and Post-Nicene Fathers, First Series；Vol. 6.；Edited by Philip Schaff.；Buffalo, NY: Christian Literature Publishing Co.,；1888.；<http://www.newadvent.org/fathers/1602221.htm>.}

% structure: p0001 source=h1 target=section reason=page-title

\section{福音之谐和，卷二}

% structure: p0002 source=h2 target=subsection reason=relative-heading-rank; base=h2

\subsection{第二十二章 论本段之后所记事件的顺序，以及\Person{玛窦}、\Person{马尔谷}和\Person{路加}在这些事上是否彼此一致的问题}

53. \Person{玛窦}如是继续他的叙述：\Quote{到了晚上，他们带了许多附魔的人来到祂跟前；祂用祂的话赶出了魔鬼，并治好了所有患病的人：这为应验先知\Person{依撒意亚}所说的话，\InnerQuote{祂亲自承担了我们的脆弱，负担了我们的疾病。}}他藉着随后附上的这些话，清楚表明此事属于同一日：\Quote{到了晚上。}同样，在结束医治\Person{伯多禄}岳母的叙述，并以\Quote{她就服事他们}这句话作结之后，\Person{马尔谷}接着记载了以下的话：\Quote{到了傍晚，日落的时候，他们把所有患病的和附魔的人都带到祂跟前。全城的人都聚集在门口。祂治好了许多患各种疾病的人，并赶出许多魔鬼，且不许魔鬼说话，因为他们认识祂。清晨，天还很黑，祂就起身出去，到一个荒野的地方。}这里，\Person{马尔谷}似乎以这样的方式保留了顺序，即在他以\Quote{到了傍晚}这句话所作的陈述之后，他给出了这个时间说明：\Quote{清晨，天还很黑祂就起身出去。}虽然我们未必必须认为，当我们看到\Quote{到了傍晚}这句话时，它必定是指同一天的傍晚，或者当\Quote{清晨}这个短语出现时，它必定是指同一夜之后的早晨；然而，无论如何，事件中这种顺序似乎被保存下来，是为了时间的井然有序。此外，\Person{路加}在叙述\Person{伯多禄}岳母的故事之后，虽然没有明确说\Quote{到了傍晚}，但至少使用了一个表达相同意思的短语。因为他接着这样说：\Quote{日落以后，所有有各种病人的都把他们带到祂跟前；祂按手在每个人身上，治好了他们。也有魔鬼从许多人身上出来，喊着说：\InnerQuote{祢是基督，天主之子。}但祂斥责他们，不许他们说话，因为他们知道祂是基督。天一亮，祂就出去，来到一个荒野地方。}这里我们再次看到与\Person{马尔谷}中完全相同的顺序。但\Person{玛窦}似乎引入\Person{伯多禄}岳母的故事并非按照事件本身的顺序，而仅仅按照他在之前省略后心中所暗示的次序，他首先记载了同一天所发生的事，即傍晚来临；之后，他没有附上关于\Emph{清晨}的说明，而是以这些话继续他的叙述：\Quote{现在\Person{耶稣}看见许多群众围着祂，就命令渡到湖的对岸去。}那么，这就是一件新事，与\Person{马尔谷}和\Person{路加}上下文所给的记载不同，后者在提到\Emph{傍晚}之后，又提到了\Emph{清晨}。因此，关于\Person{玛窦}中的这节经文\Quote{现在\Person{耶稣}看见许多群众围着祂，就命令渡到湖的对岸去}，我们只应理解为，他在这里引入了另一件事实，即他在此时想起的事实——在某一天，当\Person{耶稣}看见许多群众围着祂时，祂指示渡到湖的对岸。\par

\SourceRecord{Translated by S.D.F. Salmond.；Nicene and Post-Nicene Fathers, First Series；Vol. 6.；Edited by Philip Schaff.；Buffalo, NY: Christian Literature Publishing Co.,；1888.；<http://www.newadvent.org/fathers/1602222.htm>.}

% structure: p0001 source=h1 target=section reason=page-title

\section{福音书和谐论，卷二}

% structure: p0002 source=h2 target=subsection reason=relative-heading-rank; base=h2

\subsection{第23章。论向主说\Quote{无论你往哪里去，我都要跟随你；}的那人，以及与此相关的其他事情，并论玛窦与路加记载这些事的次序。}

54. 他接下来说：\Quote{有一个经师来向他说：师傅，无论你往哪里去，我都要跟随你}，如此等等，一直到\Quote{任凭死人埋葬他们的死人。}在路加福音中也有类似的叙述。但他是在说了其他许多事情之后才加入这段的，没有明确的时间顺序，只是想起这件事时才写在那里。他让我们不确定他是在那里补入了一件之前遗漏的事情，还是预先提到一件在历史顺序中其实发生在后面的事情。因为他接着写道：\Quote{发生了一件事：当他们走在路上时，有一个人向他说：无论你往哪里去，我都要跟随你。}而主的回答与我们在玛窦福音中所读到的是完全相同的。现在，虽然玛窦告诉我们这事发生在主下令渡到湖的对岸的时候，而路加却提到他们在\Quote{走在路上}的时候，这并不矛盾。因为可能是他们走在路上就是为了来到湖边。再次，关于那个请求先埋葬他父亲的人，玛窦和路加完全一致。其实仅仅因为玛窦先引用了那人为父亲请求的话，然后才记下主的话\Quote{跟随我}，而路加则先记下主的命令\Quote{跟随我}，然后才是请求者的声明，这对意义本身并无影响。路加还提到了另一个人，他说：\Quote{主，我要跟随你，但请容许我先回家告别我家里的人。}玛窦没有提到这个人。之后路加转向了一个完全不同的主题，并没有按照实际的时间顺序继续。这段经文写道：\Quote{\Emph{此后}，主又选定了另外七十二人。}而这事发生在\Quote{此后}确实明显；但主在\Quote{此后}多久才这样做并不清楚。然而，在这段时间内发生了玛窦紧接着所记述的事情。因为同一玛窦仍然保持着时间顺序，继续他的叙述，正如我们接下来将看到的。\par

\SourceRecord{Translated by S.D.F. Salmond.；Nicene and Post-Nicene Fathers, First Series；Vol. 6.；Edited by Philip Schaff.；Buffalo, NY: Christian Literature Publishing Co.,；1888.；<http://www.newadvent.org/fathers/1602223.htm>.}

% structure: p0001 source=h1 target=section reason=page-title

\section{《福音的和谐》第二卷}

% structure: p0002 source=h2 target=subsection reason=relative-heading-rank; base=h2

\subsection{第二十四章 论主在那次渡湖时在船上睡着，以及驱逐那些获准进入猪群的魔鬼；并论玛窦、马尔谷和路加对这些事所记的一切言行的一致性。}

55. \BibleQuote{祂一上了船，门徒就跟着祂。看，海里起了大风暴。}故事继续，直到\BibleQuote{祂来到了自己的城。}玛窦连续记述的两段叙事——即耶稣从睡眠中被唤醒并命令风浪后海面平静，以及那些被凶恶魔鬼附身、挣断锁链驱入荒野的人——马尔谷和路加也同样记载了。这些故事中有些部分由不同作者以不同词汇表达，但含义相同。例如，玛窦记载主说：\BibleQuote{你们为什么害怕？小信德的人哪！}而马尔谷的版本是：\BibleQuote{你们为什么害怕？难道你们没有信德吗？}因为马尔谷指的是那如同芥子般的完美信德；所以他实际上也提到了\BibleQuote{小信德}。路加则这样说：\BibleQuote{你们的信德在哪里？}因此，整句话可能如此：\BibleQuote{你们为什么害怕？你们的信德在哪里？小信德的人哪！}于是，一位记下了一部分，另一位记下了另一部分。门徒唤醒祂时所说的话也可能如此。玛窦记载：\BibleQuote{主啊，救我们，我们丧亡了！}马尔谷记载：\BibleQuote{师傅，我们丧亡，你不顾吗？}路加则简单地说：\BibleQuote{师傅，我们丧亡了！}这些不同的表达传达了唤醒主并渴望平安的同一含义。我们无需追问这些形式中哪一种更接近实际向基督所说的话。因为无论他们实际使用了这三种措辞中的哪一种，还是用了另一套未被任何福音作者记录、但同样能表达原意的不同言辞，这对事实本身有何影响？同时，也可能在多人共同唤醒祂时，所有这些表达方式分别被不同的人使用。同样，平息风浪后的感叹也是如此：按玛窦所记是\BibleQuote{这是什么样的人，连风和海也听从祂！}按马尔谷所记是\BibleQuote{你们以为这是什么人？连风和海也听从祂！}按路加所记是\BibleQuote{你们以为这是什么人？祂命令风和海，它们也听从祂。}谁看不出这些形式含义完全相同？因为\BibleQuote{你们以为这是什么人？}与\BibleQuote{这是什么样的人？}意思相同。而\BibleQuote{祂命令}一词即使被省略，至少也能理解为顺从是对发令者而言。\par

56. 此外，关于玛窦说有两个被魔鬼军团附身的人，这些魔鬼获准进入猪群，而马尔谷和路加只提到一个人，我们可以这样理解：这两人中有一位是某种特别显赫有名的人，该地区特别哀悼他的情况，并急切盼望他的解救。为了指出这一事实，两位福音作者认为只提此人即可，因与这事相关的名声更加广泛显著。也不应因不同福音作者所记的附身者话语形式不同而疑虑。因为我们可以要么将它们归于同一句话，要么认为它们全都被实际说出。同样，也不应因玛窦用复数称呼，而马尔谷和路加用单数而感到困难。因为后两者同时告诉我们，当那人被问名字时，他回答说他是\OriginalTerm{军团}{Legion}，因为魔鬼众多。此外，马尔谷说猪群在山附近，路加说它们在山坡上，这并无矛盾。因为猪群如此之大，一部分可能在山上，另一部分只在其附近。正如马尔谷明确告诉我们，约有二千头猪。\par

\SourceRecord{Translated by S.D.F. Salmond.；Nicene and Post-Nicene Fathers, First Series；Vol. 6.；Edited by Philip Schaff.；Buffalo, NY: Christian Literature Publishing Co.,；1888.；<http://www.newadvent.org/fathers/1602224.htm>.}

% structure: p0001 source=h1 target=section reason=page-title

\section{\WorkTitle{福音的和谐 卷二}}

% structure: p0002 source=h2 target=subsection reason=relative-heading-rank; base=h2

\subsection{第25章 论患瘫痪的人，主向他说：\Quote{你的罪赦了}和\Quote{拿起你的床}；特别论及玛窦和马尔谷在记载此事发生地点时是否彼此一致的问题，因为玛窦说此事发生在\Quote{自己的城}，而马尔谷说在葛法翁。}

57. 于是玛窦继续叙述，仍保持时间顺序，如下连接其叙述：\BibleQuote{祂上了船，渡过去，来到自己的城里。看，有人给祂带来一个瘫子，躺在床上；}等等，直到记载说：\BibleQuote{众人看见，就惊奇，并归光荣于天主，因为祂赐给了人这样的权柄。}马尔谷和路加也讲述了这瘫子的故事。现在，关于玛窦记载主说：\BibleQuote{孩子，放心罢，你的罪赦了；}而路加记载的称呼不是\BibleQuote{孩子}而是\BibleQuote{人}——这更有力地阐明了主的含义。因为罪被赦免的是\BibleQuote{人}，因为他是人，就不能说\Quote{我没有犯过罪；}同时，这种称呼方式也表明那赦免人罪的那位本身就是天主。马尔谷又用了和玛窦相同的措辞，但他省略了\Quote{放心罢}这几个字。的确，也有可能整句话这样说：\Quote{人，放心罢：孩子，你的罪赦了；}或这样说：\Quote{孩子，放心罢：人，你的罪赦了；}或者以其他合适的顺序说出。\par

58. 不过，当我们观察到玛窦这样记载瘫子的故事时，确实可能产生一个困难：\BibleQuote{祂上了船，渡过去，来到自己的城里。看，有人给祂带来一个瘫子，躺在床上；}而马尔谷却说此事不是发生在他的城里（即纳匝肋），而是在葛法翁。他的叙述如下：\BibleQuote{又过了些日子，祂又进了葛法翁，人听说祂在家里。就有许多人聚来，以致连门前也没有空地，祂就对他们讲道。有人带着一个瘫子来见祂，是由四个人抬来的。但因为人多，不能带到祂跟前，就拆了祂所在处的房顶，拆通了，把瘫子躺着的床缒下来。耶稣一见他们的信德；}等等。路加则没有提及事情发生的地点，而是这样叙述：\BibleQuote{有一天，祂正在施教，有法利塞人和法学士也坐在旁边，他们是从加里肋亚各乡村和犹太并耶路撒冷来的；主的德能正在催迫祂治病。看，有人用床抬着一个瘫子，要把他抬进去放在耶稣面前。但因人多，找不到抬进去的方法，就上了房顶，从瓦间把他连床缒到中间，放在耶稣面前。耶稣一见他们的信德，就说：人，你的罪赦了；}等等。因此，问题就只存在于马尔谷和玛窦之间，因为玛窦记载此事发生在主的城；而马尔谷说在葛法翁。如果玛窦明说是纳匝肋，这个问题就更难解决。但按目前情况，当我们想到加里肋亚本身就可称为基督的城，因为纳匝肋在加里肋亚，正如那由许多城组成的整个地区仍被称为罗马国家；再者，考虑到那众多民族都包括在那座城中，正如经上所写：\BibleQuote{人们论及你，啊，天主的城，说了光荣的事；}又考虑到天主的古旧子民，虽然住在许多城中，却仍被称为一个家，即以色列家——谁能怀疑耶稣是在自己的城［或\Emph{国家}］行了这件工作呢，因为祂是在葛法翁城中行的，而葛法翁是祂渡过革辣撒地区回来时所进入的加里肋亚的一座城，因此当祂来到加里肋亚时，无论祂在加里肋亚的哪个城镇，都可以正确地说祂来到了自己的城［或\Emph{国家}］？这种解释尤其可以辩护，因为葛法翁本身在加里肋亚中占有如此崇高的地位，被认为是一种都城。但即使完全不允许将\Quote{基督的城}理解为加里肋亚（纳匝肋所在地）或葛法翁（作为加里肋亚的某种首都），我们仍可断言，玛窦只是略过了耶稣来到自己的城后、抵达葛法翁之前所发生的一切，并且只是在这个点上附加了治愈瘫子的叙述；正如作者们在许多实例中所做的那样，忽略了许多中间事件，并且没有任何明确的遗漏指示，直接继续叙述，就好像他们所附加的内容确实紧接着发生一样。\par

\SourceRecord{Translated by S.D.F. Salmond.；Nicene and Post-Nicene Fathers, First Series；Vol. 6.；Edited by Philip Schaff.；Buffalo, NY: Christian Literature Publishing Co.,；1888.；<http://www.newadvent.org/fathers/1602225.htm>.}

% structure: p0001 source=h1 target=section reason=page-title

\section{\WorkTitle{福音的和谐，第二卷}}

% structure: p0002 source=h2 target=subsection reason=relative-heading-rank; base=h2

\subsection{第二十六章 论玛窦的蒙召，并论玛窦自己的记述与马尔谷和路加在谈及阿尔斐的儿子肋未时是否一致}

59. 玛窦接着以下列话语继续他的叙述：\BibleQuote{耶稣从那里前行，看见一个人名叫玛窦，坐在税关上，就对他说：跟随我！他就起来跟随了祂。} 马尔谷也记载了这件事，且保持了相同顺序，将其置于瘫痪者痊愈的记载之后。他的记载如下：\BibleQuote{祂又来到海边，所有群众都到祂那里，祂便教导他们。当祂经过时，看见阿尔斐的儿子肋未坐在税关上，就对他说：跟随我！他就起来跟随了祂。} 这里没有矛盾，因为玛窦与肋未是同一个人。路加也引述了这件事，置于同一瘫痪者痊愈的记载之后。他这样写道：\BibleQuote{此后祂出去，看见一个名叫肋未的税吏坐在税关上，就对他说：跟随我！他便放弃一切，起来跟随了祂。} 由此看来，最合理的解释是：玛窦在此以回忆先前略过之事的方式记载了这些事。因为我们确实必须相信，玛窦的蒙召发生在宣讲山中圣训之前。因为路加告诉我们，在那山上，在那次场合，从众门徒之中拣选了这十二位，耶稣也称呼他们为宗徒。\par

\SourceRecord{Translated by S.D.F. Salmond.；Nicene and Post-Nicene Fathers, First Series；Vol. 6.；Edited by Philip Schaff.；Buffalo, NY: Christian Literature Publishing Co.,；1888.；<http://www.newadvent.org/fathers/1602226.htm>.}

% structure: p0001 source=h1 target=section reason=page-title

\section{福音的和谐，卷二}

% structure: p0002 source=h2 target=subsection reason=relative-heading-rank; base=h2

\subsection{第27章 论那宴席上，人同时指责基督与罪人同席，且他的门徒不禁食；论福音书作者们对提出这些指责的人似乎给出不同记述；并论玛窦、马尔谷与路加在叙述这些人的言词及主的答覆时，彼此是否一致。}

60. 因此玛窦接着说：\BibleQuote{他坐席在屋里时，有许多税吏和罪人来，同耶稣和他的门徒一同坐席。}等等，直到我们读到：\BibleQuote{但把新酒装在新皮囊里，二者都保存了。}这里玛窦没有特别告诉我们耶稣是和税吏及罪人一同坐席是在谁的家里。这可能会使人以为他并不是严格按照时间顺序插入这条记录，而是以追忆的方式在这里叙述了一件实际上发生在不同场合的事，若不是马尔谷和路加以几乎完全相同的措辞重复了这记述，清楚地表明耶稣坐席是在肋未——即玛窦——的家里，而且所有随后的话都是在那里说的。因为马尔谷也记载了同样的事实，并保持同样的顺序，如下：\BibleQuote{当他在屋里坐席时，许多税吏和罪人也同耶稣一同坐席。}因此，当他说：\BibleQuote{在他家里}时，他显然是指他前面正在谈论的那个人，就是肋未。同样，路加在说了\BibleQuote{他对他说：跟随我；他便撇下一切，起来跟随了他。}之后，立即附上了这句话：\BibleQuote{肋未在自己家里为他设了大宴席，有许多税吏和其他人与他们同席。}这样，事情发生在谁家里就显而易见了。\par

61. 让我们接下来看看三位福音书作者都提到的、对主说的话，以及由他所作的答覆。玛窦说：\BibleQuote{法利塞人看见了，就对祂的门徒说：你们的师傅为什么同税吏和罪人一起吃饭？}这差不多以同样的措辞出现在马尔谷那里：\BibleQuote{他怎么同税吏和罪人一起吃喝呢？}只是我们发现玛窦遗漏了马尔谷插入的一件事——即加上的\BibleQuote{又喝酒}。但这又有什么关系呢？既然意思完全表达出来了，所暗示的观念是他们一同用餐。路加则似乎有些不同地记载了这场面。因为他的版本如下：\BibleQuote{但他们的经师和法利塞人向祂的门徒发怨言说：你们为什么同税吏和罪人一起吃喝呢？}但他的意图当然不是表明那次他们的师傅没有被提及，而是要暗示这指责是针对他们全体，即他本人和他的门徒一起的；然而，这控诉虽然要被认为既是指他也是指他们，却并非直接对他说，而是对他们说的。事实上，路加本人和其他人一样，也描绘主自己作答覆，说：\BibleQuote{我来不是召义人，而是召罪人悔改。}若他们的\BibleQuote{你们为什么吃喝？}不是特别针对他的，他就不会那样回答他们了。基于同样的理由，玛窦和马尔谷告诉我们，反对他的异议是直接向他的门徒提出的，因为，当指控是对门徒提出时，这控诉就更加严重地落在门徒所效法并跟随的师傅身上。因此，传达的是同一个意思；而且由于一些措辞上的变化，表达反而更好，而事实本身未受影响。同样，我们也可以处理主答覆的记载。玛窦的记载如下：\BibleQuote{健康的人不需要医生，有病的人才需要；你们去研究一下\InnerQuote{我要的是仁慈，不是祭献}是什么意思；因为我来不是召义人，而是召罪人。}马尔谷和路加也以几乎相同的措辞为我们保存了同样的意思，只是他们都未引入先知的那句引文：\BibleQuote{我要的是仁慈，不是祭献。}路加又在\BibleQuote{我来不是召义人，而是罪人}之后，加上了\BibleQuote{以悔改}一词。这一补充更充分地阐明意思，以防止任何人认为基督爱罪人纯粹是因为他们是罪人。因为这幅\Emph{病人}的比喻也清楚表明天主召唤罪人是什么意思——就像医生对待病人——其目的确是要人从罪恶中得救，如同从疾病中得救；这医治是通过悔改而实现的。\par

62. 同样地，我们也可以考察关于\Person{若翰}门徒的记载。\Person{玛窦}的记述是这样：\Quote{那时，\Person{若翰}的门徒来到祂跟前说：为什么我们和法利塞人常常禁食？} \Person{马尔谷}的版本大意相似：\Quote{\Person{若翰}的门徒和法利塞人正在禁食。他们来对祂说：为什么\Person{若翰}的门徒和法利塞人禁食，而祢的门徒却不禁食？} 其中可能被认为不一致的唯一表象，在于推测法利塞人与\Person{若翰}的门徒一同发言是\Person{马尔谷}的添加，而\Person{玛窦}只陈述了\Person{若翰}的门徒表达了上述意思。但根据\Person{马尔谷}的版本，说话者与所提及的人并非同一批人。我的意思是，那些来到\Person{耶稣}面前的，是当时在场的宾客；他们之所以前来，是因为\Person{若翰}的门徒和法利塞人正在禁食，他们便就这些人说了上述话。这样，福音作者所说的\Quote{他们来}，就不应指他刚刚插入注释\Quote{\Person{若翰}的门徒和法利塞人正在禁食}中所提到的人。而是，由于那些人正在禁食，其他一些人被此事触动，来到祂跟前，向祂提出这个问题：\Quote{为什么\Person{若翰}的门徒和法利塞人禁食，而祢的门徒却不禁食？} \Person{路加}表达得更清楚。显然，他怀着同样的想法，在陈述了主以病人为比喻论及召叫罪人的回答之后，接着这样写道：\Quote{他们对祂说：为什么\Person{若翰}的门徒常常禁食祈祷，法利塞人的门徒也是这样，而祢的门徒却吃吃喝喝？} 由此可见，和\Person{马尔谷}一样，\Person{路加}提到一方就他方说了这番话。那么，为何\Person{玛窦}说\Quote{那时，\Person{若翰}的门徒来到祂跟前说：为什么我们和法利塞人禁食？} 可能的解释是，这些人也在场，所有这些不同的群体都急于提出这个指控，各人一有机会便说。三位福音作者用不同的措辞表达了当时想要表达的意思，但并未偏离事实本身的真实陈述。\par

63. 再者，我们发现\Person{玛窦}和\Person{马尔谷}关于新郎的子女在婚宴期间不禁食的记载相似，但\Person{马尔谷}称之为\Quote{婚宴的子女}，而\Person{玛窦}称之为\Quote{新郎的子女}。这一点无关紧要。因为\Quote{婚宴的子女}自然既包括新郎方面的人，也包括新娘方面的人。因此，含义明显相同，既非不同也非矛盾。\Person{路加}则没有说\Quote{新郎的子女岂能禁食？}，而是说\Quote{当新郎和他们在一起时，你们岂能使新郎的子女禁食？} 通过这种表达方式，福音作者巧妙地揭示了相同的含义，并以一种引人深思的方式传达出来。这样便暗示，那些说话的人正试图使新郎的子女哀恸禁食，因为他们要设法置新郎于死地。此外，\Person{玛窦}所用的\Quote{哀恸}一词与\Person{马尔谷}和\Person{路加}所用的\Quote{禁食}含义相同。因为\Person{玛窦}后来也说：\Quote{那时他们就要禁食了}，而不是\Quote{那时他们就要哀恸了}。但通过使用这个词，他表明主所论及的禁食，是关乎苦难中的谦卑。同样地，主后来用关于新布和新酒的比喻，显示这种禁食对于属感官和血肉的人而言是不合宜的，他们沉溺于身体的忧虑，因此仍停留在旧心态中，借此描绘了另一种禁食，即与心神居于属灵高峰的狂喜相关，因而在某种程度上远离身体的食物。这些比喻也以类似的措辞出现在其他两位福音作者中。显而易见，不需要真正的不一致，尽管一个人可能引入某些内容，无论是属于主题本身，还是仅仅属于表达主题的措辞，而另一个人省略了，只要既没有偏离意义的真正同一性，也没有在表达同一事物的不同形式之间造成矛盾。\par

\SourceRecord{Translated by S.D.F. Salmond.；Nicene and Post-Nicene Fathers, First Series；Vol. 6.；Edited by Philip Schaff.；Buffalo, NY: Christian Literature Publishing Co.,；1888.；<http://www.newadvent.org/fathers/1602227.htm>.}

% structure: p0001 source=h1 target=section reason=page-title

\section{\WorkTitle{福音合参 卷二}}

% structure: p0002 source=h2 target=subsection reason=relative-heading-rank; base=h2

\subsection{第二十八章　论会堂主管女儿的复活与触摸祂衣边的妇人；亦论这些事件的叙述顺序是否显示任何一位记述者之间存在矛盾；尤其论会堂主管向主提出请求时所用的言词。}

64. 仍然遵循时间的顺序，玛窦接着叙述如下：\Quote{祂还在对他们说这些话的时候，看，有一位主管前来朝拜祂，说：\InnerQuote{我的女儿刚才死了；但请来，将祢的手放在她身上，她必会活过来。}}等等，直到我们读到\Quote{那女孩子就起来了。这消息传遍了那整个地方。}另外两位，即马尔谷和路加，同样记述了同一件事，只是他们没有遵循同样的顺序。因为他们将这段叙述放在不同的地方，插入到另一上下文中；即，在祂渡过湖、从革辣撒人地区回来，逐出魔鬼并准许它们进入猪群之后。马尔谷在记述了革辣撒人中间所发生的事之后，这样引入道：\Quote{耶稣又坐船渡到对岸，一大群人聚集到祂跟前，祂正靠近海边。这时，来了一个会堂主管，名叫雅依洛；他一看见祂，就俯伏在祂脚前，}等等。由此我们确实可以理解，会堂主管女儿的事是发生在耶稣再次坐船渡湖之后。但是，从这些话语本身看不出这件事发生在渡湖之后多久。然而，显然有一些时间间隔。因为如果没有间隔，就没有时间窗口容纳玛窦刚刚记述的关于他家里宴席的那些情形。他确实按福音作者的方式叙述了这些事，仿佛是在讲别人的事。但这实际上是发生在他自己身上、在他自己家里的事。在那段叙述之后，紧接着的就是这会堂主管女儿的记载。因为他编排了整个叙述，使得从一件事到另一件事的过渡本身就足够清楚地表明，紧随其后的话语所叙述的是紧接着发生的事。因为在提及前一件事中耶稣关于新布和新酒所说的话之后，他紧接着不加中断地写下了这些话语：\Quote{祂还在对他们说这些话的时候，看，有一位主管前来。}这表明，如果此人在祂说这些话时接近祂，那么祂所行或所说的事就没有任何其他事可以插入。另一方面，在马尔谷的记述中，正如我们已经指出的，那个位置很明显，其他［他未记录］的事很可能插在那里。路加的情况也大致相同，他在记述了革辣撒人中间所行的奇迹之后，接着叙述会堂主管的女儿，但他并非以与玛窦的版本矛盾的方式来衔接；玛窦的话语\Quote{祂还在说这些话的时候}使我们清楚地明白，这件事发生在关于布和酒的比喻之后。因为路加在结束陈述革辣撒人中间所发生的事之后，这样进入下一个主题：\Quote{耶稣回来的时候，民众欣然接待祂，因为他们都在等候祂。看，来了一个人，名叫雅依洛，他是会堂的主管，就俯伏在耶稣脚前，}等等。这样我们便明白，民众确实在所说时机即时接待了耶稣，因为祂就是他们等候回来的人。但是，紧接着附加的话语\Quote{看，来了一个人，名叫雅依洛，}并非字面意义上的紧接着发生。相反，正如玛窦所记载的，与税吏的宴席发生在之前。因为玛窦将当前这件事与那宴席连接起来，使我们无法设想任何其他事件顺序是正确顺序。\par

65. 那么，在我们目前着手考察的这段叙述中，这三位福音作者确实毫无异议地一致记述了患血漏病的妇人。一个人沉默不语的某些事由另一个人叙述，或者马尔谷说\Quote{谁摸了我的衣服？}而路加说\Quote{谁摸了我？}，这并不真正要紧。因为一个只是采用了惯用的说法，另一个则给出了更精确的表达。但尽管如此，两者都表达了相同的意思。因为我们更常说\Quote{你扯我}，而不说\Quote{你扯我的衣服；}尽管用词不同，我们要表达的意思却足够明显。\par

66. 然而，与此同时，仍有一个事实：玛窦记载会堂长对主谈及他的女儿，不仅像是快要死的、或正在死的、或濒临断气的，而是已经死了；而另两位福音作者却报告她当时已濒临死亡，但尚未真正死去，并且严格遵循这一版本，以至于他们告诉我们，后来有人来通报她确实已死的消息，并传话说，因此师傅现在不该再劳烦自己前来，为要按手在她身上以防止她死去——好像他并无能力使死者复活似的。因此，我们有必要调查这一事实，以免它似乎显示记述之间存在矛盾。解释的方式是假设，由于叙述的简洁，玛窦宁愿将其表达为主确实被请求做他实际所做的事，即复活死者。因为玛窦所关注的，并非父亲关于女儿所说的单纯话语，而是更重要的，他的心思和意向。因此，他给出了足以表达父亲真实想法的言辞。因为他已完全绝望于女儿的情况，认为他刚离开时她正濒死，现在绝不可能还活着，他的想法反而是她可以被复活。因此，两位福音作者记录了雅依洛实际说出的话语。但玛窦却展示了他内心隐秘的愿望和想法。因此，两个请求确实都呈给了主：要么是恢复将死的女孩，要么若她已死，就复活她。但由于玛窦的目的是简短地讲述整个故事，他让父亲在请求中直接表达了那确是他自己的真实愿望、也是基督实际所行的事。的确，如果那两位福音作者，或其中之一，告诉我们父亲自己说出了从他家来的人所说的话——即耶稣现在不应再劳烦，因为女孩已经死了——那么玛窦放入他口中的话就与他的思想不符了。但实际情况是，并没有说他同意了那些带来消息并吩咐师傅不再前来的来人。同时，我们还需注意，当主用这些话对他说：\Quote{不要怕，只管信，她必会得救}时，他并非因他不信而责备他，实际上是鼓励他要有更强的信德。因为这会堂长拥有像那说\Quote{主，我信！请帮助我的不信}之人的信德。\par

67. 既然如此，从福音作者这些多样而又不相矛盾的叙述方式中，我们显然学到了一条极其有用且非常必要的教训——即，在任何人的言辞中，我们应当严密关注的只是作者想要表达的思想，而言语应当服从于此；此外，如果有人用不同的词语传达了一个人的真实含义，而未能逐字复述其言辞，我们不应认为他在给出不正确的陈述。我们也不该让那些对字句吹毛求疵的可悲之人以为真理必然以某种方式系于字母的一撇一捺；而事实是，不仅在言语上，同样在一切表达思想的方式中，所应关注的只是思想本身，而非其他。\par

68. 此外，关于\BibleBook{玛窦}{福音}的某些抄本包含的读法：\Quote{因为那女人不是死了，而是睡着了}，而马尔谷和路加证实她是一个十二岁的女孩，我们可以认为玛窦在此遵循了希伯来语的表达方式。因为在圣经的其他地方，以及在这里，我们发现不仅那些已认识男子的女子，而且所有女性，包括未经人事的贞女，都被称为女人。例如，在论及厄娃时写道：\Quote{他把她造成一个女人}；又在\BibleBook{户籍}{纪}中，那些未与男子同寝而认识男子的女人，即贞女，被命令免于处死。保禄也沿用同样的措辞，论及基督本身说他是\Quote{生于女人}。因此，理解为这些类比远比认为这个十二岁的女孩已经结婚或认识男子更好。\par

\SourceRecord{Translated by S.D.F. Salmond.；Nicene and Post-Nicene Fathers, First Series；Vol. 6.；Edited by Philip Schaff.；Buffalo, NY: Christian Literature Publishing Co.,；1888.；<http://www.newadvent.org/fathers/1602228.htm>.}

% structure: p0001 source=h1 target=section reason=page-title

\section{\WorkTitle{福音的和谐，第二卷}}

% structure: p0002 source=h2 target=subsection reason=relative-heading-rank; base=h2

\subsection{第二十九章 论及两个瞎子和一个附魔的哑巴，其故事只由玛窦记述}

69. 玛窦继续他的叙述，如下所述：\Quote{耶稣从那里前行，有两个瞎子跟着祂，喊叫说：\InnerQuote{达味之子，可怜我们罢！}}等等，直到我们读到\Quote{法利塞人却说：\InnerQuote{祂是赖魔王驱魔。}}玛窦是唯一记载这两个瞎子和附魔的哑巴故事的人。因为其他福音书所记载的那两个瞎子，并非我们眼前的这两个。然而，这些事件如此相似，以致如果玛窦自己没有同时记录后一件事，我们可能会认为他此刻所叙述的，也已被另外两位福音作者记述了。因此，有一点我们必须牢记：有些事件彼此相似。这一点可由同一位福音作者在此提及两件事得到证明。因此，如果我们有时发现几位福音作者各自叙述了相似的事件，并发现其中存在似乎无法以对观原则解决的矛盾，那么我们可以想到，解释很简单：那（看似矛盾）的情况并未（在那一特定场合）发生，而当时发生的只是与之相似的事，或是某种以类似方式流传的事。\par

\SourceRecord{Translated by S.D.F. Salmond.；Nicene and Post-Nicene Fathers, First Series；Vol. 6.；Edited by Philip Schaff.；Buffalo, NY: Christian Literature Publishing Co.,；1888.；<http://www.newadvent.org/fathers/1602229.htm>.}

% structure: p0001 source=h1 target=section reason=page-title

\section{\WorkTitle{福音的和谐，卷二}}

% structure: p0002 source=h2 target=subsection reason=relative-heading-rank; base=h2

\subsection{第三十章 论及一段记载：祂怜悯群众，派遣门徒，赐给他们治病权柄，并以许多教训吩咐他们，指示他们如何生活；并论及玛窦在此处与玛尔谷、路加和谐一致的证据，特别是关于手杖的问题——玛窦说主告诉他们不可携带，而按玛尔谷，手杖是他们唯一可携带之物；也论及穿鞋和衣服的问题。}

70. 关于接下来记述的事件，玛窦的叙述确实没有清楚表明它们的确切顺序。因为在关于两个盲人和哑巴附魔者的两个插曲之后，他接着这样写道：\BibleQuote{耶稣周游各城各村，在他们的会堂里施教，宣讲福音的国度，治好各种疾病和灾殃。但祂看见群众，就对他们动了怜悯的心，因为他们困苦流离，像没有牧人的羊。于是祂对门徒说：庄稼确实多，但工人少；所以你们要求庄稼的主人，派工人到祂的庄稼里去。祂叫了祂十二门徒来，给他们制伏邪魔的权柄；}等等，直到\BibleQuote{我实在告诉你们，他决不会失去他的赏报。}我们现在提到的这一整段经文显明了祂如何赐给门徒许多训诲。但是，玛窦是按历史的顺序把这一段附在这里，还是仅仅根据它出现在他自己脑海中的次序来排列，正如已经说过的，并不明显。玛尔谷似乎以简洁的方式处理了这一段，以下列语句开始叙述：\BibleQuote{祂周游各村，在那些村镇施教；祂叫了十二人来，开始派遣他们两个两个地出去，并给他们制伏邪魔的权柄；}等等，直到我们读到\BibleQuote{把脚上的尘土跺下来，作为反对他们的证据。}然而，在叙述这件事之前，玛尔谷在讲述会堂长女儿复活的故事之后，立即插入了一段在主本乡发生的事：人们惊讶于主，并问祂的智慧和能力从何而来，他们觉察到祂的判断。这段记述玛窦放在了给门徒的这些训诲以及许多其他事情之后。因此，不确定本乡所发生的事是否被玛窦按照回忆的次序记录，因为他最初省略了；或者是否被玛尔谷以提前的方式引入；简言之，他们中谁保持了实际发生的顺序，谁保持了回忆的顺序。路加则紧接在雅依洛女儿复活之后，插入了这一段关于权柄和门徒训诲的内容，而且与玛尔谷一样简洁。然而，这位福音作者并不比其他几位更以这样的方式引入主题，让人感觉它也严格符合历史顺序。此外，关于门徒的名字，路加在另一处——即较早的段落中，在他们【被描述为】在山上被拣选之时——列出了他们的名字，他与玛窦在任何方面都没有分歧，除了一个例子：雅各伯的兄弟犹达斯，玛窦称他为达陡，虽然有些抄本也读作肋拜。但谁会否认一个人可能有两个或三个名字呢？\par

71. 另一个通常提出的问题是：为什么玛窦和路加记载主对门徒说，他们不可带手杖，而玛尔谷却这样说：\BibleQuote{祂吩咐他们除了手杖外，路上什么都不要带}；接着又说\BibleQuote{不要带口袋，不要带饼，不要带钱在腰带里}；这清楚表明他的叙述与那些提到不可带手杖的人属于同一场景和情况。这个问题可以通过这样的理解来解决：按照玛尔谷，要带的手杖有一种含义；而按照玛窦和路加，不可带的手杖则应作另一种解释；正如我们发现\Emph{试探}一词在\BibleQuote{天主不试探人}中是一种用法，而在\BibleQuote{主你的天主试探（考验）你们，要知道你们是否爱祂}中是另一种用法。因为前者指诱惑的试探，后者指考验的试探。另一个类似的例子是\Emph{审判}一词，在\BibleQuote{行过善的，复活进入生命；作过恶的，复活进入审判}中是一种含义，而在\BibleQuote{天主啊，求你审判我，为我辨明冤情，对抗不敬虔的民族}中是另一种含义。前者指定罪的审判，后者指辨别的审判。\par

72. 此外，还有许多词汇并非始终保留一种统一的意义，而是根据上下文的不同而被引入，因此以多种方式被理解，有时甚至伴随解释而被采用。例如，有言说：\BibleQuote{在理解上不要作孩子；但在恶事上要作婴孩，好使你们在理解上达到完全。}因为这里有一个句子，以简洁而意味深长的形式，本可以这样表达：\Quote{不要作孩子；但也要作孩子。}同样的情形也出现在这些话中：\BibleQuote{你们中间若有人自以为在今生是聪明的，就让他变成愚拙，好成为聪明的。}因为那里所说的除了\Quote{不要聪明，好成为聪明}之外，还能是什么呢？此外，句子有时也被放置得如此，以致于锻炼探究者的判断力。这一类例子见于\BibleBook{迦拉达}{书}中所说的：\BibleQuote{你们要彼此担负重担，如此就满全了基督的法律。因为人若自以为是，其实什么都不是，就欺骗了自己。但各人应当考验自己的工作，这样他就有可夸的只在自身而不在他人。因为各人要背负自己的重担。}现在，除非\Quote{重担}一词可以有不同的含义，否则无疑会有人以为同一位作者在这里所说的话自相矛盾，而且这些话在同一段落中如此相邻。因为当他刚刚说了\Quote{一人要担负他人的重担}，片刻之后又说\BibleQuote{各人要背负自己的重担}。但前者指的是在分担他人软弱时所应背负的重担，后者指的是在向天主交代自己行为时所背负的重担；前者是在与弟兄的共融中应背负的重担，后者则是我们每个人为自己背负的特有重担。同样，宗徒在\BibleQuote{我岂是带着棍棒到你们那里去吗？}这句话中所说的\BibleQuote{棍棒}是精神意义上的；而同一术语当用于棍棒打马或其他类似目的时则具有字面意义，更不用提这一短语的其他隐喻意义了。\par

73. 因此，这两种劝告都必须被接受为主对宗徒所说的；即，他们既不应携带棍杖，又应除棍杖外不带任何东西。因为祂照玛窦所记载的对他们说：\BibleQuote{不要带金、银，也不要在钱袋里带铜钱，路上不要带行囊，不要带两件内衣，也不要带鞋，连棍杖也不要带。}祂立即加上：\BibleQuote{因为工人配得他的食物。}借此祂充分表明为什么祂要他们不带这些物品。祂显示：原因不在于这些物品对于维持此生并非必要，而在于祂派遣他们的方式，要清楚宣告这些物品应由那些信从他们宣讲福音的人提供给他们；如同兵士的薪饷是应得的，葡萄园的果实是栽种者的权利，羊群的奶是牧羊人的权利。为此保禄也这样说：\BibleQuote{有谁当兵而自备粮饷呢？有谁栽种葡萄园而不会吃它的果实呢？有谁牧养羊群而不会吃羊的奶呢？}因为在这些比喻之下，他谈论的是那些为福音宣讲者所必需的事物。因此，稍后他说：\BibleQuote{如果我们把属神的事物分施给你们，那么从你们收取属血肉的事物算得大事吗？如果别人在你们身上分享这权柄，我们岂不更应分享？然而我们并没有使用这权柄。}这显明通过这些指示，主的意思并非说福音传播者除了依赖那些听福音者所提供之物外，不应寻求其他支持（否则这位宗徒本人就违反了这条诫命，因为他亲手劳动谋生，以免成为任何人的负担），而是祂赐予他们一种权柄，在执行此权柄时，他们应知道这些事物是应得的。现在，当主颁赐任何诫命时，若不遵守就有不顺服的罪；但当授予任何权柄时，任何人都可以自由选择不使用它，犹如放弃自己的权利。因此，当主对这些门徒说这些话时，祂所做的正是那位宗徒稍后更清楚地阐释的，他说：\BibleQuote{你们岂不知道在圣殿供职的人，靠圣殿之物生活吗？在祭台旁侍候的人，与祭台分享吗？同样，主也如此规定：传福音的人应靠福音生活。但我没有使用过这些权利。}所以，当他说主如此规定，而他自己并未使用这规定时，他清楚表明所赐予他的是可以使用的权柄，而非必须履行的义务。\par

74. 因此，正如我们的主所规定的，宗徒宣告祂所规定的——即那些传福音的人应靠福音生活——祂将这些劝言赐予宗徒们，使他们无需为提供或携带此生所需之物（无论是大的还是最小的）而忧虑；因此，祂引入了这句话\Quote{连杖也不要带}，以表明，对于那些忠于祂的人而言，一切均归祂的仆人所有，而他们自己也无需任何多余之物。这样，当祂加上\Quote{工人自当得他的食物}时，就清楚明确地指示了祂说这一切话的方式和原因。因此，主用\Quote{杖}一词所指的正是这种权柄，即祂说他们\Quote{不要带}什么上路，只带一根杖。因为这句话也可以简要表达为：\Quote{不要随身带任何生活必需品，连杖也不要带，只带一根杖。}所以，\Quote{连杖也不要}可被理解为\Quote{连最小的东西也不要}；而加上\Quote{只带一根杖}则可理解为，凭借他们从主领受的、以\Quote{杖}（或\Quote{棍}）之名所象征的权柄，即使那些未携带之物也不会缺乏。因此，我们的主使用了这两句话。但由于同一位福音作者并未同时记录两者，那位告诉我们\Quote{杖}在一种意义上应被带的作者，被认为与那位告诉我们\Quote{杖}在另一种意义上不应被带的作者相矛盾。然而，在对此事作出这种解释之后，不应再存有这种假设。\par

75. 同样，当玛窦告诉我们不要让鞋随身上路时，其用意在于制止那种认为必须随身携带这些东西（否则可能缺乏）的顾虑。同样，关于两件内衣所说的话的意思是，他们中无人应想到在已穿的一件之外再带一件内衣，仿佛害怕会匮乏，而同时（从主领受的）权柄已确保他能获得所需之物。同样，当玛谷说他们应穿便鞋或鞋底时，他让我们明白，这鞋之事具有某种神秘意义，要点在于脚既不被遮盖，也不裸露在地；由此可传达福音既不应被隐藏，也不应依赖于世上的美好事物。至于所禁止的不是携带或拥有两件内衣，而是更明确地穿上它们——原文是\Quote{不要穿两件内衣}——那里面所传达的劝告岂非正是：他们应行走在真诚中，而非虚伪中？\par

76. 因此，决不可怀疑我主说了所有这些话，有的按字面意义，有的按比喻意义，尽管福音作者可能只将其一部分写入他们的著作——一人插入一段，另一人给出不同的部分。同时，某些段落或由其中两位、三位、甚至所有四位以相同措辞记录。然而即使如此，我们也不能假定祂所说或所做的一切都已记载下来。此外，若有人认为主在同一个讲论中不能既使用比喻性表达又使用字面性表达，只需查阅祂的其他讲论，就会看到这种想法是多么轻率和不加考虑。因为，那时（仅举一个目前浮现在我脑海的实例），当基督劝告不要让左手知道右手所做的时，他可能会认为有必要以同样比喻的方式理解所提及的施舍本身以及当时提供的其他训诲。\par

77. 实在地，我在此必须再次重复一个读者宜牢记的劝告，以免频繁要求这种建议，即：在祂的讲论的不同段落中，我主已重复了许多祂在别处早已说过的话。确实需要提醒这一点，以免当这些段落的顺序看似与另一位福音作者的叙述不吻合时，读者臆想这构成了他们之间的矛盾；而实际上他应理解这是由于在上下文中第二次重复了已在别处表达过的内容。这一评论不仅适用于祂的话语，也适用于祂的行为。因为没有什么能阻止我们相信同一件事可能发生过不止一次。但是，一个人仅仅因为不相信某件无人能证明是不可能的事件会重复发生而抨击福音，这显露了亵渎神明的虚荣。\par

\SourceRecord{Translated by S.D.F. Salmond.；Nicene and Post-Nicene Fathers, First Series；Vol. 6.；Edited by Philip Schaff.；Buffalo, NY: Christian Literature Publishing Co.,；1888.；<http://www.newadvent.org/fathers/1602230.htm>.}

% structure: p0001 source=h1 target=section reason=page-title

\section{福音的和谐，第二卷}

% structure: p0002 source=h2 target=subsection reason=relative-heading-rank; base=h2

\subsection{关于\Person{玛窦}和\Person{路加}记载\Person{洗者若翰}在监牢中时，派遣其门徒往主那里的情况。}

78. \Person{玛窦}继续叙述如下：\BibleQuote{耶稣吩咐完了他的十二门徒，就离开那里，往他们的城里去施教宣讲。若翰在监牢里听见基督的作为，就打发他的两个门徒去问他说：\InnerQuote{你就是要来的那一位，或是我们该期待另一位呢？}}\BibleRef{玛 11:1-3}等等，直到\BibleQuote{智能藉着她的子女彰显了正义。}\BibleRef{玛 11:19}整个关于\Person{洗者若翰}的这一部分——关于他派人到\Person{耶稣}那里去的消息、他所差的人所得到的答复的内容，以及在这些离去后主论及若翰的话——\Person{路加}也记载了。但次序却不相同。然而并不清楚哪一位是按照自己的回忆所排列的次序，哪一位是依照事情本身的历史顺序。\par

\SourceRecord{Translated by S.D.F. Salmond.；Nicene and Post-Nicene Fathers, First Series；Vol. 6.；Edited by Philip Schaff.；Buffalo, NY: Christian Literature Publishing Co.,；1888.；<http://www.newadvent.org/fathers/1602231.htm>.}

% structure: p0001 source=h1 target=section reason=page-title

\section{福音的和谐 卷二}

% structure: p0002 source=h2 target=subsection reason=relative-heading-rank; base=h2

\subsection{第三十二章 论祂谴责那些城市不悔改的场合，路加与玛窦都记载了此事；以及关于玛窦与路加在顺序上一致的问题。}

79. 之后玛窦继续写道：\Quote{然后祂开始谴责那些城市，因为祂的许多大能作为行在那里，而她没有悔改；}等等，直到我们读到：\Quote{在审判之日，索多玛地所受的惩罚也要比你们容易忍受。}路加也记述了这一段，他将它作为主在一次连续讲道中所说的话来记载。这个情况似乎表明，路加是按照主所说的严格顺序记录了这些话，而玛窦则是按照他自己记忆的顺序来排列的。或者，如果我们认为玛窦的话\Quote{然后祂开始谴责那些城市}必须理解为\Quote{然后}一词旨在表达这句话被说出的精确时间点，而不是以较宽泛的方式表示许多事情被做被说的时期，那么我说，任何秉持这种想法的人同样可以相信这些话是在两个不同场合说出的。因为即便是在同一位福音作者那里，也有一些话是主说了两次的，就像这位路加在关于旅途中不要带行囊的劝告上就是一个例子；同样，我们也在其他地方发现主在两处说了类似的话——那么，如果主原本在两个不同场合的另一些话，恰好被两位不同的福音作者各自按照实际所说的顺序记录下来，从而使得顺序看起来不同，这又有什么奇怪呢？仅仅是因为这些话既是在一位福音作者所注意到的场合说的，也是在另一位所提及的场合说的。\par

\SourceRecord{Translated by S.D.F. Salmond.；Nicene and Post-Nicene Fathers, First Series；Vol. 6.；Edited by Philip Schaff.；Buffalo, NY: Christian Literature Publishing Co.,；1888.；<http://www.newadvent.org/fathers/1602232.htm>.}

% structure: p0001 source=h1 target=section reason=page-title

\section{\WorkTitle{福音的和谐，第二卷}}

% structure: p0002 source=h2 target=subsection reason=relative-heading-rank; base=h2

\subsection{第三十三章 论祂召唤他们负起祂的轭和重担的时机，以及关于玛窦与路加在叙述顺序上并无差异的问题}

80. 玛窦继续这样记述：\BibleQuote{那时耶稣回答说：父啊！天地的主宰！我称谢你，因为你将这些事向智慧及明达的人隐藏起来，}等等，直到我们读到：\BibleQuote{因为我的轭是柔和的，我的担子是轻松的。}这段经文路加也注意到，但只是部分。因为他没有给我们这些词句：\BibleQuote{凡劳苦和负重担的，你们都到我这里来，}等等。然而，完全可以合理地假设这一切可能是主在同一个场合所说的，而路加并没有记录那个场合所说的全部。因为玛窦的措辞是：\Quote{那时耶稣回答说：}这指的是祂斥责那些城邑之后的时间。另一方面，路加插入了一些事情，虽然不多，在那斥责城邑之后；然后他加上这句话：\BibleQuote{就在那时刻，祂因圣神而欢欣说。}同样，我们看到，即使玛窦的表达不是\Quote{那时}，而是\Quote{就在那时刻}，路加在间隔中插入的内容仍然如此之少，以至于将其视为在同一时刻所说，似乎也不是不合理的。\par

\SourceRecord{Translated by S.D.F. Salmond.；Nicene and Post-Nicene Fathers, First Series；Vol. 6.；Edited by Philip Schaff.；Buffalo, NY: Christian Literature Publishing Co.,；1888.；<http://www.newadvent.org/fathers/1602233.htm>.}

% structure: p0001 source=h1 target=section reason=page-title

\section{福音的和谐，第二卷}

% structure: p0002 source=h2 target=subsection reason=relative-heading-rank; base=h2

\subsection{第三十四章  论记载中门徒掐食麦穗的经文；以及关于玛窦、马尔谷、路加在叙事顺序上如何彼此和谐的问题。}

81. 玛窦继续叙述如下：\BibleQuote{那时，耶稣在安息日从麦田经过；祂的门徒饿了，就掐起麦穗来吃；}一直到\BibleQuote{因为人子也是安息日的主。}这件事马尔谷和路加也都记载了，且没有任何矛盾。不过，后两者没有使用\Quote{那时}这个说法。因此，这可能更表明玛窦保留了实际发生的顺序，而其他人则按他们记忆的顺序叙述；除非，确实，\Quote{那时}这个短语应从更广泛的意义上理解，即指这些众多不同事件发生的时期。\par

\SourceRecord{Translated by S.D.F. Salmond.；Nicene and Post-Nicene Fathers, First Series；Vol. 6.；Edited by Philip Schaff.；Buffalo, NY: Christian Literature Publishing Co.,；1888.；<http://www.newadvent.org/fathers/1602234.htm>.}

% structure: p0001 source=h1 target=section reason=page-title

\section{\WorkTitle{福音的和谐，第二卷}}

% structure: p0002 source=h2 target=subsection reason=relative-heading-rank; base=h2

\subsection{\Emph{第三十五章：论枯手的人在安息日获得痊愈；以及论如何使\Person{玛窦}关于此事的叙述与\Person{马尔谷}和\Person{路加}的叙述相谐和，无论就事件的顺序或就主和犹太人所说的话语的记录。}}

82. \Person{玛窦}继续叙述如下：\BibleQuote{\Person{耶稣}从那里离开，进了他们的会堂。看，那里有一个人，他的一只手枯干了；}一直到\BibleQuote{那只手就复原了，像另一只一样。}枯手的人获得痊愈，这件事也没有被\Person{马尔谷}和\Person{路加}忽略。如今，若这一天也被指明是安息日，可能会使我们猜想，摘麦穗和治愈这人是发生在同一天，若不是\Person{路加}已清楚表明，治枯手是在另一个安息日。因此，当\Person{玛窦}说\BibleQuote{\Person{耶稣}从那里离开，进了他们的会堂}时，这些话语确实表明，他是在离开先前提到的地方之后才进入会堂的；但同时，它们并未确定从他离开上述麦田到他进入会堂之间经过了多少天；也没有说明他是直接连续地去那里。这样，我们就有了空间插入\Person{路加}的叙述，他告诉我们是在另一个安息日这只手被治愈了。但是，可能有一个困难：\Person{玛窦}告诉我们，民众向主提出了这个问题：\BibleQuote{在安息日治病合法吗？}他们以此寻找控告祂的借口；主以绵羊的比喻回答他们：\BibleQuote{你们中间谁有一只羊，如果安息日掉在坑里，不抓住它拉上来呢？人比羊贵重得多了！所以，在安息日行善是合法的。}而\Person{马尔谷}和\Person{路加}则说，是主向民众提出了这个问题：\BibleQuote{在安息日行善或作恶，救命或害命，哪样是合法的？}我们通过假设来解决这个困难：起初，民众向主询问：\BibleQuote{在安息日治病合法吗？}随后，主知道这些人心里正寻找控告祂的机会，就把那位即将被治愈的人放在他们中间，向他们提出了\Person{马尔谷}和\Person{路加}所记载的那些问题；当他们沉默时，主又向他们讲了绵羊的比喻，得出结论说在安息日行善是合法的；最后，主愤怒地向四周环视他们，正如\Person{马尔谷}告诉我们，因他们心硬而悲伤，便对那人说：\BibleQuote{伸出你的手来。}\par

\SourceRecord{Translated by S.D.F. Salmond.；Nicene and Post-Nicene Fathers, First Series；Vol. 6.；Edited by Philip Schaff.；Buffalo, NY: Christian Literature Publishing Co.,；1888.；<http://www.newadvent.org/fathers/1602235.htm>.}

% structure: p0001 source=h1 target=section reason=page-title

\section{\WorkTitle{福音的和谐，第二卷}}

% structure: p0002 source=h2 target=subsection reason=relative-heading-rank; base=h2

\subsection{第三十六章 另一个需要思考的问题：即在记载枯手人复原之后，这三位福音作者在进入下一个主题时，其叙述顺序是否没有矛盾？}

83. 玛窦继续叙述，以下列方式将其与上文衔接：\Quote{但法利塞人出去，商议怎样陷害耶稣，为能除掉祂。耶稣知道了，就离开那里；有许多人跟随祂，祂治好了他们所有的人，并嘱咐他们不要张扬祂，为应验先知依撒意亚所说的话：}等等，直到以下经文：\Quote{外邦人将信靠祂的圣名。}只有他记载了这些事。另外两位已转向其他主题。诚然，马尔谷似乎在某种程度上保持了历史顺序：因为他告诉我们，耶稣察觉犹太人对他怀有恶意后，就与门徒退到海边，然后大批群众聚集到祂那里，祂治好了他们中的许多人。但同时，并不完全清楚他是在哪一点开始转入一个新主题，与严格接续的叙述不同。他并未明确说明这种过渡是否发生在讲述群众聚集在祂周围之处（因为如果当时是这种情况，那在其他时候也可能同样发生），或者是在他说\Quote{祂上山去}之处。路加似乎也注意到后一种情况，他说：\Quote{那时，耶稣出去，上山祈祷。}因为\Quote{那时}一词足够清楚地表明，所述事件并非紧接上文发生。\par

\SourceRecord{Translated by S.D.F. Salmond.；Nicene and Post-Nicene Fathers, First Series；Vol. 6.；Edited by Philip Schaff.；Buffalo, NY: Christian Literature Publishing Co.,；1888.；<http://www.newadvent.org/fathers/1602236.htm>.}

% structure: p0001 source=h1 target=section reason=page-title

\section{\WorkTitle{福音的和谐，第二卷}}

% structure: p0002 source=h2 target=subsection reason=relative-heading-rank; base=h2

\subsection{第三十七章 论\Person{玛窦}和\Person{路加}关于那个附魔的哑盲人的记述的一致性}

84. \Person{玛窦}接着按以下方式继续他的叙述：\Quote{那时，有人将一个附魔的、又瞎又哑的人带到祂面前，祂治好了他，以致那哑巴既能说话，又能看见。} \Person{路加}引入这个叙述，不是按相同的顺序，而是在许多其他事情之后。他也只提到那个人是哑的，而没有提他也是瞎的。但不能仅仅因为他在记述的某一部分或另一部分保持沉默，就推断他指的是一个完全不同的人。因为他也记录了随后发生的事（即在那次治愈之后），正如在\Person{玛窦}的记载中一样。\par

\SourceRecord{Translated by S.D.F. Salmond.；Nicene and Post-Nicene Fathers, First Series；Vol. 6.；Edited by Philip Schaff.；Buffalo, NY: Christian Literature Publishing Co.,；1888.；<http://www.newadvent.org/fathers/1602237.htm>.}

% structure: p0001 source=h1 target=section reason=page-title

\section{福音的和谐 卷二}

% structure: p0002 source=h2 target=subsection reason=relative-heading-rank; base=h2

\subsection{第三十八章  关于有人对他说他借贝耳则步驱魔的场合，以及由此情境引出的关于亵渎圣神和两棵树的言论；并讨论这些段落中玛窦与其他两位圣史，特别是玛窦与路加之间是否存在某些不一致之处。}

85. 玛窦继续叙述如下：\BibleQuote{众人都惊奇说：\InnerQuote{这不是达味之子吗？}法利塞人听见了，却说：\InnerQuote{这人驱魔，无非是仗赖贝耳则步，众魔之首领。}耶稣知道了他们的心意，便对他们说：\InnerQuote{凡是一国自相纷争，必成荒芜；}}等等，直到这句话：\BibleQuote{凭你的话，你将被成义；凭你的话，你将被定罪。}马尔谷并没有在哑巴故事之后紧接着插入对耶稣的这项指控，即祂借贝耳则步驱魔；而是在他自己所记录的其他某些事之后，才引入这一事件，或者是因为他在不同的联系中回忆起来而附加在那里，或者是因为他最初在历史中省略了一些内容，在注意到这些之后，又接续了这个叙述顺序。另一方面，路加对这些事情的记述几乎与玛窦所用的措辞相同。而路加在此处将圣神称为天主的手指，这一点并不表示对真正含义的一致认同有任何偏离；相反，它给了我们一个额外的教训，使我们知道无论在圣经何处出现\Quote{天主的手指}这一短语，应如何解释。此外，至于马尔谷和路加在本段中都未提及的其他事项，这些不会造成困难。某些由他们以略有不同的措辞叙述的其他情况也不会造成困难，因为含义仍然相同。\par

\SourceRecord{Translated by S.D.F. Salmond.；Nicene and Post-Nicene Fathers, First Series；Vol. 6.；Edited by Philip Schaff.；Buffalo, NY: Christian Literature Publishing Co.,；1888.；<http://www.newadvent.org/fathers/1602238.htm>.}

% structure: p0001 source=h1 target=section reason=page-title

\section{福音的和谐，第二卷}

% structure: p0002 source=h2 target=subsection reason=relative-heading-rank; base=h2

\subsection{第三十九章 论玛窦与路加在记载主答复某些求征兆的人时所说关于约纳先知、尼尼微人、南方女王以及邪魔离去后回来发现房子打扫干净的话语时如何一致的疑问}

86. 玛窦继续叙述如下：\BibleQuote{那时，有几个经师和法利塞人回答说：老师，我们愿意看到您的一个征兆；}一直到我们读到：\BibleQuote{同样，这邪恶的世代也要这样。}这些话路加也在此处记载，虽然顺序稍有不同。因为他在叙述中较早提到他们寻求主从天上来的一个征兆，使之紧接在他版本中关于医治哑巴的奇迹之后。但他在那里没有记载主给他们的答复。然而后来，在众人聚集之后，他说明这个答复是给那些人的，根据他的理解，那些人在之前的经文中被提到寻求从主那里来的征兆。他随后也附加了那个答复，只是在插入关于那妇人对主说：\BibleQuote{怀过你的子宫是有福的。}这段话之后。这个妇人的话，他又在讲述了主关于邪魔从人身上出去后又回来发现房子打扫干净的言论之后插入。这样，在妇人的话之后，在他陈述了关于他们寻求从主那里来的征兆的答复之后，他引入了约纳先知的比喻；然后，直接继续主的言论，他又列举了关于南方女王和尼尼微人的话。因此，他更多地是叙述了玛窦略过未提的内容，而不是省略了那位福音作者在此处叙述的任何事实。此外，谁能看不到关于这些话语由主说出的确切顺序的问题是多馀的呢？因为我们也应当从福音作者无可指责的权威中学到这个教训——即，不能认为作者违背了真理，即使他没有按照说话者可能说出的确切顺序重现某些演说者的言论；事实上，顺序的细节，无论是这样或那样，并不影响主题本身。路加在此版本中表明，主的这篇言论比我们可能认为的要长；他记录了一些讨论的主题，这些主题类似于玛窦在记载山上讲道时提到的。所以，我们认为这些话被说了两次，即，在之前的场合和这一次。但在这篇言论结束时，路加转向另一个主题，关于他是否按照实际发生的顺序记载这一点是不确定的。因为他这样连接：\BibleQuote{正说话的时候，有一个法利塞人请祂同他吃饭。}但他没有说\Quote{正说这些话的时候}，而只是说\Quote{正说话的时候}。因为如果他说了\Quote{正说这些话的时候}，那么这种表达当然会迫使我们认为所提及的事件，除了按此顺序记载外，也是按此顺序发生在主身上的。\par

\SourceRecord{Translated by S.D.F. Salmond.；Nicene and Post-Nicene Fathers, First Series；Vol. 6.；Edited by Philip Schaff.；Buffalo, NY: Christian Literature Publishing Co.,；1888.；<http://www.newadvent.org/fathers/1602239.htm>.}

% structure: p0001 source=h1 target=section reason=page-title

\section{\WorkTitle{福音合参，卷二}}

% structure: p0002 source=h2 target=subsection reason=relative-heading-rank; base=h2

\subsection{\Emph{第四十章　论有关有人通知祂母亲和兄弟来访一事，玛窦的记述与马尔谷和路加的顺序是否矛盾}}

87. 玛窦接着以这样的话继续叙述：\BibleQuote{耶稣还对众人说话的时候，不料祂的母亲和祂的兄弟站在外边，想要与祂说话。}直到\BibleQuote{凡遵行我圣父旨意的，就是我的兄弟、姊妹和母亲。}毫无疑问，我们应当明白这事是紧接在前述事件之后发生的。因为他用\Quote{耶稣还对众人说话的时候}这句话转入这段叙述；而\Quote{还}这个字指的是什么，不正是当时祂正在谈论的那件事吗？因为经文不是说：\Quote{当祂对众人说话时，不料祂的母亲和兄弟来了，}而是\Quote{祂还说话的时候，}等等。这种措辞迫使我们认为，这事正好发生在祂仍在谈论刚才所提及的那些事的时候。因为马尔谷在记述主有关亵渎圣神的言论之后，也叙述了主所说的话。他是这样记载的：\BibleQuote{耶稣的母亲和兄弟来了，}省略了与主那篇言论相关的某些内容——这些内容玛窦比马尔谷更详细地引入，而路加又比玛窦更详细。另一方面，路加在报告这一事件时并未保持历史顺序，而是提前叙述，并凭记忆在事件实际发生之前的一个时间点加以记述。而且，他将其引入的方式，使其看起来与前后文都没有紧密联系。因为他在报告了主的某些比喻之后，这样引入了有关祂母亲和兄弟的记载：\BibleQuote{那时，耶稣的母亲和兄弟来了，因为人多，不能到祂跟前。}因此，他没有说明他们是在什么具体时间来到祂面前的。再者，当他离开这个主题时，他接着这样说：\BibleQuote{有一天，耶稣上了船，同门徒们一起。}当然，当他使用\Quote{有一天}这个说法时，他足够清楚地表明，我们不必认为所说的那一天就是这事发生的那一天，或是紧接着的一天。因此，无论是主的言语方面，还是所叙述事件的历史顺序方面，玛窦关于主母亲和兄弟所发生之事的记载，与其他两位福音作者的版本并无任何不协调之处。\par

\SourceRecord{Translated by S.D.F. Salmond.；Nicene and Post-Nicene Fathers, First Series；Vol. 6.；Edited by Philip Schaff.；Buffalo, NY: Christian Literature Publishing Co.,；1888.；<http://www.newadvent.org/fathers/1602240.htm>.}

% structure: p0001 source=h1 target=section reason=page-title

\section{福音合参 卷二}

% structure: p0002 source=h2 target=subsection reason=relative-heading-rank; base=h2

\subsection{第41章 关于在船上所说的有关撒种者的话，其种子撒下后，部分落在路旁等；关于那人在麦子之外被人撒了莠子；关于芥菜子和酵母；以及关于他在屋里所说的藏在地里的宝贝、珍珠和撒在海里的网，以及那从宝库里拿出新和旧之物的家主；以及证明玛窦与马尔谷和路加在共同记述之事上和叙述顺序上和谐一致的方法。}

88. 玛窦接着这样说：\BibleQuote{在那一天，耶稣从家里出来，坐在海边。有许多群众聚集在他跟前，他只得上船坐下，群众都站在岸上。他用比喻向他们讲论了许多事，说：}等等，直到这句话：\BibleQuote{因此，凡在天国受教的经师，就像一家之主，从他的宝库里拿出新的和旧的东西。}这段叙述所记载的事紧接在关于主母亲和弟兄的事件之后，而且玛窦在其记述中也保留了这一历史顺序，这一点由以下情况表明：他在从一个话题转到另一个话题时，用这样的表达方式表述了联系：\BibleQuote{在那一天，耶稣从家里出来，坐在海边；有许多群众聚集在他跟前。}因为他使用了\Quote{在那一天}这个短语（除非\Quote{日子}一词按照圣经的习惯用法可能仅表示\Quote{时间}），他足够清楚地暗示要么现在所记述的事紧接在前面之事后发生，要么至少不可能相隔很远。这一推断被马尔谷保持相同顺序的事实所证实。另一方面，路加在记述了主母亲和弟兄所发生的事之后，转入另一个话题。但同时，在作此转换时，他并没有建立任何看起来与这一顺序不一致的关联。因此，在所有那些马尔谷和路加与玛窦共同记述的主所说的话的段落中，他们彼此之间的和谐是毋庸置疑的。此外，仅由玛窦提供的部分更是远在争论范围之外。至于叙述顺序，虽然不同圣史呈现略有不同，根据他们各自沿着历史顺序或记忆顺序进行，我看不出有什么理由可以断言任何两位作者之间存在陈述差异或矛盾。\par

\SourceRecord{Translated by S.D.F. Salmond.；Nicene and Post-Nicene Fathers, First Series；Vol. 6.；Edited by Philip Schaff.；Buffalo, NY: Christian Literature Publishing Co.,；1888.；<http://www.newadvent.org/fathers/1602241.htm>.}

% structure: p0001 source=h1 target=section reason=page-title

\section{\WorkTitle{福音的和谐，卷二}}

% structure: p0002 source=h2 target=subsection reason=relative-heading-rank; base=h2

\subsection{第四十二章　论主来到自己的家乡，及民众因轻视祂的出身而惊奇于祂的教训；论玛窦在此节与马尔谷和路加的和谐；尤其论及这些福音作者中的第一位所呈现的叙述顺序是否与其他两位有些许不一致的问题。}

89. 玛窦由此继续如下：\BibleQuote{当耶稣讲完这些比喻，就从那里离开，来到自己的家乡，在会堂里教训他们；}等等，一直到\BibleQuote{耶稣因为他们不信，就在那里不多行异能。}这样，他从上述包含比喻的讲论转到这段经文，其方式并不使我们绝对有必要认为一段紧接另一段的历史顺序。当我们看到马尔谷如何从这些比喻转到另一个主题——该主题与玛窦直接继接的主题不同，而是与之相当不同，并与路加所引入的更相符；以及他如何构建他的叙述，其方式使可信度倾向于这样一个假设，即在历史顺序中紧跟其后的应是这两位后福音作者都紧接[比喻]插入的事件——即耶稣在船中睡觉的遭遇，以及在\Place{革辣撒}地方驱逐魔鬼的奇迹——这两件事玛窦早已在其记载的更早阶段回忆并引入。因此，目前我们必须考虑[玛窦所记载的]主在自己家乡所说的话，以及别人对祂说的话，是否与另外两位即马尔谷和路加所记述的相符。因为，若望在其记载中在非常不同且不相似的段落中提到了向主说的或主说的话，这些与这三位福音作者在本段所记载的相似。\par

90. 如今，马尔谷叙述这段经文时，其措辞几乎与玛窦所载完全相同；惟一的区别是，他说主被同乡称为\Quote{木匠，玛利亚的儿子}，而非如玛窦所记的\Quote{木匠的儿子}。这一点毫不奇怪，因为主完全可以用这两个名称来称呼。当他们以为他是木匠的儿子时，自然也认为他是木匠。另一方面，路加则更广泛地记述了同一事件，并记录了当时发生的其他许多事情。他将这段叙述放在主受洗和受试探之后不久，无疑是以预叙的方式引入了实际上是在许多中间事件发生之后才发生的事。因此，每个人都可以从中看出一个原则，这个原则对于解决我们赖天主的帮助所承担的这项关于福音书和谐的最重要问题具有重大意义，即：这些作者并非因无知而有所省略，也不是由于单纯不知道事件的实际历史顺序而（有时）宁愿按照事件回忆起的顺序来安排。这个原则的正确性可以从以下事实中清楚地看出：路加在主于葛法翁行任何事之前，提前了实际日期，插入了我们目前正在考察的这段内容，其中记载了主的同乡如何既惊讶于祂身上权威的力量，又对祂家庭微贱表示轻蔑。因为他告诉我们，主对他们说了这些话：\Quote{你们必定要对我说：医生，医治你自己吧！在葛法翁所行的一切，也当在你的家乡这里行。}而就同一路加的叙述而言，我们尚未读到祂曾在葛法翁行过什么事。此外，因为这不费太多时间，而且这样做既简单又非常必要，我们在此插入完整的上下文，显示作者从哪个主题以及以何种方式得出这段内容。在他叙述了主的受洗和受试探之后，他接着这样说：\Quote{魔鬼用尽了各种试探后，就离开了祂，再等时机。耶稣因圣神的能力回到加里肋亚，祂的名声传遍了周围各地。祂在他们的会堂施教，受到众人的称扬。祂来到纳匝肋，就是祂长大的地方；照祂的习惯，在安息日祂进了会堂，站起来要诵读。有人把\BibleBook{依撒意亚}{先知书}递给祂，祂展开书卷，找到一处写着：\InnerQuote{上主的神临于我身，因为祂给我傅了油。祂派遣我向穷人传报福音，向俘虏宣布解放，向盲者宣告复明，使受压迫者获得自由，宣布上主恩慈之年和报复的日子。}祂把书卷卷起，交给侍役，就坐下了。会堂内众人的眼睛都注视着祂。祂便对他们说：\InnerQuote{你们今天听见的这段圣经，已经应验了。}众人都称赞祂，并惊奇祂口中所出的恩宠之言。他们说：\InnerQuote{这不是若瑟的儿子吗？}祂对他们说：\InnerQuote{你们必定要对我说这句俗语：医生，医治你自己吧！我们听说你在葛法翁所行的，也当在你的家乡这里行。}}随后他继续叙述其余部分，直到这段叙述全部完毕。因此，还有什么比这更明显的呢？他明知故犯地将这段内容提前到历史日期之前，因为他无疑知道并提及了主在此之前在葛法翁所行的某些大能之事，同时他也意识到自己尚未详细叙述这些事。事实上，他从叙述主的受洗开始，历史进展到此时，他不可能忘记直到此刻他还没有提到任何发生在葛法翁的事；实际上，他刚刚在受洗之后才开始给我们叙述关于主本人的事迹。\par

\SourceRecord{Translated by S.D.F. Salmond.；Nicene and Post-Nicene Fathers, First Series；Vol. 6.；Edited by Philip Schaff.；Buffalo, NY: Christian Literature Publishing Co.,；1888.；<http://www.newadvent.org/fathers/1602242.htm>.}

% structure: p0001 source=h1 target=section reason=page-title

\section{福音合参，卷二}

% structure: p0002 source=h2 target=subsection reason=relative-heading-rank; base=h2

\subsection{第四十三章——论玛窦、马尔谷及路加记载黑落德听说主奇妙作为时所说的话彼此一致，以及他们在叙述顺序上的和谐}

91. 玛窦接着记载：\Quote{那时，\OriginalTerm{分封侯}{tetrarch}黑落德听见\OriginalTerm{耶稣}{Jesus}的名声，就对他的臣仆说：\InnerQuote{这是\OriginalTerm{洗者若翰}{John the Baptist}，他从死者中复活了，因此这些异能在他身上彰显出来。}} 马尔谷也记载了同一段，方式相同，但顺序不同。因为他在记述\OriginalTerm{主}{the Lord}派遣门徒出去，嘱咐他们除了手杖外，路上什么也不要带，并结束了他所记录的那次讲论之后，才添加了这段。然而，他并非以这样的方式连接，使我们不得不认为所叙述的事确实紧接着历史中前面的内容。实际上，玛窦与他是一致的。因为玛窦的表达是\Quote{那时}，而不是\Quote{那天}或\Quote{那一刻}。只是他们之间有这样的差异：马尔谷没有指出那话是黑落德本人说的，而是百姓说的，他的陈述是：\Quote{他们说\OriginalTerm{洗者若翰}{John the Baptist}从死者中复活了}；而玛窦则让黑落德亲自说话，措辞是：\Quote{他对他的臣仆说：\InnerQuote{这是洗者若翰，他从死者中复活了，因此这些异能在他身上彰显出来。}} 路加再次保持与马尔谷相同的叙述顺序，并且也像马尔谷一样，并非以迫使我们相信他的顺序必须是实际发生的顺序的方式引入，他以下列措辞呈现同一段：\Quote{\OriginalTerm{分封侯}{tetrarch}黑落德听见主所做的一切事，就困惑起来，因为有人说是若翰从死者中复活了；又有人说是\OriginalTerm{厄里亚}{Elias}显现了；还有人说是古时的一位先知复活了。黑落德说：\InnerQuote{若翰我已经斩首了，但这个人是谁，我竟听见这样的事？}他就渴望见到祂。} 在这些话中，路加也证实了马尔谷的陈述，至少就关于不是黑落德本人，而是其他人说若翰从死者中复活这一肯定而言是如此。但就他提到黑落德困惑，以及随后引述这位王子的话：\Quote{\InnerQuote{若翰我已经斩首了，但这个人是谁，我竟听见这样的事？}}我们或者必须理解为：在那个困惑之后，他内心确信了别人所说的是真的，于是对他的臣仆说话，正如玛窦的版本所说：\Quote{他对他的臣仆说：\InnerQuote{这是洗者若翰，他从死者中复活了，因此这些异能在他身上彰显出来。}}或者必须假定这些话说出来时表现出他仍然处于困惑状态。因为假如他说：\Quote{这可能是洗者若翰吗？}或\Quote{难道这是洗者若翰？}那么就不必谈论一种表明他怀疑和困惑的表达方式了。但是，既然这些表达形式不在我们面前，他的话可以被理解为两种方式之一：我们可以假设他被别人的话说服了，因此确实相信若翰重现而说出了那些话；或者我们可以想象他仍然处于路加所提及的那种犹豫状态。我们的解释得到了这样一个事实的支持：马尔谷已经告诉我们，关于若翰从死者中复活的说法是由别人提出的，他并没有不让我们知道最后黑落德自己也说这样的话：\Quote{是我斩首的若翰，他从死者中复活了。}因为这些话也可以被理解为两种方式之一——即，要么是确认事实的宣告，要么是怀疑者的宣告。此外，虽然路加在记载了这一事件后转入了新话题，但另外两位，玛窦和马尔谷，则借机在此处告诉我们若翰是如何被黑落德处死的。\par

\SourceRecord{Translated by S.D.F. Salmond.；Nicene and Post-Nicene Fathers, First Series；Vol. 6.；Edited by Philip Schaff.；Buffalo, NY: Christian Literature Publishing Co.,；1888.；<http://www.newadvent.org/fathers/1602243.htm>.}

% structure: p0001 source=h1 target=section reason=page-title

\section{\WorkTitle{福音的和谐，第二卷}}

% structure: p0002 source=h2 target=subsection reason=relative-heading-rank; base=h2

\subsection{第44章 关于这三部圣史记载的若翰被监禁和死亡之顺序。}

92. 玛窦接着叙述如下：\BibleQuote{那时，黑落德为了他兄弟斐理伯的妻子黑落狄雅的缘故，逮捕了若翰，把他捆起来，关进监狱；}直至\BibleQuote{若翰的门徒前来，领了尸体，把它埋葬了，然后去报告给耶稣。}马尔谷以类似的方式叙述了这段事迹。路加则以不同的顺序叙述，而是在他陈述主受洗的圣洗时插入的。因此我们应理解他是预先叙述此事，并在此处录下了一件实际上发生在相当晚之后的事件。因为他先报告了若翰论主所说的话——即\BibleQuote{祂手里拿着簸箕，要扬净祂的禾场，把麦子收进祂的仓里，把糠秕用不灭的火烧掉；}然后立即附上了一件圣史若望表明并非按历史顺序发生的事件。因为这位圣史提到，耶稣受洗后，在把水变成酒的时候去了加里肋亚；在葛法翁逗留了几天后，离开该地回到犹太地区，在约但河附近给许多人施洗，那时若翰尚未被监禁。如今，哪个读者，若不曾更精通这些著作，不会理解为这里暗示的是在讲了关于簸箕和扬净禾场的话之后，黑落德才向若翰发怒并把他投进监狱？然而，这里提到的事件实际上并非按这里记载的顺序发生，我们已在别处证明；事实上，路加自己也把证据交在我们手里。因为如果[他意指]若翰的监禁是在那些话之后立即发生的，那么路加自己的叙述中耶稣的受洗是在他提到若翰被监禁之后才引入的，我们该如何理解呢？因此，显然他是想起这件事与当前场合相关，便预先在此叙述，并把它插入历史中，置于许多他有意给我们留下一些记录的事件之前，而这些事件在时间上实在于若翰所遭的这次不幸之前。然而，另外两位圣史玛窦和马尔谷也没有把若翰被监禁的事实放在他们叙述中——从他们自己的著作也可见——实际事件顺序所应有的位置。因为他们也告诉我们，主是在若翰被投入监狱后才去了加里肋亚；然后，在记录了许多主在加里肋亚所做的事之后，他们才提到黑落德关于被他斩首的若翰从死者中复活的议论或疑惑；并在后一场合中，向我们讲述了有关若翰被监禁和死亡的全部经过。\par

\SourceRecord{Translated by S.D.F. Salmond.；Nicene and Post-Nicene Fathers, First Series；Vol. 6.；Edited by Philip Schaff.；Buffalo, NY: Christian Literature Publishing Co.,；1888.；<http://www.newadvent.org/fathers/1602244.htm>.}

% structure: p0001 source=h1 target=section reason=page-title

\section{\WorkTitle{福音合参，卷二}}

% structure: p0002 source=h2 target=subsection reason=relative-heading-rank; base=h2

\subsection{第四十五章 论四位圣史叙述五饼奇蹟的次序与方法}

93. 玛窦在述及若翰死讯传到耶稣处后，继续其记述，以下列方式引入：\Quote{耶稣一听说这事，就从那里乘船退到荒野地方独自一处。民众听见了，就从各城步行跟随祂。祂出来，看见一大伙群众，就对他们动了怜悯的心，治好了他们的病人。}\BibleRef{玛 14:13-14} 因此，他提到这事发生在若翰受难后。所以，此前所记载的那些事——即令黑落德恐慌并促使他说\Quote{若翰我已经斩首了}\BibleRef{玛 14:10}的情况——是发生在之后。因为必须理解，这些事件是后来发生的，消息传到黑落德耳中，使他忧虑，困惑于他所听到的这位行如此奇事的人究竟是谁，而他自己却已将若翰处死。玛尔谷在叙述若翰受难后，提到被派遣的门徒们回到耶稣那里，将所行所教的一切报告给祂；主（这是玛尔谷独载的）吩咐他们到荒野地方去休息片刻，于是祂同他们上船离开；民众察觉后，徒步先到那地方；主怜悯他们，教导他们许多事；时辰已晚时，发生了用五个饼和两条鱼使众人吃饱的事。这奇蹟被四位圣史所记载。同样，路加在较早阶段记述了若翰之死，与我们所谈的场合相关，在此上下文中首先叙述黑落德对主是谁的困惑，接着立即附上与玛尔谷一致的陈述——即宗徒们回到祂那里，报告所行的一切；然后祂带他们退到荒野地方，群众跟随祂到那里，祂向他们谈论天主的国，并治愈了需要医治的人。随后他也提到，当日暮时，行了五饼奇蹟。\par

94. 然而若望与前三圣史大不相同，他更多记述主所讲的言论而非奇妙作为。在记载主离开犹太第二次往加里肋亚去之后——这离开被认为发生在其他圣史所说的若翰被监禁时主往加里肋亚的那个时刻——在记载这些之后，我说，若望紧接着在叙述中插入主经过撒玛黎雅时与井旁撒玛黎雅妇人相遇所说的长篇言论；然后他说两天后祂离开那里往加里肋亚去，随后到了加里肋亚的加纳，祂曾在那里变水为酒，并在那里治好了一位官员的儿子。至于其他圣史所记载的主在加里肋亚的言行，若望保持沉默。然而，他也提到其他圣史忽略的一件事——即祂在庆节日上耶路撒冷，在那里治好了患病三十八年的人，这人无人帮助他下到池中，池水能治愈各种疾病的人。与此同时，若望还叙述了祂在那场合所说许多话。他进一步告诉我们，这些事后祂渡过加里肋亚海，即提庇黎雅海，一大群人跟随祂；然后祂上山，与门徒们同坐——那时犹太人的逾越节近了；接着祂举目看见一大群人，就用五个饼和两条鱼喂饱了他们；其他圣史也给了我们这记载。这确证祂略过了那些事件，而其他圣史正是沿着这些事件引入这奇蹟的叙述。尽管看来采用了不同的叙述方法，前三圣史忽略了第四圣史所记载的某些事，但我们看到那三位一直大致保持相同进路的圣史，在这五饼奇蹟处直接相遇；而第四位圣史，他特别精通主言论的深奥教导，在叙述其他一些事（其他圣史保持沉默）时，以一种方法迂回在他们的轨迹上，并且正要短暂地离开他们的路径再次飞升至更崇高主题的领域时，与他们相遇于呈现这五饼奇蹟的叙述，这奇蹟是他们共有的。\par

\SourceRecord{Translated by S.D.F. Salmond.；Nicene and Post-Nicene Fathers, First Series；Vol. 6.；Edited by Philip Schaff.；Buffalo, NY: Christian Literature Publishing Co.,；1888.；<http://www.newadvent.org/fathers/1602245.htm>.}

% structure: p0001 source=h1 target=section reason=page-title

\section{\WorkTitle{福音的和谐 第二卷}}

% structure: p0002 source=h2 target=subsection reason=relative-heading-rank; base=h2

\subsection{第四十六章 关于四位福音作者如何在五饼奇迹这一主题上彼此谐和的问题}

95. \Person{玛窦}接着继续叙述，在适当顺序中连接到上述五饼事迹如下：\BibleQuote{到了晚上，祂的门徒来到祂跟前说，这是荒野的地方，时间已经过了；请遣散群众，叫他们往村庄里去，为自己买食物。但耶稣对他们说，他们不必去；你们给他们吃吧；} 等等，直到我们读到：\BibleQuote{吃的人数除了妇女孩子，约有五千。}因此，这个奇迹被所有四位福音作者记载，他们被认为彼此之间暴露出某些差异，必须加以审视和讨论，以便我们也从这一实例中学会一些规则，这些规则将适用于所有其他类似情况，作为规范陈述方式的原则，在这些陈述方式中，无论多么不同，仍然保留相同的含义，并在表达事实时保持相同的准确性。事实上，这项调查不应从\Person{玛窦}开始，尽管这符合福音作者的顺序，而应从\Person{若望}开始，由他极其详细地叙述了这个问题，甚至记录主与之谈论此事件的门徒的名字。因为他以下列方式叙述历史：\BibleQuote{那时耶稣举目观看，见有一大群人来到祂那里，就对\Person{斐理伯}说：我们从哪里买饼给这些人吃呢？祂说这话是要试验他，因为祂自己原知道要怎样做。\Person{斐理伯}回答说：就是二百块银币的饼，也不够他们每人拿一小块。有一个门徒，就是\Person{西满伯多禄}的兄弟\Person{安德肋}，对祂说：这里有一个孩童，带着五个大麦饼和两条鱼；但这么多人中这算什么呢？耶稣说：你们叫众人坐下。那地方有许多青草。于是人坐下了，数目约有五千。耶稣拿起饼来，祝谢了，就分给门徒，门徒又分给坐席的人；鱼也是如此，任意给他们吃。他们吃饱了，耶稣对门徒说：把剩下的零碎收拾起来，免得糟蹋了。他们就把零碎收拾起来，装满了十二篮子，是那五个大麦饼的零碎，就是众人吃了以后剩下的。}\par

96. 我们在此讨论的研究，并非针对这位福音作者所陈述的一个说法，即他具体说明了饼的种类；因为他没有忽略其他作者所忽略的事，即这些饼是大麦饼。这个问题也不涉及他未提及的事——即除了五千男人外，还有妇女和儿童，正如玛窦告诉我们的。现在，在所有这类研究中，应当确立并经常牢记一个原则：没有人会因为某件未被一位作者记载、却被另一位作者叙述的事而感到困扰。但这里的问题是，这些作者各自叙述的几件事如何能被证明全都真实，以至于其中一位在给出自己特有的版本时，不会使另一位作者的记述失效。因为，若按若望的记述，主看见眼前的群众，为试探他，便问斐理伯从哪里可得饼给他们吃，那么关于其他作者所说的话——即门徒们首先对主说，请打发群众离开，使他们可以到邻近的地方去买食物，而主回答他们说（按玛窦）：\Quote{他们不必去，你们给他们吃的吧}——的真实性就可能引起疑问。马尔谷和路加也同意这一点，只是他们省略了\Quote{他们不必去}这句话。因此，我们应当设想，在这些话之后，主看着群众，并按若望所记而其他作者省略的话对斐理伯说了。然后，按若望所记，斐理伯的答复，马尔谷却说是门徒们给出的——意思是，我们应当理解为斐理伯是代表其他人作出这个回答；尽管他们也可能按照非常常见的用法，用复数代替了单数。这里实际归于斐理伯的话——即\Quote{二百银钱的饼也不够他们每人吃一小块}——对应于马尔谷的版本：\Quote{我们去买二百银钱的饼给他们吃吗？}而同一位马尔谷所记的主说的话：\Quote{你们有多少饼？}却被其他人忽略了。另一方面，若望所记的安德肋关于五个饼和两条鱼的提议，在其他作者那里出现时，他们用复数代替了单数，作为将提议归于整个门徒群体的说明。事实上，路加将斐理伯的答复和安德肋的答案合并在一句话中。因为当他说：\Quote{我们不过有五个饼和两条鱼}时，他报告了安德肋的回答；但当他加上：\Quote{除非我们去为这所有人买食物}时，他似乎又回到了斐理伯的答复，只是他忽略了\Quote{二百银钱}。同时，那句关于去买食物的话也可以理解为包含在安德肋自己的话中。因为他在说：\Quote{这里有一个孩子，他有五个大麦饼和两条鱼}之后，又补充说：\Quote{但为这么多的人，这些算什么呢？}而这最后一句实际上与所讨论的那句话意思相同，即\Quote{除非我们去为这所有人买食物}。\par

97. 从与事实和思想有关的真实和谐中所发现的这种种不同的陈述，足以清楚地教导我们一个有益的教训：我们研究一个人的言语时，所应关注的不是别的，而是说话者的意图；所有真实的叙述者在记录任何事物时，无论是关于人、天使还是天主，都应当竭尽全力表达这种意图。因为说话者的思想和意图可以用言语来表达，这些言语应不留下任何与事实不符的表面矛盾。\par

98. 在此，我们确实不应忽略引起读者注意某些其他可能与已讨论之事类似的问题。其中之一是路加说他们被吩咐按五十人一组坐下，而马尔谷的版本是百人一组和五十人一组。然而，这个差异并不造成真正的困难。事实上，一个人只报告了部分，另一个人给出了全部。因为引入了百人和五十人的那位福音作者恰好提到了一些另一个人未提及的事。但这并不造成他们之间的矛盾。如果一个人只提到五十人，另一个人只提到百人，他们确实可能显得相互对立，并且不容易证明两个指令都被实际说出了，尽管前者只具体说了前者，后者只说了后者。然而，即使在这种情况下，谁不会承认，当事情受到更仔细的考虑时，应当能找到解决办法呢？我举此例的原因在于，这类事情常常出现，它们实际上并不包含矛盾，但在那些不够注意它们并草率下判断的人看来却似乎有矛盾。\par

\SourceRecord{Translated by S.D.F. Salmond.；Nicene and Post-Nicene Fathers, First Series；Vol. 6.；Edited by Philip Schaff.；Buffalo, NY: Christian Literature Publishing Co.,；1888.；<http://www.newadvent.org/fathers/1602246.htm>.}

% structure: p0001 source=h1 target=section reason=page-title

\section{\WorkTitle{福音的和谐，第二卷}}

% structure: p0002 source=h2 target=subsection reason=relative-heading-rank; base=h2

\subsection{第四十七章 论祂在水上行走，以及关于叙述该场景的福音作者之间和谐的问题，及他们如何从记载祂用五个饼饱饫群众的那段过渡。}

99. 玛窦继续叙述如下：\BibleQuote{祂遣散了群众以后，便独自上山祈祷：到了晚上，只有祂一人在那里。那时船已离岸很远，被波浪颠簸，因为风向不顺。夜间四更时分，耶稣踩着海面，向他们走去。门徒看见祂在海面上行走，就惊慌起来，说：\InnerQuote{是个妖怪。}}等等，直到\BibleQuote{他们前来朝拜祂，说：\InnerQuote{你真是天主子。}}同样，马尔谷在叙述五饼奇迹之后，也以如下方式叙述同一事件：\BibleQuote{到了晚上，船在海中，耶稣独自在岸上。祂看见门徒辛苦摇橹，因为风向不顺，}等等。这与玛窦的记述相似，只是没有提到伯多禄在水上行走。但这里我们必须注意，不要以为马尔谷关于主的叙述有什么难解之处，即：当祂在水上行走时，祂本想从他们旁边经过。若非祂实际上在朝与他们不同的方向前进，仿佛有意像对待陌生人一样绕过这些人，而他们又如此不认识祂，以为祂是个妖怪，他们怎能明白这事呢？然而，又有谁如此迟钝，竟看不出这事含有奥秘意义呢？同时，祂也帮助了那些惊慌失措并呼喊的人，对他们说：\BibleQuote{放心，是我；不要怕。}那么，祂为何想要绕过这些人，却又在他们恐惧时如此鼓励他们？无非是那绕行的意图是为了引出那些呼喊，以便去救助他们。\par

100. 此外，若望在这几位身上仍稍作停留。因为在他叙述了五饼奇迹之后，他也给我们讲述了一些关于那艘辛苦行船的事，以及主在水上行走的事。他以下列方式将这段记载与前文连接起来：\BibleQuote{耶稣既知道他们要来强迫祂，立祂为王，就独自又退到山上去了。到了晚上，祂的门徒下到海边；他们上了船，便向葛法翁驶去；天已黑了，耶稣还没有来到他们那里。海上因狂风大作而翻腾，}等等。其中似乎没有什么与其他福音的记载相矛盾，除非是这一情况：玛窦告诉我们，当群众被遣散后，祂上山去，为的是独自祈祷；而若望却说祂是在山上与那批用五个饼饱饫过的群众在一起。但是，既然若望也告诉我们，祂在奇迹之后退到山上，为的是避免群众抓住祂，立祂为王，那么显然，当饼分给群众时，他们已从山上下来到了较平坦的地方。因此，玛窦和若望关于祂再次上山的说法并无矛盾。唯一的区别是，玛窦用了\BibleQuote{祂上去了}，而若望用了\BibleQuote{祂退去了}。只有在退去时没有上山，这两者才会矛盾。同样，玛窦的话是\BibleQuote{祂独自上山去祈祷}，而若望说的是\BibleQuote{祂既知道他们要来立祂为王，就独自又退到山上}，这也不显明缺乏和谐。因为退去之事决不与祈祷之事相矛盾。的确，主以祂自己改变了我们卑微的身体，使它相似祂光荣的身体，以此教导我们：退去之事对我们来说也应是认真祈祷之事。此外，玛窦先告诉我们祂命令门徒上船，先到对岸去，直到祂遣散群众；然后告诉我们，群众被遣散后，祂自己上山独自祈祷；而若望先提到祂独自退到山上，然后接着说：\BibleQuote{到了晚上，祂的门徒下到海边；他们上了船，}等等。谁能看不出，若望在重述事实时，提到门徒后来才实际做的事，而耶稣早在祂上山之前就已命令他们去做？这如同讲话中常见的做法，以某种方式返回到原本可能被忽略的事情上。但由于没有特别注明这种重述，尤其是当它简短而快速时，听众有时会认为后来提到的事件也是在后来实际发生的。这样，福音作者的真实意思是，对于那些被描述为上船并渡海到葛法翁的人，主在他们辛劳于深水时，脚踏海面走向他们；而主的这次来临当然发生在较早的时刻，即他们航行前往葛法翁的行程中。\par

101. 另一方面，路加在记载五饼的奇迹后，转入另一主题，偏离了这一叙事次序。因为他完全没有提及那艘小船和主在水面上行走的事。但在陈述了\Quote{他们全都吃了，并且吃饱了，他们把剩下的碎块收集起来，装满了十二筐}之后，他又添上了以下记载：\Quote{有一次，耶稣独自祈祷，门徒们同祂在一起；祂问他们说：众人说我是谁？}因此，他这样连续叙述了一件新事，而另外三位曾记述主如何在海上行走、并来到航行中的门徒们那里的福音作者，却没有记载此事。然而，不应因此就认为，祂是在玛窦告诉我们祂独自上山祈祷的那座山上，对门徒们说：\Quote{众人说我是谁？}因为路加在这里似乎也与玛窦一致，他的话说：\Quote{当祂独自祈祷时}；而玛窦的话是：\Quote{他独自上山去祈祷}。但必须肯定的是，这个问题是在另一个场合提出的，因为这里既说祂独自祈祷，又说门徒们与祂在一起。因此，路加确实只提到了祂独自一人，但并未说祂没有门徒陪伴，不像玛窦和若望所记载的那样，即门徒们离开祂，先祂渡到海的那边去。因为路加明确地加上了\Quote{门徒们也与祂在一起}的陈述。因此，当他说祂独自一人时，他的意思是指那些没有与祂在一起的群众。\par

\SourceRecord{Translated by S.D.F. Salmond.；Nicene and Post-Nicene Fathers, First Series；Vol. 6.；Edited by Philip Schaff.；Buffalo, NY: Christian Literature Publishing Co.,；1888.；<http://www.newadvent.org/fathers/1602247.htm>.}

% structure: p0001 source=h1 target=section reason=page-title

\section{\WorkTitle{福音的和谐 卷二}}

% structure: p0002 source=h2 target=subsection reason=relative-heading-rank; base=h2

\subsection{第四十八章 论\Person{玛窦}和\Person{马尔谷}一方与\Person{若望}另一方之间并无矛盾，三者共同记载湖对岸后所发生之事}

102. \Person{玛窦}接着记述如下：\BibleQuote{他们过了海，来到革乃撒勒地方。那地方的人一认出是耶稣，就打发人到周围整个地区，把一切患病的人都带到祂跟前，恳求祂让他们只摸摸祂的衣边；凡摸着的，就都痊愈了。那时，有经师和法利塞人从耶路撒冷来到耶稣跟前说：\InnerQuote{祢的门徒为什么违犯长老的传统？因为他们吃饭时不洗手。}}等等，直到这句话：\BibleQuote{但吃饭不洗手，并不污秽人。}这事\Person{马尔谷}也记述了，其方式排除了任何关于矛盾的疑问。因为一方在此所表述的形式若与另一方有别，至少不违背意义的一致。另一方面，\Person{若望}依其习惯，专注于主的言论，从主履海到达的那只船之事，转到他们上岸后发生的事，并提及祂从吃饼这件事出发，讲授了许多道理，且主要涉及神圣事物。在这番言论之后，他的叙事又以高超的笔调，从一主题转向另一主题。同时，他这样转向不同主题，并不真正缺乏和谐，尽管他在这些事上与其他人所呈现的事件顺序有所不同。因为有什么能阻止我们立刻设想：那些\Person{玛窦}和\Person{马尔谷}所记载的人，确实被主治愈了，而\Person{若望}所记述的这篇言论，正是主向那些渡海跟随祂的群众发表的？这一假设更为合理，因为据\Person{若望}所述，他们渡过的地方葛法翁，靠近革乃撒勒湖；而据\Person{玛窦}所述，他们上岸后所进入的地区，正是那一带。\par

\SourceRecord{Translated by S.D.F. Salmond.；Nicene and Post-Nicene Fathers, First Series；Vol. 6.；Edited by Philip Schaff.；Buffalo, NY: Christian Literature Publishing Co.,；1888.；<http://www.newadvent.org/fathers/1602248.htm>.}

% structure: p0001 source=h1 target=section reason=page-title

\section{\WorkTitle{福音的和谐} 第二卷}

% structure: p0002 source=h2 target=subsection reason=relative-heading-rank; base=h2

\subsection{第四十九章 论客纳罕妇人所说的：\Quote{可是小狗也吃主人桌子上掉下来的碎屑}，并论玛窦和路加所记载的和谐。}

103. 玛窦于是继续叙述，在记载了主当着法利塞人关于不洗手的那篇讲论之后。他也保持了随后事件的次序，就各事件之间的过渡所表明的范围而言，他以如下方式将这段记述引入上下文：\BibleQuote{耶稣从那里离开，退往提洛和漆冬一带。看，有一个客纳罕妇人从那地区出来，呼喊说：\InnerQuote{主啊，达味之子，可怜我吧！我的女儿被魔鬼折磨得很苦。}祂却一句话也不回答她。}如此下去，直到这些话：\BibleQuote{妇人，你的信德真大！就照你所愿的给你成就吧。她的女儿就从那时痊愈了。}关于这客纳罕妇人的故事，马尔谷也有记载，他保持了相同的事件次序，并未引起任何关于不和谐的疑问，除非发现在这样的情形中：他告诉我们，当那妇人为她的女儿向主恳求时，主正在屋里。现在我们很容易想到，玛窦只是略过了房屋的提及，却仍然叙述了同一事件。但既然他陈述门徒们以这番话向主提议：\BibleQuote{打发她走吧，因为她在我们后面喊叫。}他似乎明确暗示，妇人是在主前行时跟在后面发出这些哀求的呼声。那么，它怎么可能是\BibleQuote{在屋里}呢？除非我们理解马尔谷暗示了这样的实情：那妇人走进了耶稣当时所在的地方，因为他在叙述的开头提到他在屋里。但是当玛窦说\BibleQuote{祂一句话也不回答她}时，他也让我们明白了两位圣史都没有明确叙述的事——即，在那沉默期间，耶稣走出了屋子。通过这种方式，所有其他细节都被联系在一起，从此以后没有任何不一致的迹象。至于马尔谷所记载的主给她的回答，即拿儿女的饼扔给小狗是不恰当的，那回答是在插入了一些玛窦没有省略的话之后才给出的。也就是说，（我们应当假定）首先有门徒们就那妇人的情况向主提出的请求，以及主给他们的回答，即祂被派遣来，无非是为以色列家迷失的羊；其次是那妇人自己的接近，换句话说，她跟在祂后面，朝拜祂说：\BibleQuote{主啊，帮助我}；此后，在所有这些事情之后，才说出了两位圣史都记载的那些话。\par

\SourceRecord{Translated by S.D.F. Salmond.；Nicene and Post-Nicene Fathers, First Series；Vol. 6.；Edited by Philip Schaff.；Buffalo, NY: Christian Literature Publishing Co.,；1888.；<http://www.newadvent.org/fathers/1602249.htm>.}

% structure: p0001 source=h1 target=section reason=page-title

\section{福音的和谐 第二卷}

% structure: p0002 source=h2 target=subsection reason=relative-heading-rank; base=h2

\subsection{第五十章 论祂以七个饼饱饫群众的场合，以及关于玛窦和马尔谷记载此奇迹时的和谐问题。}

104. 玛窦继续叙述如下：\Quote{那时耶稣从那里离开，来到加里肋亚海边，上了山，坐在那里。有许多群众来到祂跟前，带着瘸子、瞎子、哑巴、残废的，以及许多其他人，把他们放在耶稣脚前，祂就治好了他们；以致群众看见哑巴说话，残废的痊愈，瘸子行走，瞎子看见，都惊奇，就光荣了以色列的天主。于是耶稣叫祂的门徒来说：\InnerQuote{我怜悯这群众，因为他们同我一起已经三天，没有什么可吃，}} 等等，直到这些话：\Quote{吃的人，除了妇女和儿童外，约有四千人。} 这另一次用七个饼和几条小鱼行的奇迹，马尔谷也记载了，而且几乎以同样的顺序；例外的是他在此之前插入了一段其他圣史没有的叙述——即关于一个聋子的叙述，主吐唾沫说：\Quote{\OriginalTerm{开了吧}{Effeta},} 也就是说，开了。\par

105. 关于这个七个饼的奇迹，有必要提请注意一个事实：这两位圣史，玛窦和马尔谷，就是这样将其插入他们的叙述的。因为如果其中一位记载了此奇迹，而同时却没有注意到五个饼的例子，那么他就会被认为与其他圣史相矛盾。因为在这样的情况下，谁不会认为实际上只发生了一个奇迹，而那位被提及的作者，或其他人，或他们所有人，给出了一个不完整和不真实的版本；以至于我们不得不设想：要么那位圣史由于自己的错误而被引导提及七个饼而非五个；要么其他两位圣史，无论是都给出了不正确的陈述，还是因记忆失误而将五说成了七。同样，人们可能会认为十二篮和七篮之间，以及五千人和四千人之间，即被吃饱的人数之间，存在矛盾。但现在，既然记载了七个饼奇迹的圣史也没有忘记提到另一个五个饼的奇迹，那么没有人会感到困难，所有人都能看到这两个工作确实行过。因此，我们举了这个例子，以便如果在其他段落中遇到主的一些类似事迹，而一位圣史的叙述与另一位圣史的版本看起来如此完全矛盾以至于无法找到解决难题的方法时，我们可以理解解释仅仅是：这两件事确实发生了，并且分别由两位不同的作者记载。这正是我们在关于按百人和五十人安排群众就座的事情上已经建议注意的。因为，如果不是同一位圣史记下了这两个数字，我们可能会认为不同作者作了矛盾的陈述。\par

\SourceRecord{Translated by S.D.F. Salmond.；Nicene and Post-Nicene Fathers, First Series；Vol. 6.；Edited by Philip Schaff.；Buffalo, NY: Christian Literature Publishing Co.,；1888.；<http://www.newadvent.org/fathers/1602250.htm>.}

% structure: p0001 source=h1 target=section reason=page-title

\section{\WorkTitle{福音的和谐 第二卷}}

% structure: p0002 source=h2 target=subsection reason=relative-heading-rank; base=h2

\subsection{第五十一章 \Person{玛窦}的记载：离开这些地方后，祂来到\OriginalTerm{玛革丹}{Magedan}境内；以及论他与\Person{玛尔谷}在该提示上的一致，并在关于\OriginalTerm{约纳}{Jonah}的话的记述上（这话是作为对那些求征兆者的答复而再次提出的）。}

106. \Person{玛窦}接着记述：\BibleQuote{祂遣散了群众，就上船，来到\OriginalTerm{玛革丹}{Magedan}境内；}等等，直到这句话：\BibleQuote{一个邪恶和淫乱的世代寻求征兆，除了\OriginalTerm{约纳}{Jonah}先知的征兆外，再没有征兆给它。}这段记载在\Person{玛窦}福音的另一关联中已经出现过。因此我们必须一再坚持这一立场：主在不同场合重复过同样的话；因此，当祂话语的陈述之间出现任何完全不可调和的分歧时，我们应理解为这些话被说了两次。事实上，\Person{玛尔谷}也保持了相同的顺序；在他记述七个饼的奇迹之后，接着给出了与\Person{玛窦}相同的提示，只有一点差异：\Person{玛窦}对地点的表达不是\OriginalTerm{达尔玛奴达}{Dalmanutha}（某些抄本如此读），而是\OriginalTerm{玛革丹}{Magedan}。然而，没有理由怀疑两个名称所指的都是同一地点。因为大多数抄本，即使是\Person{玛尔谷}福音的抄本，也只读到\OriginalTerm{玛革丹}{Magedan}。也不应因这一事实感到困难：\Person{玛尔谷}没有像\Person{玛窦}那样，在主回答那些求征兆者时提及\OriginalTerm{约纳}{Jonah}，而只是简单地说祂这样回答：\BibleQuote{再没有征兆给它。}因为我们可以理解他们所求的是哪种征兆——即从天而来的征兆。而\Person{玛尔谷}只是省略了\Person{玛窦}所记载的关于\OriginalTerm{约纳}{Jonah}的话。\par

\SourceRecord{Translated by S.D.F. Salmond.；Nicene and Post-Nicene Fathers, First Series；Vol. 6.；Edited by Philip Schaff.；Buffalo, NY: Christian Literature Publishing Co.,；1888.；<http://www.newadvent.org/fathers/1602251.htm>.}

% structure: p0001 source=h1 target=section reason=page-title

\section{\WorkTitle{福音的和谐}，第二卷}

% structure: p0002 source=h2 target=subsection reason=relative-heading-rank; base=h2

\subsection{\Emph{第52章}。论玛窦与马尔谷关于法利塞人酵母的陈述，在主题本身与叙述次序上的一致。}

107. 玛窦接着说：\Quote{祂便离开他们走了。门徒渡到对岸，忘了带饼。耶稣对他们说：你们要谨慎，防备法利塞人和撒杜塞人的酵母；}等等，直到我们读到：\Quote{他们这才明白祂所说的不是叫他们防备饼的酵母，而是防备法利塞人和撒杜塞人的教训。}这些话马尔谷也记录了下来，而且次序相同。\par

\SourceRecord{Translated by S.D.F. Salmond.；Nicene and Post-Nicene Fathers, First Series；Vol. 6.；Edited by Philip Schaff.；Buffalo, NY: Christian Literature Publishing Co.,；1888.；<http://www.newadvent.org/fathers/1602252.htm>.}

% structure: p0001 source=h1 target=section reason=page-title

\section{\WorkTitle{福音的和谐 卷二}}

% structure: p0002 source=h2 target=subsection reason=relative-heading-rank; base=h2

\subsection{第五十三章 论祂问门徒众人说祂是谁的缘由，以及关于在内容或顺序上玛窦、马尔谷和路加之间是否有不一致的问题。}

108. 玛窦继续这样说：\BibleQuote{耶稣来到了斐理伯的凯撒勒雅境内，问祂的门徒说：\InnerQuote{人们说人子是谁？}他们说：\InnerQuote{有人说是洗者若翰；有人说是厄里亚；也有人说是耶肋米亚，或先知中的一位。}} 一直到这句话：\BibleQuote{凡你在地上所释放的，在天上也要被释放。} 马尔谷叙述此事几乎以相同的顺序。但他在此之前插入了一段唯独他记载的叙述——即关于那位对主说\BibleQuote{我看见人，好像看见树木在行走}的盲人得见光明的事。路加也记载了此事，将其放在五饼奇迹的记述之后；正如我们前面已经指出的，他所遵循的记忆顺序与另两位所采用的顺序并不冲突。然而，可能会想象到一种困难：路加记载主在独自祈祷、门徒们也同祂在一起时，向门徒提出了这个问题，即人们认为祂是谁；而马尔谷则告诉我们，主是在路上向门徒提出这个问题的。但这只对那些从未在路上祈祷过的人才会成为困难。\par

109. 我回忆已经说过，没有人应该认为伯多禄是在主对他说\BibleQuote{你是伯多禄，在这磐石上我要建立我的教会}时才首次获得这个名字。因为获得这个名字的时间是若望所提及的，当时主对他说：\BibleQuote{你要叫刻法，意即伯多禄。}因此，我们同样不应认为伯多禄是在马尔谷所指的场合得到这个称号，即当他逐一列举十二宗徒的名字时，并告诉我们雅各伯和若望被称为\Quote{雷霆之子}，仅因为在那段经文中他记载了主称他为伯多禄。因为那一情况在那里被提及，仅仅是因为在那一刻它出现在作者的记忆中，而不是因为在那个具体时间实际发生了此事。\par

\SourceRecord{Translated by S.D.F. Salmond.；Nicene and Post-Nicene Fathers, First Series；Vol. 6.；Edited by Philip Schaff.；Buffalo, NY: Christian Literature Publishing Co.,；1888.；<http://www.newadvent.org/fathers/1602253.htm>.}

% structure: p0001 source=h1 target=section reason=page-title

\section{《福音合参》卷二}

% structure: p0002 source=h2 target=subsection reason=relative-heading-rank; base=h2

\subsection{第五十四章 论主预告门徒自己将要受难，以及玛窦、马尔谷和路加关于同一事件的记述之和谐。}

110. 玛窦接着这样记载：\Quote{然后祂嘱咐门徒，不要告诉任何人祂就是耶稣基督。从那时起，耶稣开始向门徒指示，祂必须上耶路撒冷去，受长老、司祭长和经师许多的苦；}并继续下去，直到我们读到：\Quote{你所体会的不是天主的事，而是人的事。}马尔谷和路加也按同一顺序添加了这些段落。只有路加没有提及伯多禄对基督受难所表示的反对。\par

\SourceRecord{Translated by S.D.F. Salmond.；Nicene and Post-Nicene Fathers, First Series；Vol. 6.；Edited by Philip Schaff.；Buffalo, NY: Christian Literature Publishing Co.,；1888.；<http://www.newadvent.org/fathers/1602254.htm>.}

% structure: p0001 source=h1 target=section reason=page-title

\section{\WorkTitle{福音的和谐 第二卷}}

% structure: p0002 source=h2 target=subsection reason=relative-heading-rank; base=h2

\subsection{第五十五章 论三位福音圣史在记载主如何命令愿意跟随祂的人背起十字架跟随祂时彼此和谐一致}

111. 玛窦继续写道：\Quote{于是耶稣对门徒说：\InnerQuote{谁若愿意跟随我，该弃绝自己，背着自己的十字架来跟随我；}}直到这句话为止：\Quote{那时，祂要按照每人的行为予以报酬。}马尔谷也附加了同样的内容，并保持了相同的顺序。但他没有像人子带着祂的天使来临时那样，说祂要按照每人的行为予以报酬。不过，他同时提到了主说了这样的话：\Quote{谁若在这淫乱罪恶的世代中以我和我的话为耻，将来人子在自己父的光荣中与圣天使来临时，也要以他为耻。}这可以视为与玛窦所说的\Quote{祂要按照每人的行为予以报酬}具有相同的含义。路加也以同样的顺序添加了相同的陈述，表达方式虽略有变化，但在忠实传达同一思想的真实性上完全与另外两位吻合。\par

\SourceRecord{Translated by S.D.F. Salmond.；Nicene and Post-Nicene Fathers, First Series；Vol. 6.；Edited by Philip Schaff.；Buffalo, NY: Christian Literature Publishing Co.,；1888.；<http://www.newadvent.org/fathers/1602255.htm>.}

% structure: p0001 source=h1 target=section reason=page-title

\section{\WorkTitle{福音的和谐，第二卷}}

% structure: p0002 source=h2 target=subsection reason=relative-heading-rank; base=h2

\subsection{第56章：主与梅瑟和厄里亚一同在山上向门徒显示自己；以及前三部福音关于此事发生的顺序和情况的和谐问题；尤其关于天数，玛窦和马尔谷说六天后，而路加说八天后。}

112. 玛窦接着记述：\BibleQuote{我实在告诉你们：站在这里的人中，就有些人在未尝到死味以前，必要看见人子来到自己的国内。六天后，耶稣带着伯多禄、雅各伯和他的兄弟若望，领他们上了一座高山；}等等，直到我们读到：\BibleQuote{不要将所见的告诉任何人，除非人子从死者中复活。}\par

113. 至于马尔谷与玛窦一同告诉我们此事发生在六天之后，而路加却说八天之后，那些对此感到困惑的人不应被轻蔑地置之不理，而应通过解释得到启发。因为当我们用\Quote{过多少天}这样的说法来宣布一段时间时，有时我们并不把说话的那一天，或我们事先预告或承诺的事本身宣称要发生的那一天计算在内，而只计算中间的日子，这些日子真正、完全、最终地期满后，所提及的事件才会发生。玛窦和马尔谷正是这样做的。他们省去了耶稣说这些话的那一天，以及祂在山上展示那值得纪念的景象的那一天，而只考虑了中间的日子，因此用了\Quote{六天后}的说法。但路加把两端的日子都计算在内，即第一天和最后一天，因此用了\Quote{八天后}，这是一种以部分代表整体的说法。\par

114. 此外，路加关于梅瑟和厄里亚的陈述，即\BibleQuote{他们离开耶稣的时候，伯多禄对耶稣说：\InnerQuote{老师，我们在这里真好！}}等等，不应被视为与玛窦和马尔谷随后所述的内容相矛盾，仿佛他们使伯多禄在梅瑟和厄里亚仍在与主谈话时就提出了这个建议。因为他们并没有明确说是在那个时候，而只是忽略了路加所添加的事实——即当梅瑟和厄里亚离去时，伯多禄向主提出了搭建三个棚的建议。同时，路加还附加了暗示，说当他们进入云彩时，有声音从天上传来——这一情况未被其他两位肯定，但也没有被否认。\par

\SourceRecord{Translated by S.D.F. Salmond.；Nicene and Post-Nicene Fathers, First Series；Vol. 6.；Edited by Philip Schaff.；Buffalo, NY: Christian Literature Publishing Co.,；1888.；<http://www.newadvent.org/fathers/1602256.htm>.}

% structure: p0001 source=h1 target=section reason=page-title

\section{\WorkTitle{福音的和谐}，第二卷}

% structure: p0002 source=h2 target=subsection reason=relative-heading-rank; base=h2

\subsection{第57章　论\Person{玛窦}与\Person{马尔谷}在论及主对门徒述说\Person{厄里亚}来临之事上的和谐}

115. \Person{玛窦}继续这样写道：\BibleQuote{门徒问耶稣说：\InnerQuote{为什么经师说厄里亚应该先来呢？}耶稣回答说：\InnerQuote{厄里亚的确要来，并要重整一切；但我告诉你们，厄里亚已经来了，人们却不认识他，反而任意待了他；照样，人子也要受他们的磨难。}门徒这才明白耶稣给他们所说的，是指的洗者若翰。} 同样的一段经文，\Person{马尔谷}也有记载，并保持了相同的顺序；虽然他在措辞上有些差异，但他并未偏离对同一含义的真实表述。不过，他并没有加添门徒明白主在说厄里亚已经来时是指\Person{洗者若翰}的这一陈述。\par

\SourceRecord{Translated by S.D.F. Salmond.；Nicene and Post-Nicene Fathers, First Series；Vol. 6.；Edited by Philip Schaff.；Buffalo, NY: Christian Literature Publishing Co.,；1888.；<http://www.newadvent.org/fathers/1602257.htm>.}

% structure: p0001 source=h1 target=section reason=page-title

\section{\WorkTitle{福音合论·卷二}}

% structure: p0002 source=h2 target=subsection reason=relative-heading-rank; base=h2

\subsection{第五十八章　论那人将他的儿子带到主前，门徒却不能治好；并论这三位福音史家在此叙述次序上的一致性问题。}

116. \Person{玛窦}继续写道：\BibleQuote{祂来到群众那里时，有一个人来到祂跟前，跪在祂面前说：\InnerQuote{主，可怜我的儿子！因为他患癫痫病，痛苦不堪。}}等等，直到\BibleQuote{但这一类的，非藉祈祷和禁食，是不能赶出去的。}\Person{马尔谷}和\Person{路加}都记载了这件事，而且次序相同，毫无不一致之处。\par

\SourceRecord{Translated by S.D.F. Salmond.；Nicene and Post-Nicene Fathers, First Series；Vol. 6.；Edited by Philip Schaff.；Buffalo, NY: Christian Literature Publishing Co.,；1888.；<http://www.newadvent.org/fathers/1602258.htm>.}

% structure: p0001 source=h1 target=section reason=page-title

\section{\WorkTitle{福音的和谐，第二卷}}

% structure: p0002 source=h2 target=subsection reason=relative-heading-rank; base=h2

\subsection{第五十九章　论\Person{耶稣}向门徒谈论祂的苦难时，门徒极其忧伤的场合，三位福音作者以相同顺序记载此事。}

117. \Person{玛窦}接着这样记载：\Quote{他们住在\Place{加里肋亚}时，\Person{耶稣}对他们说：\InnerQuote{人子将被交在人手中，他们要杀害祂，第三日祂必要复活。}他们便极其忧伤。}\Person{马尔谷}和\Person{路加}也以相同顺序记载了这段话。\par

\SourceRecord{Translated by S.D.F. Salmond.；Nicene and Post-Nicene Fathers, First Series；Vol. 6.；Edited by Philip Schaff.；Buffalo, NY: Christian Literature Publishing Co.,；1888.；<http://www.newadvent.org/fathers/1602259.htm>.}

% structure: p0001 source=h1 target=section reason=page-title

\section{福音的和谐 卷二}

% structure: p0002 source=h2 target=subsection reason=relative-heading-rank; base=h2

\subsection{第60章 论祂从鱼口中付殿税，这是玛窦独自提及的事件}

118. 玛窦接着这样记述：\Quote{他们来到\Place{葛法翁}，收殿税的人来到\Person{伯多禄}跟前，对他说：你们的师傅不纳殿税吗？他说：是的；}直到我们读到：\Quote{你必找到一块银币，拿去给他们，作为我和你的殿税。}只有他独自记载了这件事，在此事插入之后，他又遵循了与\Person{马尔谷}和\Person{路加}一起遵循的顺序。\par

\SourceRecord{Translated by S.D.F. Salmond.；Nicene and Post-Nicene Fathers, First Series；Vol. 6.；Edited by Philip Schaff.；Buffalo, NY: Christian Literature Publishing Co.,；1888.；<http://www.newadvent.org/fathers/1602260.htm>.}

% structure: p0001 source=h1 target=section reason=page-title

\section{\WorkTitle{福音的和谐，第二卷}}

% structure: p0002 source=h2 target=subsection reason=relative-heading-rank; base=h2

\subsection{第61章 论祂放在他们面前作为他们效法的小孩，和世界的绊脚石；论身体中使人跌倒的肢体；论小孩子的天使，他们常见天父的面；论一百只羊中的一只；论私下责备弟兄；论解开和捆绑罪；论两个人的同心和三个人的聚集；论赦免罪直到七十个七次；论那仆人自己欠的大债被免了，却不肯免同伴欠他的小债；以及关于玛窦与其他福音作者在这些主题上的和谐的问题。}

119. 随后同一玛窦接着叙述如下：\BibleQuote{在那一刻，门徒们来到耶稣跟前，说：\InnerQuote{你认为在天国里谁是最大的？}耶稣叫了一个小孩子来，把他放在他们中间，说：\InnerQuote{我实在告诉你们：你们若不悔改，变成像小孩子一样，绝不能进天国；}}如此下去，直到以下的话：\BibleQuote{同样，如果你们不从心里宽恕你们的弟兄，我的天父也必这样对待你们。}在这段由主所说、较长的讲论中，马尔谷并未全部记录，只选取了部分内容，同时保持顺序一致。他也引入了一些玛窦未提及的事情。此外，在这段完整的讲论中，就我们目前所考虑的，唯一的打断是伯多禄询问应宽恕弟兄多少次。然而，主讲话的口吻清楚地表明，伯多禄提出的问题以及对他的回答，都属于同一讲论。路加则没有按此顺序记录这些事，只记录了那个小孩的插曲，即当门徒们思虑自己为大时，主将小孩放在他们中间作为效法。因为如果他也讲述了一些与这篇讲论中所插入的相似内容，那是他在其他上下文中、在不同的场合提醒的；正如若望记录了主关于赦免罪的话——即那些话的意思是：宗徒们赦免谁的罪，谁的罪就被赦免；他们保留谁的罪，谁的罪就被保留——是主在复活后说的；而玛窦提及在这篇讲论中主也做了这个声明，不过同一位福音作者同时声称这之前已经对伯多禄说过。因此，为了避免总要反复强调同一规则，我们应当记住：耶稣多次在不同地点说了同一句话——这个原则我们已经多次提醒你们注意；这样，即使语句的顺序被认为会造成一些困难，我们也不会感到困惑。\par

\SourceRecord{Translated by S.D.F. Salmond.；Nicene and Post-Nicene Fathers, First Series；Vol. 6.；Edited by Philip Schaff.；Buffalo, NY: Christian Literature Publishing Co.,；1888.；<http://www.newadvent.org/fathers/1602261.htm>.}

% structure: p0001 source=h1 target=section reason=page-title

\section{福音的和谐 第二卷}

% structure: p0002 source=h2 target=subsection reason=relative-heading-rank; base=h2

\subsection{第六十二章 论\Person{玛窦}与\Person{马尔谷}在论及休妻是否合法之时间记载上的和谐，尤其涉及\Person{主}与\Person{犹太人}之间具体问答，以及福音作者似乎略有分歧之处。}

120. \Person{玛窦}继续叙述如下：\BibleQuote{当\Person{耶稣}讲完了这些话，就离开\Place{加里肋亚}，来到\Place{约但河}对岸的\Place{犹太}境内；有许多群众跟随祂，祂就在那里治好了他们。\Person{法利塞人}也来见祂，试探祂说：人无论为了什么缘故，都可以休妻吗？}如此直到\BibleQuote{能领受的，就领受吧。}\Person{马尔谷}也记载了此事，且保持同样的顺序。同时，我们必须注意，不要以为有个矛盾：\Person{马尔谷}告诉我们，\Person{法利塞人}被\Person{主}问及\Person{梅瑟}命令他们什么，当\Person{主}这样问他们时，他们回答关于休书的事，\Person{梅瑟}准许他们写休书；而按照\Person{玛窦}的说法，是在\Person{主}说了那些话——祂从法律中向他们指出，\Person{天主}如何造男造女成为一体，因此那些（由祂这样结合的）人不应该被人分开——之后，他们才回答说：\BibleQuote{那么\Person{梅瑟}为什么吩咐给休书，就可以休她呢？}（按照\Person{玛窦}的说法）对于这个问题，\Person{主}又回答说：\BibleQuote{\Person{梅瑟}因为你们的心硬，才准许你们休妻；但起初并不是这样。}我再说一遍，这并无困难；因为\Person{马尔谷}并没有完全不提\Person{主}这样给的回答，而是在他们回答\Person{主}关于休书的问题之后，才引入这个回答。\par

121. 就此处所采用的说法次序或方法而言，我们应当明白，这绝不影响到主题本身的真理：无论是这些\Person{法利塞人}向\Person{主}提出关于那位\Person{梅瑟}（他也记载了\Person{天主}造男造女成为一体）准许写休书的问题，是在\Person{主}禁止夫妻分离并以法律的权威确认祂在这件事上的声明之时；还是这个问题是同一批人在\Person{主}问他们关于\Person{梅瑟}命令他们什么时所给回答的一部分。因为祂的意图是，在他们自己先提到这件事之前，不向他们解释\Person{梅瑟}如此准许他们的理由；\Person{马尔谷}所引入的询问就表明了祂的这一意图。另一方面，他们的愿望是利用\Person{梅瑟}关于给休书的命令的权威，来堵住祂的口——可以说——在祂禁止人休妻这件事上，因为他们相信祂无疑会禁止。因为他们接近祂，是为了说出试探祂的话。\Person{玛窦}表明了他们的这个愿望：他没有陈述他们自己先被审问，而是说他们主动提出了关于\Person{梅瑟}命令的问题，为的是借此指责\Person{主}在禁止休妻的事上做错了。因此，既然两位福音作者都表明了说话者的心意（言语应当服务于心意），那么无论二人采用的叙述方式如何不同，只要双方都没有错误地阐述主题本身，就无关紧要。\par

122. 也可以采取另一种看法，即按照\Person{马尔谷}的陈述，当这些人开始以休妻的问题质询\Person{主}时，\Person{主}反过来问他们\Person{梅瑟}命令他们什么；他们回答说\Person{梅瑟}准许他们写休书并休妻，\Person{主}就对他们作了关于\Person{梅瑟}所赐之法律的回答，提醒他们\Person{天主}如何建立了男女的结合，并用\Person{玛窦}所插入的话对他们说：\BibleQuote{难道你们没有念过，那创造者起初就造了他们一男一女？}等等。他们听了这话，就把他们在回答\Person{主}第一次质问时已经说过的话又以询问的形式重复了一遍，即：\BibleQuote{那么\Person{梅瑟}为什么吩咐给休书，就可以休她呢？}然后\Person{主}指出原因是他们的心硬；\Person{马尔谷}为了简洁，提前引入了这个解释，仿佛是作为他们先前那个回答的回复，而\Person{玛窦}省略了那个回答。他这样做，是认为无论这个解释在什么时候引入，都不会损害真理，因为促使它回答的那些话已经以同样的形式说了两次；而且\Person{主}在任何情况下都是以这样的话来回答的。\par

\SourceRecord{Translated by S.D.F. Salmond.；Nicene and Post-Nicene Fathers, First Series；Vol. 6.；Edited by Philip Schaff.；Buffalo, NY: Christian Literature Publishing Co.,；1888.；<http://www.newadvent.org/fathers/1602262.htm>.}

% structure: p0001 source=h1 target=section reason=page-title

\section{福音的谐和，第二卷}

% structure: p0002 source=h2 target=subsection reason=relative-heading-rank; base=h2

\subsection{第六三章　论祂按手在孩童身上；论祂对那富人所言：\Quote{变卖你所有的一切；}论在不同时辰雇工到葡萄园做工；以及论玛窦与另两位圣史在这几件事上并无任何矛盾之处的问题。}

一二三．玛窦继续写道：\Quote{那时，有人给祂带来一些小孩子，要祂给他们覆手祈祷；门徒却斥责他们；}一直到我们读到：\Quote{因为被召的人多，被选的人少。}马尔谷在此处与玛窦同序。但只有玛窦引入了关于雇工到葡萄园做工的那段。路加则先提到祂对那些互相询问谁将是最大的所说的话，接着立即描述关于他们所见有人驱魔那人却没有跟随祂的事；然后他在告诉我们祂如何决意面向耶路撒冷前行时，与另两位圣史分道；在插入许多题材之后，他在讲述那富人的故事时又与他们会合，那人对祂说：\Quote{变卖你所有的一切；}这个人的情况另两位圣史也在这里叙述，但仍按所有叙述都遵循的顺序。因为在路加所引的那段经文中，那位作者没有遗漏提到小孩子的故事，正如另两位圣史在提及富人之前立即记载的一样。那么，关于这位富人的记载——他询问应行什么善事才能获得永生——它们之间似乎存在一些矛盾，因为按玛窦，说的是：\Quote{你为什么问我关于善的事？}而按另两位，则是：\Quote{你为什么称我是善的？}那么\Quote{你为什么问我关于善的事？}这句话更可能是指那人所问的：\Quote{我应行什么善事？}因为在那里我们既称呼基督为\Quote{善}，又提出了问题。但\Quote{善师}这个称呼本身并不包含问题。因此，最好的处理方法是理解这两句话都被说出过：\Quote{你为什么称我是善的？}和\Quote{你为什么问我关于善的事？}\par

\SourceRecord{Translated by S.D.F. Salmond.；Nicene and Post-Nicene Fathers, First Series；Vol. 6.；Edited by Philip Schaff.；Buffalo, NY: Christian Literature Publishing Co.,；1888.；<http://www.newadvent.org/fathers/1602263.htm>.}

% structure: p0001 source=h1 target=section reason=page-title

\section{\WorkTitle{福音的和谐，第二卷}}

% structure: p0002 source=h2 target=subsection reason=relative-heading-rank; base=h2

\subsection{第六十四章 论祂私下向门徒预言受难的时机；载伯德儿子的母亲同其子前来，求主使一个坐在祂右边，一个坐在祂左边之时；以及玛窦与其他两位福音作者在这些主题上毫无矛盾之处。}

124. 玛窦继续叙述如下：\BibleQuote{耶稣上耶路撒冷去，私下把十二门徒带在一边，对他们说：\InnerQuote{看，我们上耶路撒冷去，人子要被交在司祭长和经师手中，他们要定祂死罪，把祂交给外邦人戏弄、鞭打、钉在十字架上；第三天祂要复活。}那时，载伯德儿子的母亲同她的儿子前来朝拜祂，向祂求一件事；}直至\BibleQuote{就如人子来不是为受服事，而是为服事人，并交出自己的生命为大众作赎价。}在此，马尔谷保持了与玛窦相同的顺序，只是他写道载伯德的儿子们亲自提出请求；而玛窦则陈述说是由他们的母亲代他们提出请求，而非他们本人，因为她将他们心愿呈于主前。因此马尔谷简要地暗示了当时所言是由他们而非她以他们之名说出。归根结底，按照玛窦和马尔谷的记载，主回复的对象是他们而非她。路加则在按相同顺序叙述了主关于受难与复活对十二门徒的预言后，略过了其他两位福音作者紧接着记载的内容；在这些段落插入之后，他在耶里哥事件处再次与同伴保持一致。至于玛窦和马尔谷关于外邦人的王侯管辖臣属之陈述——即门徒不应如此，而是他们中最大的应成为他人的仆人——路加也给出了类似内容，虽不在同一上下文；但这顺序本身表明主是在另一场合说了相同的话。\par

\SourceRecord{Translated by S.D.F. Salmond.；Nicene and Post-Nicene Fathers, First Series；Vol. 6.；Edited by Philip Schaff.；Buffalo, NY: Christian Literature Publishing Co.,；1888.；<http://www.newadvent.org/fathers/1602264.htm>.}

% structure: p0001 source=h1 target=section reason=page-title

\section{\WorkTitle{福音合参}，卷二}

% structure: p0002 source=h2 target=subsection reason=relative-heading-rank; base=h2

\subsection{第65章 论玛窦与马尔谷，或玛窦与路加，在记述耶里哥盲人复明一事上并无矛盾}

125. 玛窦接着这样记载：\Quote{他们从耶里哥出来时，有一大群人跟随祂。看，有两个盲人坐在路旁，听说耶稣经过，便喊叫说：\InnerQuote{主啊，达味之子，可怜我们吧！}}直至\Quote{他们立刻看见了，就跟随了祂}。马尔谷也记载了这件事，但只提到一个盲人。这个难题，正如从前在革辣撒人的地方，两个被群魔折磨的人的例子一样，可以得到解决。因为，玛窦在此单独提到的两个盲人中，有一个在城中是显赫有名的；单凭这一点就足以证明：马尔谷记录了他的名字和他父亲的名字；在主所行的许多治愈中，很少有这样的情况，除非是那位会堂长雅依洛的女儿复活的事迹，福音作者才提及了名字。在后一事例中，意图更为明显，因为这位会堂长在本地确实有地位。因此，几乎可以肯定，这位提买之子巴尔提买，曾从某种极兴旺的地位上跌落，如今被视为最臭名昭著、最悲惨可怜的对象，因为他不但瞎眼，还不得不坐着讨饭。这也就是为什么马尔谷只提了这一个人的原因，此人重见光明使该奇迹的名声传得同他自身不幸的恶名一样广远。\par

126. 路加虽然记载了一件性质完全相同的事，但应理解为：他实际上是叙述了另一盲人所经历的类似奇迹，并记录其施行方法与上述奇迹的相似之处。因为他记载这事发生在耶稣将近耶里哥时，而其他两位却说发生在耶稣离开耶里哥时。城名和事迹的相似，使人倾向认为只有一次这样的事。然而，若有人以为福音作者在此确实自相矛盾——这位说\Quote{将近耶里哥}，那位说\Quote{出了耶里哥}——除非那些宁可信福音不真实、也不信主行了两个性质类似且情况相似的奇迹的人，否则无人会接受。但凡忠于福音的儿子，都能立刻看出这两个假设哪个更可信、更应被接受为真实；而任何反对者，一旦被告知实情，或以他不得不保持的沉默，或至少以他若拒绝沉默时的思考内容，自会回答自己。\par

\SourceRecord{Translated by S.D.F. Salmond.；Nicene and Post-Nicene Fathers, First Series；Vol. 6.；Edited by Philip Schaff.；Buffalo, NY: Christian Literature Publishing Co.,；1888.；<http://www.newadvent.org/fathers/1602265.htm>.}

% structure: p0001 source=h1 target=section reason=page-title

\section{福音的和谐，第二卷}

% structure: p0002 source=h2 target=subsection reason=relative-heading-rank; base=h2

\subsection{第66章 论\Person{玛窦}所提到的驴驹，以及他与其他仅提及驴的福音作者叙述的一致性}

127. \Person{玛窦}继续叙述如下：\BibleQuote{当他们临近耶路撒冷，来到贝特法革，上了橄榄山，耶稣就派了两个门徒，对他们说：你们往对面的村庄去，立刻会看见一匹母驴拴在那里，还有一匹驴驹同她在一起；}一直到这句话：\BibleQuote{奉主名而来的当受赞美；贺三纳于至高之天。}\Person{马尔谷}也记载了这件事，并按相同顺序插入。\Person{路加}则在\Place{耶里哥}逗留了一段时间，叙述了一些其他作者省略的事情——即税吏长\Person{匝凯}的事迹，以及一些以比喻形式表达的言论。然而，举出这些事例后，这位福音作者又在有关\Person{耶稣}所骑的驴的叙述中与其他作者会合。\Person{玛窦}提到了一匹母驴和一匹驴驹，而其他作者只字不提母驴，这不应使我们困惑。因为在这里，我们必须再次记住我们之前在处理关于用五个饼饱饱喂养群众时，人们以五十人和一百人为一组就座的叙述时所引入的规则。应用这一原则后，读者不应在此情况下感到任何严重困难。事实上，即使\Person{玛窦}没有提到驴驹，正如他的同行史家没有注意到母驴一样，也不应造成任何困惑，以致认为当一位作者只提到母驴，而另一位只提到驴驹时，这两段陈述之间存在着不可克服的矛盾。那么，当我们看到一位作者提到了母驴（其他作者忽略提及）时，他同时也没有忽略其他作者谈到的驴驹，这岂不更应减少不安？总之，在假定两者都包含在事件中是可能的情况下，即使一位作者只指定一件事，另一位只指定另一件，也没有真正的对立。更何况当一位作者特别提及一件事，而另一位则同时提及两件事时，更不会有矛盾！\par

128. 再者，尽管\Person{若望}没有告诉我们主如何派遣门徒去牵这些动物，但他简短提及了这匹驴驹，并引用了\Person{玛窦}所用的先知的话。在这段先知证词中，福音作者复述的措辞虽然显示某些差异，但所表达的含义是一致的。然而，可能感到困难的是，\Person{玛窦}引用这段经文的形式似乎表示先知提到了母驴；而\Person{若望}的引文以及通行译本的教会抄本中的版本却没有这样做。在我看来，这种差异的解释在于：\Person{玛窦}被认为是用希伯来文写作他的福音。此外，名为\WorkTitle{七十贤士译本}的译本在某些细节上与熟悉希伯来文的人以及多位学者为我们提供的同一希伯来文书籍的译文中所发现的文本不同。如果问及这种不一致的原因，或者问及为何\WorkTitle{七十贤士译本}的权威在许多地方与希伯来文抄本中发现的真理的表述不同，我认为没有比以下假设更合理的解释：这七十位译者在圣神的感动下进行了翻译，而正是同一位圣神默感了他们所翻译的原文。这一观点也因他们据说具有的奇妙一致而得到确认。因此，当这些译者没有偏离这些言论所出自的天主的真实心意（言论应服务于表达这一心意），而在复述文本时对某些内容给出了不同形式时，他们并无意表明任何其他东西，而正是我们现在所惊叹的东西——即四位福音作者那种和谐的多样性。在这种多样性中，尽管一位史家可能以与另一位不同的方式描述某个主题，但并非如此不同以至于偏离了他所必须与之和谐一致的人的本意。理解这一点对品格有益，既能防范错误，又能正确判断；对信德本身也至关重要，可以防止认为真理似乎有某种神圣音调的保护，暗示天主不仅将事物本身，而且将表达所需的具体词语都交给我们；而事实是，所要表达的主题本身被认为比表达它的词语重要得多，如果我们有可能像天主知道真理那样，以及像祂的天使在祂内知道真理那样，在不用词语的情况下知道真理，那么我们就没有必要询问它们了。\par

\SourceRecord{Translated by S.D.F. Salmond.；Nicene and Post-Nicene Fathers, First Series；Vol. 6.；Edited by Philip Schaff.；Buffalo, NY: Christian Literature Publishing Co.,；1888.；<http://www.newadvent.org/fathers/1602266.htm>.}

% structure: p0001 source=h1 target=section reason=page-title

\section{福音的和谐，第二卷}

% structure: p0002 source=h2 target=subsection reason=relative-heading-rank; base=h2

\subsection{第67章 论驱逐买卖者出圣殿，以及前三部福音与若望在叙述同一事件时关联大不相同的问题}

129. 玛窦继续叙述如下：\Quote{及至祂进了耶路撒冷，全城都惊动了，说：这是谁？群众说：这是耶稣，加里肋亚纳匝肋的先知。耶稣进了天主的圣殿，把一切在圣殿内买卖的人都赶出去；}等等，直到我们读到\Quote{你们却使它成了贼窝。}众圣史都记载了被赶出圣殿的众多买卖者；但若望却以一种非常不同的顺序引入此事。原来，记录若翰洗者对耶稣的作证，提到祂在把水变酒时去了加里肋亚，并提到在葛法翁逗留了数日之后，若望继续告诉我们，在犹太人的逾越节期间祂上耶路撒冷去，用细绳做了一条鞭子，把在圣殿里做买卖的人赶了出去。由此可见，主做这事不是一次，而是两次；不过第一次只有若望记录，最后一次由其他三位记录。\par

\SourceRecord{Translated by S.D.F. Salmond.；Nicene and Post-Nicene Fathers, First Series；Vol. 6.；Edited by Philip Schaff.；Buffalo, NY: Christian Literature Publishing Co.,；1888.；<http://www.newadvent.org/fathers/1602267.htm>.}

% structure: p0001 source=h1 target=section reason=page-title

\section{\WorkTitle{福音的和谐 第二卷}}

% structure: p0002 source=h2 target=subsection reason=relative-heading-rank; base=h2

\subsection{第六十八章 论无花果树的枯萎，以及关于\Person{玛窦}与其他福音作者在记述该事件及相关其他事情上毫无矛盾的提问；尤其是关于\Person{玛窦}与\Person{马尔谷}在叙述顺序上的一致性。}

130. \Person{玛窦}继续记述说：\Quote{瞎子、瘸子来到圣殿里，祂就治好了他们。司祭长和经师看见祂所行的奇事，又看见儿童在圣殿里呼号说：\InnerQuote{贺三纳于达味之子！}就十分恼怒，对祂说：\InnerQuote{祢听见这些人说什么吗？}耶稣对他们说：\InnerQuote{是的；你们从未读过：\BibleQuote{由婴孩和吃奶者的口中，祢使赞美得以完备}吗？}祂便离开他们，出城到\Place{贝特阿尼}去，在那里过夜。早晨回城时，祂饿了。看见路旁有一棵无花果树，就走到跟前，除了叶子外，什么也没找到，就对它说：\InnerQuote{以后你永不再结果子！}那无花果树立刻枯萎了。门徒看见了，惊讶说：\InnerQuote{这无花果树怎么立刻枯萎了！}耶稣回答他们说：\InnerQuote{我实在告诉你们：你们如果有信德，毫不疑惑，不但能行这无花果树上所做的事，而且如果你们向这座山说：\InnerQuote{你移开，投到海里去！}，也必成就。凡你们在祈祷中祈求的，只要相信，就必得到。}}\par

131. \Person{马尔谷}也按着顺序记载了此事。\newline{}然而，他在叙述中并未遵循同样的顺序。首先，\Person{玛窦}所记载的事实，即\Person{耶稣}进了圣殿，赶出在那里买卖的人，\Person{马尔谷}在那一处并未提及。另一方面，\Person{马尔谷}告诉我们，祂环视了一切，到了傍晚，便与十二宗徒出去往\Place{伯达尼}去了。接着，他又告诉我们，另一天，当他们从\Place{伯达尼}回来时，祂饿了，就诅咒了那无花果树，\Person{玛窦}也暗示了这一点。然后，这位\Person{马尔谷}又补充说，祂进了\Place{耶路撒冷}，进入圣殿后，赶出在那里买卖的人，仿佛这件事不是发生在所述的第一天，而是发生在另一天。但既然\Person{玛窦}用这些话来连接：\BibleQuote{祂离开了他们，出城往\Place{伯达尼}去}，并告诉我们祂是在早晨回城时诅咒了那树，那么更合理的看法是：就赶出圣殿中买卖之人这一事件而言，\Person{玛窦}比\Person{马尔谷}更严格地保全了时间顺序。因为当他使用这句话：\BibleQuote{祂离开了他们，就出去了}，那么祂所离开的这些\Quote{他们}除了指祂先前与之交谈的那些人——即因儿童呼喊\BibleQuote{和散那归于达味之子}而极为不悦的人——还能指谁呢？由此可见，\Person{马尔谷}省略了第一天祂进圣殿所发生的事；而在提到祂在无花果树上除了叶子什么也没找到时，他引入的只是祂在那里所想起的，而这件事实际上发生在第二天，两位福音作者都作证如此。此外，他的记述表明：门徒们发现无花果树枯萎时所表现的惊讶，以及主就信德之事给他们的答复，以及移山入海，都不属于祂对树说\BibleQuote{以后永远没有人吃你的果子}的这第二天，而是属于第三天。因为关于第二天，那位\Person{马尔谷}记载了赶出圣殿买卖之人的事件，这件事他先前略去未提，以为是属于第一天的。因此，正是关于这第二天，他告诉我们\Person{耶稣}如何在天晚时出城，以及次日早晨他们经过时，门徒们如何看见那无花果树从根枯干，伯多禄如何想起这事，对祂说：\BibleQuote{老师，请看，你所诅咒的无花果树已经枯干了。}然后，他又告诉我们，祂给出了有关信德之能力的回答。另一方面，\Person{玛窦}记述这些事的方式暗示它们都发生在同一天；也就是说，对树所说的话：\BibleQuote{从今以后，永远不要有果子长在你身上}，以及随后的树迅速枯萎，还有当门徒们看到那枯萎并感到惊奇时，祂就信德之能力给他们的答复。由此我们应当理解，\Person{马尔谷}方面，他关于第二天记载了他先前略去、实际发生在第一天的事——即从圣殿中赶出买卖之人事件。另一方面，\Person{玛窦}在提及了第二天所做的事——即祂早晨从\Place{伯达尼}回城时诅咒无花果树——之后，省略了\Person{马尔谷}插入的某些事实，即祂进城和晚上出城，以及门徒们早晨经过时看到树枯干所表现的惊讶；然后，他把第二天（即树被诅咒的那天）发生的事与第三天实际发生的事——即门徒们看到树的枯萎状态的惊奇，以及他们从主那里听到关于信德之能力的宣告——连接在一起。\Person{玛窦}将这些不同的事实如此连接在一起，以致若非\Person{马尔谷}的叙述迫使我们留意此事，我们将无法认出前一位作者在叙述中是在哪一点、或就什么情况省略了某些内容。因此情况如下：\Person{玛窦}首先用这些话陈述事实：\BibleQuote{祂离开了他们，出城往\Place{伯达尼}去，并在那里过夜。第二天早晨，祂回城时饿了，看见路旁有一棵无花果树，就走到跟前，除了叶子以外，什么也没有找到，就对它说：\InnerQuote{从今以后，永远不要有果子长在你身上。}那无花果树立刻就枯萎了。}然后，他省略了属于同一天的其他事情，立即附加了这句话：\BibleQuote{门徒们看见了，便惊奇地说：\InnerQuote{这无花果树怎么立刻就枯萎了呢？}}虽然他们是在另一天看到这景象，另一天感到惊奇。但可以理解，树并不是在他们看见的时刻枯萎的，而是在它被诅咒的那一刻即刻枯萎的。因为他们看到的不是正在枯萎的树，而是已经完全枯干的树；这样他们就明白，树在主发命之后就立刻枯萎了。\par

\SourceRecord{Translated by S.D.F. Salmond.；Nicene and Post-Nicene Fathers, First Series；Vol. 6.；Edited by Philip Schaff.；Buffalo, NY: Christian Literature Publishing Co.,；1888.；<http://www.newadvent.org/fathers/1602268.htm>.}

% structure: p0001 source=h1 target=section reason=page-title

\section{福音的和谐，第二卷}

% structure: p0002 source=h2 target=subsection reason=relative-heading-rank; base=h2

\subsection{第六十九章：论前三部福音作者在记述犹太人问主凭什么权柄做这些事时的和谐一致}

132. 玛窦继续他的叙述如下：\BibleQuote{当祂进入圣殿，正在施教时，司祭长和民间长老来到祂跟前说：\InnerQuote{祢凭什么权柄做这些事？谁给了祢这权柄？}耶稣回答他们说：\InnerQuote{我也要问你们一件事，你们若告诉我，我也告诉你们我凭什么权柄做这些事。若翰的洗礼是从哪里来的？}}等等，直到这些话：\BibleQuote{我也不告诉你们我凭什么权柄做这些事。}另外两位，\Person{玛尔谷}和\Person{路加}，也几乎用同样的话陈述了整段经文。在顺序上似乎也没有任何出入，唯一的例外是我上面提到的情形——即\Person{玛窦}省略了属于另一日的某些事项，并以一种连接方式构建他的叙述，若不是我们注意到这一事实，可能会认为他仍是在叙述第二日的事，而\Person{玛尔谷}处理的是第三日。此外，\Person{路加}没有附上他对这一事件的记述，仿佛他是要按顺序逐日叙述；但在记载了把买卖人从圣殿赶出去之后，他略过了以上所述的一切——祂前往\Place{伯达尼}，祂返回城里，对无花果树所做的事，以及当门徒们惊奇时关于信德之能的回答。然后，在省略了这一切之后，他以这样的话引入了下一段叙述：\BibleQuote{祂天天在圣殿里施教。但司祭长、经师和百姓的首领想要除灭祂，却找不着可行的事，因为百姓都很留心听祂。有一天，当祂在圣殿里教训百姓，宣讲福音时，司祭长和经师们同长老们忽然来到祂跟前，对祂说：\InnerQuote{告诉我们，祢凭什么权柄做这些事？}}等等；其他两位福音作者也类似地记载了这一切。由此可见，\Person{路加}与其他两位并无矛盾，甚至在顺序上也是如此；因为他所记述的发生于\BibleQuote{那些日子中的一天}的事，可以被理解为属于他们也都记载为发生在那一特定日子的事。\par

\SourceRecord{Translated by S.D.F. Salmond.；Nicene and Post-Nicene Fathers, First Series；Vol. 6.；Edited by Philip Schaff.；Buffalo, NY: Christian Literature Publishing Co.,；1888.；<http://www.newadvent.org/fathers/1602269.htm>.}

% structure: p0001 source=h1 target=section reason=page-title

\section{\WorkTitle{福音的和谐，第二卷}}

% structure: p0002 source=h2 target=subsection reason=relative-heading-rank; base=h2

\subsection{\Emph{第七十章：关于被父亲命令去葡萄园工作的两个儿子，以及租给其他农夫的葡萄园；关于玛窦记述的这些段落与另外两位福音作者的一致性，他们保持了相同的顺序；以及特别关于三位福音作者都记述的葡萄园租给农夫的比喻的和谐；并特别涉及比喻中人物所作的答复，玛窦在此似乎与其他作者略有不同。}}

133. 玛窦接着写道：\BibleQuote{你们以为如何？一个人有两个儿子。他到第一个跟前说：孩子，你今天到我的葡萄园里去工作。他回答说：我不去；但后来后悔了，就去了。他又到第二个跟前说了同样的话。他回答说：主，我去；但却没有去。}直到\BibleQuote{凡跌在这石头上的，必被撞碎；这石头落在谁身上，必把他砸得粉碎。}马尔谷和路加没有提到被命令去葡萄园工作的两个儿子的比喻。但玛窦随后叙述的——即租给农夫的葡萄园的比喻，这些农夫迫害了派去的仆人，后来杀死了爱子，并把他扔出葡萄园——那两位作者也没有忽略。在详细叙述时，他们也都保持了相同的顺序，也就是说，在犹太人被问及若翰的洗礼时无法回答，以及主回答他们说：\BibleQuote{我也不告诉你们我凭什么权柄做这些事}之后，他们引入了这个比喻。\par

134. 这里不存在任何表明这些记述之间有矛盾的疑问，除非因为以下情况：玛窦在叙述了主如何向犹太人提出这个问题：\BibleQuote{葡萄园的主人来到时，他将怎样对待那些农夫？}之后，补充说他们回答说：\BibleQuote{他会狠狠地消灭那些恶人，并将葡萄园租给其他农夫，这些农夫会在季节里按时向他缴纳果实。}因为马尔谷并没有把这些话记载为那些人所作的答复，而是把它们引入，好像它们是主在提出问题之后立即说的，因此主在某种意义上自己回答了。因为（在这部福音中）他这样说：\BibleQuote{那么葡萄园的主人将做什么？他必定来消灭那些农夫，并将葡萄园交给别人。}但我们很容易假设，要么是那些人的话在此被附上，而没有插入说明性短语\Quote{他们说}或\Quote{他们回答}，这留待理解；要么是所说的答复归因于主自己而非那些人，因为当他们如此真实地回答时，主——祂本身就是真理——实际上也针对那些人给出了同样的答复。\par

135. 然而，更严重的困难可能来自以下事实：路加不仅没有说他们是作出那个答复的人（因为他和马尔谷一样，将这些话归给主），而且甚至记载他们给出了相反的答复，说：\BibleQuote{绝对不可！}因为他的叙述是这样继续的：\BibleQuote{那么葡萄园的主人将怎样对待他们？他必来消灭这些农夫，并将葡萄园交给别人。他们听了后说：绝对不可！祂看着他们说：那么经上写的是什么？\InnerQuote{匠人所弃的石头，竟成了屋角的基石？}}那么，根据玛窦的版本，听祂说这些话的人说：\BibleQuote{他会狠狠地消灭那些恶人，并将葡萄园租给其他农夫，这些农夫会在季节里按时向他缴纳果实；}而根据路加，他们却给出了与此相矛盾的答复，说：\BibleQuote{绝对不可！}实际上，主紧接着关于被匠人弃绝却成了屋角基石的石头所说的话，其引入方式暗示，这见证是为了驳斥那些反对比喻真正含义的人。因为玛窦和路加一样，记载了那段经文，似乎是为了反驳反对者，他说：\BibleQuote{你们没有读过圣经吗？\InnerQuote{匠人所弃的石头，竟成了屋角的基石？}}因为这个问题\BibleQuote{你们没有读过吗}暗示着他们所作的回答与（比喻的）真正意图相悖。这也由马尔谷指出，他以这样的方式给出同样的话：\BibleQuote{你们也没有读过这段经文吗？\InnerQuote{匠人所弃的石头，竟成了屋角的基石？}}因此，这句话在路加福音中似乎比在其他福音书中更占据其作为原始话语的适当位置。因为它是紧随着那些人说\BibleQuote{绝对不可！}的反对之后引入的。他以\BibleQuote{那么经上写的是什么？\InnerQuote{匠人所弃的石头，竟成了屋角的基石？}}这种形式表达，与其他表达方式在意义上等效。因为无论使用三种短语中的哪一种——\BibleQuote{你们没有读过吗？}或\BibleQuote{你们没有读过吗？}或\BibleQuote{那么经上写的是什么？}——句子的真实含义都同样得以表明。\par

136. 因此，我们应当理解，在那次听讲的群众中，有些人以玛窦所记述的方式回答，当他这样写时：\BibleQuote{他们对祂说：祂要凶恶地消灭那些恶人，并将葡萄园另租给别的佃农。}\BibleRef{玛 21:41} 也有些人的回答正如路加所指出的，即说了：\BibleQuote{但愿不如此。}\BibleRef{路 20:16} 因此，那些曾以第一种方式回答主的人，被群众中另一些人用\BibleQuote{但愿不如此。}的解释所回答。但是，这两方中第一方实际所给的回答——第二方回应他们说\BibleQuote{但愿不如此。}——被玛尔谷和路加归于主自己，其根据是，如我已暗示的，真理本身通过这些人说话，无论他们是作为不知自己是恶人的（正如祂也通过盖法说话，盖法当大司祭时预言却不知自己所说），还是作为理解并已到达知识和信仰的人。因为当时也有那群众在场，预言已经通过他们应验了：当祂临近时，他们以浩大的队伍迎接祂，高呼：\BibleQuote{奉主名而来的，当受赞美。}\BibleRef{玛 21:9}\par

137. 我们也不应被以下情况绊倒：同一位玛窦记述了司祭长和民间长老来到主前，问祂凭甚么权柄做这些事，是谁给了祂这权柄，当时祂反过来问他们关于若翰的圣洗，追问它是从天上来的，还是从人来的；当他们回答说不知道时，祂说：\BibleQuote{我也不告诉你们我凭甚么权柄做这些事。}\BibleRef{玛 21:27} 因为祂紧接着就在上下文中引入了这句话：\BibleQuote{你们以为如何？一个人有两个儿子，}\BibleRef{玛 21:28} 等等。因此，这段言论便连接起来，毫无间断，没有任何事物或人的插入，直到关于那租给佃农的葡萄园所叙述的。的确，可以设想祂对司祭长和民间长老讲了这一切话，是他们质问了祂的权柄。但是，如果这些人确实是为了试探祂、怀着敌意问祂，那他们就不能被视为那些已相信的人，也不能被视为引用了那个出自先知、有利于主的显著见证；显然，只有当他们具有相信者（而非无知者）的品格时，才可能给出这样的回答：\BibleQuote{祂要凶恶地消灭那些恶人，并将葡萄园另租给别的佃农。}\BibleRef{玛 21:41}\BibleRef{路 20:16} 然而，这一特点（玛窦记述的特点）决不应如此困扰我们，以至于让我们以为在此时听主讲比喻的群众中没有一个相信的人。因为同一位玛窦只是为了简洁而略过了路加没有忽略的事——即那个比喻并非只对质问祂权柄的人而讲，也是对群众讲的。因为后者（路加）这样记载：\BibleQuote{于是祂开始对群众讲这个比喻：有一个人栽种了一个葡萄园，}\BibleRef{路 20:9} 等等。因此，我们可以很好地理解，在当时聚集的群众中，也可能有一些人能像之前那样听祂，那些人曾说过：\BibleQuote{奉主名而来的，当受赞美。}\BibleRef{玛 21:9} 而且，这些人或他们中的某些人，正是那些以这些话回答的人：\BibleQuote{祂要凶恶地消灭这些恶人，并将葡萄园另租给别的佃农。}此外，这些人实际作出的回答也被玛尔谷和路加归于主自己，不仅因为他们的言语确实就是祂的言语——因为祂就是真理，常常甚至通过恶人和无知者说话，以一种隐秘的本能推动人的心灵，不是出于人圣洁的功绩，而是出于祂自己权能的权利——而且也因为这些人可能具有这样的品格，使他们有理由被视为已在基督真身体内的肢体，因此他们所说的完全可以归于他们所属的那位（基督）。因为此时祂已付洗比若翰更多，并有大批门徒，正如同一批福音书作者多次证明的；从这些跟随者中，祂也召选了那五百弟兄，宗徒保禄告诉我们，祂在复活后曾显现给他们。这一解释也得到了这样的事实的支持：在同一玛窦的版本中出现的短语——即\BibleQuote{他们对祂说：祂要凶恶地消灭那些恶人，}——其表述方式并不要求我们把代词\OriginalTerm{他（或他们）}{illi}理解为复数，仿佛它特意要表明那些话是那些狡猾地询问祂权柄的人的回答；相反，\BibleQuote{他们对祂说}这一分句的表达方式使得\OriginalTerm{他（或他们）}{illi}应被理解为单数代词，而非复数，并应被视为意指\Quote{对祂}，即对主自己，这在希腊文抄本中毫无歧义。\par

138. \BibleBook{若望}{福音}记载了主的一段言论，可以帮助我们更容易理解我所说的这句话。其内容如下：\BibleQuote{于是耶稣对那些信从祂的犹太人说：\Quote{你们如果坚持我的话，就确实是我的门徒；你们也会认识真理，而真理必使你们得到自由。}他们回答说：\Quote{我们是亚巴郎的后裔，从未给任何人做过奴隶；你怎么说\InnerQuote{你们必要得到自由}呢？}耶稣回答说：\Quote{我实实在在告诉你们：凡犯罪的，就是罪恶的奴隶。奴隶不能永远住在家里，圣子却永远住在家里。所以，如果圣子使你们自由，你们就的确自由了。我知道你们是亚巴郎的后裔；但你们却想杀我，因为我的话在你们心中没有容身之地。}}现在，当然不能认为祂所说的\Quote{你们想杀我}这句话是对那些已经信从祂并听到祂说\Quote{你们如果坚持我的话，就确实是我的门徒}的人说的。然而，由于祂后来对已经信从祂的人说了这些话，而且当时在场的一群人中有许多是敌视祂的，即使福音作者没有明确说明回答祂的是哪些人，但根据他们回答的性质以及祂随后恰当回应他们的话的实质，足以清楚表明当时是向哪些具体的人说话，以及对他们说了些什么话。因此，正如在\Person{若望}所指的那群人中，有些人已经信从了耶稣，也有些人想要杀祂，在我们现在讨论的另一群人中，有些人狡猾地询问主关于祂凭什么权柄做这些事，而另一些人则不是出于欺骗，而是出于信德，欢呼说：\BibleQuote{奉主名而来的，应受赞美！}同样，当时也有人会说：\BibleQuote{他要消灭那些恶人，把葡萄园另租给别的佃户。}此外，这句话也可以正确地被理解为是主自己的声音，或基于祂本身所是的真理，或基于祂身体的肢体与头的合一。当其他的人给出那种回答时，也有某些人在场对他们说：\BibleQuote{千万不要这样！}因为他们明白这个比喻是针对他们说的。\par

\SourceRecord{Translated by S.D.F. Salmond.；Nicene and Post-Nicene Fathers, First Series；Vol. 6.；Edited by Philip Schaff.；Buffalo, NY: Christian Literature Publishing Co.,；1888.；<http://www.newadvent.org/fathers/1602270.htm>.}

% structure: p0001 source=h1 target=section reason=page-title

\section{福音合参，卷二}

% structure: p0002 source=h2 target=subsection reason=relative-heading-rank; base=h2

\subsection{第七十一章　论婚宴比喻，即大众被邀赴王子婚宴；并论玛窦引入此段之顺序，以及路加在另一上下文中记载的相似叙述}

139. 玛窦接着记述：\Quote{当司祭长和法利塞人听了祂的比喻，就看出祂是指着他们说的；他们想要下手拿祂，却怕群众，因为群众以祂为先知。耶稣又用比喻对他们说：天国好比一个国王为他的儿子摆设婚宴，派仆人去叫那些被请的人来赴宴，他们却不肯来。} 如此下去，直到这些话：\Quote{因为被召的人多，被选的人少。} 这个关于被邀赴宴宾客的比喻，只有玛窦记述。路加也记载了类似的事，但那实际上是一段不同的经文，正如顺序本身充分表明的那样，尽管两者有些相似。然而，玛窦在葡萄园的比喻和杀害家主之子之后立即引入的内容——即犹太人明白这一整篇言论是针对他们的，并开始策划对祂的阴谋——同样由马尔谷和路加证实，他们也以同样的顺序插入这些内容。但在这一段之后，他们转而讨论另一个主题，并立即附上一段玛窦也确实按顺序引入的经文，只是在他的记述中，这段经文紧接在这个唯独他记载于此的婚宴比喻之后。\par

\SourceRecord{Translated by S.D.F. Salmond.；Nicene and Post-Nicene Fathers, First Series；Vol. 6.；Edited by Philip Schaff.；Buffalo, NY: Christian Literature Publishing Co.,；1888.；<http://www.newadvent.org/fathers/1602271.htm>.}

% structure: p0001 source=h1 target=section reason=page-title

\section{福音的和谐，第二卷}

% structure: p0002 source=h2 target=subsection reason=relative-heading-rank; base=h2

\subsection{第72章 论这三部福音作者叙述关于将带有凯撒像的钱币交给凯撒的义务以及关于曾先后嫁给七位兄弟的那位妇人所记叙文之间的和谐。}

140. 玛窦接着这样叙述：\Quote{那时，法利塞人就去商议，怎样在言语上陷害祂。他们派遣自己的门徒同黑落德党人到祂那里，说：\InnerQuote{师傅，我们知道你是真实的，并且真实地教导天主的道，也不顾忌任何人，因为你不看人的外貌。请告诉我们，你以为如何？纳税给凯撒，可以不可以？}}等等，直到下列这些话：\BibleQuote{民众听了，就都惊奇祂的教训。} 马尔谷和路加对这两处主的答复也有类似的记述——即关于钱币的那一处，是由纳税给凯撒的问题所引发的；以及关于复活的那另一处，是由曾先后嫁给七兄弟的那位妇人的事例所引出的。这两位福音作者在次序上也没有分歧。因为在讲了葡萄园租户的比喻——那比喻也涉及犹太人（针对他们所说）以及他们策划的恶谋（这三段是所有三位福音作者共同记述的）之后，马尔谷和路加这两部福音跳过了那宴请宾客赴婚宴的比喻（只有玛窦插入），随后他们又与第一位福音作者会合，记述这两段关于凯撒纳税和曾为七位不同丈夫之妻的妇人的经文，以完全相同的顺序插入，其一致性无可置疑。\par

\SourceRecord{Translated by S.D.F. Salmond.；Nicene and Post-Nicene Fathers, First Series；Vol. 6.；Edited by Philip Schaff.；Buffalo, NY: Christian Literature Publishing Co.,；1888.；<http://www.newadvent.org/fathers/1602272.htm>.}

% structure: p0001 source=h1 target=section reason=page-title

\section{福音的和谐，第二卷}

% structure: p0002 source=h2 target=subsection reason=relative-heading-rank; base=h2

\subsection{第73章。论被交托爱天主和爱近人两条诫命的人；并论玛窦和马尔谷所遵循的叙述顺序问题，以及他们与路加之间并无矛盾。}

141. \Person{玛窦}接着以下列词句继续叙述：\BibleQuote{法利塞人听说耶稣使撒杜塞人哑口无言，就聚集起来。其中有一个法学士，试探祂，问祂说：\InnerQuote{师傅，法律中哪条诫命是最大的？}耶稣对他说：\InnerQuote{你应全心、全灵、全意爱上主你的天主。这是第一条也是最大的诫命。第二条与此相似：你应当爱近人如你自己。全部法律和先知都系于这两条诫命。}}\Person{马尔谷}也记录了这件事，而且顺序也相同。\Person{玛窦}说向主发问的人是试探祂，这不应引起任何困难；而\Person{马尔谷}对此只字未提，但在段末告诉我们，主对那回答明智的人说：\BibleQuote{你离天主的国不远了。}因为这很有可能：那人前来试探祂，却因主的回答而醒悟。或者我们无论如何不应将这里的试探理解为贬义，仿佛它是想欺骗对手的诡计；我们宁可认为它是出于谨慎，就像一个想进一步考验不熟悉之人的人的行为。因为这节经文并非毫无意义地写着：\Quote{轻信的人轻浮，必将受损。}\par

142. 另一方面，\Person{路加}虽然不在此顺序中，而是在一个很不相同的上下文里，引入了一件类似的事。但究竟在那段经文中，他确实记录的是同一事件，还是与\Person{主}（据载）以类似方式谈论那两条诫命的另一个人，完全是无法确定的。同时，认为\Person{路加}所介绍的人是我们此处所论及的不同的人，可能更有理由，不仅因为叙述顺序的显著不同，而且因为\Person{路加}记载那人回答了主向他提出的问题，并在回答中自己提到了那两条诫命。同一看法进一步得到证实：在叙述主对他说\BibleQuote{你这样做，必得生活}——即教导他做那大事，据他自己的回答，这些事情包含在法律中——之后，福音作者接着写道：\BibleQuote{那人却愿意显示自己成义，便对\Person{耶稣}说：\InnerQuote{谁是我的近人？}}随后（据\Person{路加}记载），主讲述了那个从\Place{耶路撒冷}下到\Place{耶里哥}，落入强盗手中的人的故事。因此，考虑到那个人一开始被描述为试探\Person{基督}，并在回答中重复了那两条诫命；进一步考虑到，在主用\BibleQuote{你这样做，必得生活}的话语给予忠告之后，他并未被称赞为好人，相反，关于他写了这样一句话：\BibleQuote{那人却愿意显示自己成义}等等；而\Person{马尔谷}和\Person{路加}在平行叙述中提到的另一个人却得到了如此显著的称赞，以致主对他说：\BibleQuote{你离天主的国不远了}——那么，更可能的观点是，出现在那场合下的人是不同于我们这里所提及的另一个人。\par

\SourceRecord{Translated by S.D.F. Salmond.；Nicene and Post-Nicene Fathers, First Series；Vol. 6.；Edited by Philip Schaff.；Buffalo, NY: Christian Literature Publishing Co.,；1888.；<http://www.newadvent.org/fathers/1602273.htm>.}

% structure: p0001 source=h1 target=section reason=page-title

\section{\WorkTitle{福音合参}第二卷}

% structure: p0002 source=h2 target=subsection reason=relative-heading-rank; base=h2

\subsection{第七十四章。论犹太人被问及他们认为基督是谁之子的一段经文；以及关于玛窦与其他两位福音作者之间是否有不一致的问题，因为他将问题表述为：\BibleQuote{你们对基督怎么看？祂是谁的儿子？}并告诉我们他们回答：\BibleQuote{达味之子；}而其他作者则这样写：\BibleQuote{经师们怎么说基督是达味之子？}}

143. 玛窦接着这样记着：\BibleQuote{那时，法利塞人聚集在一起，耶稣问他们说：你们对基督怎么看？祂是谁的儿子？他们对祂说：达味之子。祂对他们说：那么，达味因圣神怎么称祂为主，说：上主对我主说：你坐在我右边，等我使你的仇敌作你的脚凳？达味既称祂为主，祂怎么又是他的儿子？没有人能回答祂一句话，从那一天起，也没有人敢再问祂什么。}这事马尔谷也按着次序记载。路加则省略了问主关于法律中哪条是第一条诫命的那个人，并且略过那件事后，也再次与其他两位相同次序记载，正如他们一样，记载了主问犹太人关于基督的问题，即祂怎么是达味之子。即使玛窦的记载是：耶稣问他们他们对基督怎么看，祂是谁的儿子，他们（法利塞人）回答说：\BibleQuote{达味之子}，然后祂继续问达味怎么又称祂为主，这一点也不影响意义。而根据其他两位——马尔谷和路加——的版本，我们既没有发现这些人直接被审问，也没有发现他们作了任何回答。因为我们应当持这样的看法：即这两位福音作者引入了主本人在那些人回答之后所说的话，并记录了祂在那些人面前所说的话，祂希望他们有用地受教于祂的权威，并远离经师的教训，而他们对基督的认识当时仅限于：祂按肉身是达味的后裔，而他们不明白祂是天主，并因此也是达味的主。因此，在这两位福音作者的记载中，主被提到的方式，好像祂是在谈论这些错谬的教师，向那些祂希望他们摆脱这些经师所陷入错误的人讲话。同样，玛窦所表达的问题：\BibleQuote{你们怎么说？}应被理解为不是直接对这些（法利塞人）说的，而是仅仅针对那些人，实际上是对祂希望教导的人说的。\par

\SourceRecord{Translated by S.D.F. Salmond.；Nicene and Post-Nicene Fathers, First Series；Vol. 6.；Edited by Philip Schaff.；Buffalo, NY: Christian Literature Publishing Co.,；1888.；<http://www.newadvent.org/fathers/1602274.htm>.}

% structure: p0001 source=h1 target=section reason=page-title

\section{\WorkTitle{福音的和谐，卷二}}

% structure: p0002 source=h2 target=subsection reason=relative-heading-rank; base=h2

\subsection{第75章 关于坐在梅瑟座位上的法利塞人，他们命令别人做自己不做的事，以及主对这些法利塞人所说的其他话语；论及\BibleBook{玛窦}{福音}的记述是否与另两位福音作者的记述一致，尤其是与\BibleBook{路加}{福音}的记述一致，因为路加引入了一段与此相似的话，虽然它不是按此顺序，而是放在另一上下文中。}

144. 玛窦继续他的记述，按以下顺序叙述：\BibleQuote{那时耶稣对群众和祂的门徒说：经师和法利塞人坐在梅瑟的座位上；因此，凡他们所吩咐你们遵守的，你们都要遵守并实行；但不要效法他们的行为，因为他们能说不能行。}一直到最后，\BibleQuote{从今以后，你们不得再见我，直到你们说：奉主名而来的那一位，是应当称颂的。}路加也提到一段类似的主反对法利塞人、经师和法学士的言论，但说它是在某法利塞人家里说的，那人曾邀请祂赴宴。为了叙述那段话，路加已从玛窦所遵循的顺序中插入了关于约纳史迹中三天三夜征兆、南方女王以及回来发现房屋打扫干净的不洁之邪魔的段落。紧接着那一段落，玛窦写道：\BibleQuote{耶稣还对众人说话的时候，看哪，祂的母亲和祂的兄弟们站在外面，想与祂说话。}然而，在第三部福音所呈现的主当时所说的那段话的版本，在叙述了玛窦未记载的主的一些话语之后，路加离开了与玛窦一致所坚持的顺序，以致他紧接着的叙述是这样的：\BibleQuote{祂说话的时候，有一个法利塞人请祂与他共餐。祂进去坐席。那法利塞人看见，诧异祂没有在饭前先行洗净。主对他说：如今你们法利塞人洗净杯盘的外面。}此后，路加报告了针对这些法利塞人、经师和法学士的其他类似言论，这些言论与玛窦在这段我们目前所讨论的段落中也叙述的言论是相似的。因此，尽管玛窦记录这些事的方式确实没有提及那法利塞人的家，但也没有指定任何完全排除祂曾在那家里的地点作为说话场所；然而，主在此时（\Emph{即按\BibleBook{玛窦}{福音}而言}）已离开加里肋亚来到耶路撒冷，并且从祂到达后直到现在所讨论的这段言论之间的上述事件，在上下文中都被如此安排，以致只有合理地理解为它们发生在耶路撒冷，而\BibleBook{路加}{福音}的叙述所涉及的是发生之时主仍在往耶路撒冷途中前往的事。这些考虑使我得出以下结论：这不是同一篇，而是两篇相似的言论，第一位福音作者报告了其中一篇，另一位报告了另一篇。\par

145. 还有一件事需要考量，那就是此处如何说：\BibleQuote{从今以后，你们不得再见我，直到你们说：\InnerQuote{奉主名而来的当受赞美}}，而按同一\Person{玛窦}的记载，他们此前已表达了此意。此外，\Person{路加}亦告诉我们，主曾将包含同样话语的答复，给予那些劝祂离开他们地区的人，因为\Person{黑落德}要杀害祂。这位福音作者记载了这些完全相同的话，即\Person{玛窦}在此处所记录的，是主在那一场合对\Place{耶路撒冷}本身所作的宣告中所用的。因为\Person{路加}的叙述继续如下：\BibleQuote{同日，有些法利塞人前来，对祂说：\InnerQuote{你离开此地，往别处去吧，因为\Person{黑落德}想要杀你。}祂对他们说：\InnerQuote{你们去告诉那只狐狸：看，我今天明天驱魔治病，第三天就成全了。但是，我今天明天后天必须前行，因为先知不能在\Place{耶路撒冷}之外丧命。\Place{耶路撒冷}，\Place{耶路撒冷}，你常杀害先知，用石头砸死那些奉派到你这里来的人；我多少次愿意聚集你的子女，如同母鸡把自己的雏鸡聚集在翅膀底下，你们却不愿意！看，你们的殿宇将留给你们荒凉。我告诉你们：从今以后，你们不得再见我，直到你们说：\InnerQuote{奉主名而来的当受赞美。}}}然而，在群众于主临近\Place{耶路撒冷}时说\BibleQuote{奉主名而来的当受赞美}的情况中，似乎与\Person{路加}所给的叙述并无矛盾。因为，按\Person{路加}所遵循的顺序，祂尚未到达所论及的场所，那话也还未被说出。但因为\Person{路加}并未告诉我们，祂当时确实离开了那地方，且不再回去，直到那话被他们说出的时刻（因为祂继续行程，直到抵达\Place{耶路撒冷}；并且那话，\BibleQuote{看，我今天明天驱魔治病，第三天就成全了}应被视为祂以奥秘和比喻意义所说：因为祂当然没有在此时之后字面意义上的第三天受难；相反，祂紧接着说：\BibleQuote{然而，我今天明天后天必须前行}），我们确实被迫也对\BibleQuote{从今以后，你们不得再见我，直到你们说：\InnerQuote{奉主名而来的当受赞美}}这句话给予奥秘的解释，并视它指向祂那要带着灿烂光辉来临的降临；由此也暗示，祂在\BibleQuote{我今天明天驱魔治病，第三天就成全了}这个声明中所表达的，关乎祂的身体，即教会。因为当外邦人放弃他们祖传的迷信并信从祂时，魔鬼就被逐出；而当人弃绝魔鬼和世俗，并按祂的诫命生活，直到复活的完成，在那时，就如同一首诗歌，将实现第三天的那种成全；也就是说，教会将在身体与灵魂的实在不朽中，达到天使圆满的尺度。因此，\Person{玛窦}所遵循的顺序绝不应被理解为涉及转向另一关联。但我们宁可认为，要么是\Person{路加}提前了发生在\Place{耶路撒冷}的事件，并将它们在此处引入，仅因它们在此刻被他的记忆所提示，而他的叙述实际上尚未将主带到\Place{耶路撒冷}；要么是主在那一场合临近同一城市时，确实以相似的话答复了那些劝祂提防\Person{黑落德}的人，正如\Person{玛窦}所记载的，祂也在祂已抵达\Place{耶路撒冷}、上述所有事件发生之后，对群众说了同样的话。\par

\SourceRecord{Translated by S.D.F. Salmond.；Nicene and Post-Nicene Fathers, First Series；Vol. 6.；Edited by Philip Schaff.；Buffalo, NY: Christian Literature Publishing Co.,；1888.；<http://www.newadvent.org/fathers/1602275.htm>.}

% structure: p0001 source=h1 target=section reason=page-title

\section{\WorkTitle{福音的和谐，第二卷}}

% structure: p0002 source=h2 target=subsection reason=relative-heading-rank; base=h2

\subsection{第七十六章：论在叙述顺序上玛窦与其它两位福音作者记载主预言圣殿毁灭事件时的和谐}

146. 玛窦接着以下列词句继续他的记叙：\BibleQuote{耶稣出了圣殿，正走的时候，祂的门徒前来，把圣殿的建筑指给祂看。耶稣对他们说：\InnerQuote{你们不是看见这一切吗？我实在告诉你们：将来这里决没有一块石头留在另一块石头上，而不被拆毁的。}} 这件事玛尔谷也记载了，而且顺序大致相同。但他是在一段较小的插叙之后引入的，这段插叙是为了提及那位向银库投了两个小钱的寡妇的事迹——这件事只有玛尔谷和路加记载。因为（为证明玛尔谷的顺序本质上与玛窦相同，我们只需注意）在玛尔谷的版本中，在记述主与犹太人讨论祂如何问他们怎样以为基督是达味之子后，接着叙述了祂警告他们提防法利塞人及其伪善的话——这部分玛窦已作了最详尽的呈现，引入了主在那场合的大量言论。然后，在这段被玛尔谷简短处理、被玛窦极为详细地处理过的段落之后，正如我所说，玛尔谷引进了关于那位极其贫穷却又如此丰盛的寡妇的段落。最后，他没有插入任何别的东西，就续接了一段与玛窦再次一致的记载——即关于圣殿毁灭的记载。同样，路加先陈述了关于基督如何是达味之子的提问，然后提到了几句警告他们提防法利塞人伪善的话。此后，他像玛尔谷一样，讲述了那位向银库投了两个小钱的寡妇的故事。最后，他附上了关于圣殿注定倾覆的陈述，这也出现在玛窦和玛尔谷的记载中。\par

\SourceRecord{Translated by S.D.F. Salmond.；Nicene and Post-Nicene Fathers, First Series；Vol. 6.；Edited by Philip Schaff.；Buffalo, NY: Christian Literature Publishing Co.,；1888.；<http://www.newadvent.org/fathers/1602276.htm>.}

% structure: p0001 source=h1 target=section reason=page-title

\section{福音的和谐，第二卷}

% structure: p0002 source=h2 target=subsection reason=relative-heading-rank; base=h2

\subsection{第77章 论三位福音圣史关于主在橄榄山上与门徒论及何时终结的叙述中的和谐一致}

147. 玛窦接着叙述道：\BibleQuote{耶稣在橄榄山上坐着，门徒们私下前来对他说：\InnerQuote{请告诉我们，什么时候要发生这些事？又什么是你来临和世界终结的记号？}耶稣回答他们说：\InnerQuote{你们要谨慎，免得有人欺骗你们；因为将有许多人假冒我的名来说：\InnerQuote{我是基督}，并要欺骗许多人；}}等等，直到我们读到：\BibleQuote{这些人要进入永罚，而义人却要进入永生。}我们现在必须研究这篇长篇言论，正如它在三位圣史玛窦、马尔谷和路加中所呈现的那样。因为他们都在自己的叙述中引入了这篇言论，并且保持了同样的顺序。此处，与别处一样，每位作者都讲述了一些他特有的内容，但我们在其中不必担心任何不一致之处。但我们要确认的是，在那些完全平行的段落中，他们绝不能被看作相互矛盾。因为如果在这里遇到任何看似矛盾的内容，简单地断言这是主在类似措辞但完全不同场合所说的完全不同的话，在这种情况下不能被视为合理的解释方式，因为三位圣史所给出的叙述在主题和时间顺序上都保持一致。此外，作者们在报告主所说的相同思想时并非都遵循同样的顺序，这一事实绝不会影响对主题本身的理解或传达，只要他们代表主所说的话彼此之间没有矛盾。\par

148. 同样，玛窦以这种形式陈述的：\BibleQuote{这天国的福音必先传遍天下，为给万民作证，然后终结才来到。}，马尔谷在同一段中以以下方式记载：\BibleQuote{福音必须先传给万民。}马尔谷没有加上\Quote{然后终结才来到}这句话，但当他在句子中使用\Quote{先}这个词时，就暗示了其含义：\BibleQuote{福音必须先传给万民。}因为他们曾问他关于终结的事。因此，当他这样对他们说：\BibleQuote{福音必须先传给万民}时，\Quote{先}这个词清楚地暗示了在终结来临之前必须完成某事。\par

149. 同样，玛窦这样记载的：\BibleQuote{当你们看见达尼尔先知所说的\InnerQuote{招致荒凉的可憎之物}站在圣地时——诵读的应明了——}，马尔谷则以以下形式表达：\BibleQuote{但当你们看见那招致荒凉的可憎之物站在不应该站的地方时——诵读的应明了。}尽管措辞有所改变，但所传达的意义是相同的。因为\Quote{不该在的地方}其含义就是那招致荒凉的可憎之物不应站在圣地。路加的表达方式则既不是\Quote{当你们看见那招致荒凉的可憎之物站在圣地}，也不是\Quote{不该在的地方}，而是\BibleQuote{当你们看见耶路撒冷被军队围困时，那时你们就知道它的荒凉近了。}因此，那时那招致荒凉的可憎之物将站在圣地。\par

150. 再者，\Person{玛窦}以下列语句所记述的：\BibleQuote{那时，在犹大的，应当逃到山上；在屋顶上的，不要下来取家里的东西；在田里的，也不要回去取衣裳。} \Person{玛尔谷}也几乎用同样的话作了报道。另一方面，\Person{路加}的版本则这样进行：\BibleQuote{那时，在犹大的，应当逃到山上。} 到此为止，他与另外两位一致。但他以不同的形式呈现了随后的内容。因为他接着说：\BibleQuote{在城中的，应当出来；在乡下的，不要进城；因为这是报复的日子，为要应验所写的一切。} 这些陈述彼此之间似乎显示出足够的差异。因为正如前两位福音作者所记载的那样：\BibleQuote{在屋顶上的，不要下来取家里的东西；} 而第三位福音作者所给出的却是：\BibleQuote{在城中的，应当出来。} 然而，其含义可能是：在如此巨大的迫在眉睫的危险面前，所引发的巨大骚动中，那些被困在围城状态中的人（即用\BibleQuote{在城中的}这个短语所表达的）会出现在屋顶（或\Quote{城墙}）上，惊惶不安地想知道什么恐怖笼罩着他们，或者有什么逃脱的方法可能出现。但仍然存在一个问题：这第三位福音作者在这里说\BibleQuote{应当出来}，而他之前已经用了这些词语：\BibleQuote{当你们看见耶路撒冷被军队围困的时候}？因为此后引入的——即这句话：\BibleQuote{在乡下的，不要进城}——似乎构成了一致劝告的一部分；我们可以理解城外的人如何不应进城；但困难在于，当城市已被军队围困时，城中的人如何出来。或许，这个\BibleQuote{在城中}的表达指示了一个时刻，危险如此紧迫，以致在现世的手段方面，没有机会为保存此世的身体生命留有余地；而灵魂应该准备好、自由，不被肉身的欲望牵挂或束缚，这个事实是由前两位作者所用的短语——即\BibleQuote{在屋顶上}或\Quote{在城墙上}——传达的？这样，第三位福音作者的措辞\BibleQuote{应当出来}（实际上意味着，不再沉迷于此生的欲望，而是准备好进入另一生命）在意义上等同于其他两位所用的语句：不要下来取家里的东西（实际上意味着，\InnerQuote{不要让他的情感转向肉身，好像肉身那时能给他任何好处一样}）。同样，第三位所采用的短语\InnerQuote{在乡下的，不要进城}（也就是说，\InnerQuote{不要让那些已怀着良好心意置身于城外的人，再沉溺于任何肉欲或对它的渴望}）正表示前两位福音作者在句子\InnerQuote{在田里的，也不要回去取衣裳}中所体现的内容，这差不多等于说，他不应再次卷入他已摆脱的忧虑之中。\par

151. 此外，玛窦继续写道：\BibleQuote{但你们要祈祷，叫你们的逃走不是在冬天，也不是在安息日。} 马尔谷给出了其中的一部分，省略了另一部分，他说：\BibleQuote{你们要祈祷，叫你们的逃走不是在冬天。} 路加却完全省略了这一点，取而代之的是他特有的一段话，在我看来，他借此光照了这句由前两位作者以稍显晦涩的方式摆在我们面前的经文。因为他的版本是这样：\BibleQuote{你们要谨慎，免得你们的心被宴饮、醉酒及今生的挂虑所累，那日子就忽然临到你们。因为那日子要如同网罗般地临到全地面上一切居住的人。所以你们要时时醒寤祈祷，好使你们能配得逃脱将要发生的一切事。} 这应被理解为与玛窦所提的同一个逃亡，不应在冬天或安息日进行。此外，那个\Quote{冬天}是指路加直接指明了的这些\Quote{今生的挂虑}；同样地，\Quote{安息日}是指\Quote{宴饮和醉酒}。因为悲伤的挂虑好比冬天；而宴饮和醉酒将心灵淹没并葬送于肉体的欢愉和奢侈之中——这一罪恶用\Quote{安息日}一词来表达，是因为古时，正如他们现在仍如此，犹太人在那日子有放纵享乐的非常有害的习惯，那时他们并不了解属神的安息。或者，若在玛窦和马尔谷中出现的这些话是指别的事，路加的话也可能被当作指向别的事，尽管如此，不需要提出任何暗示它们矛盾的问题。然而，目前我们并未承担诠释福音的任务，而只是为它们辩护，以驳斥对其虚假和欺骗的无端指控。此外，玛窦插入此次讲论中的其他事，即他与马尔谷共享的内容，并不带来困难。另一方面，关于那些他与路加共享的段落，值得注意的是，路加并未将它们引入当前的讲论中，尽管在叙事次序上这里是一致的。但他记录了类似语句的其他联系，要么是凭记忆复述它们，从而提前引入，以便在较早的节点讲述实际上属于较晚的主所说的话；要么是让我们明白，这些话是主两次说出的，一次是在玛窦所指的场合，第二次是路加自己所处理的场合。\par

\SourceRecord{Translated by S.D.F. Salmond.；Nicene and Post-Nicene Fathers, First Series；Vol. 6.；Edited by Philip Schaff.；Buffalo, NY: Christian Literature Publishing Co.,；1888.；<http://www.newadvent.org/fathers/1602277.htm>.}

% structure: p0001 source=h1 target=section reason=page-title

\section{福音的和谐，第二卷}

% structure: p0002 source=h2 target=subsection reason=relative-heading-rank; base=h2

\subsection{第七十八章。论玛窦和马尔谷一方面与若望另一方面之间是否有矛盾：前者说两天后将是逾越节，随后告诉我们他在伯达尼，而后者给出在伯达尼发生的平行叙述，但提到那是逾越节前六天。}

152. 玛窦继续写道：\BibleQuote{耶稣讲完了这一切话，就对门徒说：你们知道，两天后就是逾越节，人子将要被交出去，钉在十字架上。} 另外两位——即马尔谷和路加——也同样证实了这一点，而且在叙述顺序上完全一致。然而，他们并没有将这句话当作主亲口所说的。他们没有这样声明。同时，马尔谷以自己的口吻告诉我们：\BibleQuote{两天后就是逾越节和无酵饼节。} 路加也这样肯定：\BibleQuote{无酵节近了，名叫逾越节；} 也就是说，它\Quote{近了}的意思，是指再过两天就要发生，正如其他两位更明显地一致表达的那样。另一方面，若望在三处提到了同一个节日的临近。前两次提及是在他记载某些不同性质的事情时。但第三次，他的叙述显然涉及那些时间——也就是其他三位福音作者也提到这一主题的时间——即主的受难实际迫在眉睫的时候。\par

153. 但对于那些不够仔细考察此事的人而言，似乎存在矛盾：玛窦和马尔谷在述说过两天就是逾越节之后，立刻告诉我们耶稣当时在伯达尼，并引出了珍贵香膏的记述；然而若望，在正要给我们关于香膏的相同叙述时，却一开始就告诉我们耶稣在逾越节前六天来到伯达尼。问题是，那两位福音作者如何能说过两天就要庆祝逾越节，而我们却发现他们在做出这一陈述后，紧接着又与若望一道，给我们记述了伯达尼的香膏场景；而在那同一记述中，后一位作者告诉我们，逾越节庆典在六天后才举行。然而，被这一难题困扰的人，只是没有看出玛窦和马尔谷引入伯达尼场景的记述，实际上是采用了\OriginalTerm{复述}{recapitulation}的形式，仿佛其发生的时间是在他们作出两日间隔的宣告之后，而事实上，那是在逾越节前六天时已发生的事件。因为他们两人中，没有一人将伯达尼所发生之事的记述，以这样的措辞续写在关于两日间隔后庆祝逾越节的陈述之后：\BibleQuote{这些事以后，当祂在伯达尼时。} 玛窦的表述是：\BibleQuote{当耶稣在伯达尼时。} 马尔谷的版本则是：\BibleQuote{且当祂在伯达尼时，}等等；这种表达方式，无疑可以被理解为指涉逾越节前两日所说之话之前的时间。因此，情况如下：正如我们从若望的叙述中得知，耶稣在逾越节前六天来到伯达尼；那里举行了晚餐，与此相关我们得到了珍贵香膏的记述；祂离开那里后，骑着一匹驴驹来到耶路撒冷；此后发生了那些事，就是他们述说的祂到达耶路撒冷之后的事。因此，尽管福音作者没有提及这一事实，我们明白：从祂来到伯达尼并目睹香膏场景的那一天，到目前这些言行所属的那一天之间，相隔了四天，以至于在那一点上，可以插入那两位福音作者用他们关于两日后庆祝逾越节的陈述所定义的那一天。此外，当路加说：\BibleQuote{无酵节近了}，他确实没有明确提及两日间隔；但他所指的临近，仍应视为由这二日间隔所证实。同样，当若望说：\BibleQuote{犹太人的逾越节近了}，他并非意指两日间隔，而是指逾越节前还有六天的时间。因此，他在此肯定之后立即记录某些事情，旨在指明他称逾越节为近时所意指的临近程度，接着他以如下言语继续：\BibleQuote{逾越节前六天，耶稣来到伯达尼，那里有拉匝禄，就是祂从死者中复活的；在那里他们为祂摆设了晚餐。} 这就是玛窦和马尔谷在陈述过两天就是逾越节之后，以复述形式引入的事件。在他们的复述中，他们回到了伯达尼的那一天，那时距逾越节还有六天，并向我们提供了若望也提到的晚餐和香膏的记述。在那场景之后，我们应假设祂来到耶路撒冷，然后，在记述的其他事情发生后，到达这一天，这一天距逾越节还有两天；这些福音作者正是从这一天做了这个题外回顾，目的是为了对伯达尼的香膏事件做一个复述式的说明。在那段叙述完成后，他们再次回到做出题外回顾的起点；也就是说，他们现在继续记录主在逾越节前两天所说的话。因为如果我们移除关于伯达尼事件的记述——他们以题外方式引入此记述，背离了时间顺序，并以回忆和复述的形式插在实际历史位置之后——然后我们将叙述按常规顺序排列，那么叙述将如下进行——根据玛窦，主的话这样接续：\BibleQuote{你们知道，两天后就是逾越节，人子要被交付，钉在十字架上。那时，司祭长和民间长老聚集在大司祭名叫盖法的庭院里，商议如何用诡计捉拿耶稣，并杀害祂。但他们说：不可在庆节期间，以免民间发生骚动。然后，十二人中的一个，名叫犹达斯依斯加略的，去见司祭长，}等等。因为在\BibleQuote{以免民间发生骚动}这句话和\BibleQuote{那时，一个门徒名叫犹达斯的，去了}等语之间，插入着那伯达尼场景的记述，那是他们以复述方式引入的。因此，通过省略它，我们在叙述中建立了这样一种关联，足以使我们的结论令人满意，即在此处时间顺序上不存在矛盾。同样，如果我们以类似方式处理马尔谷福音，并省略同样发生在伯达尼的晚餐记述——他也以复述形式引入了它——那么他的叙述将按以下顺序进行：\BibleQuote{两天后就是逾越节和无酵节：司祭长和经师们设法如何用诡计捉拿祂，把祂处死。因为他们说：不可在庆节期间，以免民间发生骚动。犹达斯依斯加略，十二人中的一个，去见司祭长，要把祂交出。} 在这里，这些福音作者以复述方式插入的伯达尼事件，又被置于\BibleQuote{以免民间发生骚动}这一分句与我们紧接着附上的\BibleQuote{犹达斯依斯加略，十二人中的一个}这一节之间。另一方面，路加却完全省略了上述伯达尼事件。这就是我们针对逾越节前六天所作的说明——那是若望在叙述伯达尼所发生之事时提及的间隔；以及针对逾越节前两天所作的说明——那是玛窦和马尔谷在陈述之后，紧接着呈现他们对伯达尼同一场景（也为若望所记载）的记述时所指定的时期。\par

\SourceRecord{Translated by S.D.F. Salmond.；Nicene and Post-Nicene Fathers, First Series；Vol. 6.；Edited by Philip Schaff.；Buffalo, NY: Christian Literature Publishing Co.,；1888.；<http://www.newadvent.org/fathers/1602278.htm>.}

% structure: p0001 source=h1 target=section reason=page-title

\section{福音的和谐，第二卷}

% structure: p0002 source=h2 target=subsection reason=relative-heading-rank; base=h2

\subsection{第七十九章 论玛窦、马尔谷与若望在叙述伯达尼晚餐中妇人以珍贵香液倾注于主的事件上如何协调一致，以及这些记载与路加在另一时期记载的类似事件如何调和的方法。}

154. 玛窦接着从他完成审查的那个点继续叙述，用了以下的话：\BibleQuote{那时，司祭长和民间的长老聚集在大司祭的院子里，那大司祭名叫盖法，商议要用诡计捉拿耶稣，并杀害祂；但他们说：不要在过节的日子，恐怕民间生乱。当耶稣在伯达尼癞病人西满家里坐席的时候，有一个女人拿着一玉瓶贵重的香液来到祂跟前，倒在祂的头上。} 一直到这句话：\BibleQuote{这女人所做的，也要传扬，作为对她的纪念。} 我们现在要考虑伯达尼那个妇人和贵重香液的事件，就如它被详细记载的那样。因为虽然路加也记载了一个类似的事件，而且他指定耶稣晚餐所在之人的名字也可能暗示两个叙述是同一件事（因为路加也称主人为\Quote{西满}），但是，无论是自然情况还是人的习俗，都没有使这件事变得不可信——就像一个人可以有两个名字，两个人更有可能有同一个名字一样——所以更合理的看法是：在伯达尼发生此事的这个西满的家（按照路加的版本，人们是这样猜测的），是与玛窦所提到的西满不同的一个人。此外，路加并没有指明他所记载的事件发生在伯达尼。尽管他确实没有具体说明那件事发生在哪个城镇或村庄，但他的叙述似乎并不涉及同一个地点。因此，我的意见是：对这件事只能有一种解释。然而，这并不是要认为玛窦记载的那个妇人，与那天以罪妇身份来到耶稣脚前，亲吻祂的脚，用眼泪洗祂的脚，用头发擦干，并抹上香液的那个妇人，是完全不同的一个人。关于这件事，耶稣也曾用两个债户的比喻，说她那许多的罪得了赦免，因为她爱得多。我的看法是：同一个玛利亚在两次不同的场合做了这事——一次是路加所记载的，那时她初次以极大的谦卑和眼泪来到祂面前，获得了罪的赦免。因为若望虽然不像路加那样详细叙述了那件事的情形，但至少提到了这个事实，在开始讲述拉匝禄复活的故事、叙述还未将主带到伯达尼本身之前，就称赞了这位玛利亚。他对那件事的记载如下：\BibleQuote{有一个病人，名叫拉匝禄，住在伯达尼，就是玛利亚和她的姐姐玛尔大所住的村庄。这玛利亚就是用香液抹了主，又用头发擦祂的脚的那个女人；她的兄弟拉匝禄病了。} 若望以此证实了路加所告诉我们的，即路加记载了在一个名叫西满的法利塞人家里发生的这类事件。因此，我们看到玛利亚在此之前曾这样做过。她在伯达尼第二次所做的，是另一件事，不属于路加的叙述，而是由三位福音作者共同记载的，即若望、玛窦和马尔谷。\par

155. 因此，让我们注意这三福音作者——\Person{玛窦}、\Person{马尔谷}和\Person{若望}——之间如何保持和谐；无疑他们都记载了在\Place{伯达尼}发生的同一事件，当时门徒也正如三人所记，表面上因那极珍贵的香膏被浪费而对那女人抱怨。如今，\Person{玛窦}和\Person{马尔谷}告诉我们香膏是倒在主的头上，而\Person{若望}说是倒在祂的脚上，这一点可证明并无矛盾，只要我们应用我们在处理五饼饱众场景时已阐述的原则。因为正如有一位作者在记述那件事时，并未忽略说明众人立即按五十人或一百人坐下，而另一位只提到五十人，这并不能认为出现矛盾。确实，如果一位福音作者只提到一百人，另一位只提到五十人，那似乎有些困难；但即使如此，正确的结论应是他们既按五十人也按一百人坐下。这个例子应当使我们明白，正如我在讨论那段经文时向读者强调的，即使几位福音作者各自只引入一个事实，我们也应认为实际情况是两者都发生在事件中。同样，关于我们现在这段经文，我们的结论应是：那女人不仅把香膏倒在主的头上，也倒在祂的脚上。诚然，可能会有人荒谬而狡诈地争辩说，因为\Person{马尔谷}说香膏是在打破玉瓶后才倒出的，破碎的器皿中不可能剩下任何东西用来抹祂的脚。但这类人为了否认福音的真实性，会争辩说玉瓶被打破的方式使其中的任何部分都不可留；而一个完全不同的人，其立场则更好、更合虔敬，他的目的是维护福音的真实性，因此可以争辩说，玉瓶并未以使香膏全部倒出的方式被打破！再者，如果那诽谤者如此固执盲目，以致试图在打碎玉瓶这件事上破坏福音作者的和谐，他倒不如接受另一种可能：主的脚在那玉瓶本身被打碎之前已被抹膏，如此玉瓶仍完整，且盛满足够的香膏以涂抹头部，然后通过所说的打破，全部内容物被倒出。因为我们承认，当行为从头部开始时，对我们本性的各部分有适当照顾，但（我们也可以说）当我们从脚上升到头部时，同样自然的顺序也被保存。\par

156. 属于这一事件的其他事项，在我看来并不提出任何实际涉及困难的问题。其他福音作者提到门徒因珍贵香膏的浪费而抱怨，而\Person{若望}说\Person{犹达斯}是那个如此发言的人，并解释事实说\Quote{他是个贼}。但我认为很明显，这个\Person{犹达斯}就是门徒这（一般）名称所指的那一位，这里用复数代替单数，符合我们已在\Person{斐理伯}和五饼奇迹的例子中解释过的说话方式。也可以这样理解：其他门徒要么与\Person{犹达斯}感受相同，要么像他那样说话，要么因\Person{犹达斯}所说而被带到那种观点，因此\Person{玛窦}和\Person{马尔谷}用言语表达了整个群体的真实想法；但\Person{犹达斯}那样说话，只是因为他是个贼，而其余门徒的动机是对穷人的关心；此外，\Person{若望}选择只记录门徒中的那一个人发出此类怨言，因为他认为有必要在此场合揭露他惯于偷窃的习性。\par

\SourceRecord{Translated by S.D.F. Salmond.；Nicene and Post-Nicene Fathers, First Series；Vol. 6.；Edited by Philip Schaff.；Buffalo, NY: Christian Literature Publishing Co.,；1888.；<http://www.newadvent.org/fathers/1602279.htm>.}

% structure: p0001 source=h1 target=section reason=page-title

\section{《\WorkTitle{福音的和谐}》第二卷}

% structure: p0002 source=h2 target=subsection reason=relative-heading-rank; base=h2

\subsection{第80章　论\Person{玛窦}、\Person{马尔谷}和\Person{路加}所记载的关于祂派遣门徒为祂吃逾越节做准备之事的和谐。}

157. \Person{玛窦}继续写道：\BibleQuote{那时，十二人中的一个，名叫犹达斯·依斯加略的，去见司祭长，对他们说：你们愿意给我什么，我就把祂交给你们？他们议定给他三十块银钱；}一直到这些话：\BibleQuote{门徒就照耶稣所吩咐的做了，预备了逾越节。}在本段中，没有任何内容被认为与\Person{马尔谷}和\Person{路加}的记述相矛盾，他们以类似的方式记录了这段经文。因为关于\Person{玛窦}以这些措辞所作的陈述：\BibleQuote{你们进城去，到某人那里，对他说：师傅说：我的时候近了，我要在你家里守逾越节，并与我的门徒同席。}它只是指明那个人，就是\Person{马尔谷}和\Person{路加}所称的\Quote{家主}或\Quote{屋主}，在他那里他们被指示了餐厅，他们要在那里预备逾越节。\Person{玛窦}通过简单地引入短语\Quote{到某人那里}，作为他自己所作的简短解释，以便简洁地让我们明白所指的那个人是谁。因为如果他曾说过主以这样的话对他们说：\Quote{你们进城去，对他说〔或\InnerQuote{对它}〕：师傅说，我的时候近了，我要在你家里守逾越节}，那么人们可能会以为这些话打算是对着城市本身说的。因此，为了这个缘故，\Person{玛窦}插入了这句话，即主吩咐他们去\Quote{到某人那里}，但这并不是主所说的话（他正在记录主的指示），而只是他自己主动添加的话，目的是为了避免详细叙述全部内容，因为在他看来，只需要提到这些，就足以准确地表达发出命令的人的真实意图。因为谁能看不到，没有人会自然地以这样不确定的方式对别人说：\Quote{到某人那里去}？再者，如果话是：\Quote{到随便哪个人那里去}，或\Quote{到你愿意的任何一个人那里去}，这样的表达方式也许是足够正确的，但门徒被派去的人就会不确定；而\Person{马尔谷}和\Person{路加}则把他呈现为一个明确指出的个体，尽管他们略过他的名字未提。我们确信，主自己知道他差遣他们去的是哪一个人。并且为了使他这样差遣的人也能发现所指的那个人，他在他们出发前给了他们一个要遵循的特殊标记——即一个拿着一壶或一瓶水的人的形像——并告诉他们，如果他们跟着他，就会到达他所指的房屋。因此，鉴于在这里使用措辞\Quote{到你愿意的任何一个人那里去}是不适当的——尽管就语言规范的要求而言，它确实足够正确，但此处所处理的准确陈述使得它在本段中不可接受——那么，像\Quote{到某人那里去}这样的表达，正确语言用法根本不能承认的，在这里（由说话者自己）使用，又有什么理由呢？但明显的是，门徒被主差遣，显然不是到他们愿意的任何人那里，而是到\Quote{某人}那里，也就是说，到一个明确的个体那里。而这是福音作者以他自己身份说话时，完全可以这样告诉我们的，通过这样表达：\Quote{他派他们到某人那里，为了对他说：\InnerQuote{我要在你家里守逾越节}。}他也可以这样表达：\Quote{他派他们到某人那里，说：去，对他说：\InnerQuote{我要在你家里守逾越节}。}就这样，在给了我们主自己实际说的话之后，即\BibleQuote{你们进城去}，他加上了他自己的补充\Quote{到某人那里}，他这样做，并不是好像主这样表达了自己，而只是为了让们明白，尽管名字未记录，但在城中有特定一个人，主的门徒被派到他那里，以预备逾越节。同样，在这样作为他自己的解释引入的两个〔或三个〕词之后，他又重新接上主自己所说的话的顺序，即\BibleQuote{对他说：师傅说：}现在如果你问他们\Quote{对谁}要说这话，正确的回答〔立刻〕以这些措辞给出：就是那个特定的人，福音作者以他自己身份说话时，引入了短语\Quote{到某人那里}，从而使们明白主差遣他们到那人那里。这样插入的从句确实可能包含相当不寻常的表达方式，但这样理解时它仍然是一个完全合法的措辞。或者，在希伯来语中（据记载\Person{玛窦}是用希伯来语写的），可能有一些特殊用法，使得即使整个句子被视为主自己所说的话，也完全符合正确表达法则。是否如此，懂那语言的人可以决定。即使在拉丁语本身，这种表达也可能被使用，以这样的话为例：\Quote{你们进城去，到那个将会遇到你们、拿着一壶水的人所指示的人那里。}如果指示以这样的话语传达，它们就可以毫无歧义地执行。或者，如果措辞像是这样：\Quote{你们进城去，到某人那里，他住在某个地方、某某房屋}，那么这样给出的地点和房屋的指定就使得理解并执行所交付的使命成为可能。但当这些指示以及其他所有类似指示完全未被说明时，在这种情况下使用这种说法\Quote{到某人那里，并对他说}的人，不可能被理解地听从，原因很明显：当他使用\Quote{到某人那里}这些词时，他意图让们理解某个特定的个体，却没有提供任何线索使我们能认出他。但如果我们假设所指的从句是福音作者自己插入作为解释，〔我们可能会发现〕简练的要求会使表达有些晦涩，但不会使之处不正确。此外，关于这个事实，即\Person{马尔谷}说到一壶水，而\Person{路加}提到一个器皿，简单的解释是：一个用了指示器皿种类的词，另一个用了指示其容量的词，而两位福音作者都保留了实际意图的真正含义。\par

158. 玛窦接着叙述：\Quote{到了晚上，耶稣与十二门徒坐席。他们正吃晚餐的时候，耶稣说：\InnerQuote{我实在告诉你们：你们中有人要出卖我。}他们非常忧闷，开始各自对祂说：\InnerQuote{主，是我吗？}}等等，直到我们读到：\Quote{背叛祂的犹达斯也问说：\InnerQuote{辣彼，是我吗？}耶稣对他说：\InnerQuote{你说的是。}}在我們现在所提交考虑的内容中，其他三位圣史，他们也记录了这些事，没有提供任何可能引起严重困难的问题。\par

\SourceRecord{Translated by S.D.F. Salmond.；Nicene and Post-Nicene Fathers, First Series；Vol. 6.；Edited by Philip Schaff.；Buffalo, NY: Christian Literature Publishing Co.,；1888.；<http://www.newadvent.org/fathers/1602280.htm>.}

% structure: p0001 source=h1 target=section reason=page-title

\section{\WorkTitle{福音的和谐，卷三}}

\EditorialNote{本书从晚餐记述直到福音结尾证明福音作者的和谐，各作者的叙述被校勘，整体按有序连接排列。}\par

% structure: p0003 source=h2 target=subsection reason=relative-heading-rank; base=h2

\subsection{序言}

1. 鉴于我们现在已经到达历史中的那个点，即所有四位福音作者必须一起共同前行直到结尾，彼此之间没有任何严重的分歧，如果在某处其中一位提到某件事而另一位忽略，在我看来，如果我们从现在开始将所有作者的所有陈述一起纳入一个联结，并将整体安排成一个单一的叙述，且在一个视图中，我们就可以更快速地证明各位福音作者所保持的一致性。我认为，通过这种方式，我们所承担的任务可以比原本更便利、更容易地完成。因此，我们目前要做的是尝试构建一个单一的叙述，在其中我们将包括所有细节，并且我们将拥有那些福音作者的证词，他们（从全部事实中各自选择自己有能力或愿望叙述的）为我们准备了这些记录：此外，以这样的方式，所有这些我们需证明完全没有矛盾的陈述，都被当作所有福音作者共同作出的。\par

\SourceRecord{Translated by S.D.F. Salmond.；Nicene and Post-Nicene Fathers, First Series；Vol. 6.；Edited by Philip Schaff.；Buffalo, NY: Christian Literature Publishing Co.,；1888.；<http://www.newadvent.org/fathers/1602300.htm>.}

% structure: p0001 source=h1 target=section reason=page-title

\section{\WorkTitle{福音书的和谐，卷三}}

% structure: p0002 source=h2 target=subsection reason=relative-heading-rank; base=h2

\subsection{第一章 论四位福音传道者在记述主的晚餐及指出出卖者一事上的一致性}

2. 让我们据此开始，先看玛窦所记的记载，[其文如下]：\Quote{当他们正吃的时候，耶稣拿起饼来，祝福了，掰开，递给祂的门徒，说：\InnerQuote{你们拿去吃罢，这是我的身体。}} 马尔谷和路加也记了这一段。确实，路加两次提到杯：第一次是在祂给饼之前；第二次是在给饼之后。但事实是，前面所记的是按这位作者的习惯提前插入的，而他在这里按顺序插入的词语在前面的经节中并未记录，将这两处合起来正好构成其他两位福音传道者所表达的内容。若望则在此上下文中没有提及主的体和血；但他清楚地证明主在另一场合以更详细的叙述说了类似的话。然而此时，在记述了主从晚餐起来为门徒洗脚，并告诉我们主如此行的原因——在表达中祂已经暗示了，尽管仍隐晦，并引用了一段圣经的证言，表明祂正被那要吃祂面包的人出卖——之后，若望来到了这段经文，即其他三位福音传道者也共同引入的。他这样记载：\Quote{耶稣说了这话，心神烦乱，就明说：\InnerQuote{我实实在在告诉你们：你们中有一人要出卖我。}门徒们（如同一若望所加）就面面相觑，猜疑祂说的是谁。} \Quote{（如玛窦和马尔谷所记）他们就非常忧愁，一个一个地开始对祂说：\InnerQuote{是我吗？}祂回答说（如玛窦接着所述）：\InnerQuote{那同我一起把手蘸在盘子里的，那人才要出卖我。}} 玛窦接着又加了以下的话：\Quote{人子固然要照所记载的离去，但出卖人子的那人却是有祸的！那人若没有生下来，为他更好。} 马尔谷在此也与他一致，无论言词还是叙述顺序。然后玛窦接着说：\Quote{那出卖祂的犹达斯回答说：\InnerQuote{辣彼，是我吗？}祂对他说：\InnerQuote{你说的是。}} 这些话并没有明确说他自己就是那人。因为这句话仍可理解为它的意思是：\Quote{说这话的人不是我。} 这一切也完全可能是犹达斯所说的，主回答了他，而其他人都没有注意到。\par

3. 此后，玛窦接着插入祂的体和血的奥迹，即当时主所托付给门徒的。马尔谷和路加也相应如此行。但祂将杯递给它们之后，〔我们发现〕祂再次谈论祂的出卖者，按路加所记的话，他说：\Quote{但看，那出卖我之人的手与我一同在桌子上。人子固然要照所预定的而去，但那人出卖祂的，却有祸了。} 此时我们必须认为，若望所记的应插入此处，而其他三者省略了它，正如若望也略过了他们所详述的某些事情一样。与此相应，在递杯之后，以及主后来所说的话——即路加所引入的：\Quote{但看，那出卖我之人的手与我一同在桌子上}等等——之后，若望的叙述〔应被视为〕紧接着加入。其文如下：\Quote{那时，有一位门徒斜靠在耶稣的怀里，就是耶稣所爱的那位。西满伯多禄向他示意，并对他说：\InnerQuote{祂说的是谁？}他便紧靠在耶稣的胸膛上，对祂说：\InnerQuote{主，是谁？}耶稣回答说：\InnerQuote{我蘸了一块饼，递给谁，谁就是。}祂便蘸了饼，递给依斯加略人西满的儿子犹达斯。他吃了饼后，撒殚就进入了他。}\par

4. 在此我们必须小心，不要使若望显得既与路加相矛盾——路加在此之前已说过，撒殚在犹达斯与犹太人商议、为得到一笔钱而出卖祂时，已进入了他的心——又与他自己相矛盾。因为早前，在〔他记载〕接受这块饼之前，他已用了这些词句：\Quote{晚餐后，魔鬼已将出卖祂的念头放入犹达斯心中。} 魔鬼如何进入人的心，难道不是将不义的劝说放入不义之人的思想吗？然而，解释是这样的：我们应当认为犹达斯此时更完全地被魔鬼所占据，正如另一方面，在善人身上，那些在复活后、祂向他们嘘气说：\Quote{你们领受圣神罢！}时已经领受了圣神的人，后来在圣神从上降下的五旬节那天，也获得了更圆满的圣神恩赐。同样，撒殚在这块饼之后进入了这个人。而且（如若望自己在紧接着的上下文中提到的）\Quote{耶稣对他说：\InnerQuote{你所要做的，快去做罢！}那时在座的人没有一个知道祂为什么向他说这话；有些人因为犹达斯管钱囊，就以为耶稣对他说：\InnerQuote{你去买我们过节所需用的东西}，或是叫他拿些东西给穷人。他接了饼，就立刻出去，那时是夜间。他出去以后，耶稣就说：\InnerQuote{现在人子获得了光荣，天主也在祂身上获得了光荣。天主既然在祂身上获得了光荣，天主也要在自己内光荣祂，并且立刻就要光荣祂。}}\par

\SourceRecord{Translated by S.D.F. Salmond.；Nicene and Post-Nicene Fathers, First Series；Vol. 6.；Edited by Philip Schaff.；Buffalo, NY: Christian Literature Publishing Co.,；1888.；<http://www.newadvent.org/fathers/1602301.htm>.}

% structure: p0001 source=h1 target=section reason=page-title

\section{\WorkTitle{福音的和谐，卷三}}

% structure: p0002 source=h2 target=subsection reason=relative-heading-rank; base=h2

\subsection{第二章 证明在预言伯多禄否认的记载中无任何矛盾}

5. \BibleQuote{孩子们！我同你们在一起的时候不多了。你们将要寻求我，正如我曾对犹太人说过：我所去的地方，你们不能去；现在我也给你们说。我给你们一条新命令：你们该彼此相爱，如同我爱了你们，你们也该照样彼此相爱。如果你们彼此相亲相爱，众人因此就可认出你们是我的门徒。西满伯多禄问耶稣说：主！你往哪里去？耶稣回答说：我所去的地方，你如今不能跟我去，但后来却要跟我去。伯多禄向他说：主！为什么现在我不能跟你去？我要为你舍掉我的性命。耶稣回答说：你要为我舍掉你的性命吗？我实实在在告诉你：鸡叫以前，你要三次不认我。} 以上引文出自\BibleBook{若望}{福音}，并非只有若望记录了这一预告伯多禄否认的事件。其他三位圣史也记载了同样的事。然而，他们并非都以基督言论中同一具体要点作为叙述此事的契机。因为玛窦和马尔谷按完全平行的顺序，并在叙述的同一阶段引入此事，即主离开他们吃逾越节晚餐的房子之后；另一方面，路加和若望则在他离开那个场景之前就引入此事。不过，我们很容易假设，要么是前两位圣史以追述的方式插入，要么是后两位圣史以预述的方式插入；但这种假设可能因以下情况而更加可疑：不仅主的话有显著差异，甚至引入此事的那些情感——伯多禄正是因这些情感而冒失地断言他愿与主同死或为主舍命——也大不相同。这些考虑可能迫使我们理解为：此人实际上三次作出那冒失的宣告，每次都是在基督系列言论中不同的场合被激发出来的，而且主也三次回答他，预示鸡叫以前他会三次否认主。\par

6. 当然，设想伯多禄在几次彼此相隔一定时间的场合被引向如此妄断的行动，正如他实际被煽动多次否认主一样，这丝毫不足为奇。同样，认为主在三个连续的时刻给予他类似的答复，也并非不合理；尤其当我们看到，基督接连三次问他是否爱祂，彼此紧接，毫无任何其他事实或言语插入其间；而当伯多禄三次给出同样的回答时，主也三次给他同样的嘱咐，要他牧养祂的羊群。正如我们所看到的，福音书作者所采用的措辞记载了主以不同形式、不同含义所说的话语，这进一步证明：更合理的推测是，伯多禄在三个不同的场合表现了他的妄断，并三次从主那里领受关于他三重否认的警告。在此让我们再次注意我稍早从\BibleBook{若望}{福音}中引出的那段经文。在那里我们确实发现祂这样说：\BibleQuote{孩子们！我同你们在一起的时候不多了。你们将要寻找我；就如我曾对犹太人说：我去的地方，你们不能去；现在我也对你们说。我给你们一条新命令：你们该彼此相爱；如同我爱了你们，你们也该照样彼此相爱。如果你们之间彼此相亲相爱，世人因此就可认出你们是我的门徒。西满伯多禄问祂说：主！祢往哪里去？}显然，这里促使伯多禄发出\BibleQuote{主！祢往哪里去？}这一问题的，正是主自己所说的话。因为他已听见祂说：\BibleQuote{我去的地方，你们不能去。}于是耶稣对伯多禄回答说：\BibleQuote{我去的地方，你现在不能跟随我，但后来却要跟随我。}伯多禄随即这样说：\BibleQuote{主！为什么现在我不能跟随祢？我要为祢舍掉我的性命。}对这妄断的宣告，主以预言他的否认来回应。\Person{路加}则首先提到主如何说：\BibleQuote{西满，西满，看，撒殚求得了许可，要筛你们像筛麦子一样；但我已为你祈求了，为叫你的信德不致失落；待你回头后，要坚固你的兄弟。}紧接着他立即告诉我们伯多禄如此回答：\BibleQuote{主！我准备同祢一起下狱，同受死。}然后他继续写道：\BibleQuote{祂说：伯多禄，我告诉你：今天鸡叫以前，你要三次否认认识我。}谁看不出这是一个单独的场合，而在此场合中伯多禄被激发作出上述妄断宣告的境遇是全然不同的呢？但再看\Person{玛窦}，他给我们提供以下经文：\BibleQuote{他们唱了圣咏后，他说，就出来往橄榄山去。那时耶稣对他们说：今夜你们都要为我的缘故跌倒；因为经上记载：我要打击牧人，羊群就四散。但在我复活以后，我要在你们以先往加里肋亚去。}\Person{马尔谷}以完全相同的形式记载了同一段。然而，这些话或其中表达的思想，与若望所记伯多禄妄断宣告的措辞，或路加所描绘他发出如此断言的方式，有什么相似之处呢？因此我们发现，在\Person{玛窦}的叙述中，紧接着就是这样：\BibleQuote{伯多禄回答祂说：即便众人都为祢的缘故跌倒，我永远不会跌倒。耶稣对他说：我实在告诉你：今夜鸡叫以前，你要三次否认我。伯多禄对祂说：即便我必须同祢一起死，我也决不会否认祢。众门徒也都这样说了。}\par

7. 这一切也几乎以相同的语言由\Person{玛尔谷}记载下来，只是他并未如此一般地记录主关于伯多禄的跌倒事件如何发生所说的话，而是给予了更具体的表述。因为他的版本是这样的：\BibleQuote{我实在告诉你们：就在今天，就在这一夜，鸡叫两遍以前，你们要三次否认我。}由此可见，所有圣史都告诉我们主预言伯多禄在鸡叫以前要否认他，但他们并未都提及鸡叫的次数，只有\Person{玛尔谷}在叙述中对此事给出了更明确的提示。因此，有人认为\Person{玛尔谷}的陈述与其他圣史不一致。但这仅仅是因为他们没有充分留意事实的真相，尤其是因为他们带着偏见的心态来审视这个问题，从而受到对福音怀有敌意的态度所支配。事实上，伯多禄的否认，就其整体而言，是一次三次的否认。因为他一直处于同样的心神不安状态，并怀有同样的说谎意图，直到关于他的预言被提醒，并通过痛苦的哭泣和内心的忧伤得到了医治。然而，显然如果这完整的否认——也就是说，三次否认——被认为是在第一次鸡叫之后才开始，那么三位圣史就会显得对此事的记载有误。因为\Person{玛窦}的版本是这样的：\BibleQuote{我实在告诉你：今夜鸡叫以前，你要三次否认我。}\Person{路加}这样表述：\BibleQuote{我告诉你，\Person{伯多禄}，今天鸡叫以前，你要三次说不认识我。}\Person{若望}则这样呈现：\BibleQuote{我实实在在告诉你：鸡叫以前，你要三次否认我。}因此，用不同的措辞和不同的语序，这三位圣史表达了主所说的话的同一种含义——即在鸡叫以前，伯多禄要三次否认他。另一方面，如果我们假设他在鸡开始叫之前就完成了全部三次否认，那么\Person{玛尔谷}就会被指控给出了多余的陈述，因为他把以下这些话放在主的口中：\BibleQuote{我实在告诉你们：就在今天，鸡叫两遍以前，你们要三次否认我。}因为说\BibleQuote{在鸡叫两遍以前}有什么必要呢？如果按照假设，这全部三次否认在第一次鸡叫之前就已经完成了，那么显然，在第一次鸡叫之前就已经完成的否认，理所当然地被认为是在第二次鸡叫之前、第三次鸡叫之前，简言之，在那个晚上的所有鸡叫之前都已经完全说出了。但是，既然这三次否认是在第一次鸡叫之前开始的，那么那三位圣史所关注的并不是伯多禄完成否认的时间，而是否认的程度和开始的时间；也就是说，他们的目的是要说明否认将会重复三次，并且会在鸡叫之前开始。同时，就人自己的内心而言，我们也很容易理解，这否认作为一个整体是在第一次鸡叫之前就已在心中酝酿了。因为尽管就实际发出否认之言的人而言，这三次否认是在第一次鸡叫之前才开始，并在第二次鸡叫之前才真正完成，但就\Person{伯多禄}内心的倾向和恐惧而言，这否认作为一个整体是在第一次鸡叫之前就已构思。如果他在第一次鸡叫之前就已经完全心中存有否认——实际上是由于他充满了如此卑怯的恐惧，以至于不仅一次，而且第二次、甚至第三次被问及主时都能否认他——那么，在各次否认之间所间隔的时间长短并不重要。因此，对事情更正确和仔细的考虑可能会表明，正如经上所说，\BibleQuote{凡注视妇女而贪恋她的，已经在心里与她犯了奸淫}，同样，在这里，由于\Person{伯多禄}所说的话仅仅表达了他心中早已构思并强烈到足以持续三次否认主的恐惧，这三次否认应被归因于那个特定时刻——即足以驱使他一再否认的恐惧占据他内心的时刻。这样的话，即使否认的言辞是在第一次鸡叫之后才开始从他口中发出，因为他的内心被询问所打击，那么说他在鸡叫以前就否认了三次也不会包含任何荒谬或不真实，因为无论如何，在鸡叫以前，他的内心已被如此强烈的恐惧所侵袭，足以使他陷入第三次否认。因此，我们更不应感到困难，如果这三次否认，就否认者所发出的三次言论而言，是在鸡叫之前开始的，尽管不是在第一次鸡叫之前完成的。我们可以举一个平行的例子：假设有人被告知：\Quote{今夜，鸡叫以前，你要写一封信给我，在信中你要三次辱骂我。}那么，如果这个人是在鸡叫之前开始写信，而在第一次鸡叫之后才写完，这并不能说明之前的告知是假的。因此，事实上，\Person{玛尔谷}将以下这些话放在主的口中：\BibleQuote{在鸡叫两遍以前，你们要三次否认我，}这更清楚地表明了各次否认之间的时间间隔。当我们到达福音叙述的那一段时，我们会看到情景的呈现方式也在那个上下文中展示了圣史们之间的一致性。\par

8. 然而，如果要求得到主对\Person{伯多禄}所说的原话，完全照字面意义，那么我们必须回答，这是不可能发现的；此外，追问这些话也是多余的，因为说话者的含义——他说这些话的目的——尽管福音作者使用了不同的措辞，却可以极其清楚地被理解。那么，无论是\Person{伯多禄}在主的言论过程中，因不同场合的刺激，三次作出他那冒昧的宣告，并三次被主预言他将否认，正如我们的研究所指向的更可能的结果那样；还是所有福音作者所给的记述，在采用某种叙述顺序后，都可以归结为一种单一的陈述，从而可以证明，主只在一次场合下，在\Person{伯多禄}表现出他那冒昧的精神时，预言他将否认祂——无论哪种情况，都无法发现福音作者之间存在任何矛盾，因为实际上并不存在这样的矛盾。\par

\SourceRecord{Translated by S.D.F. Salmond.；Nicene and Post-Nicene Fathers, First Series；Vol. 6.；Edited by Philip Schaff.；Buffalo, NY: Christian Literature Publishing Co.,；1888.；<http://www.newadvent.org/fathers/1602302.htm>.}

% structure: p0001 source=h1 target=section reason=page-title

\section{福音的和谐 第三卷}

% structure: p0002 source=h2 target=subsection reason=relative-heading-rank; base=h2

\subsection{第三章 论如何能证明福音书在记载主从用完晚餐后离屋之前所说的话时，彼此之间并无矛盾。}

9. 因此，现在我们可以尽可能遵循从所有福音书作者共同收集的叙述顺序。这样，在根据若望的记载，向伯多禄作出上述预言之后，同一若望继续他的陈述，并在此处引入了主的言论，内容如下：\BibleQuote{你们心里不要忧愁；你们信天主，也当信我。在我父的家里有许多住处；}他详细叙述了在那段讲话中主所发表的如此难忘且极其崇高的言论，直到他适当地来到主说以下话的地方：\BibleQuote{公义的父啊，世界没有认识你，我却认识你；这些人也知道是你差了我来。我已将你的名指示他们，还要指示他们，使你所爱我的爱在他们里面，我也在他们里面。}此外，根据路加的记载，我们发现他们中间\BibleQuote{起了争论，他们中间哪一个可算为大。耶稣对他们说：外邦人有君王为主治理他们，那掌权管他们的称为恩主。但你们不可这样；你们里头为大的，倒要像年幼的；为首领的，倒要像服事人的。是谁为大？是坐席的呢？是服事人的呢？不是坐席的大吗？然而，我在你们中间如同服事人的。你们是在我的试炼中始终与我同在的；我将国赐给你们，正如我父赐给我一样，叫你们在我国里，在我的席上吃喝，并且坐在宝座上，审判以色列十二个支派。}接着，上述路加立即将以下段落附在这些话之后：\BibleQuote{主又说：西满，西满，撒但想要得着你们，好筛你们像筛麦子一样；但我已经为你祈求，叫你的信德不至失丧；你回头以后，要坚固你的弟兄。伯多禄说：主啊，我就是同你下监，同你受死，也是甘心。耶稣说：伯多禄，我告诉你，今日鸡还没有叫，你要三次不认我。耶稣又对他们说：我差你们出去的时候，没有钱囊，没有口袋，没有鞋，你们缺少什么没有？他们说：没有。耶稣说：但如今有钱囊的可以带着，有口袋的也可以带着，没有剑的，要卖衣服买一把剑。我告诉你们，经上写着说：他被列在罪犯之中。这句话必应验在我身上，因为那关系我的事必然成就。他们说：主啊，请看！这里有两把剑。耶稣说：够了。}接下来是玛窦和马尔谷共同记载的段落：\BibleQuote{他们唱了诗，就出来往橄榄山去。那时，耶稣对他们说：今夜，你们为我的缘故都要跌倒。因为经上记着说：我要击打牧人，羊就分散了。但我复活以后，要在你们以先往加利利去。伯多禄回答说：众人虽然为你的缘故跌倒，我却永不跌倒。耶稣说：我实在告诉你，今夜鸡叫以先，你要三次否认我。伯多禄说：我就是必须和你同死，也总不能不认你。众门徒都是这样说。}我们按照玛窦的记载介绍了前面的段落。但马尔谷的记载也几乎与这些相同，除了我们在上面已经澄清的关于鸡叫的明显矛盾之处。\par

\SourceRecord{Translated by S.D.F. Salmond.；Nicene and Post-Nicene Fathers, First Series；Vol. 6.；Edited by Philip Schaff.；Buffalo, NY: Christian Literature Publishing Co.,；1888.；<http://www.newadvent.org/fathers/1602303.htm>.}

% structure: p0001 source=h1 target=section reason=page-title

\section{福音合论 第三卷}

% structure: p0002 source=h2 target=subsection reason=relative-heading-rank; base=h2

\subsection{第四章　离开晚餐厅后，来到园地或花园中所发生的事；以及在若望对此沉默的情况下，如何在其他三位福音作者——即玛窦、马尔谷和路加——之间证明真正的和谐一致。}

10. 玛窦接着按同一顺序叙述如下：\BibleQuote{那时，耶稣同他们来到一个名叫革责玛尼的地方。}马尔谷也提到这一点。路加也提及此地，但他没有提到那园地的名称。因为他这样说：\BibleQuote{耶稣出来，照常往橄榄山去，门徒也跟随着他。到了那地方，便对他们说：你们要祈祷，免得陷于诱惑。}这就是其他两人以革责玛尼之名所举出的地点。我们知道，那里就是若望在以下叙述中所提到的那个园子：\BibleQuote{耶稣说了这些话，就和门徒出去，过了克德龙溪，在那里有一个园子，他和门徒都进去了。}接着，按玛窦的记载，下一句话是：\BibleQuote{他对门徒说：你们坐在这里，等我到那边去祈祷。他带了伯多禄和载伯德的两个儿子同去，开始忧闷悲恸起来。然后对他们说：我的心灵忧闷得要死；你们留在这里，同我一起醒寤。他稍往前行，俯伏在地祈祷说：我圣父，若是可能，就让这杯离开我吧！但不要照我所愿意的，而要照你所愿意的。他来到门徒那里，见他们睡着了，便对伯多禄说：怎么，你们竟不能同我醒寤一个时辰？醒寤祈祷吧！免得陷于诱惑；心神固然愿意，肉体却是软弱的。他第二次再去祈祷说：我圣父，如果这杯不能离开我，非要我喝不可，就愿你的旨意成就吧！他又回来，见他们仍然睡着，因为他们的眼睛很沉重。他离开他们，再去祈祷，第三次说了同样的话。然后他来到门徒那里，对他们说：现在你们还睡下去，休息下去吧！看，时候到了，人子要被交在罪人手中。起来，我们走吧！看，那出卖我的人近了。}\par

11. 马尔谷也记载了这些段落，以完全相同的方式和顺序引入。不过，他有些句子写得更简短，另一些则稍加解释。主所说的这些话，在玛窦的版本中，有一部分似乎彼此有些矛盾。我指的是：在第三次祈祷后，他来到门徒那里，对他们说：\BibleQuote{现在你们还睡下去，休息下去吧！看，时候到了，人子要被交在罪人手中。起来，我们走吧！看，那出卖我的人近了。}因为当随后紧接另一句话\BibleQuote{看，时候到了}，以及之后的吩咐\BibleQuote{起来，我们走吧}时，该如何理解上面所给的指示\BibleQuote{现在你们还睡下去，休息下去吧}？那些在这里觉得有矛盾的人，试图以一种表示责备而非允许的语调来念\BibleQuote{现在你们还睡下去，休息下去吧}。如果有必要，这个办法完全可以采用。然而，马尔谷在复述这些话时，暗示主在说完\BibleQuote{现在你们还睡下去，休息下去吧}之后，又加了\BibleQuote{够了}，然后附上进一步的陈述\BibleQuote{时候到了；看，人子要被出卖了}。因此我们可以得出结论：实际情况是这样：主在说完\BibleQuote{现在你们还睡下去，休息下去吧}之后，沉默了片刻，以便他允许他们做的事能够真正实行；然后他才说了另一句话\BibleQuote{看，时候到了}。因此，在马尔谷福音中，我们看到这些话（关于睡觉）之后紧接着\BibleQuote{够了}，也就是说，\BibleQuote{你们现在休息得够了}。但由于没有明确提到主当时所插入的沉默，这段经文就以牵强的方式被理解，并被认为必须以特别的发音来念这些话。\par

12. 另一方面，路加省略了他祈祷的次数。但他告诉我们其他福音作者没有记载的一件事：即当他祈祷时，有一位天使在加强他；当他更恳切地祈祷时，他流出血汗，汗珠滴落到地上。因此，当他做如下陈述\BibleQuote{他从祈祷中起来，来到门徒那里}时，并未说明那时他已祈祷了多少次。然而，他这样做并不与其他两人相矛盾。此外，若望确实提到他和门徒进入园子，却没有叙述在叛徒与犹太人前来捉拿他之前，他在那里做了什么。\par

13. 因此，这三位福音作者以这种方式叙述了同一事件，正如另一方面，一个人可能对同一件事给出三种不同的陈述，措辞上有所差异，但没有任何真正的矛盾。例如，路加更明确地指出了祂离开门徒向前走的距离——也就是说，当祂离开他们去祈祷时——因为他说是\Quote{约有一掷石之远}。\newline{}马尔谷则先以自己的话记述主祈祷说，\InnerQuote{若是可能，就让这时辰离开祂}，指的是祂受难的时辰，随后他也用\InnerQuote{杯}这个词来表达。然后他以如下方式复述了主自己的话：\InnerQuote{阿爸，圣父，一切为祢都可能：请将这杯从我撤去。}如果我们把另外两位福音作者所给的从句，以及马尔谷本人已经引入的一个清晰的平行句（以他自己代替主陈述的形式）与这些词连接起来，那么整个句子将以这种形式呈现：\InnerQuote{圣父，若是可能（因为一切为祢都可能），请将这杯从我撤去。}这样表述正是为了防止任何人以为当祂说\InnerQuote{若是可能}时，祂贬低了圣父的能力。因为祂的话不是\InnerQuote{若祢能做}，而是\InnerQuote{若是可能}。而凡是祂愿意的事，都是可能的。因此，这里的\InnerQuote{若是可能}与\InnerQuote{若祢愿意}具有相同的含义。因为马尔谷在说\InnerQuote{一切为祢都可能}时，已使\InnerQuote{若是可能}这一短语的涵义十分清楚。此外，这些作者记载了祂说\InnerQuote{然而，不要照我所愿的，而要照祢所愿的}（这个表述与另一种形式\InnerQuote{然而，不要照我的意愿，而要照祢的意愿成就}完全等同），这清楚地表明，\InnerQuote{若是可能}这句话所指的，并非圣父方面有任何绝对的不可能，而只是指祂的意愿。路加所陈述的更为清晰的说法使这一点更加明显。因为他的版本不是\InnerQuote{若是可能}，而是\InnerQuote{若祢愿意}。对于这一更清晰的真实含义的表达，我们可以再加上马尔谷插入的从句，从而使整个句子这样呈现：\InnerQuote{若祢愿意（因为一切为祢都可能），请将这杯从我撤去。}\par

14. 再次，关于马尔谷提到主不仅说\Quote{圣父}，而且说\Quote{阿爸，圣父}，解释很简单：\Quote{阿爸}在希伯来文中与\Quote{父}在拉丁文中完全等同。也许主使用这两个词带有某种象征意义，意在表明，在承受这痛苦时，祂承担了祂的身体——教会——的角色，祂已成了教会建筑的角石，而教会部分来自希伯来人（祂用\Quote{阿爸}指称他们），部分来自外邦人（祂用\Quote{父}指称他们）。宗徒保禄也使用了同样有意义的表达。因为他说：\BibleQuote{我们藉着祂呼唤：阿爸，圣父}，在另一处又说：\BibleQuote{天主派遣祂的圣神到你们心中，呼唤：阿爸，圣父}。因为好导师和真救主，通过分担比较软弱者的痛苦，应在自己身上说明一个真理：祂的见证人不应当绝望，尽管可能由于人性的软弱，在受苦之时悲伤会潜入他们的心；因为他们若能铭记天主知道什么对祂所关心的人最好，并将祂的意愿置于自己的意愿之上，就能克服悲伤。然而，关于这个问题，整体上现在不是进行更详细讨论的时候。因为我们只处理关于福音作者和谐性的问题；从他们多样化的叙述方式中，我们学到了有益的教训：为了得到真理，在处理措辞时唯一要紧的是说话者使用它们的意图。因为\Quote{圣父}一词与\Quote{阿爸，圣父}这个短语含义相同。但为了揭示神秘意义，\Quote{阿爸，圣父}这一表达更为清晰；而为了指明合一，\Quote{圣父}一词已经足够。主的确使用了\Quote{阿爸，圣父}这种称呼方式，这必须被接受为事实。但祂的意图不会非常明显，除非有方法（既然别人只使用\Quote{圣父}一词）表明在这种表达形式下，那由犹太人构成和由外邦人构成的两个教会也真正是一体。这样，我们可以认为\Quote{阿爸，圣父}这个短语被采用是为了传达主在另一场合所说的同样观念：\BibleQuote{我还有其他羊，不属于这羊栈。}这句话无疑指外邦人，因为祂在以色列人中也有羊。但在那段经文中，祂紧接着补充说：\BibleQuote{我必须也把他们领来，好使成为一牧一栈。}因此我们可以说，正如\Quote{阿爸，圣父}这个短语包含了（两个种族）以色列人和外邦人的观念，单独使用\Quote{圣父}一词则指向这两个种族合而为一的羊群。\par

\SourceRecord{Translated by S.D.F. Salmond.；Nicene and Post-Nicene Fathers, First Series；Vol. 6.；Edited by Philip Schaff.；Buffalo, NY: Christian Literature Publishing Co.,；1888.；<http://www.newadvent.org/fathers/1602304.htm>.}

% structure: p0001 source=h1 target=section reason=page-title

\section{福音的和谐，第三卷}

% structure: p0002 source=h2 target=subsection reason=relative-heading-rank; base=h2

\subsection{第五章 论四位福音传道者在祂被捕时所行所说之事上的记载，并证明这些不同叙述并无真正矛盾。}

15. 当我们依据玛窦和马尔谷的版本时，我们发现历史叙述如此进行：\BibleQuote{当祂还在说话时，看，十二人之一的犹达斯来了，并同他一起有一大群人，带着刀剑棍棒，从司祭长和民间长老那里来。那出卖祂的人给了他们一个记号，说：我亲吻谁，谁就是祂；你们把祂拿住。他随即来到耶稣跟前，说：老师，你好！并亲吻了祂。} 然而，首先，正如我们从路加的叙述中得知的，祂对叛徒说：\BibleQuote{犹达斯，你用亲吻来出卖人子吗？} 接着，从玛窦我们学到，祂这样说：\BibleQuote{朋友，你为什么来？} 此后，祂又说了某些话，这些出现在若望的叙述中，其内容如下：\BibleQuote{你们找谁？他们回答说：纳匝肋人耶稣。耶稣对他们说：我就是。出卖祂的犹达斯也同他们站在一起。祂一说我就是，他们就倒退，跌倒在地上。于是祂又问他们说：你们找谁？他们说：纳匝肋人耶稣。耶稣回答说：我已告诉你们我就是；如果你们找我，就让这些人去吧。这应验了祂所说的话：你所赐给我的人，我没有失落一个。}\par

16. 接着出现一段路加记载的话：\BibleQuote{当祂周围的人看见将要发生的事时，他们对祂说：主，我们可用剑砍吗？其中一人击打了大司祭的仆人，} 正如四位史家都注意到的，\BibleQuote{削掉了他的耳朵，} 而根据路加和若望，那是他的\BibleQuote{右耳。} 此外，我们还从若望得知，击打仆人的人是伯多禄，他击打的那个人的名字叫玛尔苛。接下来我们看路加提到的：\BibleQuote{耶稣回答说：且到此为止。} 与之相连的是玛窦补充的话：\BibleQuote{把你的剑放回原处：因为凡动剑的，必死于剑下。你想我不能现在祈求我的父，祂会立刻给我十二营以上的天使吗？但圣经怎能应验，说这事必须如此？} 同时还可加上若望告诉我们在同一场合祂所说的问题：\BibleQuote{我父所给我的杯，我岂可不喝？} 然后，如路加所记载的，祂触摸了那被打之人的耳朵，治好了他。\par

17. 同样，我们不应让这样的想法困扰我们，即\Person{路加}告诉我们，当门徒问主是否可以用剑击打时，主用这些话回答说：\BibleQuote{你们暂且容忍至此，}这似乎暗示主在击打之后如此表达，表示他对至此所做的事满意，但不愿再有进一步行动；而\Person{玛窦}所用的措辞则让我们理解为，\Person{伯多禄}动刀这件事完全令主不悦。因为更正确的推测是，当他们向主提问说：\BibleQuote{主，我们可以用剑击打吗？}时，他当时回答说：\BibleQuote{你们暂且容忍至此；}他的意思是：\BibleQuote{不要因即将发生的事而惊慌。这些人应被容许至此；也就是说，容许他们来捉拿我，从而成就那关于我所写的一切。}然而，我们还需进一步推测，在他们向主提问和主对他们回答的对话过程中，\Person{伯多禄}因渴望作为护卫者，并因对主更强烈的激动而动手击打。但这两件事可能同时轻易发生，同一人不可能同时说出两句不同的话。因为作者不会使用\BibleQuote{\Person{耶稣}回答说}这样的表述，除非这些话是对周围之人所提问题的回答，而不是对\Person{伯多禄}行为的评述。因为只有\Person{玛窦}记录了\Person{耶稣}对\Person{伯多禄}行为的评判。在那段经文中，\Person{玛窦}所用的措辞也不是\BibleQuote{\Person{耶稣}这样回答\Person{伯多禄}：把剑收起来；}而是以这样的措辞：\BibleQuote{\Person{耶稣}于是对他说：把剑收起来；}由此看来，主是在行动之后才如此表明的。至于\Person{路加}所用的措辞，即\BibleQuote{\Person{耶稣}回答说：你们暂且容忍至此，}必须被认为是回答那些向他提问之人的。然而，按照我们先前的解释，仆人挨的那一剑恰恰是在这些人与主之间进行上述问答的时刻发生的，因此圣史认为按同样的顺序记录这一行为是合理的，因此它被插入在问话和答话的措辞之间。因此，这与\Person{玛窦}引入的陈述并无矛盾，即\BibleQuote{凡动剑的，必死于剑下，}——也就是说，那些动剑的人。但如果认为主的回答表明他赞同这次主动使用剑，即使只造成一处非致命的伤口，那么这里似乎有些不一致。然而，向\Person{伯多禄}说的话可以被理解为整体上与其他部分完全和谐；因此，将\Person{路加}和\Person{玛窦}所记载的结合起来，如上所述，我们得到以下连贯的语句：\BibleQuote{你们暂且容忍至此。把你的剑收回原处；因为凡动剑的，必死于剑下，}等等。此外，如何理解\BibleQuote{你们暂且容忍至此}这句话，我已经解释过了。如果有更好的解释方法，也可以。只是无论如何都要维护圣史的诚实。\par

18. 此后，\Person{玛窦}继续叙述，并提及在那时刻，\Person{耶稣}向群众说了如下的话：\BibleQuote{你们带着刀剑棍棒出来拿我，如同对付强盗一样吗？我天天坐在圣殿里教导人，你们并没有捉拿我。}然后他又加了一些话，\Person{路加}如此记载：\BibleQuote{但现在是你们的时候，黑暗的权势。}接下来是\Person{玛窦}给出的句子：\BibleQuote{这一切事的发生，是为应验先知书上的记载。于是，众门徒都撇下他逃跑了。}最后一件事实也被\Person{马尔谷}记载。同一位圣史还补充了以下内容：\BibleQuote{有一个少年人，身上只披着一块麻布，跟随他；当他们捉拿他时，他就撇下麻布，赤身从他们逃走了。}\par

\SourceRecord{Translated by S.D.F. Salmond.；Nicene and Post-Nicene Fathers, First Series；Vol. 6.；Edited by Philip Schaff.；Buffalo, NY: Christian Literature Publishing Co.,；1888.；<http://www.newadvent.org/fathers/1602305.htm>.}

% structure: p0001 source=h1 target=section reason=page-title

\section{福音的和谐 卷三}

% structure: p0002 source=h2 target=subsection reason=relative-heading-rank; base=h2

\subsection{第六章 论这些福音作者在叙述主被带到大司祭府邸时所发生之事的和谐，以及在半夜被带到该府邸后所发生的事件，特别是伯多禄否认的事}

19. 继玛窦的叙述，我们读到这样一句话：\BibleQuote{他们拿住耶稣，把祂带到大司祭盖法那里去；经师和长老已在那里聚集。}然而，我们从若望得知，祂先被带到亚纳斯那里，他是盖法的岳父。另一方面，马尔谷和路加都未提大司祭的名字。此外，我们发现祂是被捆绑着带去的。因为若望告诉我们，在那人群中，有千夫长和一营兵，以及犹太人的仆役。接着玛窦记载：\BibleQuote{伯多禄远远地跟着耶稣，直到大司祭的庭院，并且进去同仆役坐着，要看结局。}马尔谷对此补充说：\BibleQuote{他在火旁取暖。}路加也说了类似的话：\BibleQuote{伯多禄远远地跟着。他们在庭院中间生了火，一同坐下，伯多禄也坐在他们中间。}若望这样继续：\BibleQuote{西满伯多禄跟着耶稣，另一个门徒也跟着。那个门徒（即那个另一个门徒）是大司祭认识的，就同耶稣一起进了大司祭的庭院。但伯多禄（如同一若望所述）站在门外。后来那大司祭认识的门徒出去，对看门的使女说了一声，就领伯多禄进去。}最后这一事实归功于若望的叙述。这样我们看出伯多禄如何得以进去，并如其他福音作者所述，留在了庭院内。\par

20. 玛窦接着这样记载：\BibleQuote{那时，司祭长、长老和整个公议会，寻找假见证来反对耶稣，为要处死祂；但找不到。尽管有许多假见证前来，还是找不到。}马尔谷在这里解释：\BibleQuote{他们的见证互不相合。}但如玛窦继续所说：\BibleQuote{最后来了两个假见证人，说：这个人说过：我能拆毁天主的圣殿，三天内重建起来。}马尔谷提到还有其他的人说：\BibleQuote{我们听见祂说：我要拆毁这座人手所造的圣殿，三天内另建一座非人手所造的。因此（如马尔谷在同一段所注）他们的见证互不相合。}接着玛窦叙述：\BibleQuote{大司祭站起来对祂说：你什么都不回答吗？这些人作证反对你的是什么？耶稣却默不作声。大司祭又对祂说：我因生活的天主命令你起誓，告诉我们你是不是基督，天主子。耶稣对他说：你说的是。}马尔谷以不同的措辞记载同一事件，只是他没有提到大司祭令祂起誓。但他明确指出，耶稣回答大司祭的两句话——即\InnerQuote{你说的是}和\InnerQuote{我是}——实际意思相同。因为正如上面提到的马尔谷所述，叙述继续如下：\BibleQuote{耶稣说：我是；你们将要看见人子坐在大能者的右边，乘着天云降来。}这正如玛窦也展示的段落，唯一区别是玛窦没有说耶稣回答的是\InnerQuote{我是}。玛窦继续这样叙述：\BibleQuote{大司祭撕破自己的衣服说：祂说了亵渎的话；何必还需要见证人呢？你们已听见了祂的亵渎。你们以为怎样？他们回答说：祂该被处死。}马尔谷的记载完全一致。于是玛窦继续：\BibleQuote{那时他们向祂脸上吐唾沫，用拳头打祂；另有人用手掌打祂，说：给我们说预言，基督，打你的是谁？}马尔谷同样记录了这些事。他还提到另一件事，即他们蒙住了祂的脸。路加也为这些事件作证。\par

21. 这些事主经历了，发生在最初被带去的那个大司祭的府邸中，从半夜直到清晨，伯多禄也在此受试探。不过，关于伯多禄的这次试探——发生在主承受这些凌辱期间——各位福音作者在叙述顺序上并不一致。因为玛窦和马尔谷先叙述对主的凌辱，然后才讲伯多禄的试探。路加则先描述伯多禄的试探，之后才写主所受的侮辱；而若望则先叙述伯多禄试探的一部分，然后插入几节记录主所承受的事，接着补充说主从亚纳斯那里被解送到大司祭盖法那里（\Emph{i.e.} from Annas），随后在此处接续并总括他之前所开始的对伯多禄在最初被带去的府邸中受试探的叙述，完整描述那件事，此后又回到主所经历的事上，告诉我们祂如何被带到盖法面前。\par

22. 据此，\Person{玛窦}继续写道：\BibleQuote{那时，\Person{伯多禄}坐在外面庭院里；一个使女来到他跟前，说：\InnerQuote{你也是同加里肋亚人\Person{耶稣}一起的。}他却在众人面前否认说：\InnerQuote{我不知道你在说什么。}当他出去到了门廊，另一个使女看见他，就对那里的人说：\InnerQuote{这个人也是同纳匝肋人\Person{耶稣}一起的。}他又起誓否认说：\InnerQuote{我不认识那个人。}过了一会儿，那些站着的人前来对\Person{伯多禄}说：\InnerQuote{你确是他们中的一个，因为你的口音把你暴露了。}于是他就开始诅咒发誓说他不认识那个人。立刻鸡就叫了。}这就是\Person{玛窦}的叙述。但我们也得知，在他出去之后，当他已否认主一次时，第一遍鸡叫了——\Person{玛窦}没有指明这点，但\Person{玛尔谷}暗示了。\par

23. 但他第二次否认主，并非在他出门外之时。那是在他回到火堆旁之后发生的。然而，没有必要提及他返回的具体时间。因此，\Person{玛尔谷}继续叙述这件事：\BibleQuote{他出去到了门廊，鸡就叫了。一个使女又看见他，就开始对站着的人说：\InnerQuote{这是他们中的一个。}他又否认了。}但这使女并非先前那个，而是另一个，正如\Person{玛窦}告诉我们的。不仅如此，我们还进一步得知，在第二次否认时，有两方对他说话：即\Person{玛窦}和\Person{玛尔谷}所提到的使女，以及\Person{路加}所提到的另一个人。因为\Person{路加}的叙述是这样的：\BibleQuote{\Person{伯多禄}远远地跟着。他们在庭院中间生了火，一同坐下，\Person{伯多禄}也坐在他们中间。有一个使女看见他坐在火光中，定睛看他，说：\InnerQuote{这个人也是同他一起的。}他却否认说：\InnerQuote{女子，我不认识他。}过了一会儿，另一个人看见他，说：\InnerQuote{你也是他们中的。}}现在，\Person{路加}所插入的\Quote{过了一会儿}，覆盖了（我们可以推测）\Person{伯多禄}出去和第一遍鸡叫的那段时间。但到此时，他已经又进来了；因此我们可以理解\Person{若望}叙述的一致性，他告诉我们\Person{伯多禄}站在火旁时第二次否认了主。因为在\Person{若望}对\Person{伯多禄}第一次否认的叙述中，他不仅没有提到鸡的第一次啼叫（其他福音作者除了\Person{玛尔谷}外也同样没有提到），而且也没有注意到是当\Person{伯多禄}坐在火旁时使女认出了他。因为\Person{若望}在那里只说了这些：\BibleQuote{于是那个看门的使女对\Person{伯多禄}说：\InnerQuote{你不也是这个人的门徒中的一个吗？}他说：\InnerQuote{我不是。}}然后他插入了他认为应当叙述的同一天在耶稣面前所发生的事。他的记录如下：\BibleQuote{那时，仆人和差役因为天冷，生了炭火，站在那里取暖；\Person{伯多禄}也同他们站在一起取暖。}因此，在这里我们可以设想\Person{伯多禄}出去过，到这个时候又进来了。因为起初他是坐在火旁；过了一段时间，如我们所知，他回来了，开始站着（在火边）。\par

24. 然而，或許有人會對我們說：伯多祿實際上還沒有出去，只是起身想要出去。這可能是那些認為第二次審問和否認發生在伯多祿在外門口時的人的說法。我們來看看若望接著的敘述。其內容如下：\Quote{大司祭就問耶穌有關祂的門徒和祂的教訓。耶穌回答他說：\InnerQuote{我向世界公開講話；我常在會堂和聖殿裏教導，那裏是猶太人常聚集的地方；我暗地裏沒有說過甚麼。你為甚麼問我呢？問那些聽過我的人，我對他們說了甚麼；看，他們知道我所說的。}當祂說了這些，旁邊站着的一個差役用手掌打了耶穌，說：\InnerQuote{你這樣回答大司祭嗎？}耶穌回答他說：\InnerQuote{如果我說得不好，你可以作證不好；如果說得好，你為甚麼打我？}亞納斯就把祂捆綁着解到大司祭蓋法那裏。}這確實表明亞納斯是大司祭。因為當祂受這樣的質問：\Quote{\InnerQuote{你這樣回答大司祭嗎？}}時，耶穌還沒有被解到蓋法那裏。路加也在他福音的開頭提到亞納斯和蓋法是大司祭。這些話之後，若望回到他先前開始的關於伯多祿否認的敘述。這樣他把我們帶回他所記錄的那些事件發生的房屋中，耶穌就是從那裏被解到蓋法那裏，正如瑪竇告訴我們的，在這一幕開始時祂正被帶往那裏。此外，若望是通過重述的方式記錄他在此處引入的有關伯多祿的事。他回顧那件事的敘述，以完整描述三次否認，於是這樣繼續：\Quote{西滿站着烤火。他們就對他說：\InnerQuote{你也是祂的門徒嗎？}他否認說：\InnerQuote{我不是。}}因此，我們發現伯多祿的第二次否認不是發生在門口，而是當他站在火旁的時候。然而，若不是他此時已從外面返回，這是不可能發生的。因為不是說這次他確實出去了，而另一位提到的使女在外面看見了他；而是說她是在他出去時看見他的；換句話說，就在他起身要出去時，她注意到了他，並對那裏的人——即那些在庭院內火旁聚集的人——說：\Quote{這個人也是同納匝肋人耶穌一起的。}然後我們可以設想，那個已經出去的人聽到這話後，又進來，對那些懷有敵意的人發誓說：\Quote{我不認識那個人。}同樣，馬爾谷也論到同一個使女說：\Quote{她就開始對旁邊站着的人說：\InnerQuote{這個人是他們中的一個。}}因為這個使女不是對伯多祿說的，而是對那些在他出去時仍留在那裏的人說的。與此同時，她說話的聲音讓他聽見了；於是他回來，又站在火旁，以否認回應他們的話。然後若望這樣陳述：\Quote{他們說：\InnerQuote{你也不是祂的門徒嗎？}}我們理解這個問題是在他返回站在那裏時向他提出的；我們也認識到這個敘述與以下情況一致：此時伯多祿不僅要應對瑪竇和馬爾谷提到的與第二次否認相關的另一個使女，還要應對路加介紹的另一個人。這就是若望使用複數\Quote{\InnerQuote{他們說}}的原因。解釋可能是：當使女在他出去時對庭院內與她在一起的人說：\Quote{\InnerQuote{這個人是他們中的一個。}}他聽到了她的話就返回，企圖藉否認來為自己辯護。或者，按照更可能的推測，我們可以設想他出去時沒有聽到關於他的話，在他回來時，使女和路加所介紹的另一個人這樣對他說：\Quote{\InnerQuote{你也不是祂的門徒嗎？}}他以否認回應，\Quote{\InnerQuote{說：我不是；}}並且，當路加所說的另一個人更執着地追問說：\Quote{\InnerQuote{你確實是他們中的一個。}}伯多祿就回答：\Quote{\InnerQuote{朋友，我不是。}}儘管如此，當我們比較所有福音作者關於此事的所有陳述時，我們清楚地得出結論：伯多祿的第二次否認不是發生在門口，而是在庭院內的火旁。因此顯而易見，瑪竇和馬爾谷告訴了我們他如何出去，但為了簡潔，卻沒有提及他回來的事實。\par

25. 因此，现在让我们考察在第三次否认这一点上诸位福音作者的一致性，这一点我们先前仅在玛窦的记述中举例说明。马尔谷接着以这样的措辞继续他的叙述：\Quote{过了一会儿，旁边站着的人又对伯多禄说：你确是他们中的一个，因为你也是加里肋亚人。他就开始诅咒发誓说：我不认识你们所说的这个人。立刻鸡叫了第二遍。}路加则继续他的叙述，以这种方式记述同一事件：\Quote{大约隔了一个时辰，另一个人肯定地说：这个人实在同他在一起，因为他也是加里肋亚人。伯多禄说：人啊，我不知道你说什么。他说话的时候，立刻鸡就叫了。}若望接着记述伯多禄的第三次否认，如此写道：\Quote{大司祭的一个仆役，即被伯多禄砍掉耳朵那人的亲戚，说：我不是看见你同他在花园里吗？伯多禄又否认了，立刻鸡就叫了。}玛窦和马尔谷所用的短语\Quote{过了一会儿}所指的具体时间段，路加在说\Quote{大约隔了一个时辰}时阐明了。然而若望没有提示这一时间段。至于玛窦和马尔谷使用复数而非单数，称那些与伯多禄说话的人为\Quote{他们}，而路加只提到一个个体，若望也只指一个人，并进一步指明他是被伯多禄砍掉耳朵那人的亲戚；我们容易解释：或者理解为玛窦和马尔谷在这里采用了常见的表达方式，用复数代替单数；或者设想在场的一个人——一个认识伯多禄并见过他的人——带头说出那话，其余人效仿他的自信，同他一起向伯多禄强调这个说法。如果是这样，那么两位福音作者给出了泛泛的表述，仅用了复数；而另外两位则宁愿具体指明那一位在事件中扮演主要角色的特殊个体。但是，玛窦断言\Quote{你确是他们中的一个，因为你的口音暴露了你}这句话是对伯多禄本人说的。同样，若望告诉我们\Quote{我不是看见你同他在花园里吗？}这个问题是直接对伯多禄提出的。另一方面，马尔谷让我们理解\Quote{他确是他们中的一个，因为他也是加里肋亚人}这句话是旁边站着的人彼此谈论伯多禄的。同样，路加表示另一个人所讲的\Quote{这个人实在同他在一起，因为他也是加里肋亚人}不是对伯多禄说的，而是关于伯多禄的。然而这些差异可以这样解释：或者理解那些说伯多禄是直接受话人的福音作者恰当地转述了大意，因为当着一个人的面谈论他同直接对他说话几乎一样；或者设想这两种说话方式确实都使用了，前者被前几位福音作者注意到，后者被后几位注意到。此外，我们认为第三次否认之后发生了第二次鸡叫，正如马尔谷明确告知我们的。\par

26. \Person{玛窦}接着叙述说：\BibleQuote{\Person{伯多禄}想起了\Person{耶稣}对他所说的话：\InnerQuote{鸡叫以前，你要三次否认我。}他就出去痛哭。}\Person{马尔谷}也是这样记载：\BibleQuote{\Person{伯多禄}想起了\Person{耶稣}对他说的话：\InnerQuote{鸡叫两遍以前，你要三次否认我。}他就开始哭泣。}\Person{路加}的记载如下：\BibleQuote{主转过身来，看了看\Person{伯多禄}。\Person{伯多禄}就想起主的话，主怎样对他说：\InnerQuote{鸡叫以前，你要三次否认我。}\Person{伯多禄}就出去痛哭。}\Person{若望}没有提及\Person{伯多禄}的回忆和哭泣。\Person{路加}在此的陈述，即\BibleQuote{主转过身来，看了看\Person{伯多禄}}，需要更加仔细地考虑，以便正确理解。因为，虽然有所谓的\Quote{内厅（或庭院）}，但\Person{伯多禄}这次是出现在仆役们中间的外庭（或大厅），他们连同他一起在火边取暖。\Person{耶稣}在这个地方被犹太人听见，这也是不可信的假设，这样我们也可以理解所说的\Emph{注视}是肉眼注视。因为\Person{玛窦}首先给我们这样叙述：\BibleQuote{随后，他们就向祂脸上吐唾沫，拳打祂；另有人用手掌打祂，说：\InnerQuote{基督，向我们说预言，打你的是谁？}}然后他立即接着关于\Person{伯多禄}的一段：\BibleQuote{\Person{伯多禄}却坐在庭院外面。}他若不是先前的事都是在屋内向主做的，就不会用后面这个说法。而且，根据\Person{马尔谷}的记载，这些事不仅发生在屋内，还发生在房屋的上层。因为，在记录了上述情况之后，\Person{马尔谷}接着说：\BibleQuote{\Person{伯多禄}在庭院下面。}因此，正如\Person{玛窦}的话\BibleQuote{\Person{伯多禄}却坐在庭院外面}向我们表明先前提及的事是在屋内发生的，\Person{马尔谷}的话\BibleQuote{\Person{伯多禄}在庭院下面}指出这些事不仅发生在屋内，而且发生在房屋的上层。但若是这样，主怎能用肉眼实际的一瞥注视\Person{伯多禄}呢？这些考虑使我得出这样的结论：所提到的注视是从天堂投向\Person{伯多禄}的，其效果是让他在脑海中回想起他现在否认的次数，以及主曾预言性地对他说的话，这样（主如此仁慈地注视着他），引导他悔改，并流出有益救恩的泪水。因此，这个说法将与我们日常使用的其他说话方式类似，比如当我们这样祈祷时：\Quote{主，求祢垂顾我；}或者当提到某人藉天主的仁慈从某种危险或困境中被救出时，我们说\Quote{主垂顾了他。}在圣经中，我们也发现这样的话：\Quote{求祢垂顾我，俯听我；}以及\Quote{主，求祢回来，拯救我的灵魂。}依我判断，这里所用的说法也应当作类似的理解，即说\BibleQuote{主转过身来，看了看\Person{伯多禄}；\Person{伯多禄}就想起主的话。}最后，我们需要注意，虽然福音作者通常更习惯使用\Quote{\Person{耶稣}}这个名号，而不是\Quote{主}这个字，但\Person{路加}在上述句子中专门使用了后者，明确地说：\BibleQuote{\Emph{主}转过身来，看了看\Person{伯多禄}；\Person{伯多禄}就想起\Emph{主}的话。}而\Person{玛窦}和\Person{马尔谷}对这次\Emph{注视}默不作声，因此他们只说\Person{伯多禄}想起的不是\Quote{主}的话，而是\Quote{\Person{耶稣}}的话。由此，我们可以推断，从\Person{耶稣}发出的这次\Emph{注视}不是用肉身眼睛的注视，而是从天堂发出的注视。\par

\SourceRecord{Translated by S.D.F. Salmond.；Nicene and Post-Nicene Fathers, First Series；Vol. 6.；Edited by Philip Schaff.；Buffalo, NY: Christian Literature Publishing Co.,；1888.；<http://www.newadvent.org/fathers/1602306.htm>.}

% structure: p0001 source=h1 target=section reason=page-title

\section{福音的和谐 第三卷}

% structure: p0002 source=h2 target=subsection reason=relative-heading-rank; base=h2

\subsection{第七章　论福音作者在记载清晨耶稣被解送到\Person{比辣多}前所发生之事时，彼此完全和谐一致；并论及关于主被定价的引文的问题，该引文\Person{玛窦}归诸\Person{耶肋米亚}先知，但在该先知著作中却找不到这样的段落。}

27. \Person{玛窦}接着记载如下：\BibleQuote{到了早晨，众司祭长和民间的长老商议要处死耶稣；他们把祂捆绑后，押解去交给总督\Person{般雀比辣多}。} \Person{马尔谷}的记载与此类似：\BibleQuote{清早，司祭长与长老、经师和全公会议决后，就捆绑耶稣，押解去交给\Person{比辣多}。} \Person{路加}则在其记述\Person{伯多禄}背主之后，似乎在天将破晓时回顾了耶稣所受的苦，并接续叙述如下：\BibleQuote{那些看守耶稣的人戏弄祂，击打祂；又蒙住祂的眼睛，问祂说：\InnerQuote{说说看，打你的是谁？}他们还说了许多亵渎的话反对祂。天一亮，民间的长老、司祭长和经师都聚齐了，把祂带到他们的公议会里，说：\InnerQuote{祢是基督吗？告诉我们。}耶稣对他们说：\InnerQuote{我若告诉你们，你们也不信；我若问你们，你们也不回答我，也不放我走。从今以后，人子要坐在天主权能的右边。}众人便都说：\InnerQuote{那么，祢是天主之子吗？}耶稣对他们说：\InnerQuote{你们说我便是。}他们说：\InnerQuote{我们何需其它见证？我们亲自从祂口中听见了。}于是全体起来，把祂押送到\Person{比辣多}那里。} \Person{路加}记载了这一切。他的叙述包含了某些\Person{玛窦}和\Person{马尔谷}也记载的事实，即主被问是否是天主之子，祂回答说：\BibleQuote{我告诉你们，从此以后，你们要看见人子坐在权能者的右边，驾着天上的云彩降来。}我们又推断这些事发生在天将破晓时，因为\Person{路加}的说法是：\BibleQuote{天一亮。}因此，\Person{路加}的叙述与其它二者相似，尽管他也提到了别人未曾提及的事。我们进一步推断，在夜间，主已面对假见证的磨难——\Person{玛窦}和\Person{马尔谷}简要记载了此事，\Person{路加}却略过不提，但他讲述了黎明时所发生的事。前两位，即\Person{玛窦}和\Person{马尔谷}，连贯地叙述了主直到清晨所受的一切。但之后，他们又回头记载\Person{伯多禄}背主的事；叙述完此事后，他们又回到清晨的事件，引入了为完成对主遭遇的记述而仍需讲述的其他情况。但之前，他们未曾专门描述清晨所发生的事。同样，\Person{若望}在按自己认为需要的方式详细记述了主所行的事，并讲述了\Person{伯多禄}背主的全部故事后，接续叙述如下：\BibleQuote{于是他们把耶稣带到\Person{盖法}那里，进了总督府。那时是清早。}在这里，我们或者可以假设，有某种急需\Person{盖法}在场于总督府的事，或在其他司祭长审问主时他缺席；或者总督府是在他的家中；而且，从这一幕一开始，他们只是把耶稣押送到最终要交给的那个人面前。但由于他们把被告当作已定罪的人带来，且此前\Person{盖法}已判定耶稣该死，于是不再耽搁，将祂交给\Person{比辣多}，以处死祂。因此，\Person{玛窦}在此讲述了\Person{比辣多}与主之间发生的事。\par

28. 然而，他首先岔开话题，讲述\Person{犹达斯}的结局，此事只有他记载。他的叙述如下：\BibleQuote{那时，出卖耶稣的\Person{犹达斯}看见祂已被定罪，就后悔了，把那三十块银钱退还给司祭长和长老，说：\InnerQuote{我犯了罪，因我出卖了无辜的血。}他们说：\InnerQuote{这与我们何干？你自己承当吧。}于是他把银钱扔在圣殿里，就出去上吊死了。司祭长拾起银钱说：\InnerQuote{不可放入献仪箱内，因为这是血价。}他们商议后，用那钱买了陶匠的田，作为埋葬外乡人之处。因此那块田直到今日还叫做\InnerQuote{血田}。这就应验了\Person{耶肋米亚}先知所说的话：\InnerQuote{他们取了那三十块银钱，即以色列子民为祂所估定的价钱，给了陶匠的田，如同主所吩咐我的。} }\par

29. 现今，若有人因这段话未见于耶肋米亚先知的作品而感到困惑，并认为这有损于福音作者的真实性，请他首先注意一个事实：并非所有福音抄本都包含这段话出自耶肋米亚的归属，有些抄本仅记载这段话是\Quote{\Emph{由先知。}}因此，可以断言，那些未包含耶肋米亚名字的抄本更值得遵循。因为这些话确实是由一位先知所说，只是那位先知是匝加利亚。这样一来，可以假设那些包含耶肋米亚名字的抄本是有缺陷的，因为它们本应给出匝加利亚的名字，或者像某些抄本那样根本不提名字，仅说明是\Quote{\Emph{由先知说，}}这先知无疑会被理解为匝加利亚。然而，若有人愿意，可采纳这种辩护方式。但就我而言，我并不满意，原因是大多数抄本都包含耶肋米亚的名字，并且那些在希腊文抄本中更仔细研究福音的评注者报告说，他们在更古老的希腊文样本中发现它也是这样记载的。我还考虑到另一点：即没有理由在原本正确的文本中加上这个名字从而造成讹误；相反，却有明显的理由从许多抄本中删去这个名字。因为当面对这段话在耶肋米亚书中找不到的难题时，鲁莽无知者很容易这样做。\par

30. 那么，如何解释此事呢？只有假设这一切是按照天主眷顾的更隐秘旨意而行的，正是这眷顾引导了圣史们的心思。因为情况可能是这样：当\Person{玛窦}正在编撰他的福音时，根据一种熟悉的经验，他心里冒出了\Person{耶肋米亚}这个名字，而不是\Person{匝加利亚}。然而，这样的不准确之处，他必定会加以改正（因为当他还活在肉身时，肯定会有一些读了这作品的人提醒他注意），如果他当时没有意识到，在回忆情况的过程中（这回忆的过程也由圣神引导），一个先知的名字被提示出来代替另一个，这并非没有目的，而且若非主的旨意要如此记载，这事也不会发生在他身上。然而，若有人问为什么主要如此决定，那么有一个首要且最有益的理由，值得我们立刻思考，即：借此传达了一个概念，关于所有圣先知因同一圣神说话，在他们的话语中持续完全一致的奇妙方式——这情况无疑比所有先知的话都由一个人口述出来更能打动人心。同样的思考也可合适地提示我们有责任毫不迟疑地接受圣神借着这些先知所表达的一切，并把每一位先知个人的传述也视为全体先知的传述，将全体先知的集体传述也视为每一位先知个人的传述。那么，如果实际情况是，\Person{耶肋米亚}所说的话真的同样属于\Person{匝加利亚}，也属于\Person{耶肋米亚}；另一方面，\Person{匝加利亚}所说的话真的同样属于\Person{耶肋米亚}，也属于\Person{匝加利亚}，那么当\Person{玛窦}重读他所写的，发现一个名字代替了另一个名字出现在他心中时，他有何必要改正他的文本呢？那么，他正确的做法难道不是以顺服圣神的权威（他肯定比我们更真切地感到自己的心思处于圣神的引导之下），因此，按照主的旨意和安排，对他如此写下的内容不予改动，目的是让我们明白，先知们在他们的言论中彼此保持如此完全的和谐，以至于我们把原本由\Person{匝加利亚}所说的一句话归给\Person{耶肋米亚}，就绝非荒谬之举，而事实上是一件最合宜的事？因为如果在我们的时代，有人想要提醒我们注意某个人的话，却偶然提到了另一个人的名字，而此人实际上并未说出那些话，但同时也是那真正说话者最亲密的朋友和同伴；如果他随即想起自己给出了一个名字代替另一个，他回过神来更正了这个错误，但却是以这样的方式来做，说\Quote{毕竟，我所说的话并没有错；}那么，我们以为这意指什么呢？岂不正是说明在这两者之间存在着如此完美的感情共鸣——也就是说，这个试图重复某人话语的人所意指的是那个真正说话者，而当时出现在他脑海中的是另一个人的名字代替了前者——以至于将那些话归于前者与归于后者几乎是一回事？那么，在圣先知们的事例上，这种做法岂不更应被理解并特别推荐给我们注意，好让我们接受他们整个系列所编成的书卷，仿佛它们只是由一位作者写成的一本书，在其中，关于所处理的主题，不应认为存在任何不一致之处（因为实际也找不到），而在其中，会有一个比由单一个人（即使是最博学的人）陈述所有这些言论更为显著的一致性、真实性的范例？因此，虽然有些人（无论是无信仰者还是无知之人）试图在这里找到一个论据，来帮助他们证明圣史们之间缺乏和谐，但另一方面，有信德与学问的人却应当将此服务于证明圣先知们所特有的合一性。\par

31. 我还有另一个理由来解释为何圣神的权威允许或更确切地说指引在此处用耶肋米亚的名字代替匝加利亚的名字（对此理由的更充分讨论，我认为必须留待另一机会，以免当前论述超出我们完成此著作所需限定的范围）。耶肋米亚记载他买了同族兄弟儿子的田地，并付了钱。那笔钱的确没有以匝加利亚中所见的具体价银即三十块银钱的名义给出；但另一方面，匝加利亚中并未提及买田之事。显然，福音作者将预言提到三十块银钱解释为只有在主身上才得应验，因而以此作为为祂所定的价钱。然而，耶肋米亚关于买田的话也与同一件事有关，这种关联可能由有意选择引入福音叙述的耶肋米亚之名（他说买田之事）代替匝加利亚之名（我们借此得知三十块银钱）奥秘地表明。这样，读者先读福音见到耶肋米亚之名，再读耶肋米亚未发现关于三十块银钱的段落，却看到买田的章节，就会被教导比较这两段，理解预言的真正含义，并知道它如何也与在主身上显示的预言应验相关。因为玛窦还在所引段落中加了这段话：\BibleQuote{以色列子民给祂估了价，给了那陶工的田地，如同主所吩咐我的。}这些话无论在匝加利亚还是耶肋米亚中都找不到。因此，我们宁可认为这些是福音作者在主的启示下凭己意巧妙而奥秘地插入的——主让他明白，上述预言的真正含义与基督被定价这件事有关。此外，在耶肋米亚中，买田的证据要投入瓦器里。同样，我们在福音中看到，为所付的价钱买了陶工的田地，这田地要用作埋葬外乡人的坟地。这可能都是象征那些在现世如客旅寄居、藉圣洗与基督同葬者安息的恒久。因为主也对耶肋米亚说过，买田表明在那片土地上将有遗民从被掳中得释放。我认为在提请你们注意这些先知见证互相协调并与福音叙述比较时，真正细心刻苦研究应寻求何种意义后，对上述内容作一番勾勒是合适的。这些就是玛窦关于叛徒犹达斯的记述。\par

\SourceRecord{Translated by S.D.F. Salmond.；Nicene and Post-Nicene Fathers, First Series；Vol. 6.；Edited by Philip Schaff.；Buffalo, NY: Christian Literature Publishing Co.,；1888.；<http://www.newadvent.org/fathers/1602307.htm>.}

% structure: p0001 source=h1 target=section reason=page-title

\section{福音的和谐，第三卷}

% structure: p0002 source=h2 target=subsection reason=relative-heading-rank; base=h2

\subsection{第8章 论福音史家对发生在比拉多面前的事件的记述毫无矛盾。}

32. 祂接着如下进行：\BibleQuote{耶稣站在总督面前，总督问祂说：\InnerQuote{祢是犹太人的君王吗？}耶稣对他说：\InnerQuote{你说的是。}当司祭长和长老控告祂时，祂什么也不回答。那时比拉多对祂说：\InnerQuote{祢没有听见他们控告祢多少事吗？}耶稣连一句话也不回答他，以致总督大为惊异。每逢节日，总督照例给民众释放一个他们所愿意的囚犯。那时，有一个出名的囚犯，名叫巴拉巴。当他们聚集在一起时，比拉多对他们说：\InnerQuote{你们愿意我给你们释放哪一个？巴拉巴，还是那称为基督的耶稣？}原来他知道他们是出于嫉妒才把祂解送来的。当比拉多坐在审判座上时，他的妻子派人来对他说：\InnerQuote{你千万不要干涉那义人的事，因为今天我因为祂在梦中受了许多苦。}但是司祭长和长老说服了民众，叫他们要求巴拉巴，而除掉耶稣。总督又回答他们说：\InnerQuote{这两个当中，你们愿意我给你们释放哪一个？}他们说：\InnerQuote{巴拉巴。}比拉多对他们说：\InnerQuote{那么，对于那称为基督的耶稣，我该怎么办？}众人都说：\InnerQuote{该钉祂在十字架上。}总督对他们说：\InnerQuote{祂究竟做了什么恶事？}他们却越发喊叫说：\InnerQuote{该钉祂在十字架上。}比拉多见事情毫无成效，反而更为骚乱，就拿水，当着民众洗手说：\InnerQuote{对这义人的血，我是无罪的，你们自己负责吧。}全体百姓回答说：\InnerQuote{祂的血归在我们和我们的子孙身上。}于是比拉多给他们释放了巴拉巴；并因耶稣加以鞭打后，把祂交给人钉在十字架上。}以上是玛窦所记载的比拉多对主所行的事。\par

33. 马尔谷所记载的，无论在用语还是内容上，也与上述几乎完全相同。不过，当民众要求按节期惯例释放一名囚犯时，比拉多回答他们的话，这位福音史家记载如下：\BibleQuote{比拉多回答他们说：\InnerQuote{你们愿意我给你们释放犹太人的君王吗？}}另一方面，玛窦的记载是这样的：\BibleQuote{你们愿意我给你们释放哪一个？巴拉巴，还是那称为基督的耶稣？}玛窦没有提及民众要求释放一名囚犯，这并不构成难题。但值得一问的是，比拉多实际所说的话究竟是玛窦所记的，还是马尔谷所记的。因为这两种表达形式似乎有差异，即\BibleQuote{你们愿意我给你们释放哪一个？巴拉巴，还是那称为基督的耶稣？}和\BibleQuote{你们愿意我给你们释放犹太人的君王吗？}然而，由于他们习惯称自己的君王为\Quote{受傅者}，有人可能用这个说法，有人可能用那个说法，所以很明显，比拉多问他们是否愿意释放犹太人的君王，即基督。马尔谷只关心与主有关的事，所以这里没有提到巴拉巴，但这无碍于意义的一致性。因为他在记载他们的回答时，已足够清楚地表明他们所要求释放的人是谁。他的版本是：\BibleQuote{但是司祭长煽动民众，宁愿给他释放巴拉巴。}然后他接着补充说：\BibleQuote{比拉多又回答他们说：\InnerQuote{那么，你们愿意我怎样对待你们所称的犹太人的君王呢？}}这就足够清楚地表明，马尔谷所说的\Quote{犹太人的君王}正是玛窦所说的\Quote{基督}。因为只有犹太人才称君王为\Quote{受傅者}；而玛窦记载比拉多的话是：\BibleQuote{比拉多对他们说：\InnerQuote{那么，对于那称为基督的耶稣，我该怎么办？}}于是马尔谷继续：\BibleQuote{他们又喊叫说：\InnerQuote{钉祂在十字架上。}}玛窦的记载是：\BibleQuote{众人都对他说：\InnerQuote{把祂钉在十字架上。}}马尔谷接着：\BibleQuote{比拉多对他们说：\InnerQuote{祂做了什么恶事？}他们却越发喊叫说：\InnerQuote{钉祂在十字架上。}}玛窦没有记录这段；但他记载了\BibleQuote{比拉多见事情毫无成效，反而更为骚乱}，并告诉我们他如何当着民众洗手以表明自己对这义人的血无罪（马尔谷和其他福音史家没有记载此事）。这样他也足够清楚地表明总督如何设法让民众同意释放祂。马尔谷简短地提到这一点，他说比拉多说了\BibleQuote{\InnerQuote{祂做了什么恶事？}}然后马尔谷也以此结束他对比拉多与主之间发生的事情的记载：\BibleQuote{比拉多愿意满足民众，就给他们释放了巴拉巴，并把耶稣鞭打后，交给他们钉在十字架上。}以上是马尔谷对发生在总督面前之事的叙述。\par

34. \Person{路加}记载了发生在\Person{比拉多}面前的经过如下：\BibleQuote{他们便开始控告祂，说：\InnerQuote{我们查出这人煽惑我们的民族，禁止纳税给\Person{凯撒}，并说他自己是\Person{基督}，是君王。}} 前两位圣史没有记录这些话，虽然他们确实提及这些人控告了祂。因此，\Person{路加}是具体说明对祂提出的虚假控告内容的那一位。另一方面，他没有记载\Person{比拉多}对祂说：\BibleQuote{\InnerQuote{你什么都不回答吗？你看，他们提出多少证据控告你。}} \Person{路加}没有引入这些话，而是继续叙述其他事情，这些事情也由那两位圣史报告了。于是他继续写道：\BibleQuote{\Person{比拉多}问祂说：\InnerQuote{你是犹太人的君王吗？}祂回答他说：\InnerQuote{你说的是。}} \Person{玛窦}和\Person{马尔谷}也同样记载了这件事，是在叙述耶稣因不回答控告者而被责问之前。然而，\Person{路加}叙述这些事情的顺序丝毫不影响事实；同样，某一位圣史略过不提而另一位明确指明的事件，也无关紧要。接下来的事就是一例：\BibleQuote{于是\Person{比拉多}对司祭长及群众说：\InnerQuote{我在这人身上查不出什么罪状来。}他们却更加激烈地说：\InnerQuote{祂在犹太全境施教，从\Place{加里肋亚}开始直到这里，煽动民众。}\Person{比拉多}一听说\Place{加里肋亚}，便问这人是否是\Place{加里肋亚}人。他既知祂属\Person{黑落德}管辖，就把他送到\Person{黑落德}那里，那时\Person{黑落德}也正在\Place{耶路撒冷}。\Person{黑落德}见了耶稣，非常高兴；因为他原已久闻耶稣的事，渴望见祂，并且希望看到祂行些奇迹。于是他问祂许多话，但祂什么也不回答。司祭长和经师们站在那里，极力控告祂。\Person{黑落德}同他的兵士们藐视并戏弄祂，给祂换上华丽的衣服，将祂送回\Person{比拉多}那里。\Person{黑落德}和\Person{比拉多}就在那一天成了朋友：因为他们彼此以前有仇。} 这些都是\Person{路加}独自叙述的，即主被\Person{比拉多}送到\Person{黑落德}那里，以及当时所发生的事。与此同时，他在这段话中所作的某些陈述，与其他圣史在叙述的不同部分中所报告的事情有相似之处。然而，那些圣史的即时目的是简单地叙述从主被交给钉十字架之前，在\Person{比拉多}面前所做的一切事。而\Person{路加}则按照自己的计划，为了讲述与\Person{黑落德}发生的事作了上述插叙；之后他又回到在总督面前所发生的历史。于是他继续叙述如下：\BibleQuote{\Person{比拉多}召集了司祭长、首领和民众，对他们说：\InnerQuote{你们把这人带到我这里，说祂煽惑民众；看，我在你们面前审问了祂，关于你们控告祂的这些事情，我在祂身上查不出什么罪状来。}} 这里我们注意到，他省略了\Person{比拉多}问主要回答控告之事。之后他继续写道：\BibleQuote{\Person{黑落德}也没有查出什么：因为我将你们送到了他那里，看，祂并没有做什么该死的事。因此，我要责打祂，然后将祂释放。因为在庆节上，他必须给他们释放一个人。于是他们齐声喊叫说：\InnerQuote{除掉这人，给我们释放\Person{巴辣巴}。}\Person{巴辣巴}是因为在城里作乱和杀人而被下狱的。\Person{比拉多}愿意释放耶稣，就又对他们说话。但他们喊叫说：\InnerQuote{钉死祂，钉死祂！}\Person{比拉多}第三次对他们说：\InnerQuote{祂做了什么恶事？我在祂身上没查出死罪来。因此，我要责打祂，然后将祂释放。}他们却更大声地喊叫，要求钉死祂；他们的喊声占了上风。} \Person{比拉多}一心想要释放耶稣，多次设法争取民众的同意，这由\Person{玛窦}用寥寥数语就充分证明了，他说：\BibleQuote{\Person{比拉多}见事情无法挽回，反而发生骚动。} 他之所以这样写，是因为\Person{比拉多}确实在这方面下了很大功夫，尽管他没有详述他多少次试图从他们的愤怒中解救耶稣。因此，\Person{路加}对总督面前所发生的事的叙述作结如下：\BibleQuote{于是\Person{比拉多}决定满足他们的要求。他释放了他们所要求的那个因作乱和杀人入狱的人，但把耶稣交给了他们的意愿。}\par

35. 现在让我们来看若望所记载的同一事件——即那些涉及比拉多的事件——的记述。他如此写道：\Quote{他们自己却不进入总督府，免得受玷污，但为能吃逾越节。于是比拉多出来到他们那里，说：\InnerQuote{你们对这个人提出什么控告？}他们回答说：\InnerQuote{如果这个人不是作恶的，我们绝不会把他交给你。}}我们必须查看这段经文，以表明它与路加所记载的版本没有任何不一致之处。路加的记载指出，有人对祂提出了某些指控，并且也说明了这些指控的内容。因为路加的话是这样：\Quote{他们就开始控告祂，说：\InnerQuote{我们发觉这个人煽惑我们的民族，阻止给凯撒纳税，并说他自己是基督，是君王。}}另一方面，根据我刚才从若望引用的段落，当比拉多问他们\Quote{你们对这个人提出什么控告？}时，犹太人似乎不愿提出任何具体的指控。因为他们回答说：\Quote{如果这个人不是作恶的，我们绝不会把他交给你；}其含义是，他应当接受他们的权威，不再追问对祂提出的过失是什么，并仅仅因为祂已被他们认为配得由他们交给他而相信祂有罪。既然如此，我们应当认为这两份记述都报告了实际说过的话——目前我们面前的这份和路加的那份都是。因为在双方众多的言论和回答中，这些作者根据自己的判断尽可能进行了选择，每个人都在他的叙述中引入了他认为足够的内容。另外，若望自己也提到了对祂提出的某些指控，我们将在适当的地方看到这些指控。在这里，他接着这样写道：\Quote{于是比拉多对他们说：\InnerQuote{你们自己把祂带去，按照你们的法律审判祂吧。}犹太人对他说：\InnerQuote{我们不可以处死任何人；}这是为应验耶稣所说的话，表明祂将怎样死去。然后比拉多又进了总督府，叫耶稣来，对祂说：\InnerQuote{你是犹太人的君王吗？}耶稣回答说：\InnerQuote{你自己说这话，还是别人对你论及我的？}}这一点似乎与其他人所记载的——即\Quote{耶稣回答说：\InnerQuote{你说的是。}}——不一致，除非在后续内容中明确说明那一句话与另一句话都被说过。因此，他让我们明白，他接下来记录的事情[不应被视为]主从未实际说过的话，而是应被视为其他福音作者所略过的事情。因此，请看他叙述的其余部分。它这样写道：\Quote{比拉多回答说：\InnerQuote{我难道是犹太人吗？你本国的人和司祭长把你交给了我：你做了什么事？}耶稣回答说：\InnerQuote{我的王国不属于这世界；如果我的王国属于这世界，我的臣仆就会为我争战，使我不致被交给犹太人；但现在我的王国不是从这里来的。}比拉多于是对祂说：\InnerQuote{那么，你是君王吗？}耶稣回答说：\InnerQuote{你说我是君王。}}看，这就是他到达其他福音作者所报告之点的所在。然后他继续下去，主仍然是发言者，叙述其余没有记载的其他事情。他的话是这样：\Quote{我为此而诞生，也为此而来到世界，为给真理作证。凡属于真理的，都听我的声音。比拉多对他说：\InnerQuote{什么是真理？}说了这话，他又出去到犹太人那里，对他们说：\InnerQuote{我在这人身上查不出什么罪来。但你们有个惯例，在逾越节我该给你们释放一个人：那么，你们愿意我给你们释放这犹太人的君王吗？}他们又都喊叫说：\InnerQuote{不要这人，而要巴拉巴。}巴拉巴是个强盗。于是比拉多把耶稣带去鞭打了。士兵们用荆棘编了一个冠冕，戴在祂头上，给祂披上一件紫红袍；他们来到祂面前说：\InnerQuote{犹太人的君王，万岁！}并用手掌打祂。比拉多又出去对他们说：\InnerQuote{看，我给你们带出祂来，为叫你们知道我查不出祂有什么罪。}耶稣就出来了，戴着荆棘冠冕，穿着紫红袍。比拉多对他们说：\InnerQuote{看，这个人！}司祭长和差役们一看见祂，就喊叫说：\InnerQuote{钉祂在十字架上！钉祂在十字架上！}比拉多对他们说：\InnerQuote{你们自己把祂带去钉十字架吧；因为我查不出祂有什么罪。}犹太人回答他说：\InnerQuote{我们有法律，按那法律祂应当死，因为祂自命为天主子。}}这可以与路加所报告的犹太人在控告中提出的说法相吻合——即\Quote{我们发觉这个人煽惑我们的民族}——这样我们就可以在这里附加它所给出的理由：\Quote{因为祂自命为天主子。}若望接着这样写道：\Quote{比拉多一听这话，越发害怕，又进了总督府，对耶稣说：\InnerQuote{你从哪里来？}耶稣没有回答他。于是比拉多对祂说：\InnerQuote{你不跟我说话吗？你不知道我有权释放你，也有权钉你在十字架上吗？}耶稣回答说：\InnerQuote{除非赐给你从上头来的，你对我没有任何权柄；所以把我交给你的那个人有更大的罪。}从此比拉多想要释放祂；但犹太人喊叫说：\InnerQuote{如果你释放这个人，你就不是凯撒的朋友；凡是自命为王的，就是反对凯撒。}}这完全符合路加在涉及犹太人所提出的上述控诉时所记载的内容。因为在\Quote{我们发觉这个人煽惑我们的民族}这句话之后，他又加上了\Quote{并且阻止给凯撒纳税，还说他自己是基督，是君王。}这也为前面提到的难题提供了解决之道，即可能似乎有理由认为若望暗示犹太人没有对主提出任何具体的指控，因为他们回答他说：\Quote{如果这个人不是作恶的，我们绝不会把他交给你。}若望接着这样写道：\Quote{比拉多听了这话，就把耶稣带出来，在一个名叫石铺地的地方——希伯来语叫加巴萨——坐上了审判席。那是逾越节预备日，大约第六时辰；他对犹太人说：\InnerQuote{看，你们的君王！}但他们喊叫说：\InnerQuote{除掉祂！钉祂在十字架上！}比拉多对他们说：\InnerQuote{我可以把你们的君王钉十字架吗？}司祭长回答说：\InnerQuote{除了凯撒，我们没有君王。}于是他把耶稣交给他们去钉十字架。}以上是若望关于比拉多所行事迹的记述。\par

\SourceRecord{Translated by S.D.F. Salmond.；Nicene and Post-Nicene Fathers, First Series；Vol. 6.；Edited by Philip Schaff.；Buffalo, NY: Christian Literature Publishing Co.,；1888.；<http://www.newadvent.org/fathers/1602308.htm>.}

% structure: p0001 source=h1 target=section reason=page-title

\section{\WorkTitle{福音的和谐，卷三}}

% structure: p0002 source=h2 target=subsection reason=relative-heading-rank; base=h2

\subsection{第九章 论祂在\Person{比拉多}卫队手中所受的嘲弄，以及三位报道该场景的福音作者——即\Person{玛窦}、\Person{马尔谷}和\Person{若望}——之间的谐和}

36. 现在我们到了可以研读主严格意义上的苦难的阶段，正如这四位福音作者在叙述中所呈现的那样。\Person{玛窦}这样开始他的记述：\BibleQuote{那时，总督的兵士把耶稣带进总督府，召集了全队兵士围着祂。他们脱去祂的衣服，给祂披上一件朱红袍，又用荆棘编了一个冠冕，戴在祂头上，把一根芦苇放在祂右手里；然后跪在祂面前，戏弄祂说：\InnerQuote{犹太人的君王，万岁！}}在同一叙述阶段，\Person{马尔谷}这样写道：\BibleQuote{兵士把耶稣带进庭院，即总督府；他们召集了全队兵士，给祂穿上紫袍，又编了一个荆棘冠冕，戴在祂头上，开始向祂致敬说：\InnerQuote{犹太人的君王，万岁！}然后用芦苇打祂的头，并向祂吐唾沫，屈膝朝拜祂。}因此，我们在此看到，\Person{玛窦}说他们给祂披上\Quote{朱红袍}，而\Person{马尔谷}却提到紫袍，祂被穿上紫袍。解释可能是，这些戏弄者是用那件朱红袍来代替王室的紫袍。此外，有一种偏红的紫色，与朱红色非常接近。也可能\Person{马尔谷}注意到了袍子所含的紫色，尽管它本质上是朱红色。\Person{路加}对此略而不提。另一方面，在叙述\Person{比拉多}把耶稣交出去钉十字架之前，\Person{若望}插入了以下段落：\BibleQuote{于是\Person{比拉多}把耶稣带去鞭打了。兵士们用荆棘编了一个冠冕，戴在祂头上，给祂披上一件紫袍，来到祂跟前说：\InnerQuote{犹太人的君王，万岁！}并用手掌打祂。}这清楚地表明，\Person{玛窦}和\Person{马尔谷}是以回顾性重述的方式报道了这件事，而它实际上并非发生在\Person{比拉多}把祂交出钉十字架之后。因为\Person{若望}足够明确地告诉我们，这些事情发生在祂仍在\Person{比拉多}面前的时候。因此我们断定，其他福音作者在那一特定点引入此事，只是因为之前略过了，此时记起而插入。\Person{玛窦}接着叙述的内容也证实了这一点。他继续写道：\BibleQuote{他们遂向祂吐唾沫，拿芦苇打祂的头。戏弄完了，就给祂脱去袍子，穿上祂自己的衣服，带祂出去，要钉死祂。}这里我们得知，脱下袍子、穿上祂自己的衣服是在结束时，当祂被带走时进行的。\Person{马尔谷}这样记述：\BibleQuote{戏弄完了，就给祂脱去紫袍，穿上祂自己的衣服。}\par

\SourceRecord{Translated by S.D.F. Salmond.；Nicene and Post-Nicene Fathers, First Series；Vol. 6.；Edited by Philip Schaff.；Buffalo, NY: Christian Literature Publishing Co.,；1888.；<http://www.newadvent.org/fathers/1602309.htm>.}

% structure: p0001 source=h1 target=section reason=page-title

\section{福音的和谐 第三卷}

% structure: p0002 source=h2 target=subsection reason=relative-heading-rank; base=h2

\subsection{第十章 论如何协调\Person{玛窦}、\Person{马尔谷}和\Person{路加}所述另一个人被迫背负\Person{耶稣}的十字架与\Person{若望}所说\Person{耶稣}亲自背负十字架的说法。}

37. 因此，\Person{玛窦}继续叙述如下：\BibleQuote{他们出来时，遇见一个\Place{基勒乃}人，名叫\Person{西满}，就强迫他背祂的十字架。} 同样，\Person{马尔谷}说：\BibleQuote{他们把祂带出去要钉十字架。他们强迫一个从乡下来、路过的\Place{基勒乃}人\Person{西满}——他是\Person{亚历山大}和\Person{鲁富斯}的父亲——背祂的十字架。} \Person{路加}的记载也如此：\BibleQuote{他们带走祂时，抓住一个从乡下来的\Place{基勒乃}人\Person{西满}，把十字架放在他身上，叫他背着跟随\Person{耶稣}。} 另一方面，\Person{若望}记载如下：\BibleQuote{他们抓住\Person{耶稣}，把祂带走。祂自己背着十字架出来，到了一个地方，叫髑髅地，希伯来语叫\Place{哥耳哥达}，他们就在那里把祂钉在十字架上。} 由此我们明白，\Person{耶稣}在前往所说的地方时，是亲自背着十字架的。但在路上，上文提到的\Person{西满}——其他三位圣史所提及的——被强迫服役，接过了十字架，背着走完余下的路程，直到到达地点。因此我们发现，两种情形确实都发生了：先是\Person{若望}所记载的那种，然后是其他三位所举例的那种。\par

\SourceRecord{Translated by S.D.F. Salmond.；Nicene and Post-Nicene Fathers, First Series；Vol. 6.；Edited by Philip Schaff.；Buffalo, NY: Christian Literature Publishing Co.,；1888.；<http://www.newadvent.org/fathers/1602310.htm>.}

% structure: p0001 source=h1 target=section reason=page-title

\section{\WorkTitle{福音的和谐，第三卷}}

% structure: p0002 source=h2 target=subsection reason=relative-heading-rank; base=h2

\subsection{第十一章 关于\Person{玛窦}的记载与\Person{马尔谷}的记载在给祂饮用的饮料一事上的一致性——该叙述是在祂被钉十字架的叙事之前引入的}

38. \Person{玛窦}接着这样说：\BibleQuote{他们来到了一个名叫\Place{哥耳哥达}的地方，即\InnerQuote{髑髅}的地方。} 就地点而言，他们显然是一致的。同一\Person{玛窦}接着又说：\BibleQuote{他们拿苦胆调和的酒给祂喝；祂只尝了尝，便不愿意喝。} \Person{马尔谷}的记载如下：\BibleQuote{他们将没药调和的酒给祂喝，祂却没有接受。} 在此，我们可以理解，当\Person{玛窦}说酒是\Quote{苦胆调和}时，他表达了与\Person{马尔谷}相同的含义。因为提到苦胆是为了强调饮料的苦味。而没药调和的酒以其苦味著称。事实上，也可能是苦胆和没药一起使酒变得极其苦涩。此外，当\Person{马尔谷}说\Quote{祂没有接受}时，我们理解这句话的意思是祂没有接受以致实际喝下。然而，祂确实尝了，正如\Person{玛窦}所证实的。因此，\Person{马尔谷}的话\Quote{祂没有接受}表达了与\Person{玛窦}的版本\Quote{祂不愿意喝}相同的含义。但前者（\Person{马尔谷}）却没有提及祂尝了这饮料。\par

\SourceRecord{Translated by S.D.F. Salmond.；Nicene and Post-Nicene Fathers, First Series；Vol. 6.；Edited by Philip Schaff.；Buffalo, NY: Christian Literature Publishing Co.,；1888.；<http://www.newadvent.org/fathers/1602311.htm>.}

% structure: p0001 source=h1 target=section reason=page-title

\section{福音的和谐 第三卷}

% structure: p0002 source=h2 target=subsection reason=relative-heading-rank; base=h2

\subsection{第十二章 论四位福音圣史关于分祂衣服一事所保存的一致性}

39. \Person{玛窦}接着写道：\BibleQuote{他们钉死祂以后，分取了祂的衣服，抽签拈阄；然后坐下，看守着祂。} \Person{马尔谷}记载同一事件如下：\BibleQuote{他们钉祂在十字架上以后，分取了祂的衣服，抽签拈阄，看谁得着什么。} 同样，\Person{路加}说：\BibleQuote{他们分取了祂的衣服，抽签拈阄。民众站着观望。} 前三部福音如此简略地记载了这件事。但\Person{若望}给了我们关于此事如何进行的更详细叙述。他的版本如下：\BibleQuote{那时，兵士钉死耶稣以后，拿走了祂的衣服，分成四份，每个兵士一份；还有祂的里衣：那件里衣是无缝的，从上到下浑然织成的。于是他们彼此说，我们不要撕破它，我们抽签，看是谁的。这就应验了经上的话：\InnerQuote{他们分了我的衣服，为我的里衣拈阄。}}\par

\SourceRecord{Translated by S.D.F. Salmond.；Nicene and Post-Nicene Fathers, First Series；Vol. 6.；Edited by Philip Schaff.；Buffalo, NY: Christian Literature Publishing Co.,；1888.；<http://www.newadvent.org/fathers/1602312.htm>.}

% structure: p0001 source=h1 target=section reason=page-title

\section{\WorkTitle{福音的和谐 第三卷}}

% structure: p0002 source=h2 target=subsection reason=relative-heading-rank; base=h2

\subsection{第十三章 论主受难的时辰，以及马尔谷与若望在\Quote{第三时辰}和\Quote{第六时辰}记载中无矛盾的问题}

40. 玛窦接着写道：\Quote{他们在祂头上安放了祂的罪状，写着：\InnerQuote{这是耶稣，犹太人的君王。}} 另一方面，马尔谷在作此陈述之前，插入了这句话：\BibleQuote{那时是第三时辰，他们就钉了祂在十字架上。} 他是在述及分衣之事后立即接上此语的。因此，我们必须特别谨慎地考虑此事，以免出现严重错误。因为有人认为主必定是在第三时辰被钉十字架；此后，从第六时辰到第九时辰，黑暗笼罩大地。按照这种说法，我们就得理解为从祂被钉到黑暗降临之间经过了三个时辰。若非若望记载比拉多坐在审判台上（那地方叫石铺地，希伯来语则叫加巴达）时大约是第六时辰，这种观点或许可以有理有据地坚持。因为他的叙述如下：\BibleQuote{正是逾越节的预备日，大约第六时辰；他对犹太人说：看，你们的君王！但他们喊叫说：除掉，除掉！钉死祂！比拉多对他们说：我可以钉死你们的君王吗？司祭长回答说：除了凯撒，我们没有君王。于是他将耶稣交给他们，去钉十字架。} 因此，如果耶稣是在大约第六时辰、比拉多正坐审判台时被交给犹太人钉十字架，那么祂怎会在第三时辰被钉呢？正如有些人因误解马尔谷的话而认为的那样。\par

41. 首先，让我们思考祂实际被钉的时辰；然后我们再看马尔谷为何记载祂在第三时辰被钉。当时大约是第六时辰，比拉多正如所说正坐在审判台上，将祂交出钉十字架。经文并非说完全第六时辰，只说\Quote{大约第六时辰}；也就是说，第五时辰完全过去，第六时辰也已开始。然而，这些作者自然不能使用诸如第五时辰又一刻、或第五时辰又三分之一、或第五时辰又一半之类的说法。因为圣经习惯简单地使用整数，不提及分数，尤其是在时间方面。我们有一个例子：\Quote{八天}之后，正如他们所说，祂上山去了——玛窦和马尔谷则记载为\Quote{六天以后}，因为他们只考虑从开始计算的那一天到结束的那一天之间的天数。这一点特别要注意，因为我们看到若望在此用词多么有分寸。因为他说的不是\Quote{第六时辰}，而是\Quote{大约第六时辰}。即使他没有这样表达，而是仅说第六时辰，我们仍可以根据我在圣经中熟悉的说话方式（即使用整数）来解释这个短语。这样，我们完全可以将意思合理地理解为：在第五时辰结束、第六时辰开始时，那些与主被钉相关的事正在进行，直到第六时辰结束、祂挂在十字架上时，发生了三位福音传道者（玛窦、马尔谷和路加）所证实的黑暗。\par

42. 按适当的顺序，让我们现在探究：马尔谷在告诉我们，当他们钉死祂时，分了祂的衣服，拈阄决定各人得什么之后，又加上了这句话：\Quote{到了第三时辰，他们就钉死了祂。} 他在此已经说过：\Quote{钉死祂时，他们分了祂的衣服}；其他福音作者也证实，当祂被钉死时，他们分了祂的衣服。因此，如果马尔谷的目的是要指明事件发生的时间，他只需简单地说\Quote{到了第三时辰}就够了。那么，他加上这些话\Quote{他们钉死了祂}，除了在他的概括陈述中，有意暗示一件值得思考的事——当这段圣经被诵读时，整个教会完全知道主被悬在木头上是哪个时辰，并且有方法纠正作者的错误或驳斥他的不真实——还有什么理由呢？但是，由于他十分清楚主是被士兵悬在十字架上的，而不是被犹太人，正如若望明确肯定的，他隐藏的目的（引出所述条文）是传达这样的观念：那些呼喊祂应该被钉死的人，才是主真正的钉死徒，而非那些仅仅遵照职责为长官服役的人。因此我们明白，当犹太人呼喊主应该被钉死时，正是第三时辰。因此，这最真实地暗示，这些人在他们呼喊的时候，确实钉死了基督。这一点更值得注意，因为他们不愿意显出是自己做了这事，因此将祂交给比拉多，正如若望所记载的他们的话足够清楚地表明的那样。因为，在叙述比拉多如何对他们说：\Quote{你们控告这个人什么？}之后，他的记载继续道：\Quote{他们回答说：如果这人不是作恶的，我们便不会把他交给你。比拉多于是对他们说：你们带他去，按你们的法律审判他。犹太人于是对他说：我们不可杀死任何人。} 因此，他们特别不愿意显出做了的事，正是马尔谷在此表明他们在第三时辰实际上做了的。因为他最真实地判断，杀害主的不是士兵的手，而是犹太人的舌头。\par

43. 此外，如果有人声称，当犹太人第一次以所述的话呼喊时，并非第三时辰，他只是极其疯狂地显示自己是福音的敌人；除非他或许能证明自己能够提出某种新的解决办法。因为他不可能确立所指的时候不是第三时辰的立场。因此，我们当然更应相信一位诚实的福音作者，而非人们好争辩的猜疑。但你会问：你怎么证明是第三时辰？我回答：因为我相信福音作者；如果你也相信他们，请告诉我，主如何能在第六时辰和第三时辰都被钉死？因为，坦诚地说，我们无法绕过若望叙述中的第六时辰；而马尔谷记载了第三时辰；因此，如果我们两人都接受这些作者的话，请告诉我任何其他方式，使这两个时间记载都能按字面正确理解。如果你能做到，我将非常乐意地接受。因为我宝贵的不是自己的意见，而是福音的真理。我确实希望别人能发现更多澄清这个问题的方法。然而，在那之前，如果你愿意，请和我一起利用我已经提出的解决办法。因为如果找不到任何解释，这个本身就足够了。但如果可以设计出另一个，当它被阐明时，我们将做出选择。只是不要认为这是一个不可避免的结论：四位福音作者中的任何一位说了假话，或是在如此崇高和神圣的权威职位上陷入了错误。\par

44. 再者，如果有人声称他能证明，群众在上述的呼喊中并非发生在第三时辰，因为根据马尔谷的记载：\BibleQuote{比拉多又回答他们说：那么，你们要我给你们处置这称为犹太人之王的人呢？他们又喊叫说：钉死他。}此后在同一位圣史叙述中我们找不到其他细节，而是立刻读到主被比拉多交出去钉十字架——若望记载这件事大约发生在第六时辰——我再说，如果有人提出这样的论点，就让他明白，这里有许多事情被略过未记，其间比拉多正设法寻找解救耶稣脱离犹太人之手的方法，并且极尽全力用一切手段抵制他们疯狂的要求。因为玛窦说：\BibleQuote{比拉多对他们说：那么，我该怎么办那称为基督的耶稣呢？众人都说：该钉死他。}于是我们断言那时是第三时辰。同一玛窦接着添加了这一句：\BibleQuote{比拉多见事无济，反而更为骚动，}我们理解已经过了两个时辰，这段时间比拉多试图释放耶稣，犹太人则骚动破坏他的努力，而第六时辰已经开始，在这时辰结束之前，那些记述中在比拉多交出主到黑暗来临之间发生的事就发生了。再者，关于玛窦以上记载的——即\BibleQuote{当他坐在审判座上时，他的妻子派人来对他说：你切勿干涉那义人的事，因为我今天在梦中为他受了许多苦}——我们注意到，比拉多实际上是在稍后才坐上审判席的，但在玛窦叙述较早发生的事件时，他想起了比拉多妻子的报告，决定在此处插入，以便让我们明白比拉多有一个特别急迫的理由，即使到最后也不愿将主交给犹太人。\par

45. 路加在提到比拉多说\BibleQuote{那么，我惩戒他以后，便释放他}之后，告诉我们全体群众都喊叫说：\BibleQuote{除掉这个人，给我们释放巴拉巴。}但也许他们还没有喊叫\BibleQuote{钉死他！}因为路加接着这样说：\BibleQuote{比拉多愿意释放耶稣，便再次向他们讲话。但他们喊叫说：钉死他，钉死他！}据理解，这是在第三时辰。路加接着以这些话继续：\BibleQuote{他第三次对他们说：这人做了什么恶事？我在他身上找不到死罪；那么，我惩戒他以后，便释放他。他们却大声喊叫，要求把他钉死，他们的喊声占了优势。}因此，这位圣史也清楚地表明当时有大的骚动。为了我研究真理的目的，我们可以足够准确地收集究竟过了多少时间之后他才对他们说了这些话：\BibleQuote{这人做了什么恶事？}而当他随后补充说\BibleQuote{他们却大声喊叫，要求把他钉死，他们的喊声占了优势}时，谁能看不出这喊声正是因为他们看到比拉多不愿将主交出来而发出的呢？由于他极不情愿交出他，他当然不会立刻屈服，而是实际上犹豫了两个多时辰。\par

46. 同样询问\Person{若望}，看看\Person{比拉多}这种犹豫有多强烈，以及他如何回避这耻辱的差事。因为这位福音作者更详细地记录了这些事件，虽然他当然也没有提到从第二时辰到第六时辰一部分的所有事情。在叙述了\Person{比拉多}鞭打\Person{耶稣}，任凭士兵戏弄并给他穿上袍子，忍受虐待和许多嘲弄之后（我相信\Person{比拉多}所有这些许可，确实是为了缓解他们的愤怒，使他们不再执意要疯狂地杀死他），\Person{若望}继续他的叙述如下：\BibleQuote{\Person{比拉多}又出来，对他们说：看，我带他出来见你们，叫你们知道我查不出他有什么罪。\Person{耶稣}出来，戴着荆棘冠冕，穿着紫红袍。\Person{比拉多}对他们说：你们看这个人！}目的是让他们看到这耻辱的景象而平息下来。但福音作者继续：\BibleQuote{司祭长和差役看见他，就喊着说：钉死他，钉死他！}那时是第三时辰，正如我们所坚持的。注意随后的话：\BibleQuote{\Person{比拉多}对他们说：你们自己带他去钉十字架吧，因为我查不出他有什么罪。\Person{犹太人}回答说：我们有法律，按那法律他应该死，因为他自命为天主子。\Person{比拉多}听了这话，更加害怕，又进了总督府，对\Person{耶稣}说：你是从哪里来的？\Person{耶稣}却没有回答。\Person{比拉多}对他说：你不对我说话吗？你不知道我有权释放你，也有权钉死你吗？\Person{耶稣}回答说：若不是从上头赐给你的，你根本没有权柄办我；所以把我交给你的那人罪更重。从此\Person{比拉多}想要释放他。}当这里说\BibleQuote{\Person{比拉多}想要释放他}时，我们可以推测他在这项努力中花了多长时间，这里省略了多少\Person{比拉多}所说的话，以及\Person{犹太人}提出的反对，直到这些\Person{犹太人}说出打动他并使他屈服的话。因为作者继续这样说：\BibleQuote{但\Person{犹太人}喊着说：你若释放这个人，就不是凯撒的朋友；凡自命为王的，就是反对凯撒。\Person{比拉多}听了这话，就带\Person{耶稣}出来，到了一个地方叫铺华石处，希伯来话叫加巴塔，就在那里坐堂。那天是逾越节的预备日，约在第六时辰。}这样，从\Person{犹太人}第一次喊叫\BibleQuote{钉死他}那时是第三时辰，到他坐堂的此刻，两个小时过去了，这期间\Person{比拉多}试图拖延，\Person{犹太人}骚动；此时第五时辰已经过完，第六时辰也过去了一些。然后叙述继续：\BibleQuote{他对\Person{犹太人}说：看，你们的王！但他们喊着说：除掉他，除掉他！钉死他！}但即使现在，\Person{比拉多}也没有因为担心他们控告自己而如此屈服。因为他的妻子在他此时坐堂时派人来告诉他——这是\Person{玛窦}所记录的，他是唯一记载此事的人，并提前叙述，在他认为合适的时机插入。如此，\Person{比拉多}仍然继续努力阻止事态进一步恶化，便对他们说：\BibleQuote{我要钉死你们的王吗？}于是\BibleQuote{司祭长回答说：除了凯撒，我们没有王。于是\Person{比拉多}将\Person{耶稣}交给他们去钉十字架。}在他上路、与两个强盗同钉十字架、分衣服并拈阄决定他的里衣、以及各种侮辱加在他身上的时候（因为在这些事进行的同时，也向他投以嘲弄），第六时辰完全过去，黑暗降临，这是\Person{玛窦}、\Person{马尔谷}和\Person{路加}所提到的。\par

47. 因此，让这种不虔敬的固执终止吧，并且让我们相信，主耶稣基督既在第三时辰被犹太人的呼声钉死，又在第六时辰被士兵的手钉死。因为在这些犹太人方面的骚动和比拉多方面的激动中，从他们爆发出喊声\Quote{钉死祂}时起，已经过去了两个多小时。但是，即使是在所有圣史中最追求简洁的\Person{马尔谷}，也乐意简明地指出\Person{比拉多}的意愿和拯救主生命的努力。因为，在给出这句话\Quote{他们又喊说：\InnerQuote{钉死祂}}（在此他使我们明白他们在此之前已经喊过，当时他们要求释放\Person{巴辣巴}给他们）之后，他接着附上这些话：\Quote{于是\Person{比拉多}继续对他们说：\InnerQuote{祂作了什么恶事？}}因此，通过这一简短的句子，他让我们了解了那些需要很长时间才能完成的事情。然而，同时考虑到正确理解他的意思，他没有说\Quote{于是\Person{比拉多}\Emph{说}对他们}，而是这样表达：\Quote{于是\Person{比拉多}\Emph{继续说}对他们：\InnerQuote{祂作了什么恶事？}}因为，如果他的措辞是\Quote{\Emph{说}}，我们可能会理解为这些话只说了一次。但是，通过采用\Quote{\Emph{继续说}}这个说法，他向有智慧的人清楚地表明，\Person{比拉多}反复说了很多次。因此，让我们思考一下，与\Person{玛窦}相比，\Person{马尔谷}表达此事是多么简洁；与\Person{路加}相比，\Person{玛窦}是多么简洁；与\Person{若望}相比，\Person{路加}是多么简洁；同时，这些作者中的每一位都引入了一些自己特有的内容，有时是这一件，有时是那一件。最后，让我们也思考一下，与所发生的众多事情以及它们发生所占据的时间相比，即使是\Person{若望}本人的叙述也是多么简洁。让我们放弃反对的疯狂，并相信在所指的间隔内可能已经过去了两个多小时。\par

48. 然而，如果有人断言，如果情况确实如此，\Person{马尔谷}本应在真正是第三时辰的时刻明确提及第三时辰，即当犹太人的声音响起要求主被钉死时；并且，他还可能在那里清楚地告诉我们，那些叫嚷者确实在那时钉死了祂——这样的推理者只是以自己傲慢的骄傲对真理的历史家施加法则。因为他同样可以主张，如果他本人是这些事件的叙述者，那么所有其他作者都应该以与他自己记录完全相同的方式和顺序来记录它们。因此，让他满足于认为自己的见解低于圣史\Person{马尔谷}的见解，\Person{马尔谷}认为正确的做法是在由天主默感所提示的那个时刻插入陈述。因为那些历史家的回忆是由统治众水的那位的手所掌管，正如经上所载，按照祂自己的美意。因为人的记忆穿越各种思想，任何人都无法控制进入他脑海的主题或其提示的时间。因此，鉴于那些神圣而真实的人，在这叙述顺序的事情上，将他们回忆的偶然性（如果可以这样说）交托给天主的隐藏权能的引导，对祂而言没有什么是偶然的，那么任何卑微的人，远离天主的视线，遥远地寄居在祂之外，就不应该说：\Quote{这应该放在这里}；因为他完全不知道使天主愿意将其放在它所占据位置的原因。一位宗徒的话是这样说的：\BibleQuote{但如果我们的福音被蒙蔽，那是被蒙蔽在丧亡的人身上。}他又说：\BibleQuote{对一些人，我们是生命的芬芳引人入生；对另一些人，是死亡的芬芳引人死亡；}并立即加上：\BibleQuote{谁能胜任这些事呢？}——也就是说，谁能足够理解这是如何公义地成就的？主亲自表达了同样的事，祂说：\BibleQuote{我来了，是为叫那看不见的看得见，叫那看得见的变成瞎子。}因为天主知识和智慧的财富深不可测，同一团泥土，一个器皿成为贵重的，另一个成为卑贱的。而对于血肉之躯，它说：\BibleQuote{人啊！你是谁，竟敢向天主抗辩？}那么，在目前考虑的事上，谁知道主的心意？或者谁作过祂的谋士，祂以这样的方式掌管了这些圣史在回忆中的心，并将他们提升到祂教会崇高建筑中如此权威的位置，以至于那些在他们看来可能呈现矛盾的同一事物，成为许多人心眼盲目的手段，他们理应被交给自己内心的情欲和败坏的心智；并且也成为许多人在虔诚理解上受到彻底锻炼的原因，按照全能者隐藏的正义？因为一位先知向主说话的话是这样：\BibleQuote{你的意念何其深奥。愚拙人不会知道，愚昧人不会明白这些事。}\par

49. 此外，我请求并劝诫那些阅读我们依靠主的助佑所精心撰写的这部论述的人，要牢记这篇我在此处认为有必要引入的言论，以便在类似查询中提出的每一个相似的困难中都能由此获益，而无需反复重复同样的事。此外，任何愿意洗脱自己不虔敬的顽固之心并专注于此事的人，都能轻易看出，\Person{马尔谷}在此插入第三时辰这一记述是多么恰当，以便每个人都能由此想到一个时辰，即犹太人真正将主钉死在十字架上的时辰，尽管他们企图把这罪责推给罗马人，无论是推给他们的首领还是士兵，当我们在此看到士兵执行职责的记录时便可明了。因为这位作者在此说道：\BibleQuote{他们钉死了祂，就分开祂的衣服，拈阄分取，看谁得什么。}若不是指士兵，还能指谁呢？这在\Person{若望}的叙述中已清楚显明。因此，为了不让人忽略犹太人，而将如此大罪的观念归咎于那些士兵，\Person{马尔谷}在此给了我们这样的记述：\BibleQuote{这是第三时辰，他们钉死了祂}——他的目的是要让人发现那些犹太人才是真正的钉死主的人，细心的研究者会发现，他们完全有可能在第三时辰就已呼喊要求钉死主，同时他观察到士兵所做的事是在第六时辰发生的。\par

50. 然而，同时也不乏一些人，他们认为约翰所说的预备日的时间——他说：\BibleQuote{那正是逾越节的预备日，约在第六时辰}——应理解为这一天的第三时辰，也就是比拉多坐审判席的那个时辰。这样，所说的第三时辰结束之时，似乎就是他被钉十字架、悬在木架上的时候。这样，必须认为又过了三个时辰，在结束时他交出了灵魂。根据这一想法，黑暗也就会从他死亡的时辰——即白天的第六时辰——开始，持续到第九时辰。因为这些人断言，犹太人的逾越节预备日确实是在安息日次日的那一天，因为无酵节始于那安息日；然而，真正的逾越节——即正在主的受难中实现的逾越节，不是犹太人的逾越节，而是基督徒的逾越节——从夜间第九时辰开始被预备，即有了它的\OriginalTerm{预备日}{parasceue}，因为主那时正被犹太人准备处死。因为\OriginalTerm{预备日}{parasceue}这个词，按解释意为\Quote{预备}。因此，在所说的夜间第九时辰与他被钉十字架之间，有一段时间，约翰称之为\OriginalTerm{预备日}{parasceue}的第六时辰，马尔谷称之为白天的第三时辰；这样，按照这一看法，马尔谷并不是以回溯方式在他的记述中引入犹太人喊叫的时辰：\BibleQuote{钉死他，钉死他}，而是明确提到第三时辰为主被钉在木架上的时辰。哪个信徒不会欣然接受这个问题的解决方案呢，只要能找到与所说的第九时辰相关的某一点（在事件的叙述中），我们可以合理且与其他情况一致地推测，我们逾越节的\OriginalTerm{预备日}{parasceue}——即基督之死的预备——从此开始。因为，如果我们说它始于主被犹太人逮捕的时候，那仍是初夜时分。如果我们认为它始于他被带到盖法岳父的家的时候，在那里他也被司祭长们审问，那时鸡还没有叫，这我们从伯多禄的否认中得知，他的否认是在鸡叫时才发生的。同样，如果我们假定它始于他被交给比拉多的时候，我们就有圣经最明确的陈述，大意是那时已经是早晨。因此，我们只能理解，这逾越节的\OriginalTerm{预备日}{parasceue}——即主之死的预备——始于所有司祭长（在他面前首次受审的那些司祭长）回答并说\BibleQuote{他该死的}的时刻，这句话我们见玛窦和马尔谷都有记载；这样，他们被认为是以回溯的方式，在较后的阶段引入了关于伯多禄否认的事实，而这些事实在历史顺序上发生在较早的时候。而且，推测如下毫无不合理之处：如我所说，他们宣判他该死的时刻很可能就是夜间的第九时辰，从那个时辰到比拉多坐审判席的时辰之间，插入了这所谓的第六时辰——然而不是白天的第六时辰，而是\OriginalTerm{预备日}{parasceue}的第六时辰——也就是说，主祭献的预备，这才是真正的逾越节。按照这一理论，当同一\OriginalTerm{预备日}{parasceue}的第六时辰结束时，主被悬于木架，这发生在白天第三时辰结束时。因此，我们可以在这一看法和另一种看法之间做出选择，后者认为马尔谷是以回忆的方式引入了第三时辰，并且在提到那个时辰时特别着眼于表明：在主的钉十字架一事上，犹太人所受的谴责，即他们被认为处于这样的地位，以至于呼喊将祂钉死，其效果是我们认为他们是实际钉死祂的人，而不是那些用双手将祂悬在木架上的人；正如前面提到的百夫长，比主所派遣的那些朋友（在事实的使命上）更能以真切的意义接近主。但无论我们采纳这两种看法中的哪一个，无疑都为主受难时辰的这个难题找到了一个解决方案，这个方案极其恰当地既能激发争辩者的无耻，又能搅动弱者的不成熟。\par

\SourceRecord{Translated by S.D.F. Salmond.；Nicene and Post-Nicene Fathers, First Series；Vol. 6.；Edited by Philip Schaff.；Buffalo, NY: Christian Literature Publishing Co.,；1888.；<http://www.newadvent.org/fathers/1602313.htm>.}

% structure: p0001 source=h1 target=section reason=page-title

\section{\WorkTitle{福音书的和谐 第三卷}}

% structure: p0002 source=h2 target=subsection reason=relative-heading-rank; base=h2

\subsection{第14章 论所有圣史在与他一同被钉的两个强盗之事上的一致}

51. 玛窦继续叙述如下：\Quote{那时，有两个强盗与他一同被钉，一个在右边，一个在左边。} 马尔谷和路加也以类似的形式记载了此事。若望也没有提出难题，尽管他没有提及那些强盗。因为他说道：\Quote{另有两人和他一起，两侧各一个，耶稣在中间。} 但如果若望说这些其他人是无罪的，而其他圣史称他们为强盗，那就会有矛盾。\par

\SourceRecord{Translated by S.D.F. Salmond.；Nicene and Post-Nicene Fathers, First Series；Vol. 6.；Edited by Philip Schaff.；Buffalo, NY: Christian Literature Publishing Co.,；1888.；<http://www.newadvent.org/fathers/1602314.htm>.}

% structure: p0001 source=h1 target=section reason=page-title

\section{\WorkTitle{福音的和谐，第三卷}}

% structure: p0002 source=h2 target=subsection reason=relative-heading-rank; base=h2

\subsection{第十五章 论\Person{玛窦}、\Person{马尔谷}和\Person{路加}关于侮辱主之人群的记述的一致性}

52. \Person{玛窦}接着这样说：\Quote{那些路过的人讥诮祂，摇着头说：\InnerQuote{你这拆毁圣殿、三日内把它重建起来的，救你自己吧！你如果是天主子，就从十字架上下来吧！}}\Person{马尔谷}的记述几乎与之一字不差。然后\Person{玛窦}接着说：\Quote{同样，司祭长、经师和长老们也讥笑祂说：\InnerQuote{他救了别人，却救不了自己。他如果是以色列王，现在就从十字架上下来，我们就信他。他信赖天主，如果祂愿意，就让祂现在救他，因为他曾说：\InnerQuote{我是天主子。}}}\Person{马尔谷}和\Person{路加}虽然措辞不同，但传达的意义相同，尽管一人忽略了另一人提及的某些内容。因为他们确实在司祭长这件事上一致，让我们明白他们在主被钉十字架时侮辱了祂。唯一的区别是：\Person{马尔谷}没有特别提到长老，而\Person{路加}举了首领们为例，但没有加上\Quote{司祭的}称呼，因此他用一个总称涵盖了全体首领；所以我们可以认为他的描述中包括了经师和长老。\par

\SourceRecord{Translated by S.D.F. Salmond.；Nicene and Post-Nicene Fathers, First Series；Vol. 6.；Edited by Philip Schaff.；Buffalo, NY: Christian Literature Publishing Co.,；1888.；<http://www.newadvent.org/fathers/1602315.htm>.}

% structure: p0001 source=h1 target=section reason=page-title

\section{福音的和谐 第三卷}

% structure: p0002 source=h2 target=subsection reason=relative-heading-rank; base=h2

\subsection{第十六章 论归于强盗的讥讽，以及玛窦与马尔谷一方和路加另一方之间毫无矛盾的疑问，当时最后一位圣史叙述说，两人中有一个讥笑祂，另一个信了祂。}

53. 玛窦继续叙述说：\BibleQuote{同祂一起被钉的强盗，也这样讥诮祂。} 马尔谷在此与玛窦完全一致，以不同措辞表述了同一件事。另一方面，路加似乎与之矛盾，除非我们留意到一种相当常见的表达方式。因为路加的记述如下：\BibleQuote{悬挂着的凶犯中，有一个侮辱祂说：\InnerQuote{如果你是基督，救你自己和我们吧！}} 然后同一位圣史接着在同一上下文中叙述道：\BibleQuote{另一个应声责斥他说：\InnerQuote{你既然受同样的刑罚，连天主你都不怕吗？对我们来说，这是理所当然的，因为我们所受的，正配我们所行的；但这个人没有做过什么不正当的事。}他又对耶稣说：\InnerQuote{主，当你进入你的王国时，请你记念我。}耶稣对他说：\InnerQuote{我实在告诉你，今天你就要同我在乐园里。}} 于是问题在于：我们如何能使玛窦的报导\BibleQuote{同祂一起被钉的强盗也这样讥诮祂}或马尔谷的报导\BibleQuote{同祂一起被钉的人辱骂祂}与路加的证词相一致？路加的证词大意是，其中一人辱骂基督，但另一人制止了他并信了主。解释是，玛窦和马尔谷在概述这段经文时，使用了复数而非单数；正如在\WorkTitle{希伯来书}中，我们看到以复数形式陈述说：\BibleQuote{他们堵住了狮子的口}，而所指的只有达尼尔一人。同样，在说到他们\BibleQuote{被锯死}时也用了复数，而那种死亡方式只记载于依撒意亚身上。同样，在\WorkTitle{圣咏}中说：\BibleQuote{世上的君王起来，首领聚在一起}等，也使用了复数而非单数，按照\WorkTitle{宗徒大事录}中对这段经文的解释。因为在那部书中引用该圣咏的人，认为君王指黑落德，首领指比拉多。此外，既然外邦人惯于对福音提出此类诽谤，我愿请他们思考，他们自己的作家是如何谈论\Person{淮德拉}、\Person{美狄亚}和\Person{克吕泰涅斯特拉}等人的，而实际上每个名字之下只有一个闻名的人物。例如，有一个人说：\Quote{乡下人也对我无礼}，即便只有一个人行为粗鲁，这难道不是更常见的事吗？总之，即使路加的记载只限于两个强盗中的一个，也不会造成真正的矛盾，除非其他圣史明确说过\Quote{两者}凶犯都辱骂了主；因为若那样，我们就不能假设复数指的是只有一个人了。但既然所用的措辞是\Quote{强盗}或\Quote{同祂一起被钉的人}，并且没有加上\Quote{两者}一词，那么这种表达既可以用于两人都参与的情况，也同样可用于其中一人涉及其中——此时通过使用复数来传达该事实，这是常见的说法方式。\par

\SourceRecord{Translated by S.D.F. Salmond.；Nicene and Post-Nicene Fathers, First Series；Vol. 6.；Edited by Philip Schaff.；Buffalo, NY: Christian Literature Publishing Co.,；1888.；<http://www.newadvent.org/fathers/1602316.htm>.}

% structure: p0001 source=h1 target=section reason=page-title

\section{\WorkTitle{福音的和谐，第三卷}}

% structure: p0002 source=h2 target=subsection reason=relative-heading-rank; base=h2

\subsection{第十七章   论四位福音圣史关于醋饮记载的和谐}

54. 玛窦接着这样叙述：\BibleQuote{从第六时辰起，直到第九时辰，遍地都黑暗了。}另有两位福音圣史也证实了同一事实。路加还附加了黑暗的原因，即\BibleQuote{太阳失了光。}玛窦又继续写道：\BibleQuote{约在第九时辰，耶稣大声呼喊说：厄里！厄里！肋玛撒巴黑塔尼！就是说：我的天主，我的天主，你为什么舍弃了我？站在那里的人，有几个听见了，就说：这人呼唤厄里亚。}玛尔谷的记述与此几乎完全一致，就词语而言几乎，就意义而论则完全一致。接着玛窦这样说：\BibleQuote{其中一个人立即跑去，拿了海绵，浸满了醋，绑在芦杆上，递给他喝。}玛尔谷以类似的形式呈现：\BibleQuote{有一个人跑去，把海绵浸满了醋，绑在芦杆上，递给他喝，说：且等着，看厄里亚是否来把他取下。}然而，玛窦却把关于厄里亚的这些话，不是当作递醋海绵的人说的，而是其余的人说的。因为他这样写道：\BibleQuote{其余的人却说：等着，看厄里亚是否来救他；}——由此我们推断，所特指的那个人和那里的其他人都用了这些词语。路加又在叙述强盗的蛮横之前提到了醋。他这样记载：\BibleQuote{兵士也戏弄他，上前来，递给他醋喝，说：你若是犹太人的君王，就救你自己吧！}路加的目的是将兵士所做所说的事综合在一句话里。而他没有明确说是其中\BibleQuote{一个}人递醋，这并不成为难题。因为采用我们前面讨论过的一种表达方式，他只是以复数代替了单数。若望也记载了醋的事，他说：\BibleQuote{此后，耶稣知道一切都完成了，为应验经上的话，就说：我渴。那里放着一个器皿盛满了醋；他们就用海绵浸满了醋，绑在牛膝草上，送到他口边。}虽然这位若望这样告诉我们，耶稣说了\BibleQuote{我渴}，并且提到那里放着一个器皿盛满了醋，而其他福音圣史没有详细说明这些，这并不足为奇。\par

\SourceRecord{Translated by S.D.F. Salmond.；Nicene and Post-Nicene Fathers, First Series；Vol. 6.；Edited by Philip Schaff.；Buffalo, NY: Christian Literature Publishing Co.,；1888.；<http://www.newadvent.org/fathers/1602317.htm>.}

% structure: p0001 source=h1 target=section reason=page-title

\section{\WorkTitle{福音合参·第三卷}}

% structure: p0002 source=h2 target=subsection reason=relative-heading-rank; base=h2

\subsection{第十八章　论主临死时的连续话语；并论玛窦和马尔谷在其记载中是否与路加一致，以及这三位福音作者是否与若望一致。}

55. \Person{玛窦}继续写道：\BibleQuote{耶稣又大喊一声，交付了灵魂。}同样，\Person{马尔谷}说：\BibleQuote{耶稣大声喊叫，就交付了灵魂。}\Person{路加}则告诉我们祂发出那大声喊叫时所说的话。他的记载是这样的：\BibleQuote{耶稣大声喊着说：\InnerQuote{父啊，我把我的灵魂交在祢手中}；说了这话，便交付了灵魂。}另一方面，\Person{若望}没有记载\Person{玛窦}和\Person{马尔谷}所报告的第一个声音——即\BibleQuote{厄里，厄里}——也略过了只有\Person{路加}叙述而其他两位称之为\Quote{大声}的那一句话。我指的是\BibleQuote{父啊，我把我的灵魂交在祢手中}这一呼喊。\Person{路加}也证实了这一呼喊是大声发出的；因此我们可以理解这特定的呼喊就是\Person{玛窦}和\Person{马尔谷}所提及的大声喊叫。但\Person{若望}记载了其他三位都没有注意到的一件事，即祂在接受了醋以后说\BibleQuote{完成了}。我们认为这一呼喊是在那大声喊叫之前发出的。因为\Person{若望}的话是这样：\BibleQuote{耶稣一尝了那醋，便说：\InnerQuote{完成了}；就低下头，交付了灵魂。}在这\BibleQuote{完成了}的呼喊与后面所说的\BibleQuote{就低下头，交付了灵魂}之间，发出了\Person{若望}本人没有记载而其他三位所注意到的那个声音。因为准确的顺序似乎是：祂首先说\BibleQuote{完成了}，那时关于祂所预言的都已应验；此后——好像祂一直在等待这一刻，确实祂是自愿死去的——祂将灵魂交托（给圣父），并放弃了它。但无论一个人认为这些话是按什么顺序说的，他最应当警惕的是，不要认为任何一位福音作者与另一位相矛盾，当一个作者略过了另一个作者重复的内容，或者详述了另一个作者略过的事情时。\par

\SourceRecord{Translated by S.D.F. Salmond.；Nicene and Post-Nicene Fathers, First Series；Vol. 6.；Edited by Philip Schaff.；Buffalo, NY: Christian Literature Publishing Co.,；1888.；<http://www.newadvent.org/fathers/1602318.htm>.}

% structure: p0001 source=h1 target=section reason=page-title

\section{福音的和谐 第三卷}

% structure: p0002 source=h2 target=subsection reason=relative-heading-rank; base=h2

\subsection{第十九章 论圣殿帐幔的撕裂，以及玛窦和马尔谷在叙述此事发生的顺序上是否与路加真正相符的问题。}

56. 玛窦接着写道：\BibleQuote{看，圣殿帐幔从上到下裂为两半。} 马尔谷的记载也如下：\BibleQuote{圣殿帐幔从上到下裂为两半。} 路加同样以类似的话作了陈述：\BibleQuote{圣殿帐幔从中间裂开。} 然而，他并非按相同的顺序引入这段话。因为，为了将奇迹与奇迹相连，他先告诉我们：\BibleQuote{太阳变暗了。} 然后认为应当紧接着附加上面的句子，即：\BibleQuote{圣殿帐幔从中间裂开。} 因此，看来他是较早引入了这个事件——实际上该事件发生在主断气的时候——以便在那里给我们一个关于喝醋、大声呼喊以及死亡本身情形的概括性描述，而这些被认为发生在帐幔撕裂之前、黑暗降临之后。因为玛窦将这句话：\BibleQuote{看，圣殿帐幔裂了。} 紧接着放在叙述：\BibleQuote{耶稣又大声喊叫，交付了灵魂。} 之后，从而清楚地告诉我们，帐幔撕裂的时间是在耶稣交出祂的灵魂之后。然而，如果他未曾加上\Quote{看哪}二字，而只是简单地说\BibleQuote{圣殿帐幔裂了。}，那么就不确定玛窦和马尔谷是否以综述形式叙述了此事，而路加保持了准确的顺序，或者路加是否给出了对这些人在正确历史顺序中所引入的事件的概要。\par

\SourceRecord{Translated by S.D.F. Salmond.；Nicene and Post-Nicene Fathers, First Series；Vol. 6.；Edited by Philip Schaff.；Buffalo, NY: Christian Literature Publishing Co.,；1888.；<http://www.newadvent.org/fathers/1602319.htm>.}

% structure: p0001 source=h1 target=section reason=page-title

\section{\WorkTitle{福音的和谐，卷三}}

% structure: p0002 source=h2 target=subsection reason=relative-heading-rank; base=h2

\subsection{第二十章 论\Person{玛窦}、\Person{马尔谷}和\Person{路加}关于百夫长及与他同在一起之人的惊讶一事所提供之几处记载的一致性}

57. 玛窦接着这样记载：\BibleQuote{大地震动，岩石崩裂，坟墓自开，许多已睡的圣人之身体复活了，在耶稣复活后，他们由坟墓出来，进入圣城，发显给许多人。} 没有理由担心这些只有玛窦记载的事实会与其他人所述的有任何不符。同一位福音作者接着这样记载：\BibleQuote{百夫长和与他一同看守耶稣的人，看见地震及所发生的事，就非常害怕，说：\Quote{这人真是天主子。}} 马尔谷的记载是这样的：\BibleQuote{对面站着的百夫长，看见耶稣这样喊叫而断气，就说：\Quote{这人真是天主子。}} 路加的记载如下：\BibleQuote{百夫长看见所发生的事，就光荣天主说：\Quote{这人实在是个义人。}} 在此玛窦说是百夫长和他同在一起的人看见地震而如此惊讶，而路加表示那人的惊奇是因耶稣发出那样的呼声并断了气而引发的；这便清楚地表明耶稣如何能够自己掌握自己死亡的时间。但这并不矛盾。因为上述玛窦不仅告诉我们百夫长\Quote{看见了地震}，还附加了\Quote{及所发生的事}这句话，他已经指出路加有足够的空间将主的死亡本身视为引发百夫长惊奇的事件。因为那个事件也是当时以如此奇妙方式发生的事之一。同时，即使玛窦没有加上这样的声明，我们仍然完全可以合理地认为，当时发生了许多令人惊奇的事，百夫长和与他在一起的人可能都以惊愕的眼光看待所有这一切，史家们可以自由选择叙述任何一个他们各自愿意举例作为那人惊奇对象的事件。如果仅仅因为其中一人可能指出了某一事件作为百夫长惊奇的直接原因，而另一人介绍了不同的细节，就指责他们的记载不一致，那是不公平的。因为这些事件全部都是那人确实感到惊奇的对象。再者，一位福音作者告诉我们百夫长说：\Quote{这真是天主子}，而另一位说：\Quote{这人真是天主子}，对于任何还记得前面已论述过多次的类似情形的人来说，这一点也不会造成困难。因为这些不同的表述都传达了完全相同的含义，而且尽管一位作者加上了\Quote{这人}一词而另一位没有，这也绝不构成矛盾。可能更明显的矛盾是路加所报道的百夫长的话不是\Quote{这是天主子}，而是\Quote{这是个义人}。但我们应该认为，要么百夫长实际说了这两句话，两位福音作者记录了一句，第三位记录了另一句；要么或许路加有意表达百夫长在说耶稣是天主子时所持有的确切想法。因为可能百夫长并不真正理解他是与父同等的独生子，而只是称他为天主子，因为他相信他是个义人，正如许多义人都被称为天主的儿子一样。再者，当路加说\Quote{百夫长看见所发生的事}时，他实际上是用一个术语涵盖了当时发生的一切奇妙之事，纪念了一件奇事，而所有这些奇迹事件可以说是这个奇事的组成部分和部分。此外，关于玛窦也提到了与百夫长在一起的人，而其他作者没有提及这些人，对于任何明白仅仅因为一位作者记载了另一位未曾提及的事并不一定构成矛盾这一公认原则的人来说，这难道不是不言自明的吗？最后，关于玛窦告诉我们\Quote{他们非常害怕}，而路加没有说那人害怕，却告诉我们\Quote{他光荣了天主}，谁不能理解他正是因所表现的恐惧而光荣了天主呢？\par

\SourceRecord{Translated by S.D.F. Salmond.；Nicene and Post-Nicene Fathers, First Series；Vol. 6.；Edited by Philip Schaff.；Buffalo, NY: Christian Literature Publishing Co.,；1888.；<http://www.newadvent.org/fathers/1602320.htm>.}

% structure: p0001 source=h1 target=section reason=page-title

\section{\WorkTitle{福音的和谐，第三卷}}

% structure: p0002 source=h2 target=subsection reason=relative-heading-rank; base=h2

\subsection{第二十一章。论当时站在那里的妇女，以及玛窦、马尔谷和路加说她们远远站着，是否与若望说其中一人站在十字架旁相矛盾。}

58. 玛窦这样叙述：\BibleQuote{有许多妇女在那里远远地观看，她们是从加里肋亚跟随耶稣来的；其中有玛利亚玛达肋纳，雅各伯和若瑟的母亲玛利亚，以及载伯德儿子的母亲。} 马尔谷这样记载：\BibleQuote{也有些妇女远远地观看：其中有玛利亚玛达肋纳，小雅各伯和若瑟的母亲玛利亚，以及撒罗默（她们当耶稣在加里肋亚时就跟随祂，并服事祂）；还有许多其他与祂一同上耶路撒冷的妇女。} 我看不出这些作者之间有任何可视为不一致之处。因为有些妇女在两个名单中都被提名，而有些只在一个名单中提到，这怎么会影响真理呢？路加也同样这样叙述：\BibleQuote{所有聚集来看这景象的群众，看见所发生的事，都捶着胸膛回去了。所有祂的熟人和从加里肋亚跟随祂的妇女，都远远地站着观看这些事。} 这里我们看出，就妇女在场这一点而言，他与前两位十分一致，尽管他没有提任何人的名字。关于也在场的群众，他们看见所发生的事，都捶着胸膛回去了，在这方面他也同样与玛窦一致，尽管那位福音作者在上下文中特别提到：\BibleQuote{百夫长和与他在一起的人。} 因此，显然只有路加明确提到了祂的\Quote{熟人}远远站着。因为若望也注意到在主断气前妇女们在场的状况。他的叙述如下：\BibleQuote{在耶稣的十字架旁，站着祂的母亲、祂母亲的姐妹、克罗帕的妻子玛利亚和玛利亚玛达肋纳。耶稣看见母亲和祂所爱的门徒站在旁边，就对母亲说：\InnerQuote{女人，看，你的儿子！}然后对门徒说：\InnerQuote{看，你的母亲！}从那时起，那门徒就把她接到自己家里。} 关于这段话，若非玛窦和马尔谷同时明确提到了玛利亚玛达肋纳的名字，我们或许可以说，有一群妇女远远站着，另一群在十字架旁。因为在这些作者中，除了若望本人，没有人提到主的母亲。因此，现在的问题是这样：我们如何理解同一位玛利亚玛达肋纳既与别的妇女远远站着，如玛窦和马尔谷所记载的，又站在十字架旁，如若望告诉我们的？除非这些妇女所站的距离，既可以合理地说她们是近的，因为她们就在祂眼前，但比起那些与百夫长和士兵们更近地围在周围的群众来说，又是远的。因此，我们可以设想，那些与主的母亲一同在场目睹此景的妇女，在祂将母亲托付给门徒之后，就开始退后，以便从人群密集处脱身，并从更远的地方观看剩下要发生的事。这样，其他福音作者只在主死后才提到这些妇女，他们正确地报导说，那时她们已远远站着。\par

\SourceRecord{Translated by S.D.F. Salmond.；Nicene and Post-Nicene Fathers, First Series；Vol. 6.；Edited by Philip Schaff.；Buffalo, NY: Christian Literature Publishing Co.,；1888.；<http://www.newadvent.org/fathers/1602321.htm>.}

% structure: p0001 source=h1 target=section reason=page-title

\section{\WorkTitle{福音的和谐}，第三卷}

% structure: p0002 source=h2 target=subsection reason=relative-heading-rank; base=h2

\subsection{第二十二章 论福音作者关于若瑟向比拉多求主的身体之叙述是否一致，以及若望的记载有无自相矛盾之处的问题}

59. 玛窦接着记述：\BibleQuote{到了傍晚，来了一个阿黎玛特雅的富人，名叫若瑟，他也是耶稣的门徒；他去见比拉多，求耶稣的身体。比拉多就下令把身体交给他。}\newline{}马尔谷的记述如下：\BibleQuote{到了傍晚，因为是预备日，即安息日的前一天，阿黎玛特雅人若瑟，一位受人尊敬的议员，他本人也等待天主的国，前来，大胆地进入比拉多那里，求耶稣的身体。比拉多惊奇祂是否已经死了；便叫来百夫长，问他是否已经死了。从百夫长那里得知后，就把身体赐给了若瑟。}\newline{}路加的记述是这样的：\BibleQuote{看，有一个人名叫若瑟，是个议员，是个善人，又是个义人（他没有赞同他们的计谋和行事）。他是阿黎玛特雅人，一个犹太人的城；他也等待天主的国。这人去见比拉多，求耶稣的身体。}\newline{}然而若望先叙述了与主同钉之人的腿被打断，以及主的肋旁被长枪刺透（这一段只有他记录），然后加上与其他福音作者相同的叙述。其文如下：\BibleQuote{此后，阿黎玛特雅人若瑟，因为怕犹太人，暗暗地作耶稣的门徒，来求比拉多，要取耶稣的身体；比拉多允准了他。他就来，把耶稣的身体取去。}\newline{}这里没有任何使任何人显得彼此矛盾之处。但有人或许会问，若望是否自相矛盾，因为他一方面与其他人一致地告诉我们若瑟如何求耶稣的身体，另一方面又独自陈述若瑟因怕犹太人而暗暗地作耶稣的门徒。因为可以合理地提出疑问：一个因害怕而暗暗作门徒的人，怎么会有勇气去求祂的身体——这件事连那些公开跟随祂的人都不敢做。然而我们应理解，这个人仗着他尊贵的地位所给与的胆量而行此事，这地位使他能够熟悉地进入比拉多面前。我们还须设想，在履行与埋葬有关的最后服务时，他较少顾忌犹太人，尽管他在平时听主讲道时，总是尽力避免招惹他们的敌意。\par

\SourceRecord{Translated by S.D.F. Salmond.；Nicene and Post-Nicene Fathers, First Series；Vol. 6.；Edited by Philip Schaff.；Buffalo, NY: Christian Literature Publishing Co.,；1888.；<http://www.newadvent.org/fathers/1602322.htm>.}

% structure: p0001 source=h1 target=section reason=page-title

\section{福音的和谐 第三卷}

% structure: p0002 source=h2 target=subsection reason=relative-heading-rank; base=h2

\subsection{第二十三章 论前三部福音作者在记载主的安葬时是否与若望完全和谐的问题}

60. 玛窦继续写道：\BibleQuote{若瑟取了遗体，用洁净的细麻布裹好，安放在自己的新坟墓里，就是他曾在岩石中凿出的；又把一块大石头滚到墓门口，就离开了。} 马尔谷的记载如下：\BibleQuote{他买了细麻布，把祂取下来，用细麻布裹好，安放在岩石中凿成的坟墓里，又把一块石头滚到墓门口。} 路加是这样记述的：\BibleQuote{他把祂取下来，用细麻布裹好，安放在岩石中凿成的坟墓里，那里从来没有安放过人。} 就这三部记述而言，不可能提出任何不和谐的指控。然而，若望告诉我们，主的安葬不仅有若瑟参与，还有尼苛德摩。因为他从尼苛德摩开始，与上文适当衔接，继续叙述如下：\BibleQuote{又有尼苛德摩，就是当初夜间来见耶稣的，带着约一百斤的没药和沉香调和的香料来。} 然后他再次引入若瑟，接着这样说：\BibleQuote{他们取了耶稣的遗体，用香料和细麻布裹好，这是犹太人安葬的惯例。在祂被钉十字架的地方有一个园子，园子里有一座新坟墓，从来没有安放过人。因为那天是犹太人的预备日，坟墓又近，他们就把耶稣安放在那里。} 然而，如果我们正确理解这段话，实在没有理由像之前那样认为这里有矛盾。因为那些没有提及尼苛德摩的福音作者，并没有声称主是仅由若瑟安葬的，尽管若瑟是他们记载中唯一出现的人。而且，这三位福音作者一致告诉我们主是由若瑟用细麻布裹好的事实，并不妨碍我们设想尼苛德摩可能带来了其他麻布衣物，加在若瑟提供的上面，这样若望的叙述就可以完全正确，尤其因为他告诉我们主不是裹在一条细麻布里，而是裹在细麻布衣物中。同时，考虑到包头的手巾和缠裹全身的绷带，所有这些都是细麻布做的，我们可以看到，即使那里实际上只有一条细麻布（如前三福音作者所指的那种），仍然可以极其真实地说\BibleQuote{他们用细麻布衣物裹好祂}。因为\BibleQuote{细麻布衣物}这个短语一般适用于所有用亚麻织成的织物。\par

\SourceRecord{Translated by S.D.F. Salmond.；Nicene and Post-Nicene Fathers, First Series；Vol. 6.；Edited by Philip Schaff.；Buffalo, NY: Christian Literature Publishing Co.,；1888.；<http://www.newadvent.org/fathers/1602323.htm>.}

% structure: p0001 source=h1 target=section reason=page-title

\section{\WorkTitle{福音的和谐，第三卷}}

% structure: p0002 source=h2 target=subsection reason=relative-heading-rank; base=h2

\subsection{第二十四章 论四位福音作者关于主复活前后事件的叙述毫无矛盾之处}

61. 玛窦继续写道：\Quote{在那里有玛利亚玛达肋纳和另一位玛利亚，对着坟墓坐着。} 马尔谷的记载如下：\Quote{玛利亚玛达肋纳和若瑟的母亲玛利亚，看见了祂安放的地方。} 至此，显然这些记述之间没有任何形式的不一致。\par

62. 玛窦接着这样记述：\Quote{次日，就是预备日之后的那一天，司祭长和法利塞人一同来见比拉多，说：大人，我们记得那骗子还活着的时候曾说：三天后我要复活。所以请下令将坟墓把守妥当，直到第三天，免得祂的门徒夜间来偷走祂，并向百姓说：祂从死人中复活了；这样，后一个骗局比前一个更坏。比拉多对他们说：你们有守卫；去吧，尽你们所能把守妥当。他们就去，封了石头，派守卫把守坟墓。} 这段叙述只有玛窦记载。然而其他作者没有提到任何可能看似矛盾之处。\par

63. 同一位玛窦接着叙述如下：\Quote{安息日过后，在向一周的第一天黎明的时候，玛利亚玛达肋纳和另一位玛利亚来看坟墓。看，忽然大地震，因为主的天使从天降下，来到石头前，把它滚开，坐在上面。他的容貌像闪电，衣服洁白如雪；看守的人因怕他，就战抖，变得像死人一样。天使回答妇女说：不要害怕！我知道你们寻找被钉十字架的耶稣。祂不在这里；祂已照祂所说的复活了。来看主安放的地方。快去告诉祂的门徒：祂已从死人中复活了；并且祂在你们以先往加利利去，在那里你们要看见祂；看，我已经告诉你们了。} 马尔谷与这相合。然而，有人可能会觉得有困难，因为根据玛窦的版本，石头已经从坟墓滚开，天使坐在上面。因为马尔谷告诉我们，妇女们进入坟墓，看见一个少年人坐在右边，穿着白袍，她们就惊惶。但解释可能是：玛窦只字未提她们进入坟墓后看见的天使，而马尔谷未提她们看见坐在外面石头上的那位。这样，她们就看见了两位天使，并分别听了关于耶稣的天使报告——首先是从坐在外面石头上的天使那里，然后是从进入坟墓后看见坐在右边的天使那里。这样，坐在外面的天使所说的话：\Quote{来看主安放的地方}，就鼓励了她们走进坟墓；到了那里，如前所述，她们鼓起勇气进去，我们推测她们看见了玛窦没有提及、而马尔谷所描述的坐在右边的天使，并从祂那里听到了与之前类似的话。如果这个解释不能令人满意，我们当然应该接受另一种观点：当她们进入坟墓时，她们进入了一片区域，按理说，当时已经安全地围起了一块地方，在凿开岩石建造坟墓之处稍前方延伸了一段距离；因此，根据这种观点，她们实际看见的是坐在右边的唯一天使，而玛窦也将这位天使描述为在地震发生时坐在滚开墓口（即从岩石中挖出的坟墓）的石头上。\par

64. 可能还有人会问，马尔谷说：\Quote{她们就出来，从坟墓逃走，因为她们战栗惊惶；什么也不告诉人，因为她们害怕。} 而玛窦的记载是这样：\Quote{她们就急忙离开坟墓，怀着恐惧和极大的喜乐，跑去报告祂的门徒。} 然而解释可以是：妇女们不敢告诉天使们——即她们没有足够的勇气对听到天使的话做出回应。或者，她们也不敢对看到躺在那里的守卫说话；因为玛窦提到的喜乐与马尔谷提到的恐惧并不矛盾。实际上，我们本应认为两种情绪都占据了她们的心，即使玛窦本人没有提到恐惧。但现在，这位福音作者也特别提到恐惧，说：\Quote{她们就急忙离开坟墓，怀着恐惧和极大的喜乐}，他就没有留下任何能引起疑惑的问题。\par

65. 与此同时，这里确实出现了一个不容轻率处理的问题，涉及妇女们来到坟墓的确切时辰。因为当玛窦说：\BibleQuote{安息日傍晚，天正要破晓一周的第一天时，玛利亚玛达肋纳和另一个玛利亚来看坟墓}，我们该如何理解马尔谷的陈述，他这样说：\BibleQuote{非常清早，一周的第一天，他们在日出时来到坟墓}？应当注意，在此马尔谷所说的，与另外两位圣史，即路加和若望的记载并无矛盾。因为路加说：\BibleQuote{非常清早，} 若望则说：\BibleQuote{清早，天还黑的时候，} 他们传达了与马尔谷所说的 \BibleQuote{非常清早，日出之时} 相同的意思；也就是说，他们都指向东方天空开始发亮的时刻，这当然只有在日出临近时才会发生。因为日出所散发的光明通常被称为黎明。因此，马尔谷与使用 \BibleQuote{天还黑的时候} 这一短语的圣史并无矛盾；因为破晓时，黑夜的残余黑暗随着太阳升起而逐渐消逝。而 \Quote{非常清早} 这一短语不必理解为此时太阳本身已在地面上被实际看见；它更应被视作我们惯于使用的一种表达，当我们向人们暗示某事应比平时更早完成时。因为当我们用了 \Quote{清早} 一词，如果希望避免让听者以为我们直接指太阳已明显可见于地上的时刻，我们通常加上 \Quote{非常} 一词，说 \Quote{非常清早}，好让他们明白我们所指的时间也称为破晓。与此同时，人们惯常的做法是，在多次听到鸡鸣后，开始猜测白天即将来临时，便说：\Quote{现在是清早}；然后当他们斟酌言辞，观察到太阳正在升起——即它立即要出现在这些地方——天空开始泛红或发亮时，那些说过 \Quote{现在是清早} 的人便扩大其表达，说：\Quote{现在是非常清早}。但这有什么关系呢？只要无论我们偏爱哪种解释，都能理解马尔谷使用 \BibleQuote{清早} 一词时所表达的意思，与路加采用 \BibleQuote{清早} 短语时所指的相同；并且前者使用的完整表达——即 \BibleQuote{非常清早}——相当于我们在路加中所见的 \BibleQuote{非常黎明}，以及若望所选的 \BibleQuote{清早，天还黑的时候}？此外，当马尔谷说到 \BibleQuote{日出} 时，他仅仅意味着太阳升起开始将光照到天空。但现在的问题是：玛窦既没有说 \BibleQuote{清早} 也没有说 \BibleQuote{非常清早}，而是说 \BibleQuote{安息日傍晚，天正要破晓一周的第一天时}，他如何与这三位一致？这是必须仔细研究的问题。现在，玛窦用夜晚的第一部分（即此处称为\Quote{傍晚}）来指这一特定的夜晚，妇女们在其结束时来到坟墓。我们理解他如此指代该夜晚的原因是：到傍晚时分，她们可以合法地携带香料，因为安息日确实已经结束。因此，由于她们先前因安息日受阻无法这么做，他给夜晚的命名取自她们开始可以合法做某事的时间，而她们在同一夜晚的任何喜欢的时刻做了此事。所以，\Quote{安息日傍晚} 这一短语被用来表达，就好像说的是 \Quote{安息日夜晚}，换句话说，就是安息日白天之后的夜晚。他使用的明确措辞充分清楚地表明了这一点。因为他的措辞是：\BibleQuote{安息日傍晚，天正要破晓一周的第一天时}；如果我们需要理解的 \Quote{傍晚} 仅仅指夜晚最初的短暂时间，换句话说，只是夜晚的开始，那么情况就不会如此。因为所谓 \Quote{开始破晓一周的第一天} 所指的并非夜之初，而是夜本身，即它开始被光亮的推进带向结束。因为夜晚第一部分的终点正是第二部分的起点，但整个夜晚的终点是光。因此，我们只能说傍晚破晓一周的第一天，除非在 \Quote{傍晚} 一词下我们理解它指的是夜本身，它作为一个整体被光带向结束。在圣经中，以部分表达整体也是一种常见的说话方式；因此，在 \Quote{傍晚} 一词下，圣史在此指代整个夜晚，其终点在黎明。因为那些妇女是在黎明时来到坟墓；这样，她们确实是在夜晚来的，这里用 \Quote{傍晚} 一词来指代。因为，如我已说，整个夜晚由这个词表示；因此，无论她们在该夜晚的任何时刻来到，她们确实是在该夜晚来的。因此，即使她们是在该夜晚最晚的时刻来的，这仍然是毫无疑问的。但除非整个夜晚能由该表达理解，否则不能说她们是在 \BibleQuote{安息日傍晚，天正要破晓一周的第一天时} 来的。因此，在所提及的夜晚来的妇女，正是在指定的傍晚来的。如果她们在任何时刻来，甚至该夜晚最晚的时刻，她们肯定是在夜晚本身来的。\par

66. 主死亡与复活之间所经历的三日时间，除非按照以部分代整体的表达方式，否则无法正确理解。因为祂亲自说：\Quote{正如约纳在鲸腹中三天三夜，人子也要在地心中三天三夜。}\BibleRef{玛 12:40} 现在，无论我们如何计算时间，是从祂断气之时还是从祂被埋葬之日算起，除非我们将中间那一日，即安息日，视为完整的一日——也就是说，连同其夜间作为完整的一日——并且另一方面，理解那介于其间的两日，即预备日和一週的第一日（我们称为主日），是以部分代整体的原则来处理的，否则总数无法明确。因为有些人在这些难题面前感到窘迫，且不了解圣经在解决这些难题时极为常见的以部分代整体的表达方式，便想出将那三个时辰（即从第六时辰到第九时辰，太阳变暗期间）算作一个独立的夜间，而将另外三个时辰（即从第九时辰到日落，太阳重新照耀大地期间）算作一个独立的白天。但接下来的安息日之夜与其白昼相连，若我们将其与白昼一同计算，则会有两日两夜。然而，安息日之后又是一週第一日（即主日破晓，主复活之时）的夜间。因此，这种计算方式的结果将只有两日两夜，再加上一个夜间——即便我们假定最后一个夜间可以算作完整的夜晚，并且前提是我们不指出那所谓的破晓实际上是该夜的最后一刻。由此可见，即使我们以这种方式计算这六个时辰（其中三个时辰太阳变暗，另外三个时辰重新照耀），我们仍然无法得出三日三夜的满意计算。因此，根据圣经语言中如此频繁出现的用法（即以部分代整体），我们必须将预备日视为一端之日，主在这一天被钉十字架并埋葬，并从这一界限算出一个完整的白昼与其夜间完全度过。同样，我们必须将中间部分，即安息日，视为不仅是部分计算，而是真正完整的一日。第三日则必须从其第一部分开始计算，即从夜间算起，并将其白昼部分与之相连，视为构成完整的一日。这样我们将得到三日的空间，类比于一个已经考虑过的案例，即主上山后八日；关于这段时间，我们发现玛窦和马尔谷只注意中间完整的几日，记载为\Quote{六天以后}\BibleRef{玛 17:1}，而路加则记载为\Quote{八天以后}\BibleRef{路 9:28}。\par

67. 因此，现在我们继续考察这段经文的其余部分，看看这些叙述在其他方面如何与\Person{玛窦}的记载完全一致。因为\Person{路加}以极其明确的方式告诉我们，那些来到坟墓的妇女看见了两位天使。我们认为，前两位福音书作者各自提及了其中一位天使：就是说，\Person{玛窦}提到了一位，即坐在外面石头上的那位；\Person{马尔谷}提到了另一位，即坐在坟墓里面右边的那位。但\Person{路加}对这一幕的记述如下：\BibleQuote{那天是预备日，安息日也快到了。那些从\Place{加里肋亚}与祂同来的妇女，跟在后面，看见了坟墓，也看见了祂的身体是怎样安放的。她们就回去，预备了香料和香膏。安息日，她们遵命安息。一周的第一天，大清早，她们带着预备好的香料，来到坟墓前，见石头已由墓穴滚开了，就进去，却找不着主耶稣的身体。她们正为此事疑虑的时候，忽然有两个人，穿着耀眼的衣服，站在她们身边。她们都害怕，把脸伏在地上，那两个人对她们说：你们为什么在死人中找活人呢？祂不在这里，已经复活了！你们应该记得：祂还在\Place{加里肋亚}时，怎样告诉你们说：人子必须被交付于罪人之手，被钉在十字架上，并在第三日复活。她们就想起了祂的话。从坟墓回去，把这一切事报告给那十一位宗徒和所有其余的人。}因此问题在于：如果\Person{路加}在同一事件中记载那两位天使站在那些妇女旁边，而\Person{玛窦}和\Person{马尔谷}却说她们各自看见一位天使坐着——\Person{玛窦}说一位坐在外面的石头上，\Person{马尔谷}说一位坐在坟墓里面右边——尽管天使们所说的话相似，这怎么可能呢？然而，我们仍可以设想：妇女们看见了一位天使，按照\Person{玛窦}所指的位置，以及\Person{马尔谷}所指的情况，正如我们已经解释过的。这样，我们可以理解这些妇女进入了坟墓，即进入了一个在某种围墙内围起来的空间，当她们来到建造坟墓的岩石区域前面时，可以说她们进入了；在那里，我们可以认为她们看见了天使坐在从墓穴滚开的石头上，如\Person{玛窦}所说，换句话说，天使坐在右边，如\Person{马尔谷}所表达。然后我们进一步推测，这些妇女进去后，在观看主身体所躺的地方时，看见了另外两位站着的天使，如\Person{路加}告诉我们，他们用类似的话对她们说话，以激励她们的心灵，建立她们的信德。\par

68. 但我们也来考察\Person{若望}的版本，看它是否、或以何种方式与其他版本一致。于是\Person{若望}叙述这些事如下：\BibleQuote{在一周的第一天，\Person{玛利亚玛达肋纳}清早，天还黑的时候，来到坟墓前，看见石头已从坟墓挪开了。她就跑去见\Person{西满伯多禄}和\Person{耶稣}所爱的另一位门徒，对他们说：他们从坟墓里把主搬走了，我们不知道他们把祂放在哪里。\Person{伯多禄}和那位门徒就出来，往坟墓那里去。两人一起跑，那位门徒跑得比\Person{伯多禄}快，先到了坟墓。他弯下身，看见麻布放倒在地上，却没有进去。\Person{西满伯多禄}随后赶到，进了坟墓，看见麻布放倒在地上，也看见裹头巾没有和麻布放在一起，而是卷着放在另一处。然后先到坟墓的那位门徒也进去了，他看见了，就相信了。因为他们还不明白圣经，必须从死者中复活。于是两个门徒又各自回家去了。\Person{玛利亚}却站在坟墓外面哭泣；她哭泣的时候，弯身往坟墓里看，看见两位白衣的天使，一位坐在头，一位坐在脚，就是\Person{耶稣}身体曾放置的地方。他们对她说：女人，你为什么哭？她对他们说：因为我主被搬走了，我不知道他们把祂放在哪里。她说了这话，就转过身，看见\Person{耶稣}站在那里，却不知道那就是\Person{耶稣}。\Person{耶稣}对她说：女人，你为什么哭？你找谁？她以为祂是园丁，就对祂说：先生，若是你把祂搬走了，请告诉我你把祂放在哪里，我必去取回来。\Person{耶稣}对她说：\Person{玛利亚}。她转过身，用希伯来语对祂说：\OriginalTerm{辣波尼}{Rabboni}；意思是师傅。\Person{耶稣}对她说：不要摸我，因为我还没有升到父那里去；你到我弟兄们那里去，告诉他们：我要升到我的父，也是你们的父那里去；升到我的天主，也是你们的天主那里去。\Person{玛利亚玛达肋纳}就去告诉门徒们，她见了主，并主对她说了这些话。}在\Person{若望}的叙述中，关于坟墓的日子或时间的陈述与其他记载一致。再者，在关于看见两位天使的报道上，他也与\Person{路加}一致。但当我们观察到一位福音书作者（\Person{若望}）告诉我们这些天使被看见是坐着，而另一位（\Person{路加}）说她们站着；当我们还注意到还有其他一些事情是这两位作者没有记载的；而且，当我们进一步思考这些问题如何被提出，涉及能否证明这些史书之间在这些主题上的一致性，以及确定所说之事发生的顺序时，我们看到，除非我们仔细审查全部，否则这些叙述之间很容易出现矛盾。\par

69. 既然如此，就让我们尽主所能帮助我们的程度，将福音作者们一同陈述中发生在主复活期间的所有这些事件，放在一起，按它们实际可能发生的顺序，以一段连贯的叙述加以编排。是在一周的第一天清晨，正如所有福音作者一致见证的，妇女们来到坟墓那里。那时，只有玛窦所记载的一切已经发生了；也就是说，关于地震、石头滚开、守卫的恐惧，他们被吓得像死人一样。然后，如若望告诉我们的，玛利亚玛达肋纳来了，她无疑在爱主方面比其他服侍主的妇女更加炽热；所以若望只提她一人，不提名其他妇女，是合理的，但她们与她同在，我们从其他福音作者的记述中可以知道。她于是来了；当她看见石头从坟墓挪开，没有停下来作更仔细的检查，且毫不怀疑耶稣的身体已被从坟墓移走，她跑去，如那位若望所说，告诉伯多禄和若望本人这事。因为若望就是耶稣所爱的那位门徒。他们于是跑去坟墓；若望先到，弯腰看见细麻布放在那里，但没有进去。伯多禄跟着到了，进入坟墓，看见细麻布放着，也看见裹头巾，不是与细麻布一同放着，而是卷在另一处地方。然后若望也进去，看见了，就相信了玛利亚告诉他的话，即主已从坟墓被取走了。\BibleQuote{因为他们还不明白圣经，就是说，祂必须从死人中复活。于是两个门徒又回到自己的家里去了。但玛利亚却站在坟墓外面哭泣}——也就是说，在岩石中凿成的坟墓本身所在的地点之前，但同时也在他们已进入的园子里；因为那里有一个园子，如那位若望所提到的。然后他们看见天使坐在右边，在石头上面，那石头是从坟墓滚开的；这个天使玛窦和马尔谷都提到。\BibleQuote{天使对妇女说：你们不要害怕！我知道你们是寻找那被钉死的耶稣。祂不在这里，照祂所说的，已经复活了。你们来看安放主的地方。快去告诉祂的门徒，祂从死人中复活了，并且看，祂在你们以先往加里肋亚去，在那里你们要看见祂。看，我已经告诉了你们。}在马尔谷那里也有类似的话。听了这些话，玛利亚仍然哭泣，弯腰往坟墓里看，看见两个天使，在若望的叙述中被介绍，穿着白衣，坐在安放耶稣身体的地方，一个在头，一个在脚。\BibleQuote{他们对她说：女人，你为什么哭？她对他们说：因为他们把我的主搬走了，我不知道他们把祂放在哪里。}这里我们该认为天使已经站起来，以便他们能被看见站着，如路加所说那样，然后，按同一位路加的叙述，对妇女们说话，因为她们害怕而把脸伏在地上。话是这样：\BibleQuote{为什么在死人中找活人呢？祂不在这里，已经复活了。你们应当记得祂还在加里肋亚时怎样告诉你们，说：人子必须被交在罪人手里，被钉在十字架上，第三日复活。她们就想起祂的话来。}此后，如我们从若望所学，\BibleQuote{玛利亚转过身来，看见耶稣站着，却不知道那是耶稣。耶稣对她说：女人，你为什么哭？你找谁？她以为是园丁，就对祂说：先生，若是你把祂搬走了，请告诉我你把祂放在哪里，我好把祂取回。耶稣对她说：玛利亚！她便转过身，用希伯来话对祂说：\OriginalTerm{辣步尼}{Rabboni}！就是说\Quote{老师}。耶稣对她说：不要摸我，因为我还没有升到父那里去；但你去我弟兄那里，对他们说：我升到我的父和你们的父那里去，到我的天主和你们的天主那里去。}然后她离开坟墓，也就是离开石头面前有园子空间的地方。和她一同的还有那些其他妇女，如马尔谷告诉我们，她们惊惧战栗，什么都不告诉任何人。至此我们接着取玛窦所记载的以下段落：\BibleQuote{看，耶稣迎上她们说：愿你们平安！她们就上前抱住祂的脚朝拜了祂。}因为这样我们推断，她们来到坟墓时，被天使两次说话；同样，她们也被主亲自两次说话，一次是当玛利亚以为祂是园丁时，第二次是当祂在路上遇见她们时，为了通过这样的重复来坚固她们，带她们脱离恐惧状态。\BibleQuote{于是耶稣对她们说：不要害怕！去告诉我的弟兄，叫他们往加里肋亚去，在那里他们要看见我。}\BibleQuote{玛利亚玛达肋纳就去告诉门徒说：她看见了主，并且主对她说了这些话；}——然而不是她独自一人，而是与她同去的那些妇女，就是路加所指的，\BibleQuote{她们就把这些事告诉了十一个门徒和其余的人。她们的话，在他们看来好像是无稽之谈，他们就不相信。}马尔谷也证实了这些事；因为他在告诉我们妇女怎样战栗惊奇地从坟墓出来，什么也不告诉任何人之后，接着陈述：主在一周的第一天清晨复活，首先显现给玛利亚玛达肋纳，她曾是他从她身上赶出七个魔鬼的，她去告诉那些曾与祂同在的人，他们正在哀恸哭泣，他们听见祂活着，且被玛利亚看见了，却不相信。还应注意到，玛窦也插入了一段话，说那些看见并听了所有这些事的妇女正离开时，有些曾像死人一样躺着的守卫也进城去，把这些所发生的一切事，即他们自己也能观察到的事，报告给司祭长。他还告诉我们，他们和长老聚集商议后，给了兵士许多钱，吩咐他们说：祂的门徒趁他们睡觉时来把祂偷去了，同时答应他们，他们会向总督说好话，因总督曾派那些守卫。最后，他补充说：他们拿了钱，就照他们所教的去做了，这话在犹太人中间流传直到今日。\par

\SourceRecord{Translated by S.D.F. Salmond.；Nicene and Post-Nicene Fathers, First Series；Vol. 6.；Edited by Philip Schaff.；Buffalo, NY: Christian Literature Publishing Co.,；1888.；<http://www.newadvent.org/fathers/1602324.htm>.}

% structure: p0001 source=h1 target=section reason=page-title

\section{\WorkTitle{福音的和谐，第三卷}}

% structure: p0002 source=h2 target=subsection reason=relative-heading-rank; base=h2

\subsection{第25章 论\Person{基督}后来向门徒显现自身，以及当比较四位圣史以及\Person{保禄}宗徒和\WorkTitle{宗徒大事录}所给出的记载时，是否能建立彻底和谐的问题。}

70. 我们必须开始考察主在复活后如何向门徒显示自己的方式，这不仅是为了清楚显明四位福音作者在这些主题上的一致性，也为了展示他们与宗徒保禄的一致，保禄在《致格林多人前书》中论及此事。保禄的陈述如下：\Quote{因为我当日所领受而传交给你们的，其中首要的是：基督照圣经所载，为我们的罪死了，被埋葬了，第三天照圣经所载复活了，并显现给刻法，然后显现给那十二位；此后又同时显现给五百多位弟兄，其中大多数至今仍在，但有些已经长眠了。此后，显现给雅各伯，然后显现给众宗徒。最后也显现给了我，如同一个流产儿。}然而，这一系列显现是任何一位福音作者都没有记载的。因此我们必须考察他们记录的顺序是否与保禄的记载相矛盾。因为保禄并未讲述全部，福音作者也没有在他们的报告中包罗一切。因此，我们真正要研究的乃是这样一个问题：在这些不同叙述中我们所遇到的那些事件中，是否有任何东西足以建立作者之间的不一致。在四位福音作者中，惟有路加省略了主如何被妇女看见的记载，而只限于讲述天使的显现。玛窦则告诉我们，主在她们从坟墓返回时遇见了她们。马尔谷也提到主首先显现给玛利亚玛达肋纳；若望同样如此。只有马尔谷没有说明主如何显现给她，而若望却给了我们解释。此外，路加不仅默然不语主曾向妇女显现的事实，正如我已指出的，而且记载了两个门徒——其中之一是克罗帕——在认出主之前与主交谈，他们的语气似乎暗示妇女们除了看见天使告诉她们主活着之外，没有讲述任何其他显现。因为路加的叙述如下：\Quote{就在那一天，他们中有两个人往一个村庄去，村名厄玛乌，离耶路撒冷约有六十斯塔狄。他们彼此谈论所发生的一切事。正谈论讨论时，耶稣亲自走近，与他们同行。但他们的眼睛被蒙蔽，以致认不出祂。祂对他们说：你们走路彼此谈论的是些什么事？你们面带愁容。那一个名叫克罗帕的回答祂说：独有你在耶路撒冷作客，不知道这几天在那里所发生的事吗？祂问他们说：什么事？他们对祂说：关于纳匝肋人耶稣，祂在天主和众百姓面前，是位言行有能力的先知。我们的司祭长和首领竟把祂交出，判了死罪，钉在十字架上。我们原本指望祂就是要拯救以色列的那位。此外，这些事发生到今天已是第三天。还有我们中的几个妇女惊吓了我们：她们清早到了坟墓那里，没有找到祂的遗体，就回来说看见了天使显现，天使说祂活着。我们中的几个人也到了坟墓那里，发现正如妇女所说，但没有看见祂。}所有这些事，根据路加的记载，他们讲述的情形正如他们能回忆并想起妇女们所报告给他们的，或是门徒们在得知遗体被移走的消息后跑到坟墓时所报告的。同时，路加自己只记载了伯多禄跑到坟墓，弯腰看见细麻布单独放在那里，就回去了，心里希奇所发生的事。关于伯多禄的这条记录，是在讲述这两个在路上遇见主的门徒之前，也是在妇女们看见天使并听到天使说耶稣复活的故事之后插入的；因此这个位置似乎标志了伯多禄跑到坟墓的时间。但我们仍必须认为路加在这里以总结性重述的形式插入了关于伯多禄的段落。因为伯多禄跑到坟墓的时候也是若望跑到坟墓的时候；那时他们听到的仅仅是妇女们，特别是玛利亚玛达肋纳传达给他们的消息，说遗体被搬走了。而且，那妇人带来这消息的时候，正是她看见石头从坟墓被滚开的时候。后来才发生了其他事情，与天使的显现以及主自己的显现有关：主向妇女们显示了两次，一次在坟墓旁，另一次在她们从坟墓返回时遇见她们。然而这发生在他被那两个在路上的人看见之前，其中一人是克罗帕。因为当这个克罗帕与主谈话时，在他认出主之前，他没有明确说伯多禄去过坟墓。但他的话说：\Quote{我们中的几个人也到了坟墓那里，发现正如妇女所说；}这句话也应理解为以总结性重述的形式引入的。因为它指的是妇女们首先向伯多禄和若望关于遗体被搬走的报告。因此，当路加在这里告诉我们伯多禄跑到坟墓，并且也记载克罗帕提到他们中的一些人去过坟墓时，应视为证实若望的记载，即两个人去了坟墓。但路加第一次只特别提到伯多禄，正是因为玛利亚最早把消息带给了他。然而还有一个困难：同一路加没有说伯多禄进去，只说弯下腰看见细麻布单独放着，就回去了，心里希奇；而若望暗示是本人（即耶稣所爱的那个门徒）这样观看，没有进入坟墓（他是第一个到达的），只是弯腰看见细麻布放在原处；尽管他也补充说后来进了坟墓。因此解释很简单：伯多禄起初确实弯腰观看，如路加所描述的那样，若望未提及；伯多禄稍后进去了，但仍是在若望进去之前。这样，我们将发现所有这些作者都真实地记述了所发生的事，用词上并无矛盾。\par

71. 因此，我们不仅要采纳四位福音作者所呈现的记述，也要采纳宗徒保禄的陈述，将整体连成一个连贯的叙述，展示所有这些事件可能发生的次序，涵盖主向男性门徒的所有显现，而省略祂先前对妇女们的宣告。现在，在所有男性中，伯多禄被认为是基督首先显现的那一位。至少，就四位福音作者和宗徒保禄实际提及的所有个体而言，这一点成立。但同时，谁敢肯定或否认祂可能在向伯多禄显现之前，已显给他们中的某人，尽管所有作者都对此保持沉默？因为保禄所给的陈述并非\BibleQuote{祂首先显现给刻法。}而是这样：\BibleQuote{祂首先显现给刻法，然后显现给十二位；此后又同时显现给五百多位弟兄们。}这样，就不清楚这十二位是谁，正如我们不知道那五百位是谁一样。的确，很可能此处所列举的十二位是属于众门徒中的某十二位不知名者。因为那时宗徒保禄可能将主所指定的宗徒们称为十一位，而非十二位。事实上，一些抄本确实包含这一读法。然而，我认为这是那些对经文感到困惑的人所作的修改，他们以为这指的是那十二位宗徒，而在犹达斯消失时，他们实际上只有十一位。那么，可能正是那些包含\Quote{十一位}读法的抄本更为正确；也可能保禄意指用这句话来理解某些其他的十二位门徒；或者，事实可能是他意指那个被祝圣的数字保持如前，尽管群体已缩减为十一位：因为十二这个数字，用于宗徒时，具有如此神秘的重要性，以至于为了保持同一数字的精神象征，只能选举一个人，即玛弟亚，来填补犹达斯的位置。但无论采纳这些观点中的哪一个，都不一定会产生与真理不一致或与其中任何一位最可信的史家相矛盾的结果。尽管如此，仍然有可能的假设：在祂显现给伯多禄之后，祂接着显现给那两人，其中一个是克罗帕，关于他们，路加给了我们一个完整的叙述，而玛尔谷只给了我们一个非常简短的记载。后者以这些简洁的措辞报告了同一件事：\BibleQuote{此后，祂以另一种形像显现给他们中的两个人，当他们走路并前往一个乡间宅邸时。}因为我们假定所提及的居所也可能被称为乡间宅邸，并非不合理；正如白冷本身，从前被称为城，即使在现今也被称为村庄，尽管自从在这城中出生的主的名如此广泛地在万国教会中被宣告以来，它的尊荣已经变得如此之大。在希腊抄本中，我们发现的读法与其说是\Quote{乡间宅邸}，不如说是\Quote{庄园}。但那个术语不仅用于住所，也用于城外的自由市镇和殖民地，而城市是其余地方的元首和母亲，因此被称为大都会。\par

72. 再者，如果马尔谷告诉我们主以另一种形像显现给这些人，路加则提及同一件事，说他们的眼睛被\Quote{阻止}，以致不认识祂。因为他们的眼睛被某种东西遮住，这遮住一直持续到擘饼时，为的是指向一个众所周知的奥迹，直到那时祂身上不同的形像才向他们显现，而他们直到擘饼时才认出祂，如路加的叙述所示。因此，与他们仍不明白默西亚必须受死并复活的心智状态相应，他们的眼睛也承受了类似的情况；这并不是说真理本身误导人，而是说他们自己没有能力认识真理，把这事想成了别的东西。这一切的深层含义是：任何人若不在祂的身体内——也就是说，在祂的教会内——就不应以为自己已获得对基督的认识；这身体的合一由一位宗徒在饼的圣事象征下提醒我们注意，他说：\BibleQuote{我们众人都是一个饼、一个身体。}正是如此，当祂递给他们祂所祝福的饼时，他们的眼睛开了，认出了祂，也就是说，他们的眼睛开了，得以认识祂，因为那阻止他们认出祂的障碍如今被移除了。因为他们肯定不是闭着眼睛走路。但他们身上有某种东西阻碍他们正确看见眼前的事物——这种状况通常由黑暗或某种体液引起。然而，这并不是说主不能改变祂肉身的形状，使祂的外表真实地、实际上是不同的，而不是他们习惯看到的那一个。因为，甚至在被捕以前，祂在山上显圣容，祂的面容\BibleQuote{发光如同太阳。}祂曾把真正的水变成真正的酒，也能毫无疑义地把任何形体变成祂喜欢的任何其他形体。但这里的意思是，当祂以另一种形像显现给那两个人时，祂并没有那样做。因为祂对这些人显现的并不是祂所是的，因为他们的眼睛被阻止，以致不认出祂。此外，我们也可以合宜地认为，他们眼睛的这种障碍来自撒殚，目的是阻止他们认出耶稣。然而，基督允许这种情况一直持续到接受饼的奥迹为止。这样，教训可能是：当我们成为祂身体合一的参与者时，我们就理解到敌人的障碍被移除，并得到自由去认识基督。\par

73. 此外，有必要相信这些就是马尔谷所提及的同一批人。因为他告诉我们，他们去把这事报告给其余的人；正如路加所述，这些人在同一时刻起来，返回耶路撒冷，找到那十一位门徒和跟他们在一起的人，他们说：\BibleQuote{主果然复活了，而且显现给西满。}然后他补充说，这两个人也报告了路上发生的事，以及怎样在擘饼时被他们认出来。因此，此时耶稣复活的报告已由那些妇女和西满伯多禄传达，因为主已向他显现过。这两个门徒发现他们到耶路撒冷所遇见的人正在谈论同一件事。因此，可能是恐惧使他们在路上先前不愿提及他们已听说祂复活了，而只限于陈述妇女们如何看见了天使。因为他们不知道自己在跟谁说话，所以有理由担心不要漏出任何有关基督复活的话，以免落入犹太人手中。但再次，我们必须指出，马尔谷说\BibleQuote{他们就去报告给其余的人，但那些人也不相信他们。}然而路加告诉我们，那些人已在说主果然复活，并显现给西满。但解释不是很简单吗：他们当中有些人拒绝相信所报告的事。此外，谁不明白马尔谷只是省略了路加叙述中详细说明的某些事情——也就是说，耶稣认出他们之前与他们谈话的内容，以及他们在擘饼时认出祂的方式？因为，在记载祂如何以另一种形像向他们显现，当他们走向一个乡间住处时，马尔谷立即加上这句话：\BibleQuote{他们就去报告给其余的人，但那些人也不相信他们。}好像人们能报告一个不认识的人，或者那些只见过祂另一种形像的人能认识祂！因此，毫无疑问，马尔谷没有给我们解释他们认出祂的方式，以便他们能向别人报告。因此，这件事值得我们铭记，以便我们习惯于注意到这些福音作者的习惯：他们省略那些没有记录的事情，并且把他们叙述的事实以这样的方式连接起来，以致于那些没有适当考虑上述习惯的人，没有什么比这更容易引起误解，导致他们以为圣史们之间存在矛盾。\par

74. 路加接着按以下所述继续叙述：\BibleQuote{他们正谈论这些事的时候，耶稣亲自站在他们中间，对他们说：\Quote{愿你们平安！是我，不要害怕！}他们却惊慌害怕，以为看见了鬼魂。祂对他们说：\Quote{你们为什么惊慌？为什么心里起疑念？看看我的手和脚，实在是我！摸摸我，看看！因为鬼魂没有骨肉，如同你们看我所有的。}说了这话，就把手和脚给他们看。} 主复活后显现自己，若望的下述言论也被认为是指同一行为：\BibleQuote{正是那周的第一天晚上，门徒所在的地方，因为怕犹太人，门都关着，耶稣来了，站在中间对他们说：\Quote{愿你们平安！}说了这话，就把手和肋旁指给他们看。} 如此，我们也可将路加所记而若望略去的某些事情，联系若望的这些话。因为路加继续这样说：\BibleQuote{他们由于欢喜，还是不敢相信，只是惊讶；耶稣对他们说：\Quote{你们这里有什么吃的没有？}他们便给了祂一片烤鱼和一块蜜饼。祂在他们面前吃了，把剩下的给了他们。} 再者，路加略去、若望记载的一段，可以接在这些话之后。内容如下：\BibleQuote{门徒看见主，就欢欣踊跃。耶稣又对他们说：\Quote{愿你们平安！就如我圣父派遣了我，我也同样派遣你们。}说了这话，就向他们嘘了一口气，说：\Quote{你们领受圣神吧！你们赦免谁的罪，谁的罪就得赦免；你们保留谁的罪，谁的罪就被保留。}} 又一次，我们可以附上若望略去、路加插入的另一段。它这样说：\BibleQuote{耶稣对他们说：\Quote{我以前还同你们在一起的时候，就对你们说过这话：凡梅瑟法律、先知和圣咏上指着我所记载的话，都必须应验。}于是祂开启了他们的明悟，叫他们理解圣经；又对他们说：\Quote{经上曾这样记载：基督必须受苦，第三天从死者中复活；并且必须从耶路撒冷开始，因祂的名向万邦宣讲悔改和罪过的赦免。你们就是这些事的见证人。我要把我圣父所恩许的，派遣到你们身上；至于你们，你们要留在城中，直到佩戴上自高天而来的能力。}} 那么，请注意路加在此如何提到了圣神的恩许，这恩许除了在\BibleBook{若望}{福音}中，我们在别处并未看到主曾赐予。这一点值得特别注意，以便我们牢记：众福音作者如何彼此证实对方的真理，即使在有些主题上，他们中有些人并未亲自记载，但他们知道这些事已被其他作者记述。 在这些事之后，路加对其它所发生的一切缄默不言，在他的叙述中只引入耶稣升天这一事件。同时他附加了这【升天的叙述】，仿佛它紧接在主在一周的第一天（即主复活的那一天）所说的这些话以及其他事迹之后；然而，在\BibleBook{宗徒}{大事录}中，同一位路加告诉我们，这事件实际发生在主复活后的第四十天。 最后，关于若望记载宗徒多默在所述场合并不在场，而根据路加，两位门徒（其中之一是克罗帕）返回耶路撒冷，发现那十一人和同他们一起的人聚集在一起，那么几乎无需怀疑，我们必须认为多默只是在主按上述方式向弟兄们显现他们正交谈时之前离开了他们。\par

75. 既如此，若望现在记载主第二次显示自己，是在八天后赐予门徒的，那次多默也在场，他此前一直未见过主。叙述接着这样写道：\BibleQuote{过了八天，门徒又在屋里，多默也和他们在一起。那时，耶稣来了，门都关着，站在他们中间说：\InnerQuote{愿你们平安！}然后他对多默说：\InnerQuote{伸过你的手指来，看看我的手；伸过你的手来，探入我的肋旁；不要作不信的人，而要作相信的人。}多默回答他说：\InnerQuote{我主！我天主！}耶稣对他说：\InnerQuote{多默，因为你看见了我，才相信；那些没有看见而相信的，才是有福的。}}主第二次向门徒显现——也就是说，若望在第二次所记载的显现——我们也可以认出马尔谷以简要的方式提及，这符合这位福音作者的惯例。然而，他的措辞\BibleQuote{\Emph{末后}，祂显现给那十一位，他们正在坐席}这一事实造成了困难。困难不在于若望没有提及他们坐席——他很可能省略了这一点——而在于使用了\Quote{末后}这个词，因为这似乎表明主在那次之后没有显示给他们，而若望接着还记载了主第三次在提庇黎雅海边显现。此外，我们还须注意到，同一马尔谷告诉我们，耶稣\BibleQuote{责斥他们的不信和心硬，因为他们没有相信那些在祂复活后看见过祂的人}。这些话指的是祂复活后显现给的两个门徒——他们正走向一个村庄——以及伯多禄；路加叙事的考察已向我们显示，主首先（在宗徒中）显现给伯多禄——也许还有玛利亚玛达肋纳，以及与她一起的其他妇女，当时祂在坟墓旁被她们看见，又当祂在路上遇见她们返回时。因为上述马尔谷以这样一种方式构建了他的记载：他首先插入关于两个门徒的简短说明——祂显现给正在走向村庄的他们——以及他们报告给其余人却未获相信，然后紧接着在相同的上下文中附上这一陈述：\BibleQuote{末后，祂显现给那十一位，他们正在坐席，并责斥他们的不信和心硬，因为他们没有相信那些在祂复活后看见过祂的人。}那么，这个\Quote{\Emph{末后}}一词怎么被使用，仿佛主在此之后就没有被他们看见呢？因为宗徒们最后一次在世上看见主，实际上是祂升天的时候，那事件发生在祂复活后的第四十天。难道在那个时候祂会因他们没有相信那些在祂复活后看见过祂的人而责斥他们吗？那时他们自己已经多次在祂复活后看见祂，尤其是在祂复活的那天——也就是说，一周的第一天，那时将近夜晚，如路加和若望所记载的。因此，我们只能假定，在所讨论的这段经文中，马尔谷的意图是以他简洁的方式，仅就主的复活那天作出陈述；也就是说，那一周的第一天，玛利亚和与她一起的其他妇女在拂晓后看见祂，那天伯多禄也看见祂，那天祂也显现给两个门徒，其中一个是克罗帕，马尔谷自己也似乎提到了他们；那天，当将近夜晚时，祂显现给那十一位（但多默除外）以及与他们在一起的人；那天，最后，上述那些人向门徒报告了他们所看见的事。因此，他使用了\Quote{\Emph{末后}}这个词，因为提到的事件是这同一天发生的最后一件。因为当两个门徒从他们擘饼时认出祂的地方返回，并进入耶路撒冷，找到那十一位以及与他们在一起的人——如路加告诉我们——他们正彼此谈论主的复活以及祂已显现给伯多禄；这两个人也述说了路上发生的事，以及他们怎样在擘饼时认出了祂。但是，肯定也有在场的人不相信。因此，我们看出马尔谷的话\BibleQuote{他们也不相信他们}是真实的。因此，正如马尔谷告诉我们的，当这些人现在坐席时，正如路加告诉我们的，他们正在谈论这些事，主站在他们中间，对他们说：\BibleQuote{愿你们平安。}正如路加和若望都记载的。此外，当祂进入他们中间时，门是关着的，若望独自提到这点。这样，在路加和若望所报告的主在那次场合对门徒所说的话中，也包括了马尔谷所举出的斥责，其中主责备他们不相信那些在祂复活后看见过祂的人。\par

76. 然而，若按路加和若望所记，这个主日之夜初时，则马尔谷说主向十一人显现时他们正坐席，也可能产生理解上的困难。因为若望分明告诉我们，那次多默宗徒不在场；我们相信，当主进入他们中间之前，多默已经离开，但这是在两个从村庄返回的门徒与十一人交谈之后，正如我们从路加所发现的。诚然，路加在他的叙述中提出了一个时间点，我们可以合理地假定，首先多默在他们谈论这些事时出去了，然后主进来了。然而马尔谷说：\Quote{最后，他向那十一人显现，他们正坐席}，这迫使我们承认多默也在那里。但或许他特意称他们为十一人，尽管其中一人缺席，因为犹达斯的位置尚未由玛弟亚补上之前，这个宗徒团体就是以这个数目称呼的。或者，若难以接受这个解释，我们仍可假定，在四十天期间主多次显示给门徒之后，他也在最后一次向十一人显现时他们正坐席——即是在第四十天本身；而且，因他即将离开他们升天，他在那值得纪念的一天特意责斥他们不信那些在他复活后见过他的人，直到他们亲眼见到为止；特别是因为，当他升天后他们传福音时，连外邦人也愿意相信他们所未见之事。因为，马尔谷在提及这次责斥之后，立刻接着引用这段话：\Quote{他又对他们说：\InnerQuote{你们往普天下去，向一切受造物宣讲福音。信而受洗的，必要得救；但不信的，必被定罪。}}因此，若他们受命宣讲不信者将受审判，而他们所不信的正是他们未曾看见的，难道他们自己不应首先被责备，因他们拒绝相信那些主更早显示给他们的人，直到他们亲眼看到主为止吗？\par

77. 在接下来的部分，我们有进一步的提示，认为这是主最后一次以身体形态向宗徒显现。因为同一位马尔谷继续写道：\BibleQuote{这些记号将伴随着信的人：因我的名驱逐魔鬼，说新语言，手拿毒蛇，若喝了什么致命的毒物，也决不受害，按手在病人身上，病人就会痊愈。}然后他补充了这一陈述：\BibleQuote{主耶稣给他们说了这些话以后，就被接上升天，坐在天主的右边。他们出去，到处宣讲，主与他们合作，并以随后发生的奇迹来证实所传的道理。}当他写道\BibleQuote{主给他们说了这些话以后，就被接上升天}，这似乎足以表明这是他在世上对他们说的最后一番话。同时，这些话似乎并不绝对限于此意。因为他说的不是\BibleQuote{他给他们说了这些话以后}，而只是\BibleQuote{他给他们说了话以后}；因此，若有需要作此解释，完全可以假定这不是最后一次讲论，也不是他在世上与他们同在的最后一天，但所有他在这些日子里对他们所说的事，都可以包括在\BibleQuote{他给他们说了话以后，就被接上升天}这一句里。然而，鉴于我们上面详细论述的理由，更使我们断定这是最后一天，而不是专指那十一人，当时因多默缺席他们只有十人，所以我们认为，在马尔谷所记的这一讲论之后——我们需按正确顺序将其与宗徒大事录所记的其他话语，无论是门徒的还是主自己的话，联系起来——我们必须相信主被接升天，即在他复活后第四十天。\par

78. 再者，若望虽然明确说他略过了耶稣所行许多事，但他仍愿意给我们叙述主复活后第三次向门徒显现，即在提庇黎雅海边，向七个门徒——即伯多禄、多默、纳塔乃耳、载伯德的两个儿子，以及另外两个未提名者。那一次他们正在捕鱼；遵从他的命令，将网撒在右边，拉上许多大鱼，共一百五十三条；他又三次问伯多禄是否爱他，嘱托他牧养自己的羊，预言他将受的苦，并论及若望说：\BibleQuote{我愿意他存留直到我来。}若望就此结束了他的福音。\par

79. 现在我们须考虑\Person{耶稣}在\Place{加里肋亚}首次向门徒显现的场合。因为\Person{若望}记述为第三件的事件，发生在\Place{加里肋亚}的\Place{提庇黎雅}海边。若有人想起五饼的奇迹，同一\Person{若望}以这些话开始记述：\BibleQuote{此后，\Person{耶稣}渡过\Place{加里肋亚}海，即\Place{提庇黎雅}海。}，便可看出场景是在那个地区。那么，除了\Place{加里肋亚}，还有什么地方应自然地被认定为祂复活后向门徒初次显现的合适地点呢？当我们想起天使的话时，似乎就得出了这一结论。按\WorkTitle{玛窦福音}记载，天使对来到坟墓的妇女说：\BibleQuote{你们不要害怕；我知道你们寻找那被钉的纳匝肋\Person{耶稣}。祂不在这里，因为祂已照祂所说的复活了。你们来看看安放主的地方；快去告诉祂的门徒：祂已由死者中复活了；看，祂在你们以先往\Place{加里肋亚}去，在那里你们要看见祂。看，我已经告诉了你们。}\Person{马尔谷}也有类似的记述，无论他所提到的天使是否同一位。其版本如下：\BibleQuote{不要惊慌！你们寻找那被钉的纳匝肋\Person{耶稣}。祂复活了，不在这里；请看安放祂的地方。但你们去，告诉祂的门徒和\Person{伯多禄}：祂在你们以先往\Place{加里肋亚}去，在那里你们要看见祂，正如祂以前对你们说过的。}现在这些话给人的印象似乎是：\Person{耶稣}复活后并不向门徒显现，而是在\Place{加里肋亚}。然而，连\Person{马尔谷}本人也没有记载这所指的显现；他告诉我们祂如何在每周第一日的清晨先显现给\Person{玛利亚玛达肋纳}；她怎样去告诉那些与祂一起的人，他们正在哀悼哭泣；这些人怎样拒绝相信她；此后，祂又怎样被那两个往乡下住所去的门徒看见；这两个人怎样将所发生的事报告给其余的人——按\Person{路加}和\Person{若望}一致证实，这事发生在\Place{耶路撒冷}，正在主复活的那一天，且夜幕将至。之后，同一福音作者接着讲到祂称为最后一次的显现，即那十一位坐席时蒙受的显现；他叙述了那一幕之后，告诉我们祂如何被接升天，如我们所知，这事发生在离\Place{耶路撒冷}不远的\Place{橄榄山}。因此，\Person{马尔谷}没有一处记载他所宣称的天使预告的实际应验。另一方面，\Person{玛窦}将自己的记述局限于一件事件，并且无论之前或之后，都没有提到门徒看见复活的主的任何其他地点，除了天使预言中所指定的\Place{加里肋亚}。简言之，这位福音作者先介绍了天使对妇女说话的话；然后加上她们去时发生的事，以及守卫如何被收买而散布假报告；接着插入他的记述[关于加里肋亚的显现]，仿佛那正是紧接着他所叙述的事。因为天使的话，\BibleQuote{祂复活了；看，祂在你们以先往\Place{加里肋亚}去。}，的确使人似乎有理由认为，在那次显现之前不会有其他事情发生。因此，\Person{玛窦}的记述如下：\BibleQuote{于是十一位门徒往\Place{加里肋亚}去，到了\Person{耶稣}给他们指定的山上。他们看见祂，就朝拜了祂，但有些人怀疑。\Person{耶稣}前来对他们说：天上地下的一切权柄都赐给了我。所以你们要去使万民成为门徒，因父及子及圣神之名给他们授洗，教训他们遵守我所吩咐你们的一切。看，我同你们天天在一起，直到今世的终结。}\Person{玛窦}就是用这些话结束了他的\WorkTitle{玛窦福音}。\par

80. 因此，若不是其他圣史的叙述促使我们必须更仔细地考察全部问题，我们或许会认为主在复活后初次显现给门徒的地点，只能是加里肋亚，别无他处。同样，如果玛尔谷对天使的宣告略而不提，任何人都可能认为玛窦之所以告诉我们门徒如何前往加里肋亚的一座山上，并在那里朝拜了主，是出于渴望表明那命令，以及他也记录由天使传达的预言，确实得到了应验。然而，就实际情况而论，路加和若望都足够清楚地证实，主在复活当天就在耶路撒冷被祂的门徒看见；而耶路撒冷距离加里肋亚甚远，因此祂不可能在同一天内被同一批人在两地看见。同样，玛尔谷虽以相似的措辞报告了天使的宣告，却从未提及主复活后确实在加里肋亚被门徒看见。因此，这些考虑迫使我们探究这句话的真实含义：\BibleQuote{看，祂要在你们以先往加里肋亚去，在那里你们要看见祂。}因为如果连玛窦本人也没有说明十一位门徒前往加里肋亚、到耶稣指定的山上，并在那里看见祂并朝拜祂，我们或许会认为那预言并没有字面上的应验，而整个宣告是旨在传达一种比喻意义。我们便会从路加记载的话中找到一种对应，即：\BibleQuote{看，我今天明天驱魔治病，第三天就要完毕。}这个预言确实没有在字面上应验。同样，如果天使曾说：\Quote{祂要在你们以先往加里肋亚去，在那里你们将首先看见祂；}或\Quote{唯有在那里你们要看见祂；}或\Quote{除那里以外别无他处你们要看见祂；}那么毫无疑问，在这种情况下，玛窦就会与其他圣史相矛盾。然而，实际情况是，话语就是如此：\BibleQuote{看，祂要在你们以先往加里肋亚去，在那里你们要看见祂。}而且没有说明那会面发生的精确时间——是在最早的时机、在其他地方看见祂之前，还是在较晚的时期、在他们也在加里肋亚以外的地方看见祂之后；此外，尽管玛窦记载门徒前往加里肋亚的一座山上，他既没有指定那趟行程的日子，也没有按某种顺序来叙事，迫使我们非认为这一事件必定是首次显现不可。因此，我们可以得出结论：玛窦与其他圣史的叙述并无矛盾，而且他使我们能够在与其自身报告一致的前提下，理解并接受这些其他记载的含义与真实性。同时，由于主如此指示——不是指向祂打算初次显现的地点，而是指向加里肋亚地区，祂无疑后来在那里显现了；并且祂通过天使（他说：\BibleQuote{看，祂要在你们以先往加里肋亚去，在那里你们要看见祂。}）和祂自己的话（\BibleQuote{你们去报告我的弟兄，叫他们往加里肋亚去，在那里必要见我。}）传达这些关于看见祂的指示——在这些事实中，我们发现一些考虑，使每个信徒都渴望探究：这一切可能被理解为具有怎样的神秘意义。\par

81. 然而，首先我们也必须考虑祂可能如此以身体形态在\Place{加里肋亚}显现的时间问题，根据\Person{玛窦}以这些话所作的陈述：\Quote{于是那十一位门徒往\Place{加里肋亚}去了，到了\Person{耶稣}给他们指定的山上；他们看见祂，就朝拜了祂，但有些人怀疑。}显然，这并非祂复活的那一天。因为\Person{路加}和\Person{若望}一致告诉我们，非常清楚地，祂在那一天当夜幕降临时在\Place{耶路撒冷}被人看见；而\Person{马尔谷}在这件事上却不那么清楚。那么，他们是在什么时候在\Place{加里肋亚}看见了主呢？我指的不是\Person{若望}所提及的、在\Place{提庇黎雅海}边的那次显现；因为在那次，只有他们七人在场，且被发现正在捕鱼。但我指的是\Person{玛窦}所详细描述的那次显现，那时那十一人在山上，\Person{耶稣}按天使的预告先他们而去了那里。因为\Person{玛窦}陈述的要旨似乎是：他们就在那里找到了祂，正是因为祂按约定先他们而去。那么，这事既没有发生在祂复活的那一天，也没有发生在随后的八天里；过了这八天，\Person{若望}记载主向门徒们显现，那时\Person{多默}——他在祂复活的那一天没有看见祂——第一次看见了祂。因为，确实，假如那十一人在这八天期间真的在\Place{加里肋亚}的山上看见了祂，那么很可以问：\Person{多默}，他曾是这十一人中的一员，怎能说他在这八天结束时是第一次看见祂呢？对于那个问题，没有答案，除非真的有人说：他们不是那十一人（那些人那时已带有宗徒的特定称号），而是从祂众多的追随者中选出的另外十一位门徒。因为那十一人确实是仅有的尚被称为宗徒的人，但他们并非仅有的门徒。因此，情况或许是：那次确实指的是宗徒们；但并非全部，只有他们中的一些人在场；还有其他门徒与他们在一起，以致在场的人数凑足了十一人；而在那八天结束时第一次看见主的多默，那次却不在场。因为当\Person{马尔谷}提及那十一人时，他不用一般的表达\Quote{十一人}，而是明确地说：\Quote{祂向\Emph{那}十一人显现了。}\Person{路加}也同样记载说：\Quote{他们回到\Place{耶路撒冷}，发现那十一人聚集在一起，还有与他们在一起的人。}在那里，他让我们明白，这些人就是\Emph{那}十一人——也就是说，宗徒们。因为当他加上\Quote{以及与他们在一起的那些人}时，他确实足够清楚地表明，那些其他人与之同在的人，是在某种卓越意义上被称为\Quote{那十一人}的；这使我们明白，这里指的是那些现在被特别称为宗徒的人。因此，很有可能，从宗徒和其他门徒的团体中，凑足了十一位门徒的数目，他们在八天之内于\Place{加里肋亚}的山上看见了\Person{耶稣}。\par

82. 但在这一解决方法的道路上，又出现了另一个困难。因为，当\Person{若望}记载主如何被看见——不是那十一人在山上，而是他们中的七人在\Place{提庇黎雅海}捕鱼时——他附上了这样的陈述：\Quote{这是\Person{耶稣}从死者中复活后，第三次向门徒们显现。}现在，如果我们接受这样的说法：主在这八天期间被那十一位门徒的团体看见，且在祂被\Person{多默}看见之前，那么\Place{提庇黎雅海}边的这次显现就不会是第三次，而是祂第四次显现自己。的确，在此我们必须小心，不要让任何人以为，当\Person{若望}说\Quote{第三次}时，他的意思是主总共只有三次显现。相反，我们必须理解他所指的是天数，而非显现本身的数量；而且，还需注意到，这些日子并非一个接一个紧接而来，而是有间隔之分，正如福音作者自己给出的暗示。因为，撇开祂向妇女们的显现不谈，福音中清清楚楚地记载，祂在复活后的第一天显示了三次：即一次向\Person{伯多禄}；另一次向那两位门徒，其中一位是\Person{克罗帕}；第三次向那较大的团体，当时他们正在彼此交谈，夜幕降临。但\Person{若望}鉴于这些都发生在同一天，便将它们算作一次显现。然后，他指定第二次——即在另一天的显现——与\Person{多默}也看见祂的那次同一；他又特别指出第三次是在\Place{提庇黎雅海}边，也就是说，并非字面意义上的第三次显现，而是祂自我显现的第三天。因此，结果是，在所有这些事件之后，我们不得不认为还发生了另一个场合，根据\Person{玛窦}，那十一位门徒在\Place{加里肋亚}的山上看见祂，祂已按约定先他们而去，这样，天使和祂自己所预告的一切，都得以完全应验，甚至一字不差。\par

83. 因此，在四位福音著者中，我们发现有十次不同的显现，是主在复活后向不同的人显现的。第一次，在坟墓附近向妇女们显现。第二次，在从坟墓返回的路上向同一群妇女显现。第三次，向伯多禄显现。第四次，向两个正往乡下去的人显现。第五次，在耶路撒冷向较多的人显现，当时多默不在场。第六次，在多默看见祂的那次。第七次，在提庇黎雅海边。第八次，在加里肋亚的一座山上，玛窦提及此事。第九次，在马尔谷所说的\Quote{最后，他们正坐席的时候}，由此暗示他们今后不再在地上与祂一同吃饭。第十次，在同一天，但已不是在地上，而是被云彩接去，升到天上，如马尔谷和路加所记载的。这最后一次的显现，马尔谷紧接着讲述主如何在他们坐席时向他们显现之后就直接引出的。因为他的叙述连贯如下：\Quote{主给他们说了这些话以后，就被接升天。} 而路加则略过了四十天内祂与门徒之间的所有事情，在记述了祂复活生命的第一天、在耶路撒冷向较多的人显现之后，就静默地连接了最后一天祂升天的事。他的叙述如下：\Quote{他领他们出去，直到伯大尼附近，就举手降福了他们。正降福他们的时候，就离开他们，被提升天去了。} 因此，除了在地上看见祂之外，他们也看见了祂被高举升天。所以，在祂完全升天之前，福音书中共记载了祂被不同的人看见这么多次：即九次在地上，一次在祂升天时的空中。\par

84. 但同时，并非所有的事都记录了下来，正如若望明确宣告的。因为祂在升天之前的四十天内，与门徒常有来往。但祂并非在这四十天中不间断地一直向他们显现。因为若望告诉我们，在祂复活生命的第一天之后，又过了八天，在那段时间结束时祂才再次向他们显现。被[若望]认定为第三次的显现——即在提庇黎雅海边的那次——可能发生在紧接着的第二天；因为这与任何记载并无冲突。然后，祂在认为合适的时候显现，正如祂曾与他们约定的（这约定在前面的先知预言中也已传达），要在他们之前往加里肋亚去。在这四十天里，祂按照自己的心意，在不同的时机、向不同的人、以不同的方式显现。伯多禄在科尔乃略和与他一起的人面前讲道时，提到了这些显现，他说：\Quote{就是我们在祂从死者中复活后，与祂同吃同喝的那四十天。} 但这并不意味着他们在这四十天内每天都与祂同吃同喝。因为这与若望的陈述相反，若望在其中插入了八天的时间，期间祂未被看见，并且使第三次显现发生在提庇黎雅海边。同时，即使假设祂在那之后每天都向他们显现并与他们一起生活，也不会与叙述中的任何内容产生矛盾。也许，使用\Quote{四十天}这个说法——相当于四个十，因此可能神秘地指向整个世界或整个时间时期——正是因为最初的十天（其中包含了所说的八天）可以按照圣经的惯例，以部分代表整体的原则，合宜地计入。\par

85. 因此，让我们比较宗徒保禄所说的，以判断它是否引起任何难题。他的陈述如下：\Quote{祂第三日按着圣经复活，并显现给刻法。} 他没有说\Quote{祂\Emph{首先}显现给刻法。} 因为这与福音记载祂首先显现给妇女们的事实不一致。他接着说道：\Quote{然后显现给那十二位}；不论祂当时向哪些人显现，也不论是哪个时辰，这至少是在祂复活的那一天。他又继续：\Quote{此后，祂同时显现给五百多位弟兄。} 无论这些人是与那十一位一起聚集在因怕犹太人而关着门的地方，当耶稣在多默离开团体后到来时，还是指在八天之后的另一次显现，都不会造成矛盾。他接着说：\Quote{此后，祂显现给雅各伯。} 然而，我们不应认为这是祂第一次向雅各伯显现；而是可以理解为一次单独向那位宗徒的特殊显现。接下来他补充道：\Quote{然后显现给众宗徒}，这并不表示这是祂第一次向他们显现，而是从那时起，祂与他们有更亲密的来往，直到祂升天之日。最后他说：\Quote{最后，也显现给我这个像流產儿一样的人。} 但那是在祂升天之后相当长的一段时间，从天上的一次启示。\par

86. 因此，现在让我们来探讨我们推迟了的主题，研究玛窦和马尔谷所记载的报告中可能隐藏的奥秘意义，即主复活后作了这个声明：\BibleQuote{我要在你们以先往加里肋亚去，在那里你们要看见我。}\BibleRef{玛 28:7} \BibleRef{谷 16:7} 因为这个宣布，如果它真的应验了，那肯定是在相当长的时间间隔之后才应验的；然而它的措辞似乎最自然地引导我们（尽管这样的结论并非绝对必要）期望所指的显现要么是唯一的，要么是第一次出现的。然而，我们注意到，所涉及的话不是作为圣史自己的话，以叙述过往事件的形式给出的，而是作为天使的话，天使是奉主的委派而说的，随后也是主自己的话；也就是说，这些话被圣史用在他的叙述中，但他把它们作为天使和主所说的直接陈述呈现出来。因此，这无疑迫使我们将它们当作预言性的话语来接受。现在，加里肋亚可以解释为\OriginalTerm{移居}{Transmigration}或\OriginalTerm{启示}{Revelation}。因此，如果我们采纳\OriginalTerm{移居}{Transmigration}的想法，那么对于这句话\BibleQuote{祂在你们以先往加里肋亚去，在那里你们要看见祂}，我们还能想到什么别的意义呢？无非就是基督的恩宠要从以色列子民转移到外邦人那里吗？宗徒们向外邦人宣讲福音时，除非主为他们在人心预备道路，否则不会得到接纳——这可能是对这句话\BibleQuote{祂在你们以先往加里肋亚去}的理解。而且，他们会带着喜乐和惊奇看待困难的破除和移除，以及主通过信者们的启迪为他们开了一扇门——这是对这句话\BibleQuote{在那里你们要看见祂}的理解；也就是说，你们会在那里找到祂的肢体，会在那些接待你们的人身上认出祂活泼的身体。或者，如果我们采取第二种观点，将加里肋亚解释为\OriginalTerm{启示}{Revelation}，那么想法可能是，祂现在不再处于仆人的形像，而是处于与圣父同等的形像中；正如祂在若望的记载中所许诺给爱祂之人的：\BibleQuote{我要爱他，并将我自己显示给他。}\BibleRef{若 14:21} 也就是说，祂后来要显示自己，不仅像他们以前看见的那样，也不仅像祂带着伤痕复活后，让他们可以触摸和看见的那样，而是以那不可言说的光的特征，就是祂照亮每一个来到这世界的人的光，祂在黑暗中照耀，黑暗却不理解祂。这样，祂为我们先行去到一个地方，祂从那里并不离去，虽然祂来到我们这里，并且祂先行到那里并不包含祂离开我们。那将是一个可以被称为真加里肋亚的启示，那时我们将相似祂；在那里我们必要看见祂，像祂那样真实地。那时，对于我们，也有更蒙福的移居，从这个世界进入永恒，如果我们遵守祂的诫命，配得被安置在祂的右边。因为在那里，左边的人要进入永火，而义人却要进入永生。因此他们要移居那里，并在那里看见祂，就如恶人看不见祂一样。因为恶人要被除去，以致不得看见主的荣耀；不义的人不得看见光。因为祂说：\BibleQuote{永生就是：认识你，唯一的真天主，和你所派遣来的耶稣基督。}\BibleRef{若 17:3} 正如祂将在那个永恒中被认识，祂用仆人的形像把祂的仆人带到那里，为使他们得以自由地默观主的形像。\par

\SourceRecord{Translated by S.D.F. Salmond.；Nicene and Post-Nicene Fathers, First Series；Vol. 6.；Edited by Philip Schaff.；Buffalo, NY: Christian Literature Publishing Co.,；1888.；<http://www.newadvent.org/fathers/1602325.htm>.}

% structure: p0001 source=h1 target=section reason=page-title

\section{\WorkTitle{福音的和谐 第四卷}}

\EditorialNote{本书讨论那些属于\Person{马尔谷}、\Person{路加}或\Person{若望}所独有的段落。}\par

% structure: p0003 source=h2 target=subsection reason=relative-heading-rank; base=h2

\subsection{序言}

1. 既然我们已经完整地、按着前后关联地研读了\Person{玛窦}的叙述，并且我们与其他三卷福音所进行的比较，直至其结论，已经证实了一个事实，即这些福音作者中没有一位在其自己的福音中包含了任何与其其他陈述相矛盾的内容，也没有与其他作者所呈现的记载不一致的内容，现在让我们以类似的方式审视\Person{马尔谷}。我们的计划是省略那些他与\Person{玛窦}共有的段落——这些段落我们已经按照必要的程度进行了探讨，并且现在已经处理完毕——然后处理那些剩余的段落，以便对它们进行讨论和比较，并证明它们与其他福音作者所叙述的、直至主的晚餐记述的内容完全和谐。因为从那个点直到结尾，所有四卷福音中所记载的一切事件，我们都已处理，并考虑了它们相互一致的问题。\par

\SourceRecord{Translated by S.D.F. Salmond.；Nicene and Post-Nicene Fathers, First Series；Vol. 6.；Edited by Philip Schaff.；Buffalo, NY: Christian Literature Publishing Co.,；1888.；<http://www.newadvent.org/fathers/1602400.htm>.}

% structure: p0001 source=h1 target=section reason=page-title

\section{福音的和谐，卷四}

% structure: p0002 source=h2 target=subsection reason=relative-heading-rank; base=h2

\subsection{第一章 关于证明马尔谷福音在叙述上与其余福音和谐的问题（不考虑他与玛窦共有的段落），从开头直到记载\BibleQuote{他们进了葛法翁，一到安息日，他就教训他们}的段落；路加也记载了这件事。}

2. 马尔谷是这样开始的：\BibleQuote{耶稣基督，天主子，福音的开始，正如先知依撒意亚所记载的：}直至\BibleQuote{他们进了葛法翁；一到安息日，耶稣就进入会堂教训他们。}在这整段上下文中，一切已在论及玛窦时检查过。然而，关于他进入葛法翁会堂并在安息日教训他们的这个说法，是马尔谷与路加共有的。但这里并不引起困难。\par

\SourceRecord{Translated by S.D.F. Salmond.；Nicene and Post-Nicene Fathers, First Series；Vol. 6.；Edited by Philip Schaff.；Buffalo, NY: Christian Literature Publishing Co.,；1888.；<http://www.newadvent.org/fathers/1602401.htm>.}

% structure: p0001 source=h1 target=section reason=page-title

\section{《福音的和谐》，卷四}

% structure: p0002 source=h2 target=subsection reason=relative-heading-rank; base=h2

\subsection{第二章　论被邪魔折磨的人，并问马尔谷的叙述是否与路加一致（后者与此段记载相同）。}

3. 马尔谷继续叙述如下：\Quote{众人都惊奇他的教训，因为他教训他们正像有权威的人，不像经师们。在他们的会堂里，有一个附有邪魔的人，他喊叫说：\InnerQuote{纳匝肋的耶稣，我们与你有什麼相干？你来毁灭我们吗？}}等等，直到我们读到\Quote{他走遍了全加里肋亚，在他们的会堂里宣讲，并驱逐魔鬼}。虽然其中有些细节是马尔谷和路加共有的，但本节的全部内容在我们按顺序讲解玛窦的叙述时已经处理过，因为这些事情进入叙述顺序的方式使我不能省略。然而，路加说这个邪魔离开那人时没有伤害他；而马尔谷的陈述大意是：\Quote{邪魔就从他身上出来，撕裂他，大声喊叫}。因此，这里似乎有些矛盾。因为如果如路加所说，邪魔\Quote{没有伤害他}，那邪魔怎能\Quote{撕裂他}（或如一些\OriginalTerm{抄本}{codices}所说\Quote{折磨他}）呢？不过，路加给出了完整的说明：\Quote{魔鬼把他摔倒在中间，就从他身上出来，没有伤害他。}因此，我们要明白，马尔谷所说的\Quote{折磨他}正是指路加在\Quote{把他摔倒在中间}这句话中所表达的。而路加加上\Quote{没有伤害他}，意思仅仅是说，那人肢体的这一阵摔打和折磨并没有使他虚弱，正如魔鬼离开时常有的情况，有时甚至在除去困扰的过程中某些肢体也会受伤。\par

\SourceRecord{Translated by S.D.F. Salmond.；Nicene and Post-Nicene Fathers, First Series；Vol. 6.；Edited by Philip Schaff.；Buffalo, NY: Christian Literature Publishing Co.,；1888.；<http://www.newadvent.org/fathers/1602402.htm>.}

% structure: p0001 source=h1 target=section reason=page-title

\section{\WorkTitle{福音的和谐，卷四}}

% structure: p0002 source=h2 target=subsection reason=relative-heading-rank; base=h2

\subsection{第三章 论\Person{马尔谷}所记载的\Person{伯多禄}名字多次被突出的场合，是否与\Person{若望}所给出的宗徒获得那名字的具体时间不一致？}

4. 同一位\Person{马尔谷}接着记载：\BibleQuote{有一个癞病人来到祂跟前，求祂，跪在祂面前对祂说：\InnerQuote{你若愿意，就能洁净我。}}等等，直到经上记载：\BibleQuote{他们喊着说：\InnerQuote{你是天主子。}祂就严厉地嘱咐他们，不要将祂显露出来。}\Person{路加}也记载了与此相类似的内容。但没有任何矛盾之处。\Person{马尔谷}继续写道：\BibleQuote{祂上了山，把自己所愿意的人召来；他们便来到祂面前。祂就选定了十二人，为使他们同祂常在一起，并为派遣他们去宣讲；并赐给他们权柄，以医治疾病、驱逐魔鬼。祂还称\Person{西满}为\Person{伯多禄}。}等等，直到经上记载：\BibleQuote{那人就走了，开始在\Place{德卡波立}传扬\Person{耶稣}为他做了何等大事；众人都惊讶不已。}我知道在按照\BibleBook{玛窦}{福音}的顺序时，我已经谈过门徒的名字。因此，我在此重复提醒：不要有人以为\Person{西满}是在这次才首次得到\Person{伯多禄}的名字，也不要以为\Person{马尔谷}在此处与\Person{若望}有任何矛盾，因为\Person{若望}记载那位门徒很早以前就被这样称呼：\BibleQuote{你要叫\OriginalTerm{刻法}{Cephas}，即译作磐石。}因为\Person{若望}在那里记载了主赐给他这个名字的原话。而\Person{马尔谷}在经文中是以复述的方式引入这件事，他说：\BibleQuote{祂称\Person{西满}为\Person{伯多禄}。}因为他的意图是要在此处列出十二位宗徒的名字，并且有必要提到\Person{伯多禄}，所以他决定简要地说明：那名字并非那位门徒一直拥有的，而是主所赐的；不过，并非是在\Person{马尔谷}当时记载的时刻，而是在\Person{若望}记载主说那句话的场合。这段经文所包含的其他内容，与其他福音书没有任何不一致之处，前面也已经讨论过。\par

\SourceRecord{Translated by S.D.F. Salmond.；Nicene and Post-Nicene Fathers, First Series；Vol. 6.；Edited by Philip Schaff.；Buffalo, NY: Christian Literature Publishing Co.,；1888.；<http://www.newadvent.org/fathers/1602403.htm>.}

% structure: p0001 source=h1 target=section reason=page-title

\section{福音的和谐，卷四}

% structure: p0002 source=h2 target=subsection reason=relative-heading-rank; base=h2

\subsection{第四章：论\Quote{祂越嘱咐他们不要告诉任何人，他们越发广传}；以及关于这一说法是否与福音中向我们称赞的祂的预知不一致的问题。}

5. 马尔谷继续如此记述：\Quote{耶稣又乘船渡到对岸，有许多人聚集到祂那里；祂靠近海边；}等等，直到我们读到：\Quote{宗徒们聚集到耶稣那里，将一切，无论他们所作的还是所教的，都告诉了祂。}最后这一部分马尔谷与路加共有，二者之间并无出入。本节其余内容我们已讨论过。马尔谷继续这样记述：\Quote{祂对他们说：你们私下到荒野地方去休息一会儿；}等等，直到这些话：\Quote{但祂越嘱咐他们，他们越发广传；且极其惊讶，说：祂所做的一切都好：祂使聋子听见，使哑巴说话。}在这一切之中，没有任何东西呈现马尔谷与路加之间任何不和谐的迹象；以上所述我们已经在比较这些福音作者与玛窦时考虑过了。同时，我们必须确保没有人会认为我在此处从\BibleBook{马尔谷}{福音}引用的最后一段陈述与所有福音作者的整体相矛盾，这些福音作者在报告祂的其他大多数言行时都明确指出祂知道人心里的事；也就是说，人的思想和欲望无法向祂隐藏。因此若望在以下段落中说得非常清楚：\Quote{但耶稣并不信任他们，因为祂认识所有的人，并且不需要任何人见证人，因为祂知道人心里是什么。}但有什么奇怪呢，如果祂能识别人们当前的思想，祂预先向伯多禄宣布了他将来会有的思想——这思想在他大胆宣称自己准备为祂或与祂同死的时候肯定还没有——那又有什么奇怪呢？既然如此，那么马尔谷的陈述\Quote{祂嘱咐他们不要告诉任何人；但祂越嘱咐他们，他们越发广传}如何能不显得与祂所拥有的这种至高无上的知识和预知相矛盾呢？因为如果祂，作为一位在自己知识中掌握所有人现在和未来意图的天主，知道祂越嘱咐他们不要传扬，他们就越会传扬，那么祂给出这样的嘱咐又有何目的呢？不过，解释也许是这样：祂希望让后进的人明白，那些被赋予传教使命的人应当以更大的热情和热忱来传道，如果连被禁止的人都无法保持沉默的话。\par

\SourceRecord{Translated by S.D.F. Salmond.；Nicene and Post-Nicene Fathers, First Series；Vol. 6.；Edited by Philip Schaff.；Buffalo, NY: Christian Literature Publishing Co.,；1888.；<http://www.newadvent.org/fathers/1602404.htm>.}

% structure: p0001 source=h1 target=section reason=page-title

\section{\WorkTitle{福音的和谐，第四卷}}

% structure: p0002 source=h2 target=subsection reason=relative-heading-rank; base=h2

\subsection{第五章 论\Person{若望}关于那不属门徒圈子却驱魔的人所作的陈述；以及主的答复：\BibleQuote{不要禁止他们，因为那不反对你们的，就是帮助你们的}；并论那答复是否与另一句话相矛盾，主曾说：\BibleQuote{那不与我在一起的，就是反对我的}。}

马尔谷继续如下：\BibleQuote{在那些日子，又有一大群人，没有什么可吃的；}等等，直到这些话：\BibleQuote{若望回答祂说：师傅，我们看见一个人用你的名字驱魔，他不跟从我们，我们就禁止了他。但耶稣说：不要禁止他；因为没有一个人以我的名字行奇迹，能轻易诽谤我；因为那不反对你们的，就是帮助你们的。}路加以类似的方式叙述，只除他没有插入\BibleQuote{因为没有一个人以我的名字行奇迹，能轻易诽谤我。}这一句。因此，此处没有任何引起二者之间不一致的问题。然而，我们必须看看这句话是否应该被认为与主的另一句话相矛盾，即这意思的：\BibleQuote{那不与我在一起的，就是反对我的；那不与我聚集的，就是分散的。}因为这个人既不与主在一起（若望报告说他没有与他们联合跟从主），若主说那不与我在一起的，就是反对我的，那么这人如何不反对主呢？或者，若他反对主，主又为何对门徒说：\BibleQuote{不要禁止他；因为那不反对你们的，就是帮助你们的？}或许有人会说，重要的是要注意，这里主对门徒说：\BibleQuote{那不反对你们的，就是帮助你们的}，而在另一处，主谈自己时说：\BibleQuote{那不与我在一起的，就是反对我的}。这似乎表明，一个人可以不与主在一起，却与主的门徒——可以说是主的肢体——联合。此外，像\BibleQuote{那接待你们的，就是接待我}和\BibleQuote{你们对我这些最小兄弟中的一个所做的，就是对我做的}这样的话，其真理性如何呢？或者，一个人能反对主的门徒却不反对主吗？不，那么像\BibleQuote{那轻视你们的，就是轻视我}和\BibleQuote{你们没有对我最小兄弟中的一个做，就是没有对我做}以及\BibleQuote{扫禄，扫禄，你为什么迫害我}——虽然扫禄迫害的是主的门徒——这些话又作何解释？但是，真正的意思是这样的：一个人不与他在一起的程度，就是他反对他的程度；同样，一个人不反对他的程度，就是他与他在一起的程度。例如，就这个以基督之名行奇迹却不在基督门徒之列的人而言，就他以基督之名行奇迹这件事来说，他是在他们一边的，不反对他们。但是，既然门徒禁止那人做一件实际是在他们一边的事，主就对门徒说：\BibleQuote{不要禁止他。}因为门徒应当禁止那些在他们团契之外的事，以便将那人带回教会的合一，而不应禁止像这样与门徒一致的事，即他以主和师傅之名驱魔。这就是天主教会所遵循的原则：不谴责异端者中共同的圣事；在这些事上，他们与我们同在，不反对我们。但是，教会谴责并禁止分裂和分离，或任何反对和平与真理的情感。因为在这些事上，他们反对我们，正因为他们在那些事上不与我们在一起，并且因为他们不与我们聚集，所以分散。\par

\SourceRecord{Translated by S.D.F. Salmond.；Nicene and Post-Nicene Fathers, First Series；Vol. 6.；Edited by Philip Schaff.；Buffalo, NY: Christian Literature Publishing Co.,；1888.；<http://www.newadvent.org/fathers/1602405.htm>.}

% structure: p0001 source=h1 target=section reason=page-title

\section{福音的和谐 第四卷}

% structure: p0002 source=h2 target=subsection reason=relative-heading-rank; base=h2

\subsection{第六章 论\Person{马尔谷}所记载，主论及那奉基督之名驱魔却不跟从门徒之人所说的话比\Person{路加}更多；以及如何证明这些附加的话与基督禁止拦阻那奉祂名行奇迹之人的心意有关。}

7. \Person{马尔谷}接着叙述说：\BibleQuote{凡因你们属于基督，因这名给你们一杯水喝的，我实在告诉你们，他决不会失掉他的赏报。谁若使这些信我的小子中的一个跌倒，倒不如把一块磨石套在他的脖子上，投在海里。若你的手使你跌倒，砍掉它！你残废进入生命，比有两只手而往地狱里，到那不灭的火里去更好，那里的虫子不死，火也不灭。}以及往下直到所说的话：\BibleQuote{你们当中应当有盐，彼此和平相处。}这些话语，\Person{马尔谷}记载是主在禁止那奉祂名驱魔却不跟从门徒的人被拦阻之后，紧接着所说的。在这一段中，他也引入了一些其他福音作者没有记载的事，但也有些在\Person{玛窦}福音中出现，还有些同样在\Person{玛窦}和\Person{路加}福音中见到。然而，那些福音作者在不同上下文、不同顺序中引入这些事，而非在有人向基督报告那人不跟从门徒却奉祂名驱魔的这一点上。因此，我的意见是：主确实根据\Person{马尔谷}的见证，在这一上下文说了这些话，祂也在其他场合说过同样的话，理由很简单，这些话与祂在此表达的心意十分相关——祂禁止对奉祂名行奇迹之人施加任何阻拦，即使那人并不跟从门徒。因为\Person{马尔谷}把这两段话的关系呈现如下：\BibleQuote{因为谁不反对我们，就是向着我们；凡因你们属于基督，因这名给你们一杯水喝的，我实在告诉你们，他决不会失掉他的赏报。}这清楚地表明，即使\Person{若望}所提出的那个人，因而给主机会开始上述谈论，他也不是像异端者那样脱离门徒的团体而轻视它。但他的处境类似于常见的情况：有些人虽然尚未接收基督的圣事，却因基督徒的名号如此善待基督徒，甚至接待他们，并仅仅因为他们是基督徒而适应他们；主告诉我们，这一类型的人不会失掉他们的赏报。然而，这并不意味着他们应当立刻认为自己因着对基督徒怀有的这种善意就安全稳妥，而同时他们既未藉基督的圣洗得到洁净，也未与祂身体的合一结合。而是说，天主仁慈正在引导他们，使他们也能达到这些更高的事物，从而安全地离开现世。这样的人确实在尚未加入基督徒数目之前，就已经比那些已拥有基督徒名号并分享基督徒圣事，却只推荐那些能引诱其他人跟随他们走向永罚的途径的人更有益[仆人]。主以身体的肢体比喻这些人，并命令将他们像冒犯的手或眼一样从身体中抛出，即从合一的共融中砍掉，好让他们宁愿没有这些同伴进入生命，也不愿与他们一同进入地狱。此外，当他们邪恶的建议（即他们沉溺的绊脚石）得不到同意时，他们就与那些与他们分离的人分开了。确实，如果他们的邪恶性情被所有与他们有共融的好人所发现，他们就被完全从所有人的共融以及神圣圣事的参与中砍掉。但如果他们的这一性情只为某些人所知，而他们的邪恶性情不为大多数人所知，他们就必须被容忍，就像未簸扬前在打谷场上忍耐糠秕一样；也就是说，处置他们的方式既不涉及在不义的共融中与他们达成一致，也不致使善人因他们而离开善人的团体。这就是那些自己里面有盐、彼此和平相处的人所行的。\par

\SourceRecord{Translated by S.D.F. Salmond.；Nicene and Post-Nicene Fathers, First Series；Vol. 6.；Edited by Philip Schaff.；Buffalo, NY: Christian Literature Publishing Co.,；1888.；<http://www.newadvent.org/fathers/1602406.htm>.}

% structure: p0001 source=h1 target=section reason=page-title

\section{\WorkTitle{福音的和谐，第四卷}}

% structure: p0002 source=h2 target=subsection reason=relative-heading-rank; base=h2

\subsection{第七章 论从这一点直到主的晚餐——以晚餐为起点共同讨论四位福音书作者全部叙述的行动——其间马尔谷福音未提出任何需要特别考察的问题。}

8. 马尔谷接着写道：\BibleQuote{祂从那里起身，来到\Place{约旦河}对岸的\Place{犹太}境内；众人又聚集到祂那里，祂照常教导他们。}直到经上记载：\BibleQuote{因为众人都是拿他们多余的投入，但这寡妇却从她的匮乏中把她所有的，全部生活费都投进去了。}在这整个上下文中，当我们按照顺序将其他福音书与\Person{玛窦}进行比较时，已经为消除任何矛盾的表象而审查了上述所有内容。不过，关于那穷寡妇投入\OriginalTerm{两个小钱}{two mites}到银库的记述，只有两位福音书作者，即\Person{马尔谷}和\Person{路加}，作过报道。但他们的和谐并无问题。而且从这一点直到主的晚餐——后者构成了我们共同讨论四位福音书作者所有记录的起点——\Person{马尔谷}没有引入任何需要我们将它与其他陈述进行特别比较，或为驱散不一致的表象而进行研究的内容。\par

\SourceRecord{Translated by S.D.F. Salmond.；Nicene and Post-Nicene Fathers, First Series；Vol. 6.；Edited by Philip Schaff.；Buffalo, NY: Christian Literature Publishing Co.,；1888.；<http://www.newadvent.org/fathers/1602407.htm>.}

% structure: p0001 source=h1 target=section reason=page-title

\section{福音的和谐 第四卷}

% structure: p0002 source=h2 target=subsection reason=relative-heading-rank; base=h2

\subsection{第八章 论\BibleBook{路加}{福音}，尤其论其开端与\BibleBook{}{宗徒大事录}卷首之间的和谐。}

9. 因此，接下来让我们按顺序研读\BibleBook{路加}{福音}。然而，我们将略过他（\Person{路加}）与\Person{玛窦}和\Person{马尔谷}共有的段落，因为所有这些都已处理过了。\Person{路加}于是以如下方式开始他的记述：\BibleQuote{因为已有许多人着手整理在我们中间所成就的事迹，正如那些从起初就亲眼见过并为圣言服务的人所传给我们的；我也从起头仔细查究了一切，决意按次第写给你，最尊敬的\Person{特奥斐罗}，使你认清你所受的教导是确实的。}这个开头并不直接属于福音书的叙述。但它提示我们要注意一个事实：这位\Person{路加}也是另一本名为\BibleBook{}{宗徒大事录}的书的作者。我们持此观点的理由不仅是\Person{特奥斐罗}的名字既出现在那里也出现在这里。因为很可能有另一位同名的人；即便两次指同一个人，也可能有另一个作者向他致献另一部作品，正如\Person{路加}为他写了福音一样。然而，我们将作者同一性的观点建立在这个事实上：第二卷书以下列话语开始：\BibleQuote{\Person{特奥斐罗}，我在前一部书中已论述了耶稣开始所行所教的一切，直到祂藉着\OriginalTerm{圣神}{Holy Ghost}向祂所选、要他们宣讲福音的宗徒们颁赐命令的那一天。}这句话使我们明白，在此之前，他已写了四福音书中之一，这四福音书在教会中享有最高权威。同时，当他告诉我们他撰写了一部论述耶稣开始所行所教的一切直到祂向宗徒颁赐命令那一天的著作时，我们不应理解为他在自己的福音书中完整记述了耶稣在世上与宗徒们同在时所行所说的一切。因为这与\Person{若望}所肯定的相反，他说耶稣还做了许多别的事，倘若一一写出来，所写的书连世界也容纳不下。此外，一个公认的事实是，其他几位福音作者叙述了不少\Person{路加}在其历史中没有涉及的事。因此，其含义是，他写论述所有这些事，是在于他从全部材料中做了选择，纳入他认为适合并适宜用来完成交付给他的责任的那些事实。此外，当他提到许多人\BibleQuote{着手整理在我们中间所成就的事迹}时，他似乎指的是某些未能完成其承担任务的人。因此他也说，他也决意\BibleQuote{按次第仔细写，因为许多人已着手}等等。然而，这里所指的，我们应当理解为那些在教会中没有获得权威的作者，因为他们完全无力正确地完成所着手的事。此外，我们面前这位作者并不限于将自己的叙述延伸到主的复活和升天；他的目标也不仅仅是在四位福音书作者中取得与其劳苦相称的尊荣地位。但他也承担了记录此后由宗徒们所行的事；并且讲述了他认为对读者或听众信德的建立必需且有益的事件，他给予了我们如此忠实的叙述，以至于他的书是唯一被教会承认为值得接受的\BibleBook{}{宗徒大事录}；而其他所有试图写宗徒言行录的作者，虽缺乏第一必要的可信度，但都遭到了拒绝。而且，\Person{马尔谷}和\Person{路加}写作时，时间上完全可能使他们不仅接受基督教会，也接受仍然活在肉身中的宗徒们的检验。\par

\SourceRecord{Translated by S.D.F. Salmond.；Nicene and Post-Nicene Fathers, First Series；Vol. 6.；Edited by Philip Schaff.；Buffalo, NY: Christian Literature Publishing Co.,；1888.；<http://www.newadvent.org/fathers/1602408.htm>.}

% structure: p0001 source=h1 target=section reason=page-title

\section{\WorkTitle{福音的和谐，第四卷}}

% structure: p0002 source=h2 target=subsection reason=relative-heading-rank; base=h2

\subsection{第九章 论如何证明\Person{路加}所记载的捕鱼神迹与\Person{若望}在主复活后记载的类似事件并非同一回事；以及从这一点直到主的晚餐，从该事件直到末了，所有福音作者的记载都已联合审查，在此之后，\WorkTitle{路加福音}和\WorkTitle{马尔谷福音}都没有出现需要特别关注的难题。}

10. \WorkTitle{路加福音}是这样开始的：\BibleQuote{在犹太王\Person{黑落德}的时候，有个司铎名叫\Person{匝加利亚}，属于\Person{阿彼雅}班次；他的妻子是\Person{亚郎}的后代，名叫\Person{依撒伯尔}；}一直到这句：\BibleQuote{祂讲完话后，就对\Person{西满}说：\InnerQuote{把船划到深处去，撒下你们的网捕鱼。}}在整个这一部分，没有任何关于不一致的疑问。的确，\Person{若望}似乎记载了与最后这段类似的事。但他所记载的实在大为不同。我指的是主复活后在\Place{提庇黎雅海}边发生的事。在那件事中，不仅时间截然不同，而且情景本身也完全是另一回事。因为那时网撒在右边，捕到一百五十三条鱼。并且还提到，那都是大鱼。因此，福音作者觉得有必要说明：\BibleQuote{因为鱼虽然众多，网却没有破；}这肯定是因为他想到之前\Person{路加}所记载的那件事，在那件事中，网因鱼太多而破了。至于其余的，\Person{路加}并没有叙述\Person{若望}所记载的那些事，只有关主的受难和复活除外。而整个这一部分，即介于主的晚餐与结论之间的内容，我们已经通过联合比较所有福音作者的见证来处理过，结果证明它们之间完全没有不一致之处。\par

\SourceRecord{Translated by S.D.F. Salmond.；Nicene and Post-Nicene Fathers, First Series；Vol. 6.；Edited by Philip Schaff.；Buffalo, NY: Christian Literature Publishing Co.,；1888.；<http://www.newadvent.org/fathers/1602409.htm>.}

% structure: p0001 source=h1 target=section reason=page-title

\section{\WorkTitle{福音的和谐}，第四卷}

% structure: p0002 source=h2 target=subsection reason=relative-heading-rank; base=h2

\subsection{第十章 论福音作者\Person{若望}，以及他与其余三位的区别}

11. 若望仍保存着，在他与其他三人之间，已无任何比较可作。因为，无论每位福音作者各自记载了一些他人未记的事，但很难证明由这些事中会产生任何涉及真正矛盾的疑问。同样，这是一个明确公认的立场：前三部——即玛窦、马尔谷和路加——主要专注于我们的主耶稣基督的人性，按此人性祂既是君王又是司祭。如此，马尔谷似乎对应于那著名的神秘四活物象征中的人形，他要么看起来与玛窦更为亲近，因他与玛窦一致叙述的事比与其他两人更多，并且在此他合乎君王性格的意念（如我在第一卷所述，君王通常有随从陪伴）；要么，按更可能的解释，他与两者（另两位同观福音作者）都联合。因为尽管他在多数段落与玛窦一致，但在某些其他段落却与路加更为一致。这事实本身显示他同时与狮子和牛犊相关，即与玛窦强调的君王职份和路加引入的司祭职份相关，于此基督也特别显现为人，而马尔谷所承担的形象与这两者都相关。另一方面，基督的天主性——凭此祂与圣父同等，按此祂是圣言，与天主同在的天主，降生成人居住在我们中间的圣言，按此祂与圣父原是一体——被若望特别着手以推荐给我们心智。如同鹰一般，他停留在基督更高层次的言论中，罕有降落到地上。简言之，虽然他清楚表明自己对主母亲的认知，却既不与玛窦和路加一起记录祂的诞生，也不与三人一起叙述祂的洗礼；他所做的只是以崇高而卓绝的方式呈现若翰所作的见证，然后离开这些同伴，与主一起前往加里肋亚的加纳婚宴。在那里，尽管福音作者本人以\Quote{母亲}之名提及她，祂仍这样对她说：\BibleQuote{女人，我与祢有什么相干？}在此，[应理解为]祂并非拒绝那位由祂取得肉躯的母亲，而是要在此时（祂即将变水为酒的时刻）特别恰当地传达祂的天主性之概念；这同一的天主性，也创造了那个女人，并未在她内受造。\par

12. 然后，在提及在葛法翁度过的几天后，福音作者再次来到圣殿，在那里他陈述耶稣论及祂身体之殿的话：\BibleQuote{你们拆毁这座殿，三天内我要把它重建起来。}这句话鲜明地启示，不仅天主在那圣殿中（寓居于降生成人的圣言位格），而且祂自己真正复活了那肉躯，凭的是祂与圣父一体中所具有的特权，按此祂并非与圣父分开行动；而在其他所有段落中，圣经所用的措辞似乎是\Quote{天主复活了祂}；任何别处也找不到这样的说法，说当天主复活基督时，基督也复活了祂自己，因为祂与圣父是同一天主；这正是眼前这段经文（祂说\BibleQuote{你们拆毁这座殿，三天内我要把它重建起来}）的涵义。\par

13. 接着，与尼苛德摩所说的那些话何等伟大、何等神圣！从这些话语中，福音作者再次转向若翰的见证，让我们注意到新郎的朋友因新郎的声音而喜乐。在这句话中，祂让我们明白：人的灵魂本身没有光，也不能获得福乐，除非分享那不变的上智。此后，祂领我们到撒玛黎雅妇人那里，提到水，人若喝了这水，将永远不再渴。再者，祂又带我们回到加里肋亚的加纳，耶稣曾在那里变水为酒。在那段叙述中，祂告诉我们耶稣对那位儿子患病的贵族说：\BibleQuote{除非你们看到神迹和奇事，你们总不会信。}这话旨在将信徒的心智提升到一切可变事物之上，以至于祂甚至不愿意信徒追求神迹本身——尽管这些神迹带有神圣的印记，却仍是在可变的身体上完成的。\par

14. 接着他把我们带回耶路撒冷，讲述了医治患了三十八年疾病之人的故事。当时所说的话是何等丰富，讲论何等详尽！这里我们遇到了这句话：\BibleQuote{犹太人想要杀害祂，因为祂不但违背了安息日，而且还说天主是他的圣父，使自己与天主平等。} 在这段经文中，足够清楚地表明，祂并非像圣人们惯常使用那种说法那样，称天主为祂的父，而是意指祂与天主平等。因为在此之前不久，祂曾对那些指责祂违反安息日的人说：\BibleQuote{我的圣父一直工作，我也工作。} 于是他们的怒火爆发，不仅因为祂说天主是他的圣父，更因为祂用这句话时，希望人们理解祂与天主平等：\BibleQuote{我的圣父一直工作，我也工作。} 在这句话中，祂也理所当然地表明，既然圣父工作，圣子也应工作；因为圣父不是没有圣子而工作。这与祂在同一段落稍后，当那些人因祂的宣告而愤怒时所说的一致，即：\BibleQuote{因为圣父所做的事，圣子也同样做。}\par

15. 最后若望降下，与另外三位相伴，他们的行程与同一位主同行，但在地上，并与他们一同记录用五个饼使五千人吃饱的事。然而，在这段叙述中，唯有他提到，当众人想要立祂为王时，耶稣独自一人退到山上。在我看来，他说这话的目的，只是要向理性的灵魂传达这个真理：基督统治我们的心智和理性，纯粹是在一个祂超越我们、与人类没有本性相通、并且祂确实独自一人、如同圣父的唯一同伴的领域。然而，这是一个奥秘的真理，属血肉的人无法认识，他们的生命匍匐在这较低的地土上，正因为它是一个如此崇高的奥秘。因此，基督自己也从那些习惯于以世俗观念寻求祂国度的人那里退到山上。所以，祂在其他地方如此表达自己：\BibleQuote{我的王国不属于这个世界。} 同样，这也是唯一由若望记载的，他仿佛以太般的飞翔高高翱翔于地上，在正义之太阳的光中欢悦。然后，从与这座山相关的叙述以及五饼的奇迹出发，他与那三位同行了片刻，直到渡海的记述，以及耶稣在水面上行走的场合。但此时他立刻再次上升到主的话语的领域，并讲述了那些极其庄重、冗长、持续高深而崇高的言语，它们以饼的倍增为缘由，当时祂对群众说了以下的话：\BibleQuote{我实实在在告诉你们：你们找我，并不是因为看见了神迹，而是因为吃了饼，并且吃饱了。不要为那必坏的食物劳力，要为那存到永生的食物劳力。} 这些话之后，祂继续以类似的词语讲了很长一段时间，并且以最高昂的格调。那时，有些人从这样崇高的话语的教导中跌倒了，就是那些此后不再与祂同行的人。但也有那些依附祂的人；这些人能够理解这句话的意思：\BibleQuote{叫人活的是圣神，肉是无益的。} 因为确实如此，即使借着肉身，也是圣神有益，惟有圣神有益；而肉身没有圣神，就毫无益处。\par

16. 接下来我们来到这一段：祂的兄弟们——即祂按肉身的亲戚——催促祂上到庆节去，好使祂有机会向群众显明自己。在这里，祂的回答又是何等崇高！\BibleQuote{我的时候还没有到，你们的时候却常是准备好的。世界不能恨你们，却是恨我，因为我指证它，它的行为是邪恶的。} 所以，情况就是这样：\BibleQuote{你们的时候常是准备好的，} 因为你们渴望那样的一天，就是先知所说的：\BibleQuote{我并没有劳苦跟随祢，上主；人的日子我没有渴望过，祢知道。} 也就是说，翱翔到圣言的光中，并渴望亚巴郎所渴望看见的那一天，他看见了，就欢喜。再者，若望记载的祂在节期上到圣殿后所说的话，是多么奇妙、多么神圣、多么崇高啊！这些话是这样的：祂将要往的地方，他们不能去；他们既认识祂，也知道祂从哪里来；那派遣祂的是真实的，他们却不认识祂。这几乎等同于祂所说的：\Quote{你们既知道我从哪里来，又不知道我从哪里来。} 祂借这样的言论还想让人明白什么呢？无非是：按肉身、按族谱和地域，祂可以被他们认识；但就祂的天主性而言，祂却不被他们认识。同样，当祂谈到圣神的恩赐时，祂也向他们显示祂是谁，因为祂能够拥有赐予那至高恩惠的能力。\par

17. 再者，这位圣史记述耶稣从橄榄山返回圣殿后，在祂宽恕了由试探者带来、按律当砸死的犯奸淫妇人之后，所说的话是何等重大！那时祂用指头在地上写字，仿佛要指明，像这些人的品性，将被写在世上，而不是天上；祂也曾告诫门徒，要因他们的名字被记在天上而欢喜。或者，祂是要借此表明：祂藉着谦卑自下（弯腰示意）在地上行奇事；又或者，时候已到，祂的法律不应再像从前写在贫瘠的石版上，而应写在能结出果实的土地上。于是，在这些事以后，祂声称自己是世界的光，并宣告凡跟随祂的，不在黑暗中行走，必有生命的光。祂还说：\BibleQuote{我就是开始，就是我给你们所讲的。}藉此称号，祂清楚将自己与祂所造的光区分开来，并以那使万物被造的光自居。因此，当祂说自己是世界的光时，我们不可将这话理解为如同祂以类似措辞对门徒所说的：\BibleQuote{你们是世上的光。}因为他们只可比作点着的灯，不可放在斗底下，而要放在灯台上；正如祂论及圣若翰洗者时所说：\BibleQuote{他是一盏点着而发亮的灯。}但祂自己就是开始，关于祂也如此宣布：\BibleQuote{从他的满盈中，我们都领受了。}在所讨论的场合中，祂进一步声称祂自己，即圣子，是真理，这真理将使我们自由，没有它，无人能得自由。\par

18. 接着，若望讲述了天生瞎眼者复明的事迹后，用了相当长的篇幅记述那事件所引出的丰富讲论，论及羊、牧人、门，以及祂舍弃生命并重新取回的能力，在其中祂彰显了天主性的至高权能。此后，他记载在耶路撒冷庆祝修殿节时，犹太人对祂说：\BibleQuote{你使我们的心悬疑不定，要到几时呢？你如果是基督，就明明地告诉我们吧。}然后，他报告了主趁此讲论机会所说出的崇高言语。就在这场合，祂说：\BibleQuote{我与父原是一体。}此后，他又叙述了拉匝禄从死者中复活的事：与此奇事相关，主说：\BibleQuote{我就是复活，就是生命；信从我的，即使死了，仍要活着；凡活着而信从我的人，必永远不死。}在这些话中，我们所认出的不正是祂的天主性的崇高吗？与祂共融，我们将永远生活。此外，若望与玛窦、玛尔谷联合记述了有关伯达尼的事，那里玛利亚将珍贵的香液倒在祂的脚上和头上。然后，直到主的受难与复活，若望与其他三位圣史保持一致，但仅限于他的叙述涉及相同的地点。\par

19. 此外，就主的话语而言，他并未停止攀升至更崇高、更广泛的宣讲，主从这一点开始也亲自说出了这些话语。因为他在书中插入了一篇崇高的讲论，那是当外邦人通过斐理伯和安德肋表示渴望见到主时主所说的，这篇讲论其他福音作者都没有记载。在那里，他还记录了一些重要的话语，再次论及那光照人并使人成为光明之子的光。此后，在有关晚餐本身的内容上——其他每位福音作者都未曾忽略提及祂——若望所记载的耶稣的那些话语是多么丰富、多么崇高啊，而其他人却默然略过了！我不仅可以举出祂谦卑的赞美，当祂洗门徒的脚时，而且还有那极其感人、极其丰富的讲论，当主对那仍留在他身边的十一位门徒讲话时，这十一人是在出卖者被那块饼指出来并出去之后剩下的。正是在这篇若望久久流连的讲论中，祂说：\BibleQuote{看见我的，也就看见了圣父。}也正是在这篇讲论中，祂如此详尽地谈论了圣神，那位护慰者，祂将要派遣给他们的；以及论及祂自己的光荣，这光荣是祂在未有世界以前与圣父同享的；并论及祂使我们合而为一，如同圣父与圣子原是一体——并非祂与圣父和我们成为一体，而是我们要如同圣父与圣子那样合而为一。祂还在这篇讲论中说了许多其他奇妙崇高的事。然而，谁能看不出，用配得上这些主题的任何方式来讨论它们——即使我们有能力这样做——至少也不是我们目前所承担的任务呢？因为我们的目的是帮助那些爱慕天主圣言并研究神圣真理的人明白：在他的福音中，若望确实是同一基督的宣告者和宣讲者，这位真实而诚实的基督正是其他三位撰写福音者也为之作证的那一位，而且其余宗徒也同样为祂作证，他们虽然没有动手编写文字叙述，但至少履行了在正式宣讲中传报祂的类似服务；但同时，若望从书的开头就已被举到基督教义中远为崇高的境界，他只是很少的情况下才保持在与其他福音作者相同的水平上。这些场合特别是以下几点：首先，在约但河畔，关于若翰洗者的见证；其次，在提庇黎雅海的对岸，当主用五个饼给群众吃饱，并在水面上行走时；第三，在伯达尼，那里有一位信德虔诚的妇女将珍贵的香膏浇在祂身上。于是，他继续前行，直到他在受难时与他们相遇——这受难的内容他自然必须与他们一起来叙述。但是，即使在那部分，并且特别是在主晚餐的主题上——没有任何一位作者忽略了它——若望给了我们一个更加丰富的叙述，仿佛他直接从主胸怀（他常依靠在其中）的宝藏中汲取材料。然后，再次，[若望向我们展示]祂如何用更崇高的话语使比拉多感到惊奇，宣告祂的国不属于这世界，祂生来就是君王，祂为此来到世界，为给真理作证。[也正是在这部福音中]祂在复活后带着深邃神秘的意图从玛利亚那里退开，并对她说：\BibleQuote{不要摸我，因为我还没有升到圣父那里去。}也正是在这里，祂向门徒嘘气赐予圣神，使我们因此明白：这位与天主圣三\OriginalTerm{同质共存且同永}{consubstantialis et coaeternus}的圣神，不应被视为仅仅是圣父的神，而也应被视作圣子的神。\par

20. 最后，祂在这里将祂的羊群托付给爱祂的\Person{伯多禄}，伯多禄三次承认了这份爱。然后祂说，祂愿意这位\Person{若望}一直留在世上直到祂来。在这句话中，在我看来，祂似乎又以深邃而神秘的方式传达了一个事实：若望的这份福音使命——在此使命中，他被提升到圣言最清澈的光明中，在那里可以看见\OriginalTerm{天主圣三}{Trinity}的平等与不变，并且首先我们能看到，其他人性在本质上与那因降生成人而使圣言成为血肉的人性相距何等遥远——直到主亲自降临之前，都无法清楚地辨别和认知。因此，它将这样存留直到祂来。目前，它存留在信徒的信德中，但将来我们能够面对面地默观它，那时我们的生命基督显现，我们也与祂一同在光荣中显现。但若有人以为，人仍生活在此尘世生命中，就可能驱散并清除一切由肉身和肉体的幻想所引入的迷雾，达到不变真理的最宁静光明，并以完全脱离现世生活的头脑坚定不移地坚持它，那么此人既不理解他所求的是什么，也不明白他是谁竟如此设想。这样的人，不如接受那崇高且毫无欺骗的权威，它告诉我们：只要我们仍在肉身内，便是与主远离，我们因信德而非因目睹而行走。因此，他应恒久持守并守护自己的信德、望德和爱德，仰望那应许的景象，按我们所领受的圣神的抵押，祂将教导我们一切真理，因为那使耶稣基督从死者中复活的天主，也必藉住在我们内的圣神使我们有死的身体活过来。但在因罪而死的身体活过来之前，它无疑是可朽坏的，并压着灵魂。如果人在肉身中曾得以超越覆盖整个大地的云雾——即覆盖整个尘世生活的肉性黑暗——那只不过像被一道闪电快速触碰，随即又沉入自身的软弱，只留下渴求，使他再次激动（趋向邪恶），而纯洁并不足以坚定他（于善）。然而，任何人越能做到这一点，他就越大；越不能做到，就越小。如果一个人的心智尚未有这种体验——尽管在此心智中，基督仍藉信德居住——他应恳切地努力减少此世的私欲，并藉修德行善而终结它们，如同与这三位福音作者一起行走，与中保基督同行。并且，怀着极大希望之喜乐，让他以信德持守那位永远是圣子、但因我们而成为人子的基督，使祂永恒的能力和天主性能与我们的软弱和必死性结合，并基于我们的人性，在祂内并藉祂为我们开辟道路。为使一个人避免犯罪，他应受基督君王的统治。如果他碰巧犯了罪，他可以从基督（他也是司铎）那里获得赦免。这样，在善言善行的操练中成长，并藉双重之爱的翅膀（如同两只强劲的羽翼）脱离尘世的气氛，他就可能被同一位基督——祂也是\OriginalTerm{圣言}{Word}，太初已有的\OriginalTerm{圣言}{Word}，与天主同在的\OriginalTerm{圣言}{Word}，和天主的\OriginalTerm{圣言}{Word}——所光照；虽然这仍是藉着镜子模糊地观看，但这将是一种远比任何形体相似物崇高的光照。因此，尽管前三部福音作者以行动之德的天赋为主，而若望福音以默观之德的天赋辨识这些主题，但若望福音本身也仅是部分，它将如此存留直到那圆满的来临。事实上，圣神赐给一人智慧的言语，赐给另一人知识的言语，同一圣神运作。一人为主而看待日子；另一人从主的怀中更清地领受；还有一人被提到第三层天，听见不可言传的话语。但所有人在肉身内都与主远离。凡活在美好望德中、名字写在生命册上的所有信徒，都还保留着那应许：\BibleQuote{我要爱他，并要向我显示我自己。}然而，一个人在这远离主的时期中，在对这主题的理解和知识上进步越大，就越应谨慎提防骄傲和嫉妒这些魔鬼般的恶习。让他记住，正是这部若望福音，如此特别地激励我们默观真理，也同样显著地强调爱德甘饴的教导。让他也思考那最真实有益的诫命：\BibleQuote{你越大，在一切事上越要谦卑自己。}因为那位以比其他所有福音作者更崇高的教导方式将基督呈现给我们的福音作者，也是在其叙述中，主为门徒洗脚的那一位。\par

\SourceRecord{Translated by S.D.F. Salmond.；Nicene and Post-Nicene Fathers, First Series；Vol. 6.；Edited by Philip Schaff.；Buffalo, NY: Christian Literature Publishing Co.,；1888.；<http://www.newadvent.org/fathers/1602410.htm>.}
